% Generated mechanically from frozen input inputs/p09/ja/src/spaces-more-morphisms.tex.
% Do not edit this generated file; edit the bound locale source instead.

% OK, start here.
%

\setcounter{chapter}{75}
\chapter{代数空間の射についてさらに}



\phantomsection
\label{spaces-more-morphisms-section-phantom}


\section{はじめに}
\label{spaces-more-morphisms-section-introduction}

\noindent
本章では、代数空間の射の性質についての研究を続ける。
基本的な参考文献は \cite{Kn} である。





\section{規約}
\label{spaces-more-morphisms-section-conventions}

\noindent
すべてのスキームは大きな fppf サイト $\Sch_{fppf}$ に含まれるものと
常に仮定する。また、ここで考えるすべての環 $A$ は、$\Spec(A)$ が
この大きなサイトの対象と（同型を除いて）なるという性質をもつものとする。

\medskip\noindent
$S$ をスキームとし、$X$ を $S$ 上の代数空間とする。本章および次章では、
$X \times_S X$ と書く。この記号は $X$ とそれ自身との
（$S$ 上の代数空間の圏における）積を表し、$X \times X$ の代わりに用いる。








\section{Radicial 射}
\label{spaces-more-morphisms-section-radicial}

\noindent
スキームの射の場合とは異なり、radicial 射は普遍単射と同じものではない。
実際、radicial であることの方が少し強い。

\begin{definition}
\label{spaces-more-morphisms-definition-radicial}
$S$ をスキームとする。$f : X \to Y$ を $S$ 上の代数空間の射とする。
$f$ は、任意の射 $\Spec(K) \to Y$（ここで $K$ は体）に対し、被約化
$(\Spec(K) \times_Y X)_{red}$ が空であるか、または $K$ の純非分離
体拡大のスペクトルで表現されるとき、その射は {\it radicial} であるという。
\end{definition}

\begin{lemma}
\label{spaces-more-morphisms-lemma-radicial-implies-universally-injective}
代数空間の radicial 射は普遍単射である。
\end{lemma}

\begin{proof}
$S$ をスキームとし、$f : X \to Y$ を $S$ 上の代数空間の
radicial 射とする。定義から、射 $\Spec(K) \to Y$ が与えられたとき、
これを $X$ への射に持ち上げる方法は高々一つであることが明らかである。
したがって
Morphisms of Spaces,
Lemma \ref{spaces-morphisms-lemma-universally-injective}
により、$f$ は普遍単射である。
\end{proof}




\begin{example}
\label{spaces-more-morphisms-example-universally-injective-not-radicial}
普遍単射であることと radicial であることは、もはや同値ではない。
例えば、Spaces, Example \ref{spaces-example-Qbar} の射
$$
X = [\Spec(\overline{\mathbf{Q}})/
\text{Gal}(\overline{\mathbf{Q}}/\mathbf{Q})]
\longrightarrow
S = \Spec(\mathbf{Q})
$$
は普遍単射であるが、上の意味で radicial ではない。
\end{example}

\noindent
それでも、逆向きの含意が成り立つ場合は多い。

\begin{lemma}
\label{spaces-more-morphisms-lemma-when-universally-injective-radicial}
$S$ をスキームとする。$f : X \to Y$ を $S$ 上の代数空間の
普遍単射とする。
\begin{enumerate}
\item $f$ が decent ならば、$f$ は radicial である。
\item $f$ が準分離ならば、$f$ は radicial である。
\item $f$ が局所分離ならば、$f$ は radicial である。
\end{enumerate}
\end{lemma}


\begin{proof}
$\mathcal{P}$ を代数空間の射の性質であって、基底変換と合成で保たれ、
閉埋め込みについて成り立つものとする。$f : X \to Y$ は
$\mathcal{P}$ をもち、かつ普遍単射であると仮定する。このとき
Definition \ref{spaces-more-morphisms-definition-radicial} の状況では、射
$(\Spec(K) \times_Y X)_{red} \to \Spec(K)$
は普遍単射であり、性質 $\mathcal{P}$ をもつ。したがって
$$
\mathcal{P} + \text{普遍単射}
\Rightarrow
\text{radicial}
$$
を示す問題は、ある体上の空でない被約代数空間 $X$ で、その構造射
$X \to \Spec(K)$ が普遍単射かつ $\mathcal{P}$ をもつものは、
ある体のスペクトルで表現されることを示す問題に帰着する。% P09-ERR-1048
実際、その場合 $X \to \Spec(K)$ はスキームの射となり、
スキームの射について radicial と普遍単射が同値であること、
すなわち
Morphisms, Lemma \ref{morphisms-lemma-universally-injective}
から結論が従う。

\medskip\noindent
(1) を示す。$f$ が decent かつ普遍単射であると仮定する。
Decent Spaces,
Lemmas \ref{decent-spaces-lemma-base-change-relative-conditions},
\ref{decent-spaces-lemma-composition-relative-conditions}, and
\ref{decent-spaces-lemma-properties-trivial-implications}
（埋め込みが decent であることを確認するため）により、
第一段落の議論を適用できる。$X$ を、体 $K$ 上で普遍単射な、
空でない decent 被約代数空間とする。特に $|X|$ は一点集合である。
Decent Spaces, Lemma \ref{decent-spaces-lemma-when-field}
により、$X \cong \Spec(L)$ となるある拡大
$K \subset L$ が存在し、望む結論を得る。

\medskip\noindent
準分離射は decent である。次を参照せよ。
Decent Spaces,
Lemma \ref{decent-spaces-lemma-properties-trivial-implications}。
したがって (1) から (2) が従う。

\medskip\noindent
(3) を示す。分離公理は基底変換と合成で保たれ、閉埋め込みは分離的である。
次を参照せよ。
Morphisms of Spaces,
Lemmas \ref{spaces-morphisms-lemma-base-change-separated},
\ref{spaces-morphisms-lemma-composition-separated}, and
\ref{spaces-morphisms-lemma-immersions-monomorphisms}。
したがって、証明の第一段落の議論を適用できる。
$X$ を体 $K$ 上で普遍単射かつ局所分離な被約代数空間とする。
特に $|X|$ は一点集合なので $X$ は準コンパクトである。次を参照せよ。
Properties of Spaces, Lemma \ref{spaces-properties-lemma-quasi-compact-space}。
全射エタール射 $U \to X$ で $U$ がアフィンとなるものを取れる。次を参照せよ。
Properties of Spaces,
Lemma \ref{spaces-properties-lemma-quasi-compact-affine-cover}。
スキームの射
$$
j :
U \times_X U
\longrightarrow
U \times_{\Spec(K)} U
$$
を考える。$X \to \Spec(K)$ は普遍単射なので $j$ は全射であり、
$X \to \Spec(K)$ は局所分離なので $j$ は埋め込みである。
全射な埋め込みは閉埋め込みである。次を参照せよ。
Schemes, Lemma \ref{schemes-lemma-immersion-when-closed}。
したがって $R = U \times_X U$ はアフィンスキームの閉部分スキームとして
アフィンである。特に $R$ は準コンパクトである。
よって $X = U/R$ は準分離であり、結論は (2) から従う。
\end{proof}

\begin{remark}
\label{spaces-more-morphisms-remark-weakly-radicial}
$X \to Y$ を代数空間の射とする。radicial 射の応用によっては、
任意の $\Spec(K) \to Y$（ここで $K$ は体）に対して
\begin{enumerate}
\item 空間 $|\Spec(K) \times_Y X|$ が一点集合である、
\item モノ射 $\Spec(L) \to \Spec(K) \times_Y X$ が存在する、
\item $K \subset L$ が純非分離である、
\end{enumerate}
ことを要求すれば十分である。必要になれば、後でこのような射を
{\it 弱 radicial} と呼ぶことにする。% P09-ERR-1047
例えば $X \to Y$ が全射な弱 radicial 射ならば、
$X(k) \to Y(k)$ は任意の代数閉体 $k$ に対して全射である。Example
\ref{spaces-more-morphisms-example-universally-injective-not-radicial} の射の基底変換
$X_{\overline{\mathbf{Q}}} \to \Spec(\overline{\mathbf{Q}})$
は弱 radicial であるが、radicial ではないことに注意する。
Lemma \ref{spaces-more-morphisms-lemma-when-universally-injective-radicial}
に対応する主張は、$X \to Y$ が性質 ($\beta$) をもち普遍単射ならば
弱 radicial である、というものである（証明は省略する）。
\end{remark}

\begin{lemma}
\label{spaces-more-morphisms-lemma-check-universally-injective}
$S$ をスキームとする。$f : X \to Y$ を $S$ 上の代数空間の射とする。
次を仮定する。
\begin{enumerate}
\item $f$ は局所有限型である。
\item 任意のエタール射 $V \to Y$ に対して、写像
$|X \times_Y V| \to |V|$ は単射である。
\end{enumerate}
このとき $f$ は普遍単射である。
\end{lemma}

\begin{proof}
この問題は
Morphisms of Spaces, Lemma
\ref{spaces-morphisms-lemma-universally-injective-local}
により $Y$ 上エタール局所的である。したがって $Y$ はスキームであると
仮定してよい。このとき $Y$ は特に decent であり、
Decent Spaces, Lemma
\ref{decent-spaces-lemma-conditions-on-point-in-fibre-and-qf}
により、$f$ は局所準有限である。$y \in Y$ を点とし、
$X_y$ をスキーム論的ファイバーとする。$X_y$ が空でないと仮定する。
Spaces over Fields, Lemma
\ref{spaces-over-fields-lemma-locally-quasi-finite-over-field}
により、$X_y$ は $\kappa(y)$ 上局所準有限なスキームである。
$|X_y| \subset |X|$ は $|X| \to |Y|$ の $y$ 上のファイバーなので、
$X_y$ はただ一つの点 $x$ をもつ。同じことが
$X_y \times_{\Spec(\kappa(y))} \Spec(k)$
についても成り立つ。ここで $k/\kappa(y)$ は任意の有限分離拡大とする。
実際、$k$ は $y$ の上にある点の剰余体として実現でき、
その点は $Y$ 上のあるエタールスキームに属する。次を参照せよ。
More on Morphisms, Lemma
\ref{more-morphisms-lemma-realize-prescribed-residue-field-extension-etale}。
したがって $X_y$ は幾何学的に連結である。次を参照せよ。
Varieties, Lemma \ref{varieties-lemma-characterize-geometrically-disconnected}。
このことから有限拡大 $\kappa(x)/\kappa(y)$ は純非分離である。

\medskip\noindent
以上により（$Y$ がスキームの場合）、任意の $y \in Y$ に対して
ファイバー $X_y$ は空であるか、または
$(X_y)_{red} = \Spec(\kappa(x))$ であり、
$\kappa(y) \subset \kappa(x)$ は純非分離である。
したがって $f$ は radicial である（細部は省略する）。
ゆえに Lemma \ref{spaces-more-morphisms-lemma-radicial-implies-universally-injective}
により普遍単射である。
\end{proof}




\section{モノ射}
\label{spaces-more-morphisms-section-monomorphisms}

\noindent
本節は
Morphisms of Spaces, Section \ref{spaces-morphisms-section-monomorphisms}
の続きである。代数空間のすべてのモノ射が表現可能かどうかを知りたい。
これが真であることを証明できるか、反例があるならば、
\href{mailto:stacks.project@gmail.com}{stacks.project@gmail.com}
まで電子メールで知らせてほしい。現時点では、次の場合に知られている。
\begin{enumerate}
\item 局所有限型のモノ射の場合。より一般に、分離かつ局所準有限な任意の射は
Morphisms of Spaces, Lemma
\ref{spaces-morphisms-lemma-locally-quasi-finite-separated-representable}
により表現可能であり、局所有限型のモノ射は
Morphisms of Spaces, Lemma
\ref{spaces-morphisms-lemma-monomorphism-loc-finite-type-loc-quasi-finite}
により局所準有限である。
\item 標的が零次元局所環のスペクトルの非交和である場合
（Decent Spaces, Lemma
\ref{decent-spaces-lemma-monomorphism-toward-disjoint-union-dim-0-rings}）。
\item 平坦モノ射の場合（以下を参照）。
\end{enumerate}

\begin{lemma}[David Rydh]
\label{spaces-more-morphisms-lemma-flat-case}
代数空間の平坦モノ射はスキームで表現可能である。
\end{lemma}

\begin{proof}
$f : X \to Y$ を代数空間の平坦モノ射とする。$f$ が表現可能であることを
示すには、$X \times_Y V$ がスキームであることを示さなければならない。
ここで $V$ は $Y$ に写る任意のスキームである。
スキームであるという性質は局所的なので（Properties of Spaces,
Lemma \ref{spaces-properties-lemma-subscheme}）、$V$ はアフィンであると
仮定してよい。したがって $Y = \Spec(B)$ はアフィンスキームであると
仮定してよい。次に、$X$ を準コンパクト開部分で置き換えることにより、
$X$ は準コンパクトであると仮定できる。
空間 $X$ は、$X \to X \times_{\Spec(B)} X$ が同型なので分離的である。
Limits of Spaces, Lemma \ref{spaces-limits-lemma-enough-local}
を適用して、$B$ が局所環であり、$X \to \Spec(B)$ が平坦モノ射で、
$x \in X$ が存在し、それが $\Spec(B)$ の閉点に写る場合に帰着する。
このとき、分離的代数空間の平坦射に沿って一般化は持ち上がるので、
$X \to \Spec(B)$ は全射である。次を参照せよ。
Decent Spaces, Lemma \ref{decent-spaces-lemma-generalizations-lift-flat}。
したがって $\{X \to \Spec(B)\}$ は fpqc 被覆である。
射 $X \to \Spec(B)$ は、それ自身 $X \to \Spec(B)$ による基底変換の後に
同型になる。よって fpqc 降下により同型である。次を参照せよ。
Descent on Spaces, Lemma
\ref{spaces-descent-lemma-descending-property-isomorphism}。
\end{proof}

\noindent
以下は（ある意味で）上の補題の変種である。

\begin{lemma}
\label{spaces-more-morphisms-lemma-ui-case}
$S$ をスキームとする。$f : X \to Y$ を代数空間の準コンパクトなモノ射とし、
任意の $T \to Y$ に対して写像
$$
\mathcal{O}_T \to f_{T,*}\mathcal{O}_{X \times_Y T}
$$
が単射であると仮定する。このとき $f$ は同型である
（したがってスキームで表現可能である）。
\end{lemma}

\begin{proof}
問題は $Y$ 上エタール局所的であるから、$Y = \Spec(A)$ は
アフィンであると仮定してよい。このとき $X$ は準コンパクトであり、
アフィンスキーム $U = \Spec(B)$ と全射エタール射 $U \to X$ を選べる
（Properties of Spaces, Lemma
\ref{spaces-properties-lemma-quasi-compact-affine-cover}）。
$U \times_X U = \Spec(B \otimes_A B)$ であることに注意する。
したがって、準連接 $\mathcal{O}_X$ 加群の圏は、降下データの圏
$DD_{B/A}$、すなわち $A \to B$ に関する加群上の降下データの圏と
同値である。
Properties of Spaces, Proposition
\ref{spaces-properties-proposition-quasi-coherent},
Descent, Definition \ref{descent-definition-descent-datum-modules}、
および Descent, Subsection
\ref{descent-subsection-descent-modules-morphisms} を参照せよ。
他方、
$$
A \to B
$$
は普遍単射な環準同型である。実際、$A$ 加群 $M$ を与えると、
補題の仮定により
$A \oplus M \to B \otimes_A (A \oplus M)$
は単射である。したがって Descent, Theorem
\ref{descent-theorem-descent} により、$DD_{B/A}$ は
$A$ 加群の圏と同値である。ゆえに
$f : X \to \Spec(A)$ に沿う引き戻しは準連接加群の圏の同値を定める。
特に $f^*$ は準連接加群上完全であり、$f$ は平坦である
（細部は省略する）。さらに $f$ が全射であることは明らかである
（例えば $\Spec(B) \to \Spec(A)$ が全射だからである）。
したがって $\{X \to \Spec(A)\}$ は fpqc 被覆である。
射 $X \to \Spec(A)$ は、それ自身 $X \to \Spec(A)$ による
基底変換の後に同型になる。よって fpqc 降下により同型である。
Descent on Spaces, Lemma
\ref{spaces-descent-lemma-descending-property-isomorphism} を参照せよ。
\end{proof}

\begin{lemma}
\label{spaces-more-morphisms-lemma-flat-surjective-monomorphism}
代数空間の準コンパクト、平坦かつ全射なモノ射は同型である。
\end{lemma}

\begin{proof}
そのような射は Lemma \ref{spaces-more-morphisms-lemma-ui-case} の仮定を満たす。
\end{proof}







\section{埋め込みの余法層}
\label{spaces-more-morphisms-section-conormal-sheaf}

\noindent
$S$ をスキームとする。$i : Z \to X$ を $S$ 上の代数空間の
閉埋め込みとする。対応する準連接イデアル層を
$\mathcal{I} \subset \mathcal{O}_X$ とする。次を参照せよ。
Morphisms of Spaces,
Lemma \ref{spaces-morphisms-lemma-closed-immersion-ideals}。
短完全列
$$
0 \to \mathcal{I}^2 \to \mathcal{I} \to \mathcal{I}/\mathcal{I}^2 \to 0
$$
を考える。これは $X$ 上の準連接層の短完全列である。
層 $\mathcal{I}/\mathcal{I}^2$ は $\mathcal{I}$ により
零化されるので、Morphisms of Spaces, Lemma
\ref{spaces-morphisms-lemma-i-star-equivalence}
により $Z$ 上の層に対応する。この準連接 $\mathcal{O}_Z$ 加群を
{\it $Z$ の $X$ における余法層}といい、Morphisms of Spaces,
Section \ref{spaces-morphisms-section-closed-immersions-quasi-coherent}
で述べた記法の濫用により、しばしば
$\mathcal{I}/\mathcal{I}^2$ と表す。

\medskip\noindent
$i : Z \to X$ が（局所閉）埋め込みの場合、$i$ の余法層を
閉埋め込み $i : Z \to X \setminus \partial Z$ の余法層として定義する。
Morphisms of Spaces, Remark
\ref{spaces-morphisms-remark-immersion} を参照せよ。
これはしばしば $\mathcal{I}/\mathcal{I}^2$ と表す。ここで
$\mathcal{I}$ は閉埋め込み
$i : Z \to X \setminus \partial Z$ のイデアル層である。

\begin{definition}
\label{spaces-more-morphisms-definition-conormal-sheaf}
$i : Z \to X$ を埋め込みとする。{\it 余法層
$\mathcal{C}_{Z/X}$、すなわち $Z$ の $X$ における余法層}、
または {\it $i$ の余法層}とは、上で述べた準連接
$\mathcal{O}_Z$ 加群 $\mathcal{I}/\mathcal{I}^2$ のことである。
\end{definition}

\noindent
\cite[IV Definition 16.1.2]{EGA} では、この層を
$\mathcal{N}_{Z/X}$ と表す。本章では
$\mathcal{N}_{Z/X}$ という記号を{\it 埋め込みの法層}のために
取っておきたいので、この慣習には従わない。余法層が有限表示である場合、
法層を
$$
\mathcal{N}_{Z/X} =
\SheafHom_{\mathcal{O}_Z}(\mathcal{C}_{Z/X}, \mathcal{O}_Z) =
\SheafHom_{\mathcal{O}_Z}(\mathcal{I}/\mathcal{I}^2, \mathcal{O}_Z)
$$
と定義する（そうでなければ、法層は準連接でさえないことがある）。
法層には後で戻る（ここに将来の参照を挿入する）。

\begin{lemma}
\label{spaces-more-morphisms-lemma-etale-conormal}
$S$ をスキームとする。$i : Z \to X$ を埋め込みとする。
$\varphi : U \to X$ をエタール射とし、$U$ はスキームであるとする。
$Z_U = U \times_X Z$ とおく。これは $U$ の局所閉部分スキームである。
このとき
$$
\mathcal{C}_{Z/X}|_{Z_U} = \mathcal{C}_{Z_U/U}
$$
が標準的に、かつ $U$ に関して関手的に成り立つ。
\end{lemma}

\begin{proof}
閉部分空間 $T \subset X$ を、$i$ が $X \setminus T$ への
閉埋め込みを定めるようにとる。$\mathcal{I}$ を
$X \setminus T$ 上で $Z$ を定める準連接イデアル層とする。
このとき補題の主張は、
$\mathcal{I}|_{U \setminus \varphi^{-1}(T)}$ が埋め込み
$Z_U \to U \setminus \varphi^{-1}(T)$ のイデアル層であるということに
すぎない。これは Morphisms of Spaces, Lemma
\ref{spaces-morphisms-lemma-closed-immersion-ideals}
における $\mathcal{I}$ の構成から明らかである。
\end{proof}

\begin{lemma}
\label{spaces-more-morphisms-lemma-conormal-functorial}
$S$ をスキームとする。
$$
\xymatrix{
Z \ar[r]_i \ar[d]_f & X \ar[d]^g \\
Z' \ar[r]^{i'} & X'
}
$$
を $S$ 上の代数空間の可換図式とする。$i$, $i'$ は埋め込みであると
仮定する。このとき $\mathcal{O}_Z$ 加群の標準的な写像
$$
f^*\mathcal{C}_{Z'/X'}
\longrightarrow
\mathcal{C}_{Z/X}
$$
が存在する。
\end{lemma}

\begin{proof}
まず、開部分空間 $U' \subset X'$ と $U \subset X$ であって、
$g(U) \subset U'$ であり、$i(Z) \subset U$ および
% P09-ERR-1049: frozen source line 475 applies i to Z'; expected i'(Z').
$i(Z') \subset U'$ が閉となるものを見いだす
（存在証明は省略する）。$X$ を $U$ で、$X'$ を $U'$ で置き換えれば、
$i$ と $i'$ は閉埋め込みであると仮定してよい。
$\mathcal{I}' \subset \mathcal{O}_{X'}$ および
$\mathcal{I} \subset \mathcal{O}_X$ を、それぞれ $i'$ と $i$ に
付随する準連接イデアル層とする。次を参照せよ。
Morphisms of Spaces, Lemma
\ref{spaces-morphisms-lemma-closed-immersion-ideals}。
合成
$$
g^{-1}\mathcal{I}' \to g^{-1}\mathcal{O}_{X'}
\xrightarrow{g^\sharp} \mathcal{O}_X \to
\mathcal{O}_X/\mathcal{I} = i_*\mathcal{O}_Z
$$
を考える。$g(i(Z)) \subset Z'$ なので、この合成は零である
（Morphisms of Spaces, Lemma
\ref{spaces-morphisms-lemma-closed-immersion-ideals}
の分解に関する主張を参照せよ）。したがって可換図式
$$
\xymatrix{
0 \ar[r] &
\mathcal{I} \ar[r] &
\mathcal{O}_X \ar[r] &
i_*\mathcal{O}_Z \ar[r] &
0 \\
0 \ar[r] &
g^{-1}\mathcal{I}' \ar[r] \ar[u] &
g^{-1}\mathcal{O}_{X'} \ar[r] \ar[u] &
g^{-1}i'_*\mathcal{O}_{Z'} \ar[r] \ar[u] &
0
}
$$
を得る。$g^{-1}$ は完全関手なので下段は完全である。
完全性から
$(g^{-1}\mathcal{I}')^2 = g^{-1}((\mathcal{I}')^2)$
も従う。ゆえに、この図式は写像
$g^{-1}(\mathcal{I}'/(\mathcal{I}')^2) \to
\mathcal{I}/\mathcal{I}^2$
を誘導する。（例えば $i^{-1}$ を用いて）$Z$ に引き戻すと、
$i^{-1}g^{-1}(\mathcal{I}'/(\mathcal{I}')^2)
\to \mathcal{C}_{Z/X}$ を得る。
$i^{-1}g^{-1} = f^{-1}(i')^{-1}$ なので、これは写像
$f^{-1}\mathcal{C}_{Z'/X'} \to \mathcal{C}_{Z/X}$ を与え、
そこから所望の写像が誘導される。
\end{proof}

\begin{lemma}
\label{spaces-more-morphisms-lemma-conormal-functorial-more}
$S$ をスキームとする。
Definition \ref{spaces-more-morphisms-definition-conormal-sheaf} の余法層と
Lemma \ref{spaces-more-morphisms-lemma-conormal-functorial} の関手性は、次の性質を満たす。
\begin{enumerate}
\item $Z \to X$ が $S$ 上のスキームの埋め込みならば、その余法層は
Morphisms, Definition
\ref{morphisms-definition-conormal-sheaf} のものと一致する。
\item Lemma \ref{spaces-more-morphisms-lemma-conormal-functorial} においてすべての空間が
スキームならば、写像
$f^*\mathcal{C}_{Z'/X'} \to \mathcal{C}_{Z/X}$ は
Morphisms, Lemma
\ref{morphisms-lemma-conormal-functorial} で構成したものと同じである。
\item 可換図式
$$
\xymatrix{
Z \ar[r]_i \ar[d]_f & X \ar[d]^g \\
Z' \ar[r]^{i'} \ar[d]_{f'} & X' \ar[d]^{g'} \\
Z'' \ar[r]^{i''} & X''
}
$$
が与えられたとき、写像
$(f' \circ f)^*\mathcal{C}_{Z''/X''} \to \mathcal{C}_{Z/X}$
は、写像
$f^*\mathcal{C}_{Z'/X'} \to \mathcal{C}_{Z/X}$ と、
$f$ による次の写像の引き戻しとの合成に等しい。
$(f')^*\mathcal{C}_{Z''/X''} \to \mathcal{C}_{Z'/X'}$
\end{enumerate}
\end{lemma}

\begin{proof}
省略する。なお、(1) は Lemma \ref{spaces-more-morphisms-lemma-etale-conormal} の特別な場合である。
\end{proof}

\begin{lemma}
\label{spaces-more-morphisms-lemma-conormal-functorial-flat}
$S$ をスキームとする。
$$
\xymatrix{
Z \ar[r]_i \ar[d]_f & X \ar[d]^g \\
Z' \ar[r]^{i'} & X'
}
$$
を $S$ 上の代数空間のファイバー積図式とする。$i$, $i'$ は
埋め込みであると仮定する。このとき Lemma
\ref{spaces-more-morphisms-lemma-conormal-functorial} の標準的な写像
$f^*\mathcal{C}_{Z'/X'} \to \mathcal{C}_{Z/X}$ は全射である。
$g$ が平坦ならば、この写像は同型である。
\end{lemma}

\begin{proof}
可換図式
$$
\xymatrix{
U \ar[r] \ar[d] & X \ar[d] \\
U' \ar[r] & X'
}
$$
を選ぶ。ここで $U$, $U'$ はスキームであり、横の矢印は全射かつ
エタールである。Spaces, Lemma
\ref{spaces-lemma-lift-morphism-presentations} を参照せよ。
このとき Lemmas \ref{spaces-more-morphisms-lemma-etale-conormal} および
\ref{spaces-more-morphisms-lemma-conormal-functorial-more} により、問題はスキームの射の場合に
帰着することが分かる。スキームの場合、これは Morphisms, Lemma
\ref{morphisms-lemma-conormal-functorial-flat} である。
\end{proof}

\begin{lemma}
\label{spaces-more-morphisms-lemma-transitivity-conormal}
$S$ をスキームとする。$Z \to Y \to X$ を代数空間の埋め込みとする。
このとき標準的な完全列
$$
i^*\mathcal{C}_{Y/X} \to
\mathcal{C}_{Z/X} \to
\mathcal{C}_{Z/Y} \to 0
$$
が存在する。ここで写像は Lemma
\ref{spaces-more-morphisms-lemma-conormal-functorial} から得られ、
$i : Z \to Y$ は最初の射である。
\end{lemma}

\begin{proof}
$U$ をスキームとし、$U \to X$ を全射エタール射とする。
Lemmas \ref{spaces-more-morphisms-lemma-etale-conormal} および
\ref{spaces-more-morphisms-lemma-conormal-functorial-more} により、この列の完全性は
スキームの埋め込み
$Z \times_X U \to Y \times_X U \to U$
に対する対応する列の完全性へ直ちに移される。
したがって補題は Morphisms, Lemma
\ref{morphisms-lemma-transitivity-conormal} から従う。
\end{proof}








\section{埋め込みの法錐}
\label{spaces-more-morphisms-section-normal-cone}

\noindent
$S$ をスキームとする。$i : Z \to X$ を $S$ 上の代数空間の
閉埋め込みとする。対応する準連接イデアル層を
$\mathcal{I} \subset \mathcal{O}_X$ とする。次を参照せよ。
Morphisms of Spaces, Lemma
\ref{spaces-morphisms-lemma-closed-immersion-ideals}。
次数付き $\mathcal{O}_X$ 代数の準連接層
$\bigoplus_{n \geq 0} \mathcal{I}^n/\mathcal{I}^{n + 1}$
を考える。各層 $\mathcal{I}^n/\mathcal{I}^{n + 1}$ は
$\mathcal{I}$ により零化されるので、この次数付き代数は
Morphisms of Spaces, Lemma
\ref{spaces-morphisms-lemma-i-star-equivalence}
により、次数付き $\mathcal{O}_Z$ 代数の準連接層に対応する。
この準連接次数付き $\mathcal{O}_Z$ 代数を
{\it $Z$ の $X$ における余法代数}といい、Morphisms of Spaces,
Section \ref{spaces-morphisms-section-closed-immersions-quasi-coherent}
で述べた記法の濫用により、しばしば単に
$\bigoplus_{n \geq 0} \mathcal{I}^n/\mathcal{I}^{n + 1}$
と表す。

\medskip\noindent
$i : Z \to X$ が（局所閉）埋め込みの場合、$i$ の余法代数を
閉埋め込み $i : Z \to X \setminus \partial Z$ の余法代数として定義する。
Morphisms of Spaces, Remark
\ref{spaces-morphisms-remark-immersion} を参照せよ。
これはしばしば
$\bigoplus_{n \geq 0} \mathcal{I}^n/\mathcal{I}^{n + 1}$
と表す。ここで $\mathcal{I}$ は閉埋め込み
$i : Z \to X \setminus \partial Z$ のイデアル層である。

\begin{definition}
\label{spaces-more-morphisms-definition-conormal-algebra}
$i : Z \to X$ を埋め込みとする。{\it 余法代数
$\mathcal{C}_{Z/X, *}$、すなわち $Z$ の $X$ における余法代数}、
または {\it $i$ の余法代数}とは、上で述べた次数付き
$\mathcal{O}_Z$ 代数の準連接層
$\bigoplus_{n \geq 0} \mathcal{I}^n/\mathcal{I}^{n + 1}$ のことである。
\end{definition}

\noindent
したがって $\mathcal{C}_{Z/X, 1} = \mathcal{C}_{Z/X}$ は
この埋め込みの余法層である。また
$\mathcal{C}_{Z/X, 0} = \mathcal{O}_Z$ であり、
$\mathcal{C}_{Z/X, n}$ は次の性質で特徴づけられる準連接
$\mathcal{O}_Z$ 加群である。
\begin{equation}
\label{spaces-more-morphisms-equation-conormal-in-degree-n}
i_*\mathcal{C}_{Z/X, n} = \mathcal{I}^n/\mathcal{I}^{n + 1}
\end{equation}
ここで $i : Z \to X \setminus \partial Z$ であり、
$\mathcal{I}$ は上と同じく $i$ のイデアル層である。
最後に、準連接次数付き $\mathcal{O}_Z$ 代数の標準的全射
\begin{equation}
\label{spaces-more-morphisms-equation-conormal-algebra-quotient}
\text{Sym}^*(\mathcal{C}_{Z/X}) \longrightarrow \mathcal{C}_{Z/X, *}
\end{equation}
が存在し、次数 $0$ および $1$ では同型であることに注意する。

\begin{lemma}
\label{spaces-more-morphisms-lemma-etale-conormal-algebra}
$S$ をスキームとする。$i : Z \to X$ を $S$ 上の代数空間の
埋め込みとする。$\varphi : U \to X$ をエタール射とし、
$U$ はスキームであるとする。$Z_U = U \times_X Z$ とおく。
これは $U$ の局所閉部分スキームである。このとき
$$
\mathcal{C}_{Z/X, *}|_{Z_U} = \mathcal{C}_{Z_U/U, *}
$$
が標準的に、かつ $U$ に関して関手的に成り立つ。
\end{lemma}

\begin{proof}
閉部分空間 $T \subset X$ を、$i$ が $X \setminus T$ への
閉埋め込みを定めるようにとる。$\mathcal{I}$ を
$X \setminus T$ 上で $Z$ を定める準連接イデアル層とする。
このとき補題は、$\mathcal{I}|_{U \setminus \varphi^{-1}(T)}$ が
埋め込み $Z_U \to U \setminus \varphi^{-1}(T)$ のイデアル層であること
から従う。これは Morphisms of Spaces, Lemma
\ref{spaces-morphisms-lemma-closed-immersion-ideals}
における $\mathcal{I}$ の構成から明らかである。
\end{proof}

\begin{lemma}
\label{spaces-more-morphisms-lemma-conormal-algebra-functorial}
$S$ をスキームとする。
$$
\xymatrix{
Z \ar[r]_i \ar[d]_f & X \ar[d]^g \\
Z' \ar[r]^{i'} & X'
}
$$
を $S$ 上の代数空間の可換図式とする。$i$, $i'$ は埋め込みであると
仮定する。このとき次数付き $\mathcal{O}_Z$ 代数の標準的な写像
$$
f^*\mathcal{C}_{Z'/X', *}
\longrightarrow
\mathcal{C}_{Z/X, *}
$$
が存在する。
\end{lemma}

\begin{proof}
まず、開部分空間 $U' \subset X'$ と $U \subset X$ であって、
$g(U) \subset U'$ であり、$i(Z) \subset U$ および
% P09-ERR-1050: frozen source line 724 applies i to Z'; expected i'(Z').
$i(Z') \subset U'$ が閉となるものを見いだす
（存在証明は省略する）。$X$ を $U$ で、$X'$ を $U'$ で置き換えれば、
$i$ と $i'$ は閉埋め込みであると仮定してよい。
$\mathcal{I}' \subset \mathcal{O}_{X'}$ および
$\mathcal{I} \subset \mathcal{O}_X$ を、それぞれ $i'$ と $i$ に
付随する準連接イデアル層とする。次を参照せよ。
Morphisms of Spaces, Lemma
\ref{spaces-morphisms-lemma-closed-immersion-ideals}。
合成
$$
g^{-1}\mathcal{I}' \to g^{-1}\mathcal{O}_{X'}
\xrightarrow{g^\sharp} \mathcal{O}_X \to
\mathcal{O}_X/\mathcal{I} = i_*\mathcal{O}_Z
$$
を考える。$g(i(Z)) \subset Z'$ なので、この合成は零である
（Morphisms of Spaces, Lemma
\ref{spaces-morphisms-lemma-closed-immersion-ideals}
の分解に関する主張を参照せよ）。したがって可換図式
$$
\xymatrix{
0 \ar[r] &
\mathcal{I} \ar[r] &
\mathcal{O}_X \ar[r] &
i_*\mathcal{O}_Z \ar[r] &
0 \\
0 \ar[r] &
g^{-1}\mathcal{I}' \ar[r] \ar[u] &
g^{-1}\mathcal{O}_{X'} \ar[r] \ar[u] &
g^{-1}i'_*\mathcal{O}_{Z'} \ar[r] \ar[u] &
0
}
$$
を得る。$g^{-1}$ は完全関手なので下段は完全である。完全性から、
$(g^{-1}\mathcal{I}')^n = g^{-1}((\mathcal{I}')^n)$
がすべての $n \geq 1$ に対して成り立つことも従う。
ゆえに、この図式は写像
$g^{-1}((\mathcal{I}')^n/(\mathcal{I}')^{n + 1}) \to
\mathcal{I}^n/\mathcal{I}^{n + 1}$
を誘導する。（例えば $i^{-1}$ を用いて）$Z$ に引き戻すと、
$i^{-1}g^{-1}((\mathcal{I}')^n/(\mathcal{I}')^{n + 1}) \to
\mathcal{C}_{Z/X, n}$ を得る。
$i^{-1}g^{-1} = f^{-1}(i')^{-1}$ なので、これは写像
$f^{-1}\mathcal{C}_{Z'/X', n} \to \mathcal{C}_{Z/X, n}$ を与え、
そこから所望の写像が誘導される。
\end{proof}

\begin{lemma}
\label{spaces-more-morphisms-lemma-conormal-algebra-functorial-flat}
$S$ をスキームとする。
$$
\xymatrix{
Z \ar[r]_i \ar[d]_f & X \ar[d]^g \\
Z' \ar[r]^{i'} & X'
}
$$
を $S$ 上の代数空間のカルテジアン正方形とし、$i$, $i'$ は
埋め込みであるとする。このとき Lemma
\ref{spaces-more-morphisms-lemma-conormal-algebra-functorial} の標準的な写像
$f^*\mathcal{C}_{Z'/X', *} \to \mathcal{C}_{Z/X, *}$ は全射である。
$g$ が平坦ならば、この写像は同型である。
\end{lemma}

\begin{proof}
$X'$ 上エタール局所化した後で主張を確認してよい。
% P09-ERR-1051: frozen source line 788 reverses the displayed direction X' -> X.
この場合、$X' \to X$ はスキームの射であると仮定してよく、
したがって $Z$ と $Z'$ はスキームである。ゆえに結果は
スキームの場合から従う。次を参照せよ。
Divisors, Lemma
\ref{divisors-lemma-conormal-algebra-functorial-flat}。
\end{proof}

\noindent
代数空間上の錐とベクトル束について、スキームの場合と同じ慣習を用いる
（スキームの場合は EGA の慣習を用いる）。次を参照せよ。
Constructions, Sections \ref{constructions-section-cone} および
\ref{constructions-section-vector-bundle}。
特に、ベクトル束は非常に一般的な対象である
（アフィン空間束に局所同型であるとは限らない）。

\begin{definition}
\label{spaces-more-morphisms-definition-normal-cone}
$S$ をスキームとする。$i : Z \to X$ を $S$ 上の代数空間の
埋め込みとする。{\it 法錐 $C_ZX$、すなわち $Z$ の $X$ における法錐}とは
$$
C_ZX = \underline{\Spec}_Z(\mathcal{C}_{Z/X, *})
$$
のことである。Morphisms of Spaces, Definition
\ref{spaces-morphisms-definition-relative-spec} を参照せよ。
{\it $Z$ の $X$ における法束}とは、ベクトル束
$$
N_ZX = \underline{\Spec}_Z(\text{Sym}(\mathcal{C}_{Z/X}))
$$
のことである。
\end{definition}

\noindent
したがって $C_ZX \to Z$ は $Z$ 上の錐であり、
$N_ZX \to Z$ は $Z$ 上のベクトル束である。
さらに、次数付き代数の標準的全射
(\ref{spaces-more-morphisms-equation-conormal-algebra-quotient}) は、$Z$ 上の錐の
標準的閉埋め込み
\begin{equation}
\label{spaces-more-morphisms-equation-normal-cone-in-normal-bundle}
C_ZX \longrightarrow N_ZX
\end{equation}
を定める。










\section{射の微分層}
\label{spaces-more-morphisms-section-sheaf-differentials}

\noindent
読者には、可換代数の章の対応する節
（Algebra, Section \ref{algebra-section-differentials}）、
スキームの射の章の対応する節
（Morphisms, Section \ref{morphisms-section-sheaf-differentials}）、
ならびに Modules on Sites, Section
\ref{sites-modules-section-differentials} を見ることを勧める。
まず、スキームの射に対する微分層の概念が、小エタール（環付き）サイトの
対応する射に対するものと一致することを示す。

\medskip\noindent
次の補題を明確に述べるため、一時的に、$\mathcal{F}^a$ で
$\mathcal{O}_{X_\etale}$ 加群の層、すなわち準連接
$\mathcal{O}_X$ 加群 $\mathcal{F}$ でスキーム $X$ 上のものに
付随する層を表す記法へ戻る。
Descent, Definition \ref{descent-definition-structure-sheaf} を参照せよ。

\begin{lemma}
\label{spaces-more-morphisms-lemma-match-modules-differentials}
$f : X \to Y$ をスキームの射とする。
$f_{small} : X_\etale \to Y_\etale$ を、これに付随する小エタールサイトの
射とする。Descent, Remark
\ref{descent-remark-change-topologies-ringed} を参照せよ。
このとき、普遍導分と両立する標準的同型
$$
(\Omega_{X/Y})^a = \Omega_{X_\etale/Y_\etale}
$$
が存在する。ここで最初の加群は $X_\etale$ 上の層であり、
準連接 $\mathcal{O}_X$ 加群 $\Omega_{X/Y}$ に付随するものをいう。
Morphisms, Definition
\ref{morphisms-definition-sheaf-differentials} を参照せよ。
第二の加群は Modules on Sites, Definition
\ref{sites-modules-definition-module-differentials} のものである。
\end{lemma}

\begin{proof}
$h : U \to X$ をエタール射とする。この場合、自然な写像
$h^*\Omega_{X/Y} \to \Omega_{U/Y}$ は同型である。次を参照せよ。
More on Morphisms, Lemma
\ref{more-morphisms-lemma-sheaf-differentials-etale-localization}。
このことは、自然な $\mathcal{O}_{Y_\etale}$ 導分
$$
\text{d}^a : \mathcal{O}_{X_\etale} \longrightarrow (\Omega_{X/Y})^a
$$
が存在することを意味する。実際、$(\Omega_{X/Y})^a$ の
任意の対象 $U$（$X_\etale$ の対象）上での値が
$\Gamma(U, \Omega_{U/Y})$ と標準的に同一視されることを、いま見た。
普遍性をもつ導分
$\text{d}_{X/Y} :
\mathcal{O}_{X_\etale}
\to
\Omega_{X_\etale/Y_\etale}$
により、一意な $\mathcal{O}_{X_\etale}$ 線形写像
$c : \Omega_{X_\etale/Y_\etale} \to (\Omega_{X/Y})^a$
が存在し、等式
$\text{d}^a = c \circ \text{d}_{X/Y}$ を満たす。

\medskip\noindent
逆に、$\mathcal{F}$ を $\mathcal{O}_{X_\etale}$ 加群とし、
$D : \mathcal{O}_{X_\etale} \to \mathcal{F}$ を
$\mathcal{O}_{Y_\etale}$ 導分とする。このとき $D$ を
小ザリスキーサイト $X_{Zar}$、すなわち $X$ の小ザリスキーサイトに
制限すればよい。
$X_{Zar}$ 上の層は $X$ 上の層と一致するので
（Descent, Remark \ref{descent-remark-Zariski-site-space}）、
$D|_{X_{Zar}} : \mathcal{O}_X \to \mathcal{F}|_{X_{Zar}}$
は通常の $Y$ 導分にほかならない。したがって写像
$\psi : \Omega_{X/Y} \longrightarrow \mathcal{F}|_{X_{Zar}}$
を得て、これは
$D|_{X_{Zar}} = \psi \circ \text{d}$ を満たす。特に、これを
$\mathcal{F} = \Omega_{X_\etale/Y_\etale}$ に適用すると写像
$$
c' :
\Omega_{X/Y}
\longrightarrow
\Omega_{X_\etale/Y_\etale}|_{X_{Zar}}
$$
を得る。Descent, Remark
\ref{descent-remark-change-topologies-ringed} および Lemma
\ref{descent-lemma-compare-sites} で論じた環付きサイトの射
$\text{id}_{small, \etale, Zar} : X_\etale \to X_{Zar}$
を考える。制限関手
$\mathcal{F} \mapsto \mathcal{F}|_{X_{Zar}}$ は
$\text{id}_{small, \etale, Zar, *}$ に等しく、
$\text{id}_{small, \etale, Zar}^*$ は
$\text{id}_{small, \etale, Zar, *}$ の左随伴であり、さらに
$(\Omega_{X/Y})^a =
\text{id}_{small, \etale, Zar}^*\Omega_{X/Y}$
であるから、$c'$ は写像
$$
c'' :
(\Omega_{X/Y})^a
\longrightarrow
\Omega_{X_\etale/Y_\etale}.
$$
に随伴する。
% P09-ERR-1052: frozen source line 937 says c'' and c' are inverse despite incompatible types; expected c'' and c.
$c''$ と $c'$ は互いに逆であると主張する。
この主張により補題の証明は完了する。
これを見るには、$c''(\text{d}(f)) = \text{d}_{X/Y}(f)$ および
$c(\text{d}_{X/Y}(f)) = \text{d}(f)$ を示せば十分である。
ここで $f$ は $\mathcal{O}_X$ の、$X$ の開集合上の局所切断である。
確認は省略する。
\end{proof}

\noindent
これで次の定義への準備が整った。別の定義については
Remark \ref{spaces-more-morphisms-remark-alternative} を参照せよ。

\begin{definition}
\label{spaces-more-morphisms-definition-sheaf-differentials}
$S$ をスキームとする。$f : X \to Y$ を $S$ 上の代数空間の射とする。
{\it 微分層 $\Omega_{X/Y}$、すなわち $X$ の $Y$ 上の微分層}とは、
Modules on Sites, Definition
\ref{sites-modules-definition-sheaf-differentials} の意味での、
環付きトポスの射
$$
(f_{small}, f^\sharp) :
(X_\etale, \mathcal{O}_X)
\to
(Y_\etale, \mathcal{O}_Y)
$$
に対する微分層のことである。この射については
Properties of Spaces, Lemma
\ref{spaces-properties-lemma-morphism-ringed-topoi} を参照せよ。
{\it 普遍 $Y$ 導分}を
$\text{d}_{X/Y} : \mathcal{O}_X \to \Omega_{X/Y}$ と表す。
\end{definition}

\noindent
Lemma \ref{spaces-more-morphisms-lemma-match-modules-differentials} により、
$X$ と $Y$ が表現可能な場合、この定義は既存の概念と矛盾しない。
今後、$X$ と $Y$ が表現可能ならば、上で定義した微分層と
Morphisms, Definition
\ref{morphisms-definition-sheaf-differentials}
で定義したものを区別しない。これをスキームの射に対する通常の
微分加群と関係づけたい。鍵となる補題は次である。

\begin{lemma}
\label{spaces-more-morphisms-lemma-localize-differentials}
$S$ をスキームとする。$f : X \to Y$ を $S$ 上の代数空間の射とする。
任意の可換図式
$$
\xymatrix{
U \ar[d]_a \ar[r]_\psi & V \ar[d]^b \\
X \ar[r]^f & Y
}
$$
を考える。ここで縦の矢印は代数空間のエタール射である。このとき
$$
\Omega_{X/Y}|_{U_\etale} = \Omega_{U/V}
$$
である。特に $U$, $V$ がスキームならば、これはスキームの射
$U \to V$ の通常の微分層に等しい。
\end{lemma}

\begin{proof}
Properties of Spaces, Lemma
\ref{spaces-properties-lemma-etale-morphism-topoi} および Equation
(\ref{spaces-properties-equation-restrict}) により、
$X_\etale$ 上の層の $U_\etale$ への制限を
$a_{small}$ による引き戻しと考えてよい。$b$ についても同様である。
Modules on Sites, Lemma
\ref{sites-modules-lemma-localize-differentials} により
$$
\Omega_{X/Y}|_{U_\etale} =
\Omega_{\mathcal{O}_{U_\etale}/
a_{small}^{-1}f_{small}^{-1}\mathcal{O}_{Y_\etale}}
$$
を得る。
$a_{small}^{-1}f_{small}^{-1}\mathcal{O}_{Y_\etale}
= \psi_{small}^{-1}b_{small}^{-1}\mathcal{O}_{Y_\etale}
= \psi_{small}^{-1}\mathcal{O}_{V_\etale}$
なので、補題が成り立つ。
\end{proof}

\begin{lemma}
\label{spaces-more-morphisms-lemma-module-differentials-quasi-coherent}
$S$ をスキームとする。$f : X \to Y$ を $S$ 上の代数空間の射とする。
このとき $\Omega_{X/Y}$ は準連接 $\mathcal{O}_X$ 加群である。
\end{lemma}

\begin{proof}
Lemma \ref{spaces-more-morphisms-lemma-localize-differentials} のような図式を、
$a$ と $b$ が全射であり、$U$ と $V$ がスキームとなるように選ぶ。
すると $\Omega_{X/Y}|_U = \Omega_{U/V}$ であり、これは準連接である
（例えば Morphisms, Lemma
\ref{morphisms-lemma-differentials-diagonal} による）。
したがって Properties of Spaces, Lemma
\ref{spaces-properties-lemma-characterize-quasi-coherent}
により $\Omega_{X/Y}$ は準連接である。
\end{proof}

\begin{remark}
\label{spaces-more-morphisms-remark-alternative}
$\Omega_{X/Y}$ が準連接であることが分かったので、別の方法で構成することを
試みられる。例えば Properties of Spaces, Section
\ref{spaces-properties-section-quasi-coherent-presentation}
の結果を用いて、貼り合わせにより微分層を構成できる。
例えば $Y$ がスキームで、$U \to X$ がスキームから $X$ への
全射エタール射ならば、$\Omega_{U/Y}$ は準連接
$\mathcal{O}_U$ 加群であり、$s, t : R \to U$ はエタールなので、
$$
\alpha : s^*\Omega_{U/Y} \to \Omega_{R/Y} \to t^*\Omega_{U/Y}
$$
という同型を得る。ここでは Morphisms, Lemma
\ref{morphisms-lemma-triangle-differentials-smooth} を用いた。
これがコサイクル条件を満たすことを確認すれば完了する。
$Y$ がスキームでない場合は、$\Omega_{U/Y}$ を写像
$(U \to Y)^*\Omega_{Y/S} \to \Omega_{U/S}$
の余核として定義し、前と同じように進める。この二段階の過程は少し不格好で
ある。別の可能性は、Lemma \ref{spaces-more-morphisms-lemma-localize-differentials}
のような任意の図式について層 $\Omega_{U/V}$ を貼り合わせることだが、
これもそれほど優美ではない。それでも、どちらの方法も機能し、
微分層のやや初等的な構成を与える。
\end{remark}

\begin{lemma}
\label{spaces-more-morphisms-lemma-functoriality-differentials}
$S$ をスキームとする。
$$
\xymatrix{
X' \ar[d] \ar[r]_f & X \ar[d] \\
Y' \ar[r] & Y
}
$$
を代数空間の可換図式とする。写像
$f^\sharp : \mathcal{O}_X \to f_*\mathcal{O}_{X'}$
と写像
$f_*\text{d}_{X'/Y'} : f_*\mathcal{O}_{X'} \to f_*\Omega_{X'/Y'}$
との合成は $Y$ 導分である。したがって $\mathcal{O}_X$ 加群の
標準的な写像 $\Omega_{X/Y} \to f_*\Omega_{X'/Y'}$ を得る。
さらに $f_*$ と $f^*$ の随伴性により、標準的な
$\mathcal{O}_{X'}$ 加群準同型
$$
c_f : f^*\Omega_{X/Y} \longrightarrow \Omega_{X'/Y'}.
$$
を得る。これは、$f^*\text{d}_{X/Y}(t)$ が
$\text{d}_{X'/Y'}(f^* t)$ に写るという性質により一意に特徴づけられる。
ここで $t$ は $\mathcal{O}_X$ の任意の局所切断である。
\end{lemma}

\begin{proof}
これは Modules on Sites, Lemma
\ref{sites-modules-lemma-functoriality-differentials-ringed-topoi}
の特別な場合である。
\end{proof}

\begin{lemma}
\label{spaces-more-morphisms-lemma-check-functoriality-differentials}
$S$ をスキームとする。
$$
\xymatrix{
X'' \ar[d] \ar[r]_g & X' \ar[d] \ar[r]_f & X \ar[d] \\
Y'' \ar[r] & Y' \ar[r] & Y
}
$$
を $S$ 上の代数空間の可換図式とする。このとき
$$
c_{f \circ g} = c_g \circ g^* c_f
$$
が写像 $(f \circ g)^*\Omega_{X/Y} \to \Omega_{X''/Y''}$ として
成り立つ。
\end{lemma}

\begin{proof}
省略する。ヒント：$c_f, c_g, c_{f \circ g}$ の特徴づけを、
これらの写像が局所切断に及ぼす作用という形で用いよ。
\end{proof}

\begin{lemma}
\label{spaces-more-morphisms-lemma-triangle-differentials}
$S$ をスキームとする。$f : X \to Y$, $g : Y \to B$ を
$S$ 上の代数空間の射とする。このとき標準的な完全列
$$
f^*\Omega_{Y/B} \to \Omega_{X/B} \to \Omega_{X/Y} \to 0
$$
が存在する。ここで写像は Lemma
\ref{spaces-more-morphisms-lemma-functoriality-differentials} を適用して得られる。
\end{lemma}

\begin{proof}
スキームの場合の結果 Morphisms, Lemma
\ref{morphisms-lemma-triangle-differentials} と、Lemma
\ref{spaces-more-morphisms-lemma-localize-differentials} によるエタール局所化から従う。
\end{proof}

\begin{lemma}
\label{spaces-more-morphisms-lemma-immersion-differentials}
$S$ をスキームとする。$X \to Y$ が $S$ 上の代数空間の埋め込みならば、
% P09-ERR-1053: frozen source line 1134 states Omega_{X/S}; an immersion only forces Omega_{X/Y}=0.
$\Omega_{X/S}$ は零である。
\end{lemma}

\begin{proof}
スキームの場合の結果 Morphisms, Lemma
\ref{morphisms-lemma-immersion-differentials} と、Lemma
\ref{spaces-more-morphisms-lemma-localize-differentials} によるエタール局所化から従う。
\end{proof}

\begin{lemma}
\label{spaces-more-morphisms-lemma-differentials-relative-immersion}
$S$ をスキームとする。$B$ を $S$ 上の代数空間とする。
$i : Z \to X$ を $B$ 上の代数空間の埋め込みとする。
このとき標準的な完全列
$$
\mathcal{C}_{Z/X} \to i^*\Omega_{X/B} \to \Omega_{Z/B} \to 0
$$
が存在する。ここで最初の矢印は $\text{d}_{X/B}$ により誘導され、
第二の矢印は Lemma
\ref{spaces-more-morphisms-lemma-functoriality-differentials} から得られる。
\end{lemma}

\begin{proof}
これは Morphisms, Lemma
\ref{morphisms-lemma-differentials-relative-immersion}
の代数空間版であり、その補題と、Lemmas
\ref{spaces-more-morphisms-lemma-localize-differentials} および
\ref{spaces-more-morphisms-lemma-etale-conormal} によるエタール局所化から従う。
しかし、最初の矢印を大域的に定義できることを確認すべきである。
したがって、ここで「$\text{d}_{X/B}$ により誘導される」の意味を説明する。
$i$ は、$X$ を開部分空間で置き換えた後、閉埋め込みであると
仮定してよい。$\mathcal{I} \subset \mathcal{O}_X$ を
$Z \subset X$ に対応する準連接イデアル層とする。
% P09-ERR-1054: frozen source lines 1169-1176 switch from the asserted relative base B to S.
このとき $\text{d}_{X/S} : \mathcal{I} \to \Omega_{X/S}$ は
部分層 $\mathcal{I}^2 \subset \mathcal{I}$ を
$\mathcal{I}\Omega_{X/S}$ に写す。したがってこれは
$\mathcal{I}/\mathcal{I}^2 \to
\Omega_{X/S}/\mathcal{I}\Omega_{X/S}$ という
$\mathcal{O}_X/\mathcal{I}$ 線形写像を誘導する。
Morphisms of Spaces, Lemma
\ref{spaces-morphisms-lemma-i-star-equivalence} により、これは所望の写像
$\mathcal{C}_{Z/X} \to i^*\Omega_{X/S}$ に対応する。
\end{proof}

\begin{lemma}
\label{spaces-more-morphisms-lemma-differentials-relative-immersion-section}
$S$ をスキームとする。$B$ を $S$ 上の代数空間とする。
$i : Z \to X$ を $B$ 上の代数空間の埋め込みとし、
$i$ は（エタール局所的に）左逆をもつと仮定する。このとき
Lemma \ref{spaces-more-morphisms-lemma-differentials-relative-immersion} の標準的な列
$$
0 \to \mathcal{C}_{Z/X} \to i^*\Omega_{X/B} \to \Omega_{Z/B} \to 0
$$
は（エタール局所的に）分裂完全である。
\end{lemma}

\begin{proof}
明確化：$g : X \to Z$ が $i$ の $B$ 上の左逆ならば、
$i^*c_g$ は写像
$i^*\Omega_{X/B} \to \Omega_{Z/B}$ の右逆であると主張している。
こう述べれば、結果はスキームの射に対する対応する結果と、
Lemmas \ref{spaces-more-morphisms-lemma-localize-differentials} および
\ref{spaces-more-morphisms-lemma-etale-conormal} によるエタール局所化から従う。
\end{proof}

\begin{lemma}
\label{spaces-more-morphisms-lemma-base-change-differentials}
$S$ をスキームとする。$X \to Y$ を $S$ 上の代数空間の射とする。
$g : Y' \to Y$ を $S$ 上の代数空間の射とする。
$X' = X_{Y'}$ を $X$ の基底変換とする。
射影を $g' : X' \to X$ と表す。このとき Lemma
\ref{spaces-more-morphisms-lemma-functoriality-differentials} の写像
$$
(g')^*\Omega_{X/Y} \to \Omega_{X'/Y'}
$$
は同型である。
\end{lemma}

\begin{proof}
スキームの場合の結果 Morphisms, Lemma
\ref{morphisms-lemma-base-change-differentials} と、
Lemma \ref{spaces-more-morphisms-lemma-localize-differentials} によるエタール局所化から従う。
\end{proof}

\begin{lemma}
\label{spaces-more-morphisms-lemma-differential-product}
$S$ をスキームとする。$f : X \to B$ と $g : Y \to B$ を、
同じ標的をもつ $S$ 上の代数空間の射とする。
$p : X \times_B Y \to X$ と $q : X \times_B Y \to Y$ を射影とする。
Lemma \ref{spaces-more-morphisms-lemma-functoriality-differentials} の写像は同型
$$
p^*\Omega_{X/B} \oplus q^*\Omega_{Y/B}
\longrightarrow
\Omega_{X \times_B Y/B}
$$
を与える。
\end{lemma}

\begin{proof}
スキームの場合の結果 Morphisms, Lemma
\ref{morphisms-lemma-differential-product} と、
Lemma \ref{spaces-more-morphisms-lemma-localize-differentials} によるエタール局所化から従う。
\end{proof}

\begin{lemma}
\label{spaces-more-morphisms-lemma-finite-type-differentials}
$S$ をスキームとする。$f : X \to Y$ を $S$ 上の代数空間の射とする。
$f$ が局所有限型ならば、$\Omega_{X/Y}$ は有限型
$\mathcal{O}_X$ 加群である。
\end{lemma}

\begin{proof}
スキームの場合の結果 Morphisms, Lemma
\ref{morphisms-lemma-finite-type-differentials} と、
Lemma \ref{spaces-more-morphisms-lemma-localize-differentials} によるエタール局所化から従う。
\end{proof}

\begin{lemma}
\label{spaces-more-morphisms-lemma-finite-presentation-differentials}
$S$ をスキームとする。$f : X \to Y$ を $S$ 上の代数空間の射とする。
$f$ が局所有限表示ならば、$\Omega_{X/Y}$ は有限表示
$\mathcal{O}_X$ 加群である。
\end{lemma}

\begin{proof}
スキームの場合の結果 Morphisms, Lemma
\ref{morphisms-lemma-finite-presentation-differentials} と、
Lemma \ref{spaces-more-morphisms-lemma-localize-differentials} によるエタール局所化から従う。
\end{proof}

\begin{lemma}
\label{spaces-more-morphisms-lemma-smooth-omega-finite-locally-free}
$S$ をスキームとする。$f : X \to Y$ を $S$ 上の代数空間の
滑らかな射とする。このとき微分加群 $\Omega_{X/Y}$ は有限局所自由である。
\end{lemma}

\begin{proof}
Lemma \ref{spaces-more-morphisms-lemma-localize-differentials} により、主張は $X$ と $Y$ 上
エタール局所的である。したがってスキームの場合から従う。次を参照せよ。
Morphisms, Lemma
\ref{morphisms-lemma-smooth-omega-finite-locally-free}。
\end{proof}













\section{エタールサイトの位相的不変性}
\label{spaces-more-morphisms-section-topological-invariance}

\noindent
サイト $X_{spaces, \etale}$ が「位相的不変量」であることを示す。
すると $X_\etale$ も位相的不変量であることが従う。これは
$X_{spaces, \etale}$ の表現可能な対象からなる。Lemma
\ref{spaces-more-morphisms-lemma-topological-invariance} を参照せよ。

\begin{theorem}
\label{spaces-more-morphisms-theorem-topological-invariance}
$S$ をスキームとする。$f : X \to Y$ を $S$ 上の代数空間の射とする。
$f$ は整、普遍単射かつ全射であると仮定する。関手
$$
V \longmapsto V_X = X \times_Y V
$$
は圏の同値
$Y_{spaces, \etale} \to X_{spaces, \etale}$ を定める。
\end{theorem}

\begin{proof}
射 $f$ は表現可能な普遍同相射である。次を参照せよ。
Morphisms of Spaces, Section
\ref{spaces-morphisms-section-universal-homeomorphisms}。

\medskip\noindent
まず、この関手が忠実であることを示す。
$V', V$ を $Y_{spaces, \etale}$ の対象とし、
$a, b : V' \to V$ を $Y$ 上の異なる射とする。
$V', V$ は $Y$ 上エタールなので、
$$
E =  V' \times_{(a, b), V \times_Y V, \Delta_{V/Y}} V
$$
で与えられる $a, b$ の等化子も $Y$ 上エタールである。
したがって $E \to V'$ はエタールなモノ射
（すなわち開埋め込み）であり、全射である場合に限り同型である。
$X \to Y$ は普遍同相射なので、これが成り立つのは
$E_X = V'_X$ の場合、すなわち $a_X = b_X$ の場合に限る。

\medskip\noindent
次に、この関手が充満忠実であることを示す。
$V', V$ を $Y_{spaces, \etale}$ の対象とし、
$c : V'_X \to V_X$ を $X$ 上の射とする。
射 $a : V' \to V$ を $Y$ 上で、$a_X = c$ となるように構成したい。
$a' : V'' \to V'$ を、$V''$ が分離的代数空間となる全射エタール射とする。
射 $a'' : V'' \to V$ を、$a''_X = c \circ a'_X$ となるように
構成できれば、
二つの合成
$$
V'' \times_{V'} V'' \xrightarrow{\text{pr}_i} V'' \xrightarrow{a''} V
$$
は、第一段落で示した関手の忠実性により等しい。
したがって $a''$ は一意な射 $a : V' \to V$ を経由する。
実際、$V'$ は層として $V''$ を同値関係
$V'' \times_{V'} V''$ で割った商である。
ゆえに $V'$ は分離的であると仮定してよい。この場合、グラフ
$$
\Gamma_c \subset (V' \times_Y V)_X
$$
は開かつ閉である（細部は省略する）。$X \to Y$ は普遍同相射なので、
開かつ閉な部分空間 $\Gamma \subset V' \times_Y V$ であって
$\Gamma_X = \Gamma_c$ となるものが存在する。
射影 $\Gamma \to V'$ はエタール射であり、その $X$ への基底変換は
同型である。したがって $\Gamma \to V'$ はエタール、普遍単射かつ全射
であるから、Morphisms of Spaces, Lemma
\ref{spaces-morphisms-lemma-etale-universally-injective-open}
により同型である。ゆえに $\Gamma$ は所望の射
$a : V' \to V$ のグラフである。

\medskip\noindent
最後に、この関手が本質的全射であることを示す。
$U$ を $X_{spaces, \etale}$ の対象とする。
$V$ を $Y_{spaces, \etale}$ の対象として見いだし、
$V_X \cong U$ を満たさなければならない。
全射エタール射 $U' \to U$ を、同型
$U' \cong V'_X$ および $U' \times_U U' \cong V''_X$
が存在するようにとる。
ここで $V'', V'$ は $Y_{spaces, \etale}$ の対象である。
すると関手の充満忠実性により、射
$s, t : V'' \to V'$ であって、
$t_X = \text{pr}_0$ および $s_X = \text{pr}_1$ が
$U' \times_U U' \to U'$ の射として成り立つものを得る。
$(\text{pr}_0, \text{pr}_1) :
U' \times_U U' \to U' \times_S U'$ がエタール同値関係であること、
ならびに $U' \to V'$ と $U' \times_U U' \to V''$ が普遍単射かつ
全射であることを用いると、
$(t, s) : V'' \to V' \times_S V'$ がエタール同値関係であることが
従う。このとき商 $V = V'/V''$
（Spaces, Theorem \ref{spaces-theorem-presentation} を参照）は
代数空間 $V$ であり、$Y$ 上にある。射 $V' \to V$ が存在し、
$V'' = V' \times_V V'$ を満たす。したがって射 $V \to Y$ を得る
（Descent on Spaces, Lemma
\ref{spaces-descent-lemma-fpqc-universal-effective-epimorphisms}
を参照）。$X$ へ基底変換すると、射 $U' \to V_X$ と、それと両立する同型
$U' \times_{V_X} U' = U' \times_U U'$ を得る。
これは $V_X \cong U$ を意味する
（直前に引用した補題をもう一度用いる）。

\medskip\noindent
スキーム $W$ と全射エタール射 $W \to Y$ を選ぶ。
スキーム $U'$ と全射エタール射 $U' \to U \times_X W_X$ を選ぶ。
$U'$ と $U' \times_U U'$ は $X$ 上エタールなスキームであり、
$X$ への構造射はスキーム $W_X$ を経由することに注意する。
したがって \'Etale Cohomology, Theorem
\ref{etale-cohomology-theorem-topological-invariance}
により、スキーム $V', V''$ が存在し、これらは $W$ 上エタールで、
その $W_X$ への基底変換がそれぞれ $U'$ と $U' \times_U U'$ に
同型となる。これで証明が完了する。
\end{proof}

\begin{lemma}
\label{spaces-more-morphisms-lemma-topological-invariance}
Theorem \ref{spaces-more-morphisms-theorem-topological-invariance} と同じ仮定と記法のもとで、
圏の同値 $Y_{spaces, \etale} \to X_{spaces, \etale}$ は
圏の同値 $Y_\etale \to X_\etale$ および
$Y_{affine, \etale} \to X_{affine, \etale}$ に制限される。
\end{lemma}

\begin{proof}
これは、対象 $V \in Y_{spaces, \etale}$ に対し、$V$ が
（アフィン）スキームであることと、$V \times_Y X$ が
（アフィン）スキームであることが同値だという主張にすぎない。
$V \times_Y X \to V$ は
$X \to Y$ の基底変換として整、普遍単射かつ全射なので、これは
Limits of Spaces, Lemma
\ref{spaces-limits-lemma-integral-universally-bijective-scheme} および
Proposition \ref{spaces-limits-proposition-affine} から従う。
\end{proof}

\begin{remark}
\label{spaces-more-morphisms-remark-topological-invariance-etale-site}
\begin{reference}
2016 年 5 月 1 日付 Lenny Taelman の電子メール。
\end{reference}
代数空間の普遍同相射は表現可能であるとは限らない。次を参照せよ。
Morphisms of Spaces, Example
\ref{spaces-morphisms-example-universal-homeomorphism}。
実際、Theorem \ref{spaces-more-morphisms-theorem-topological-invariance} は
普遍同相射については成り立たない。これを見るため、$k$ を標数 $0$ の
代数閉体とし、
$$
\mathbf{A}^1 \to X \to \mathbf{A}^1
$$
を Morphisms of Spaces, Example
\ref{spaces-morphisms-example-universal-homeomorphism}
のものとする。最初の射はエタールであり、
$t$ と $-t$ を同一視する。ここで
$t \in \mathbf{A}^1_k \setminus \{0\}$ である。
第二の射は考えている普遍同相射であった。
$\mathbf{A}^1_k$ は非自明な連結有限エタール被覆をもたないので
（$k$ が標数零の代数閉体だからである。細部は省略する）、
非自明な連結有限エタール被覆 $Y \to X$ を構成すれば十分である。
そのために $Y$ を、零点を二重化したアフィン直線とする
（Schemes, Example \ref{schemes-example-affine-space-zero-doubled}）。
このとき $Y = Y_1 \cup Y_2$ であり、
$Y_i = \mathbf{A}^1_k$ を $\mathbf{A}^1_k \setminus \{0\}$ に沿って
貼り合わせたものである。射 $Y \to X$ を定義するため、射
$$
Y_1 \xrightarrow{1} \mathbf{A}^1_k \to X
\quad\text{and}\quad
Y_2 \xrightarrow{-1} \mathbf{A}^1_k \to X.
$$ 
を用いる。これらは $Y_1 \cap Y_2$ 上で、$X$ の構成により貼り合わさり、
したがって射 $Y \to X$ を定める。実際、
$$
\xymatrix{
Y \ar[d] & Y_1 \amalg Y_2 \ar[l] \ar[d] \\
X & \mathbf{A}^1_k \ar[l]
}
$$
はカルテジアン正方形であると主張する。細部は省略する。例えば
Groupoids, Lemma \ref{groupoids-lemma-criterion-fibre-product}
を用いられる。$\mathbf{A}^1_k \to X$ はエタールかつ全射なので、
これにより $Y \to X$ は次数 $2$ の有限エタール射であり、
所望の例を与えることが分かる。

\medskip\noindent
より簡単には、次のように論じられる。スキーム $Y$ には群
$G = \{+1, -1\}$ の自由作用があり、$-1$ は $Y_1$ と $Y_2$ を交換し、
座標の符号を変える。このとき $X = Y/G$
（Spaces, Definition \ref{spaces-definition-quotient} を参照）であり、
したがって $Y \to X$ は有限エタールである。また、アフィン空間上の
普遍同相射 $X \to \mathbf{A}^1_k$ が存在することも、
$t \mapsto t^2$ を用いて直接示せる。
実際、この $X$ は上の $X$ と同じである。
\end{remark}









\section{肥厚化}
\label{spaces-more-morphisms-section-thickenings}

\noindent
次の用語は完全に標準的ではないかもしれないが、便利である。

\begin{definition}
\label{spaces-more-morphisms-definition-thickening}
肥厚化。$S$ をスキームとする。
\begin{enumerate}
\item 代数空間 $X'$ が代数空間 $X$ の {\it 肥厚化}であるとは、
$X$ が $X'$ の閉部分空間であり、付随する位相空間が等しいことをいう。
\item $X'$ が $X$ の {\it 一次肥厚化}であるとは、
$X$ が $X'$ の閉部分空間であり、準連接イデアル層
$\mathcal{I} \subset \mathcal{O}_{X'}$ で $X$ を定めるものの二乗が
零であることをいう。
\item $X'$ が $X$ の {\it 有限次肥厚化}であるとは、
% P09-ERR-1055: frozen source line 1518 says X is a closed subspace of X; expected X'.
$X$ が $X$ の閉部分空間であり、準連接イデアル層
$\mathcal{I} \subset \mathcal{O}_{X'}$ で $X$ を定めるものが冪零で
あること、すなわち整数 $n \geq 0$ が存在して
$\mathcal{I}^{n + 1} = 0$ となることをいう。
\item $X'$ が {\it $n$ 次肥厚化}として $X$ を肥厚化するとは、
% P09-ERR-1056: frozen source line 1522 says X is a closed subspace of X; expected X'.
$X$ が $X$ の閉部分空間であり、$\mathcal{I}^{n + 1} = 0$ であることを
いう。ここで $\mathcal{I} \subset \mathcal{O}_{X'}$ は $X$ を定める
準連接イデアル層である。
\item 二つの肥厚化 $X \subset X'$ と $Y \subset Y'$ が与えられたとき、
{\it 肥厚化の射}とは射
% P09-ERR-1057: frozen source line 1527 uses f(X) before f is defined; expected f'(X).
$f' : X' \to Y'$ であって $f(X) \subset Y$ を満たすもの、
すなわち $f'|_X$ が閉部分空間 $Y$ を経由するものをいう。
この状況で $f = f'|_X : X \to Y$ とおき、
$(f, f') : (X \subset X') \to (Y \subset Y')$ を肥厚化の射という。
\item $B$ を代数空間とする。同様に {\it $B$ 上の肥厚化}と
{\it $B$ 上の肥厚化の射}を定義する。これは、上の空間
$X, X', Y, Y'$ が $B$ への構造射を備えた代数空間であり、射
$X \to X'$, $Y \to Y'$ および $f' : X' \to Y'$ が
$B$ 上の射であることを意味する。
\end{enumerate}
\end{definition}

\noindent
基本的な同値を述べる。$X \subset X'$ が肥厚化ならば、
$X \to X'$ は整かつ普遍全単射であることに注意する。
したがって、引き戻し関手により
\begin{equation}
\label{spaces-more-morphisms-equation-equivalence-etale-spaces}
X_{spaces, \etale} = X'_{spaces, \etale}
\end{equation}
となる。Theorem \ref{spaces-more-morphisms-theorem-topological-invariance} を参照せよ。
ゆえに $\mathcal{O}_{X'}$ を $X_{spaces, \etale}$ 上の層と
考えてよい。したがって局所環付きトポスの標準的な同値
\begin{equation}
\label{spaces-more-morphisms-equation-fundamental-equivalence}
(\Sh(X'_{spaces, \etale}), \mathcal{O}_{X'})
\cong
(\Sh(X_{spaces, \etale}), \mathcal{O}_{X'})
\end{equation}
を得る。以下では、これを Properties of Spaces, Theorem
\ref{spaces-properties-theorem-fully-faithful}
の充満忠実性の結果と頻繁に組み合わせる。例えば閉埋め込み
$i_X : X \to X'$ は全射
$i_X^\sharp : \mathcal{O}_{X'} \to \mathcal{O}_X$ に対応する。

\medskip\noindent
$S$ をスキームとし、$B$ を $S$ 上の代数空間とする。
$(f, f') : (X \subset X') \to (Y \subset Y')$ を
$B$ 上の肥厚化の射とする。連続関手の図式
$$
\xymatrix{
X_{spaces, \etale} &
Y_{spaces, \etale} \ar[l] \\
X'_{spaces, \etale} \ar[u] &
Y'_{spaces, \etale} \ar[u] \ar[l]
}
$$
は可換であり、縦の矢印は同値であることに注意する。したがって
$f_{spaces, \etale}$, $f_{small}$,
$f'_{spaces, \etale}$, $f'_{small}$ はすべて同じトポスの射を定める。
ゆえに
$$
(f')^\sharp :
f_{spaces, \etale}^{-1}\mathcal{O}_{Y'}
\longrightarrow
\mathcal{O}_{X'}
$$
を $\mathcal{O}_B$ 代数の層の写像と考えてよい。これは可換図式
$$
\xymatrix{
f_{spaces, \etale}^{-1}\mathcal{O}_Y \ar[r]_-{f^\sharp} \ar[r] &
\mathcal{O}_X \\
f_{spaces, \etale}^{-1}\mathcal{O}_{Y'} \ar[r]^-{(f')^\sharp}
\ar[u]^{i_Y^\sharp} &
\mathcal{O}_{X'} \ar[u]_{i_X^\sharp}
}
$$
に入る。
ここで $i_X : X \to X'$ および $i_Y : Y \to Y'$ は、与えられた
閉埋め込みの名称である。

\begin{lemma}
\label{spaces-more-morphisms-lemma-first-order-thickening-maps}
$S$ をスキームとする。$B$ を $S$ 上の代数空間とする。
$X \subset X'$ と $Y \subset Y'$ を $B$ 上の代数空間の肥厚化とする。
$f : X \to Y$ を $B$ 上の代数空間の射とする。
任意の $\mathcal{O}_B$ 代数の写像
$$
\alpha : f_{spaces, \etale}^{-1}\mathcal{O}_{Y'} \to \mathcal{O}_{X'}
$$
であって、図式
$$
\xymatrix{
f_{spaces, \etale}^{-1}\mathcal{O}_Y \ar[r]_-{f^\sharp} \ar[r] &
\mathcal{O}_X \\
f_{spaces, \etale}^{-1}\mathcal{O}_{Y'} \ar[r]^-\alpha
\ar[u]^{i_Y^\sharp} &
\mathcal{O}_{X'} \ar[u]_{i_X^\sharp}
}
$$
を可換にするものが与えられたとする。このとき
% P09-ERR-1058: frozen source line 1619 says 'a unique morphism of (f,f') of thickenings'.
一意な肥厚化の射 $(f, f')$ が $B$ 上で存在し、
$\alpha = (f')^\sharp$ を満たす。
\end{lemma}

\begin{proof}
$f'$ を見いだすには、Properties of Spaces, Theorem
\ref{spaces-properties-theorem-fully-faithful} により、環付きトポスの射
$$
(f_{spaces, \etale}, \alpha) :
(\Sh(X_{spaces, \etale}), \mathcal{O}_{X'})
\longrightarrow
(\Sh(Y_{spaces, \etale}), \mathcal{O}_{Y'})
$$
が局所環付きトポスの射であることを示すだけでよい。
これは、局所環付きトポスの射の定義
（Modules on Sites, Definition
\ref{sites-modules-definition-morphism-locally-ringed-topoi}）、
$(f, f^\sharp)$ が局所環付きトポスの射であること
（Properties of Spaces, Lemma
\ref{spaces-properties-lemma-morphism-locally-ringed}）、
$\alpha$ が与えられた可換図式に入ること、および
$i_X^\sharp$ と $i_Y^\sharp$ の核が局所冪零であることから直ちに従う。
最後に、$f' \circ i_X = i_Y \circ f$ であることは、この図式の可換性と
Properties of Spaces, Theorem
\ref{spaces-properties-theorem-fully-faithful}
のもう一度の適用から従う。$f'$ が $B$ 上の射であることの確認は省略する。
\end{proof}

\begin{lemma}
\label{spaces-more-morphisms-lemma-open-subspace-thickening}
$S$ をスキームとする。$X \subset X'$ を $S$ 上の代数空間の肥厚化とする。
任意の開部分空間 $U \subset X$ に対して、一意な開部分空間
$U' \subset X'$ であって
$U = X \times_{X'} U'$ を満たすものが存在する。
\end{lemma}

\begin{proof}
$U' \to X'$ を $X'_{spaces, \etale}$ の対象とし、
(\ref{spaces-more-morphisms-equation-equivalence-etale-spaces}) により
対象 $U \to X$、すなわち $X_{spaces, \etale}$ の対象に
対応するものとする。
射 $U' \to X'$ はエタールかつ
普遍単射なので開埋め込みである。次を参照せよ。
Morphisms of Spaces, Lemma
\ref{spaces-morphisms-lemma-etale-universally-injective-open}。
\end{proof}

\noindent
$i_X : X \to X'$ を代数空間の肥厚化とする。核
$\mathcal{I} = \Ker(i_X^\sharp) \subset \mathcal{O}_{X'}$
の任意の局所切断は局所冪零である。これにより
$\mathcal{O}_{X'}$ の構造層をある程度制御できるが、技術的には
有限次肥厚化 $X \subset X'$ のクラスのほうがはるかに扱いやすい。
実際、$X'$ が $n$ 次肥厚化ならば、フィルトレーション
$$
0 \subset \mathcal{I}^n \subset \mathcal{I}^{n - 1} \subset
\ldots \subset \mathcal{I} \subset \mathcal{O}_{X'}
$$
をもち、$X'$ が閉部分空間によりフィルター付けられることが分かる。
% P09-ERR-1059: frozen source line 1678 skips X_n and ends with X_{n+1}=X'.
$$
X = X_0 \subset X_1 \subset \ldots \subset X_{n - 1} \subset X_{n + 1} = X'
$$
各対 $X_i \subset X_{i + 1}$ は $B$ 上の一次肥厚化である。
% P09-ERR-1060: frozen source line 1681 invokes an undefined base B in this paragraph.
単純な帰納法を用いれば、一次肥厚化について証明された多くの結果を
有限次肥厚化についての結果として言い換えられる。

\begin{lemma}
\label{spaces-more-morphisms-lemma-first-order-thickening-surjective}
$S$ をスキームとする。$X \subset X'$ を $S$ 上の代数空間の
肥厚化とする。$U$ を $X_{spaces, \etale}$ のアフィン対象とする。
このとき
$$
\Gamma(U, \mathcal{O}_{X'}) \to \Gamma(U, \mathcal{O}_X)
$$
は全射である。ここでは (\ref{spaces-more-morphisms-equation-fundamental-equivalence}) により
$\mathcal{O}_{X'}$ を $X_{spaces, \etale}$ 上の層と考えている。
\end{lemma}

\begin{proof}
$U' \to X'$ を、
$U = X \times_{X'} U'$ を満たす代数空間のエタール射とする。
Theorem \ref{spaces-more-morphisms-theorem-topological-invariance} を参照せよ。
Limits of Spaces, Lemma \ref{spaces-limits-lemma-affine}
により $U'$ はアフィンスキームである。したがって
$\Gamma(U, \mathcal{O}_{X'}) =
\Gamma(U', \mathcal{O}_{U'}) \to \Gamma(U, \mathcal{O}_U)$
は全射である。実際、$U \to U'$ はアフィンスキームの閉埋め込みである。
以下では、実際に最もよく使われる有限次肥厚化について直接証明を与える。
\end{proof}

\begin{proof}[有限次肥厚化に対する証明]
上で説明した原理により、$X \subset X'$ は一次肥厚化であると
仮定してよい。$\mathcal{I}$ で全射
$\mathcal{O}_{X'} \to \mathcal{O}_X$ の核を表す。
$\mathcal{I}$ は準連接 $\mathcal{O}_{X'}$ 加群であり、一次肥厚化の
定義により $\mathcal{I}^2 = 0$ なので、Morphisms of Spaces, Lemma
\ref{spaces-morphisms-lemma-i-star-equivalence}
を適用して、$\mathcal{I}$ が準連接 $\mathcal{O}_X$ 加群であることが
分かる。したがって補題は、短完全列
$$
0 \to \mathcal{I} \to \mathcal{O}_{X'} \to \mathcal{O}_X \to 0
$$
に付随するコホモロジー長完全列と、
$H^1_\etale(U, \mathcal{I}) = 0$ となることから従う。
ここで $\mathcal{I}$ は準連接である。次を参照せよ。
Descent, Proposition
\ref{descent-proposition-same-cohomology-quasi-coherent} および
Cohomology of Schemes, Lemma
\ref{coherent-lemma-quasi-coherent-affine-cohomology-zero}。
\end{proof}

\begin{lemma}
\label{spaces-more-morphisms-lemma-thickening-scheme}
$S$ をスキームとする。$X \subset X'$ を $S$ 上の代数空間の
肥厚化とする。$X$ がスキーム（で表現可能）ならば、$X'$ もそうである。
\end{lemma}

\begin{proof}
$X'_{red} = X_{red}$ であることに注意する。したがって $X$ が
スキームならば、$X'_{red}$ はスキームである。ゆえに結果は
Limits of Spaces, Lemma
\ref{spaces-limits-lemma-reduction-scheme} から従う。
以下では、実際に最も頻繁に使われる有限次肥厚化について直接証明を与える。
\end{proof}

\begin{proof}[有限次肥厚化に対する証明]
$X'$ が $X$ の一次肥厚化である場合を証明すれば十分である。
Properties of Spaces, Lemma
\ref{spaces-properties-lemma-subscheme} により、$X'$ には
スキームである最大の開部分空間が存在する。したがって、任意の点
$x$、すなわち $|X'| = |X|$ の点が、スキームである $X'$ の
開部分空間に含まれることを示さなければならない。
Lemma \ref{spaces-more-morphisms-lemma-open-subspace-thickening} を用いて、
$X \subset X'$ を $U \subset U'$ で置き換えてよい。
ここで $x \in U$ であり、$U$ はアフィンスキームである。
したがって $X$ はアフィンであると仮定してよい。
こうして次段落で論じる場合に帰着する。

\medskip\noindent
$X \subset X'$ を、$X$ がアフィンスキームである一次肥厚化とする。
$A = \Gamma(X, \mathcal{O}_X)$ および
$A' = \Gamma(X', \mathcal{O}_{X'})$ とおく。
Lemma \ref{spaces-more-morphisms-lemma-first-order-thickening-surjective} により、
写像 $A' \to A$ は全射である。その核 $I$ は二乗零のイデアルである。
Properties of Spaces, Lemma
\ref{spaces-properties-lemma-morphism-to-affine-scheme} により、
次の可換図式に入る標準的な射 $f : X' \to \Spec(A')$ を得る。
$$
\xymatrix{
X \ar@{=}[d] \ar[r] &  X' \ar[d]^f \\
\Spec(A) \ar[r] & \Spec(A')
}
$$
横の矢印は肥厚化なので、$f$ が普遍単射かつ全射であることは明らかである。
したがって $f$ がエタールであることを示せば十分である。実際、そのとき
Morphisms of Spaces, Lemma
\ref{spaces-morphisms-lemma-etale-universally-injective-open}
により $f$ は同型となる。

\medskip\noindent
$f$ がエタールであることを示すため、アフィンスキーム $U'$ と
エタール射 $U' \to X'$ を選ぶ。合成
$U' \to X' \to \Spec(A')$ がエタールであることを示せば十分である。
Properties of Spaces, Definition
\ref{spaces-properties-definition-etale} を参照せよ。
$U' = \Spec(B')$ と書く。$U = X \times_{X'} U'$ とおく。
$U$ は $U'$ の閉部分空間なので閉部分スキームである。したがって
$U = \Spec(B)$ であり、$B' \to B$ は全射である。
$J = \Ker(B' \to B)$ と表す。また
$J = \Gamma(U, \mathcal{I})$ であることに注意する。
ここで $\mathcal{I} = \Ker(\mathcal{O}_{X'} \to \mathcal{O}_X)$ は、
Lemma \ref{spaces-more-morphisms-lemma-first-order-thickening-surjective} の証明と同じく
$X_{spaces, \etale}$ 上で考える。射
$U' \to X' \to \Spec(A')$ は可換図式
$$
\xymatrix{
0 \ar[r] &
J \ar[r] &
B' \ar[r] &
B \ar[r] & 0 \\
0 \ar[r] &
I \ar[r] \ar[u] &
A' \ar[r] \ar[u] &
A \ar[r] \ar[u] & 0
}
$$
を誘導する。さて、$\mathcal{I}$ は準連接 $\mathcal{O}_X$ 加群なので、
$\mathcal{I} = (\widetilde I)^a$ である。記法については
Descent, Definition \ref{descent-definition-structure-sheaf} を、
この事実については Descent, Proposition
\ref{descent-proposition-equivalence-quasi-coherent} を参照せよ。
したがって $J = I \otimes_A B$ である。最後に、
$A \to B$ はエタールであることに注意する。実際、
$U \to X$ はエタール射 $U' \to X'$ の基底変換としてエタールである。
Algebra, Lemma
\ref{algebra-lemma-lift-etale-infinitesimal} により、
$A' \to B'$ はエタールであると結論する。
\end{proof}

\begin{lemma}
\label{spaces-more-morphisms-lemma-thickening-equivalence}
$S$ をスキームとする。$X \subset X'$ を $S$ 上の代数空間の
肥厚化とする。関手
$$
V' \longmapsto V = X \times_{X'} V'
$$
は圏の同値 $X'_\etale \to X_\etale$ を定める。
\end{lemma}

\begin{proof}
関手 $V' \mapsto V$ は圏の同値
$X'_{spaces, \etale} \to X_{spaces, \etale}$ を定める。Theorem
\ref{spaces-more-morphisms-theorem-topological-invariance} を参照せよ。
したがって $V$ がスキームであることと $V'$ がスキームであることが
同値であると示せば十分である。これは Lemma
\ref{spaces-more-morphisms-lemma-thickening-scheme} の内容である。
\end{proof}

\noindent
一次肥厚化は次のように記述される。
% P09-ERR-1061: frozen source line 1837 has singular/plural disagreement in 'First order thickening are'.

\begin{lemma}
\label{spaces-more-morphisms-lemma-first-order-thickening}
$S$ をスキームとする。$f : X \to B$ を $S$ 上の代数空間の射とする。
短完全列
$$
0 \to \mathcal{I} \to \mathcal{A} \to \mathcal{O}_X \to 0
$$
を $X_\etale$ 上の層について考える。ここで $\mathcal{A}$ は
$f^{-1}\mathcal{O}_B$ 代数の層であり、
$\mathcal{A} \to \mathcal{O}_X$ は
$f^{-1}\mathcal{O}_B$ 代数の層の全射であり、
$\mathcal{I}$ はその核である。次を仮定する。
\begin{enumerate}
\item $\mathcal{I}$ は $\mathcal{A}$ の二乗零イデアルである。
\item $\mathcal{I}$ は $\mathcal{O}_X$ 加群として準連接である。
\end{enumerate}
このとき一次肥厚化 $X \subset X'$ が $B$ 上に存在し、さらに同型
$\mathcal{O}_{X'} \to \mathcal{A}$ が存在する。これは
$f^{-1}\mathcal{O}_B$ 代数の同型であって、
$\mathcal{O}_X$ への全射と両立するものが存在する。
\end{lemma}

\begin{proof}
この証明では、Lemmas
\ref{spaces-more-morphisms-lemma-first-order-thickening-surjective} および
\ref{spaces-more-morphisms-lemma-thickening-scheme} の証明で用いた議論の一部をやり直す。
まず $B = S = \Spec(\mathbf{Z})$ の場合を扱う。
$U$ をアフィンスキームとし、$U \to X$ はエタールであるとする。
このとき
$$
0 \to \mathcal{I}(U) \to \mathcal{A}(U) \to \mathcal{O}_X(U) \to 0
$$
は完全である。実際、$H^1(U_\etale, \mathcal{I}) = 0$ である。
ここで $\mathcal{I}$ は準連接である。次を参照せよ。
Descent, Proposition
\ref{descent-proposition-same-cohomology-quasi-coherent} および
Cohomology of Schemes, Lemma
\ref{coherent-lemma-quasi-coherent-affine-cohomology-zero}。
$V \to U$ が $X_{spaces, \etale}$ のアフィン対象の射ならば、
$$
\mathcal{I}(V) =
\mathcal{I}(U) \otimes_{\mathcal{O}_X(U)} \mathcal{O}_X(V)
$$
である。これは $\mathcal{I}$ が準連接 $\mathcal{O}_X$ 加群だからで
ある。Descent, Proposition
\ref{descent-proposition-equivalence-quasi-coherent} を参照せよ。
したがって $\mathcal{A}(U) \to \mathcal{A}(V)$ はエタールな環準同型で
ある。Algebra, Lemma
\ref{algebra-lemma-lift-etale-infinitesimal} を参照せよ。
ゆえに
$$
U \longmapsto U' = \Spec(\mathcal{A}(U))
$$
は $X_{affine, \etale}$ からアフィンスキームとエタール射の圏への
関手である。実際、この関手は関手 $U \mapsto U'$ として
$X_\etale$ 全体上に拡張できると主張する。
これを見るため、$U$ が $X_\etale$ の対象ならば、
$$
0 \to \mathcal{I}|_{U_{Zar}} \to \mathcal{A}|_{U_{Zar}} \to
\mathcal{O}_X|_{U_{Zar}} \to 0
$$
であり、$\mathcal{I}|_{U_{Zar}}$ は $U$ 上の準連接層であることに
注意する。Descent, Proposition
\ref{descent-proposition-equivalence-quasi-coherent-functorial}
を参照せよ。したがって More on Morphisms, Lemma
\ref{more-morphisms-lemma-first-order-thickening} により、
スキームの一次肥厚化 $U \subset U'$ であって、
$\mathcal{O}_{U'}$ が $\mathcal{A}|_{U_{Zar}}$ に同型となるものを得る。
この構成が上のアフィンの場合の構成と両立することは明らかである。

\medskip\noindent
表示 $X = U/R$ を選ぶ。Spaces, Definition
\ref{spaces-definition-presentation} を参照せよ。すると
$s, t : R \to U$ はエタール同値関係を定める。
上の関手を適用すると、スキームにおけるエタール同値関係
$s', t' : R' \to U'$ を得る。代数空間 $X' = U'/R'$ を考える
（Spaces, Theorem \ref{spaces-theorem-presentation} を参照）。
射 $X = U/R \to U'/R' = X'$ は一次肥厚化である。
$\mathcal{O}_{X'}$ を $X_\etale$ 上の層とみなして考える。
構成により同型
$$
\gamma :
\mathcal{O}_{X'}|_{U_\etale}
\longrightarrow
\mathcal{A}|_{U_\etale}
$$
を得て、$s^{-1}\gamma$ と $t^{-1}\gamma$ は
$R_\etale$ 上で一致する。したがって Properties of Spaces, Lemma
\ref{spaces-properties-lemma-descent-sheaf} により、
$\gamma$ は所望の一意な同型
$\mathcal{O}_{X'} \to \mathcal{A}$ から来る。

\medskip\noindent
一般の基底代数空間 $B$ の場合を扱うため、まず上のように
$X'$ を $\mathbf{Z}$ 上の代数空間として構成する。
次に同型 $\mathcal{O}_{X'} \to \mathcal{A}$ を用いて
$f^{-1}\mathcal{O}_B \to \mathcal{O}_{X'}$ を定義する。
Lemma \ref{spaces-more-morphisms-lemma-first-order-thickening-maps} により、これは
射 $X' \to B$ を定め、これは与えられた射 $X \to B$ と両立する。
これで完了する。
\end{proof}

\begin{lemma}
\label{spaces-more-morphisms-lemma-base-change-thickening}
$S$ をスキームとする。$Y \subset Y'$ を $S$ 上の代数空間の
肥厚化とする。$X' \to Y'$ を射とし、
$X = Y \times_{Y'} X'$ とおく。このとき
$(X \subset X') \to (Y \subset Y')$ は肥厚化の射である。
$Y \subset Y'$ が一次（それぞれ有限次）肥厚化ならば、
$X \subset X'$ は一次（それぞれ有限次）肥厚化である。
\end{lemma}

\begin{proof}
省略する。
\end{proof}


\begin{lemma}
\label{spaces-more-morphisms-lemma-composition-thickening}
$S$ をスキームとする。$X \subset X'$ と $X' \subset X''$ が
$S$ 上の代数空間の肥厚化ならば、$X \subset X''$ もそうである。
\end{lemma}

\begin{proof}
省略する。
\end{proof}

\begin{lemma}
\label{spaces-more-morphisms-lemma-descending-property-thickening}
肥厚化であるという性質は fpqc 局所的である。
一次肥厚化についても同様である。
\end{lemma}

\begin{proof}
この主張の意味は次である。$S$ をスキームとし、
$X \to X'$ を $S$ 上の代数空間の射とする。
$\{g_i : X'_i \to X'\}$ を代数空間の fpqc 被覆とし、
基底変換 $X_i \to X'_i$ がすべての $i$ に対して肥厚化であるとする。
このとき $X \to X'$ は肥厚化である。
射 $g_i$ は共同全射なので、$X \to X'$ は全射であると結論する。
Descent on Spaces, Lemma
\ref{spaces-descent-lemma-descending-property-closed-immersion}
により $X \to X'$ は閉埋め込みであると結論する。
したがって $X \to X'$ は肥厚化である。
一次肥厚化の場合の証明は省略する。
\end{proof}









\section{肥厚化の射}
\label{spaces-more-morphisms-section-morphisms-thickenings}

\noindent
$(f, f') : (X \subset X') \to (Y \subset Y')$ が代数空間の肥厚化の射ならば、
射 $f$ の性質はしばしば $f'$ に継承される。いくつかの変種がある。

\begin{lemma}
\label{spaces-more-morphisms-lemma-thicken-property-morphisms}
$S$ をスキームとする。
$(f, f') : (X \subset X') \to (Y \subset Y')$ を
$S$ 上の代数空間の肥厚化の射とする。このとき次が成り立つ。
\begin{enumerate}
\item $f$ がアフィン射であることと $f'$ がアフィン射であることは同値である。
\item $f$ が全射であることと $f'$ が全射であることは同値である。
\item $f$ が準コンパクトであることと $f'$ が準コンパクトであることは
同値である。
% P09-ERR-1062: frozen source line 2007 omits 'is' before the second quasi-compact.
\item $f$ が普遍閉であることと $f'$ が普遍閉であることは同値である。
\item $f$ が整であることと $f'$ が整であることは同値である。
\item $f$ が（準）分離的であることと $f'$ が（準）分離的であることは
同値である。
\item $f$ が普遍単射であることと $f'$ が普遍単射であることは同値である。
\item $f$ が普遍開であることと $f'$ が普遍開であることは同値である。
\item $f$ が表現可能であることと $f'$ が表現可能であることは同値である。
\item ここにさらに追加する。
\end{enumerate}
\end{lemma}

\begin{proof}
$Y \to Y'$ と $X \to X'$ は整な普遍同相射であることに注意する。
これは直ちに (2), (3), (4), (7), (8) を与える。
(1) は Limits of Spaces, Proposition
\ref{spaces-limits-proposition-affine} から従う。この命題は、
$X$ および $X'$ 上エタールなアフィンスキームの間、また
$Y$ および $Y'$ 上エタールなアフィンスキームの間に
一対一対応があることを述べる。
(5) は (1), (4) と Morphisms of Spaces, Lemma
\ref{spaces-morphisms-lemma-integral-universally-closed} から従う。
最後に、
$$
X \times_Y X = X \times_{Y'} X \to
X \times_{Y'} X' \to X' \times_{Y'} X'
$$
は肥厚化であることに注意する
（二つの矢印は Lemma \ref{spaces-more-morphisms-lemma-base-change-thickening} により
肥厚化である）。したがって射
% P09-ERR-1064: frozen source line 2037 uses an arrow, rather than inclusion, inside the target thickening pair.
$(X \subset X') \to
(X \times_Y X \to X' \times_{Y'} X')$
に (3), (4) を適用して (6) を得る。
最後に (9) は、スキームの代数空間肥厚化が再びスキームであることから
従う。Lemma \ref{spaces-more-morphisms-lemma-thickening-scheme} を参照せよ。
\end{proof}

\begin{lemma}
\label{spaces-more-morphisms-lemma-thicken-property-morphisms-cartesian}
$S$ をスキームとする。
$(f, f') : (X \subset X') \to (Y \subset Y')$ を
$S$ 上の代数空間の肥厚化の射とし、
$X = Y \times_{Y'} X'$ を満たすとする。
$X \subset X'$ が有限次肥厚化ならば、次が成り立つ。
\begin{enumerate}
\item $f$ が閉埋め込みであることと $f'$ が閉埋め込みであることは
同値である。
\item $f$ が局所有限型であることと $f'$ が局所有限型であることは
同値である。
\item $f$ が局所準有限であることと $f'$ が局所準有限であることは
同値である。
\item $f$ が相対次元 $d$ の局所有限型射であることと、
$f'$ が相対次元 $d$ の局所有限型射であることは同値である。
\item $\Omega_{X/Y} = 0$ であることと
$\Omega_{X'/Y'} = 0$ であることは同値である。
\item $f$ が非分岐であることと $f'$ が非分岐であることは同値である。
\item $f$ が固有であることと $f'$ が固有であることは同値である。
\item $f$ が有限射であることと $f'$ が有限射であることは同値である。
% P09-ERR-1063: frozen source line 2059 has the article error 'an finite morphism'.
\item $f$ がモノ射であることと $f'$ がモノ射であることは同値である。
\item $f$ が埋め込みであることと $f'$ が埋め込みであることは同値である。
\item ここにさらに追加する。
\end{enumerate}
\end{lemma}

\begin{proof}
スキーム $V'$ と全射エタール射 $V' \to Y'$ を選ぶ。
スキーム $U'$ と全射エタール射
$U' \to X' \times_{Y'} V'$ を選ぶ。
$V = Y \times_{Y'} V'$ および
$U = X \times_{X'} U'$ とおく。
すると射のエタール局所的な性質については、スキームの肥厚化の射
$(U \subset U') \to (V \subset V')$ に帰着し、
More on Morphisms, Lemma
\ref{more-morphisms-lemma-thicken-property-morphisms-cartesian}
を適用できる。これで (2), (3), (4), (5), (6) が証明される。

\medskip\noindent
(1), (7), (8), (9), (10) の射の性質は基底変換で安定なので、
$f'$ が性質 $\mathcal{P}$ をもてば $f$ ももつ。次を参照せよ。
Spaces, Lemma \ref{spaces-lemma-base-change-immersions}、および
Morphisms of Spaces, Lemmas
\ref{spaces-morphisms-lemma-base-change-proper},
\ref{spaces-morphisms-lemma-base-change-integral},
\ref{spaces-morphisms-lemma-base-change-monomorphism}。

\medskip\noindent
(1), (7), (8), (9), (10) で興味深い向きは、$f$ がその性質をもつと
仮定して、$f'$ ももつことを導く向きである。
肥厚化の次数に関する帰納法により、$Y \subset Y'$ は一次肥厚化であると
仮定してよい。上の有限次肥厚化に関する議論を参照せよ。

\medskip\noindent
(1) の証明。スキーム $V'$ と全射エタール射 $V' \to Y'$ を選ぶ。
$V = Y \times_{Y'} V'$, $U' = X' \times_{Y'} V'$,
$U = X \times_Y V$ とおく。すると $U \to V$ は閉埋め込みなので
$U$ はスキームであり、したがって $U'$ もスキームである
（Lemma \ref{spaces-more-morphisms-lemma-thickening-scheme}）。
ゆえにスキームの場合の補題
（More on Morphisms, Lemma
\ref{more-morphisms-lemma-thicken-property-morphisms-cartesian}）
を $(U \subset U') \to (V \subset V')$ に適用して結論を得る。

\medskip\noindent
(7) の証明。Lemma \ref{spaces-more-morphisms-lemma-thicken-property-morphisms} の結果 (2) と、
固有が準コンパクト $+$ 分離的 $+$ 局所有限型 $+$ 普遍閉であることを
組み合わせれば従う。

\medskip\noindent
(8) の証明。Lemma \ref{spaces-more-morphisms-lemma-thicken-property-morphisms} の結果 (2) と、
有限が整 $+$ 局所有限型であること
（Morphisms, Lemma \ref{morphisms-lemma-finite-integral}）を
組み合わせれば従う。

\medskip\noindent
(9) の証明。$f$ はモノ射なので $X = X \times_Y X$ である。
これまでに証明した結果を肥厚化の射
$(X \subset X') \to
(X \times_Y X \subset X' \times_{Y'} X')$
に適用してよい。(1) により
$X' \to X' \times_{Y'} X'$ は閉埋め込みであると結論する。
実際、閉埋め込み $X' \to X' \times_{Y'} X'$ を定めるイデアルは、
イデアル $\mathcal{I} \subset \mathcal{O}_{Y'}$ で
$Y$ を $Y'$ の中で切り出すものの引き戻しに含まれるので、
これは一次肥厚化である。実際、
$X = X \times_Y X =
(X' \times_{Y'} X') \times_{Y'} Y$ は $X'$ に含まれる。
閉埋め込み $\Delta : X' \to X' \times_{Y'} X'$ の余法層は
$\Omega_{X'/Y'}$ に等しい
（これは代数空間に対する Morphisms, Lemma
\ref{morphisms-lemma-differentials-diagonal} の類似であり、
エタール局所化により、または Lemmas
\ref{spaces-more-morphisms-lemma-differentials-relative-immersion-section} と
\ref{spaces-more-morphisms-lemma-differential-product} を組み合わせることにより従う。
細部は省略する）。したがって
$\Omega_{X'/Y'} = 0$ を示せば十分であり、これは (5) と
$X/Y$ に対する対応する主張から従う。

\medskip\noindent
(10) の証明。$f : X \to Y$ が埋め込みならば、
$X \to V \to Y$ と分解する。ここで $V \to Y$ は開部分空間であり、
$X \to V$ は閉埋め込みである。次を参照せよ。
Morphisms of Spaces, Remark
\ref{spaces-morphisms-remark-immersion}。
$V' \subset Y'$ を、台となる位相空間 $|V'|$ が
$|V| \subset |Y| = |Y'|$ と同じである開部分空間とする。
すると $X' \to Y'$ は $V'$ を経由し、(1) により
$X' \to V'$ は閉埋め込みであると結論する。これで証明が完了する。
\end{proof}

\noindent
次の補題は直前の補題の変種である。関係する肥厚化が有限次であると
仮定する代わりに（この仮定によれば、局所有限型であるという性質を
$f$ から $f'$ へ移せる）、$f$ と $f'$ がそれぞれ局所有限型であることを
仮定する。

\begin{lemma}
\label{spaces-more-morphisms-lemma-properties-that-extend-over-thickenings}
$S$ をスキームとする。
% P09-ERR-1065: frozen source line 2157 uses Y \to Y' inside a thickening pair; expected Y \subset Y'.
$(f, f') : (X \subset X') \to (Y \to Y')$ を
$S$ 上の代数空間の肥厚化の射とする。$f$ と $f'$ は局所有限型であり、
$X = Y \times_{Y'} X'$ であると仮定する。このとき次が成り立つ。
\begin{enumerate}
\item $f$ が局所準有限であることと $f'$ が局所準有限であることは
同値である。
\item $f$ が有限であることと $f'$ が有限であることは同値である。
\item $f$ が閉埋め込みであることと $f'$ が閉埋め込みであることは
同値である。
\item $\Omega_{X/Y} = 0$ であることと
$\Omega_{X'/Y'} = 0$ であることは同値である。
\item $f$ が非分岐であることと $f'$ が非分岐であることは同値である。
\item $f$ がモノ射であることと $f'$ がモノ射であることは同値である。
\item $f$ が埋め込みであることと $f'$ が埋め込みであることは同値である。
\item $f$ が固有であることと $f'$ が固有であることは同値である。
\item ここにさらに追加する。
\end{enumerate}
\end{lemma}

\begin{proof}
スキーム $V'$ と全射エタール射 $V' \to Y'$ を選ぶ。
スキーム $U'$ と全射エタール射
$U' \to X' \times_{Y'} V'$ を選ぶ。
$V = Y \times_{Y'} V'$ および
$U = X \times_{X'} U'$ とおく。
すると射のエタール局所的な性質については、スキームの肥厚化の射
$(U \subset U') \to (V \subset V')$ に帰着し、
More on Morphisms, Lemma
\ref{more-morphisms-lemma-properties-that-extend-over-thickenings}
を適用できる。これで (1), (4), (5) が証明される。

\medskip\noindent
(2), (3), (6), (7), (8) の性質は基底変換で安定なので、
$f'$ が性質 $\mathcal{P}$ をもてば $f$ ももつ。次を参照せよ。
Spaces, Lemma \ref{spaces-lemma-base-change-immersions}、および
Morphisms of Spaces, Lemmas
\ref{spaces-morphisms-lemma-base-change-proper},
\ref{spaces-morphisms-lemma-base-change-integral},
\ref{spaces-morphisms-lemma-base-change-monomorphism}。
したがって各場合に、$f$ が所望の性質をもつならば
$f'$ ももつことだけを証明すればよい。

\medskip\noindent
(2) は Lemma \ref{spaces-more-morphisms-lemma-thicken-property-morphisms} の (5) と、
有限射がちょうど局所有限型な整射であるという事実
（Morphisms of Spaces, Lemma
\ref{spaces-morphisms-lemma-finite-integral}）から従う。

\medskip\noindent
(3) の場合。これは Limits of Spaces, Lemma
\ref{spaces-limits-lemma-check-closed-infinitesimally}
から直ちに従う。

\medskip\noindent
(6) の証明。$f$ はモノ射なので $X = X \times_Y X$ である。
これまでに証明した結果を肥厚化の射
$(X \subset X') \to
(X \times_Y X \subset X' \times_{Y'} X')$
に適用してよい。(3) により
$\Delta_{X'/Y'} : X' \to X' \times_{Y'} X'$ は閉埋め込みである。
実際、$\Delta_{X'/Y'}$ は全単射
$|X'| \to |X' \times_{Y'} X'|$ を誘導するので、
$\Delta_{X'/Y'}$ は肥厚化である。他方、
$\Delta_{X'/Y'}$ は Morphisms of Spaces, Lemma
\ref{spaces-morphisms-lemma-diagonal-morphism-finite-type}
により局所有限表示である。言い換えると、
$\Delta_{X'/Y'}(X')$ は有限型の準連接イデアル層
$\mathcal{J} \subset \mathcal{O}_{X' \times_{Y'} X'}$
により切り出される。(5) により $\Omega_{X'/Y'} = 0$ なので、
閉埋め込み $X' \to X' \times_{Y'} X'$ の余法層は零である。
（閉埋め込み $\Delta_{X'/Y'}$ の余法層は
$\Omega_{X'/Y'}$ に等しい。これは代数空間に対する
Morphisms, Lemma \ref{morphisms-lemma-differentials-diagonal}
の類似であり、エタール局所化により、または Lemmas
\ref{spaces-more-morphisms-lemma-differentials-relative-immersion-section} と
\ref{spaces-more-morphisms-lemma-differential-product} を組み合わせることにより従う。
細部は省略する。）言い換えると
$\mathcal{J}/\mathcal{J}^2 = 0$ である。
例えば Algebra, Lemma
\ref{algebra-lemma-ideal-is-squared-union-connected}
により、これは $\Delta_{X'/Y'}$ が同型であることを意味する。

\medskip\noindent
(7) の証明。$f : X \to Y$ が埋め込みならば、
$X \to V \to Y$ と分解する。ここで $V \to Y$ は開部分空間であり、
$X \to V$ は閉埋め込みである。次を参照せよ。
Morphisms of Spaces, Remark
\ref{spaces-morphisms-remark-immersion}。
$V' \subset Y'$ を、台となる位相空間 $|V'|$ が
$|V| \subset |Y| = |Y'|$ と同じである開部分空間とする。
すると $X' \to Y'$ は $V'$ を経由し、(3) により
$X' \to V'$ は閉埋め込みであると結論する。

\medskip\noindent
(8) は Lemma \ref{spaces-more-morphisms-lemma-thicken-property-morphisms} と、
固有射が局所有限型、準コンパクト、普遍閉かつ分離的な射であるという
定義から従う。
\end{proof}










\section{肥厚化の Picard 群}
\label{spaces-more-morphisms-section-picard-group-thickening}

\noindent
肥厚化の Picard 群に関するいくつかの事項を述べる。

\begin{lemma}
\label{spaces-more-morphisms-lemma-picard-group-first-order-thickening}
$S$ をスキームとする。$X \subset X'$ を $S$ 上の代数空間の
一次肥厚化とし、そのイデアル層を $\mathcal{I}$ とする。
このとき、アーベル群の標準的な完全列
$$
\xymatrix{
0 \ar[r] &
H^0(X, \mathcal{I}) \ar[r] &
H^0(X', \mathcal{O}_{X'}^*) \ar[r] &
H^0(X, \mathcal{O}^*_X) \ar `r[d] `d[l] `l[llld] `d[dll] [dll] \\
& H^1(X, \mathcal{I}) \ar[r] &
\Pic(X') \ar[r] &
\Pic(X) \ar `r[d] `d[l] `l[llld] `d[dll] [dll] \\
& H^2(X, \mathcal{I}) \ar[r] & \ldots \ar[r] & \ldots
}
$$
が存在する。
\end{lemma}

\begin{proof}
$X_\etale = X'_\etale$ であることを思い出す。Lemma
\ref{spaces-more-morphisms-lemma-thickening-equivalence}、およびより一般には
Section \ref{spaces-more-morphisms-section-thickenings} の議論を参照せよ。
補題の列は、アーベル群の層の短完全列
$$
0 \to \mathcal{I} \to \mathcal{O}_{X'}^* \to \mathcal{O}_X^* \to 0
$$
に付随するコホモロジー長完全列である。これは $X_\etale$ 上で考える。
ここで最初の写像は局所切断 $f$、すなわち $\mathcal{I}$ の局所切断を
$1 + f$、すなわち $\mathcal{O}_{X'}$ の可逆切断に写す。
また、環付きサイトの Picard 群を可逆函数の層の一次コホモロジー群と
同一視する。Cohomology on Sites, Lemma
\ref{sites-cohomology-lemma-h1-invertible} を参照せよ。
\end{proof}








\section{無限小近傍}
\label{spaces-more-morphisms-section-first-order-infinitesimal-neighbourhood}

\noindent
有限次肥厚化の自然な構成は次のとおりである。
% P09-ERR-1066: frozen source line 2313 has 'Suppose that ... be' instead of 'Suppose that ... is'.
$i : Z \to X$ を代数空間の埋め込みとする。
開部分空間 $U \subset X$ を、$i$ が $Z$ を閉部分空間
$Z \subset U$ と同一視するように選ぶ
（Morphisms of Spaces, Remark
\ref{spaces-morphisms-remark-immersion} を参照）。
$\mathcal{I} \subset \mathcal{O}_U$ を、$Z$ を $U$ の中で定める
準連接イデアル層とする。次を参照せよ。
Morphisms of Spaces, Lemma
\ref{spaces-morphisms-lemma-closed-immersion-ideals}。
$n \geq 1$ に対して閉部分空間 $Z_n \subset U$ を考えられる。
これは準連接イデアル層 $\mathcal{I}^{n + 1}$ により定められる。

\begin{definition}
\label{spaces-more-morphisms-definition-first-order-infinitesimal-neighbourhood}
$i : Z \to X$ を代数空間の埋め込みとする。
\begin{enumerate}
\item $Z$ の $X$ における {\it 一次無限小近傍}とは、
上で述べた一次肥厚化 $Z \subset Z_1$、すなわち $X$ 上のものをいう。
\item {\it $n$ 次無限小近傍}、すなわち $Z$ の $X$ におけるものとは、
上で述べた $n$ 次肥厚化 $Z \subset Z_n$、すなわち $X$ 上のものをいう。
\end{enumerate}
\end{definition}

\noindent
この肥厚化は次の普遍性をもつ
（したがって、上の構成が開集合 $U$ の選択に依存するという懸念は解消される）。

\begin{lemma}
\label{spaces-more-morphisms-lemma-first-order-infinitesimal-neighbourhood}
$i : Z \to X$ を代数空間の埋め込みとする。
\begin{enumerate}
\item 一次無限小近傍 $Z'$、すなわち $Z$ の $X$ におけるものは
次の普遍性をもつ。
任意の可換図式
$$
\xymatrix{
Z \ar[d]_i & T \ar[l]^a \ar[d] \\
X & T' \ar[l]_b
}
$$
が与えられ、$T \subset T'$ が $X$ 上の一次肥厚化であるとする。
このとき一意な肥厚化の射
$(a', a) : (T \subset T') \to (Z \subset Z')$ が存在し、
$X$ 上の射である。
\item $n \geq 1$ に対して、$n$ 次無限小近傍 $Z_n$、
すなわち $Z$ の $X$ におけるものは次の普遍性をもつ。
任意の可換図式
$$
\xymatrix{
Z \ar[d]_i & T \ar[l]^a \ar[d] \\
X & T' \ar[l]_b
}
$$
が与えられ、$T \subset T'$ が $n$ 次肥厚化であり、$X$ 上にあるとする。
このとき一意な肥厚化の射
$(a', a) : (T \subset T') \to (Z \subset Z_n)$ が存在し、
$X$ 上の射である。
\end{enumerate}
\end{lemma}

\begin{proof}
(1) だけを証明する。$U \subset X$ を $Z'$ の構成で用いた
開部分空間、すなわち $Z$ が $U$ の閉部分空間と同一視され、
準連接イデアル層 $\mathcal{I}$ により切り出されるような開集合とする。
$|T| = |T'|$ なので、$|b|(|T'|) \subset |U|$ である。
したがって $b$ を $U$ への射と考えてよい。次を参照せよ。
Properties of Spaces, Lemma
\ref{spaces-properties-lemma-factor-through-open-subspace}。
$\mathcal{J} \subset \mathcal{O}_{T'}$ を、$T$ を切り出す二乗零の
準連接イデアル層とする。図式の可換性により
$b|_T = i \circ a$ である。ここで $i : Z \to U$ は閉埋め込みである。
したがって Morphisms of Spaces, Lemma
\ref{spaces-morphisms-lemma-closed-immersion-ideals} により
$b^\sharp(b^{-1}\mathcal{I}) \subset \mathcal{J}$ である。
$T'$ は $T$ の一次肥厚化なので $\mathcal{J}^2 = 0$ であり、したがって
$b^\sharp(b^{-1}(\mathcal{I}^2)) = 0$ である。
Morphisms of Spaces, Lemma
\ref{spaces-morphisms-lemma-closed-immersion-ideals}
により、これは $b$ が $Z'$ を経由することを意味する。
$a' : T' \to Z'$ をこの分解とすればよい。
\end{proof}

\begin{lemma}
\label{spaces-more-morphisms-lemma-infinitesimal-neighbourhood-conormal}
$i : Z \to X$ を代数空間の埋め込みとする。
$Z \subset Z'$ を $Z$ の $X$ における一次無限小近傍とする。
このとき図式
$$
\xymatrix{
Z \ar[r] \ar[d] & Z' \ar[d] \\
Z \ar[r] & X
}
$$
は Lemma \ref{spaces-more-morphisms-lemma-conormal-functorial} により余法層の写像
$\mathcal{C}_{Z/X} \to \mathcal{C}_{Z/Z'}$ を誘導する。
この写像は同型である。
\end{lemma}

\begin{proof}
これは上の $Z'$ の構成から明らかである。
\end{proof}





\section{形式的に滑らか、エタール、非分岐な変換}
\label{spaces-more-morphisms-section-formally-smooth-etale-unramified}

\noindent
環準同型 $R \to A$ が {\it 形式的に滑らか}、
それぞれ {\it 形式的エタール}、{\it 形式的非分岐}であるとは、
Algebra, Definition \ref{algebra-definition-formally-smooth}、
それぞれ Definition \ref{algebra-definition-formally-etale}、
Definition \ref{algebra-definition-formally-unramified} の意味で、
任意の実線部分が可換な図式
$$
\xymatrix{
A \ar[r] \ar@{-->}[rd] & B/I \\
R \ar[r] \ar[u] & B \ar[u]
}
$$
で、$I \subset B$ が二乗零イデアルであるものに対し、
図式を可換にする点線矢印が存在する、それぞれ一意に存在する、
高々一つ存在することをいう。このことが、代数空間の射、さらに一般には
関手に対する次の類似物を動機づける。

\begin{definition}
\label{spaces-more-morphisms-definition-formally-smooth-etale-unramified}
$S$ をスキームとする。$a : F \to G$ を関手
$F, G : (\Sch/S)_{fppf}^{opp} \to \textit{Sets}$ の変換とする。
次の形の実線部分が可換な図式を考える。
$$
\xymatrix{
F \ar[d]_a & T \ar[d]^i \ar[l] \\
G & T' \ar[l] \ar@{-->}[lu]
}
$$
ここで $T$ と $T'$ はアフィンスキームであり、$i$ は二乗零イデアルで
定められる閉埋め込みである。
\begin{enumerate}
\item 上の任意の実線図式に対して、図式を可換にする点線矢印が存在するとき、
$a$ は {\it 形式的に滑らか}であるという\footnote{これはここで採用できる
定義の一つにすぎない。少し弱い別の条件として、点線矢印が $T'$ 上
fppf 局所的に存在することを要求できる。この弱い概念は、ある意味で
よりよい形式的性質をもつ。}。
\item 上の任意の実線図式に対して、図式を可換にする点線矢印がちょうど一つ
存在するとき、$a$ は {\it 形式的エタール}であるという。
\item 上の任意の実線図式に対して、図式を可換にする点線矢印が高々一つ
存在するとき、$a$ は {\it 形式的非分岐}であるという。
\end{enumerate}
\end{definition}

\begin{lemma}
\label{spaces-more-morphisms-lemma-formally-etale-is-combination}
$S$ をスキームとする。$a : F \to G$ を関手
$F, G : (\Sch/S)_{fppf}^{opp} \to \textit{Sets}$ の変換とする。
このとき $a$ が形式的エタールであることと、
$a$ が形式的に滑らかかつ形式的非分岐であることは同値である。
\end{lemma}

\begin{proof}
定義から形式的に従う。
\end{proof}

\begin{lemma}
\label{spaces-more-morphisms-lemma-composition-formally-smooth-etale-unramified}
合成について次が成り立つ。
\begin{enumerate}
\item 形式的に滑らかな関手の変換の合成は形式的に滑らかである。
\item 形式的エタールな関手の変換の合成は形式的エタールである。
\item 形式的非分岐な関手の変換の合成は形式的非分岐である。
\end{enumerate}
\end{lemma}

\begin{proof}
これは形式的である。
\end{proof}

\begin{lemma}
\label{spaces-more-morphisms-lemma-base-change-formally-smooth-etale-unramified}
$S$ を $\Sch_{fppf}$ に含まれるスキームとする。
$F, G, H : (\Sch/S)_{fppf}^{opp} \to \textit{Sets}$ とする。
$a : F \to G$, $b : H \to G$ を関手の変換とする。
ファイバー積図式
$$
\xymatrix{
H \times_{b, G, a} F \ar[r]_-{b'} \ar[d]_{a'} & F \ar[d]^a \\
H \ar[r]^b & G
}
$$
を考える。
\begin{enumerate}
\item $a$ が形式的に滑らかならば、基底変換 $a'$ は
形式的に滑らかである。
\item $a$ が形式的エタールならば、基底変換 $a'$ は
形式的エタールである。
\item $a$ が形式的非分岐ならば、基底変換 $a'$ は
形式的非分岐である。
\end{enumerate}
\end{lemma}

\begin{proof}
これは形式的である。
\end{proof}

\begin{lemma}
\label{spaces-more-morphisms-lemma-representable-property-formally-property}
$S$ をスキームとする。
$F, G : (\Sch/S)_{fppf}^{opp} \to \textit{Sets}$ とする。
$a : F \to G$ を表現可能な関手の変換とする。
\begin{enumerate}
\item $a$ が滑らかならば $a$ は形式的に滑らかである。
\item $a$ がエタールならば $a$ は形式的エタールである。
\item $a$ が非分岐ならば $a$ は形式的非分岐である。
\end{enumerate}
\end{lemma}

\begin{proof}
Definition \ref{spaces-more-morphisms-definition-formally-smooth-etale-unramified} のような
実線部分が可換な図式
$$
\xymatrix{
F \ar[d]_a & T \ar[d]^i \ar[l] \\
G & T' \ar[l] \ar@{-->}[lu]
}
$$
を考える。このとき $F \times_G T'$ は $T'$ 上滑らかな
（それぞれエタールな、非分岐な）スキームである。したがって
More on Morphisms, Lemma
\ref{more-morphisms-lemma-smooth-formally-smooth}
（それぞれ Lemma \ref{more-morphisms-lemma-etale-formally-etale}、
Lemma \ref{more-morphisms-lemma-unramified-formally-unramified}）により、
図式
$$
\xymatrix{
F \times_G T' \ar[d] & T \ar[d]^i \ar[l] \\
T' & T' \ar[l] \ar@{-->}[lu]
}
$$
の点線矢印を補える
（それぞれ一意に補える、高々一通りに補える）。
% P09-ERR-1067: frozen source line 2563 has 'an hence' instead of 'and hence'.
したがって最初の図式でも対応する主張を得る。
\end{proof}

\begin{lemma}
\label{spaces-more-morphisms-lemma-etale-on-top}
$S$ を $\Sch_{fppf}$ に含まれるスキームとする。
$F, G, H : (\Sch/S)_{fppf}^{opp} \to \textit{Sets}$ とする。
$a : F \to G$, $b : G \to H$ を関手の変換とする。
$a$ は表現可能、全射かつエタールであると仮定する。
\begin{enumerate}
\item $b$ が形式的に滑らかならば、$b \circ a$ は形式的に滑らかである。
\item $b$ が形式的エタールならば、$b \circ a$ は形式的エタールである。
\item $b$ が形式的非分岐ならば、$b \circ a$ は形式的非分岐である。
\end{enumerate}
逆に、実線部分が可換な図式
$$
\xymatrix{
G \ar[d]_b & T \ar[d]^i \ar[l] \\
H & T' \ar[l] \ar@{-->}[lu]
}
$$
を考える。ここで $T'$ は $S$ 上のアフィンスキームであり、
$i : T \to T'$ は二乗零イデアルで定められる閉埋め込みである。
\begin{enumerate}
\item[(4)] $b \circ a$ が形式的に滑らかならば、任意の $t \in T$ に対し、
アフィンスキームのエタール射 $U' \to T'$ と射 $U' \to G$ であって、
$$
\xymatrix{
G \ar[d]_b & T \ar[l] & T \times_{T'} U' \ar[d] \ar[l]\\
H & T' \ar[l] & U' \ar[llu] \ar[l]
}
$$
が可換となり、かつ $t$ が $U' \to T'$ の像に属するものが存在する。
\item[(5)] $b \circ a$ が形式的非分岐ならば、上の図式の点線矢印は
高々一つしか存在しない。すなわち $b$ は形式的非分岐である。
\item[(6)] $b \circ a$ が形式的エタールならば、上の図式の点線矢印は
ちょうど一つ存在する。すなわち $b$ は形式的エタールである。
\end{enumerate}
\end{lemma}

\begin{proof}
$b$ が形式的に滑らか（それぞれ形式的エタール、形式的非分岐）であると仮定する。
エタール射は滑らかかつ非分岐なので、$a$ は表現可能かつ滑らか
（それぞれエタール、非分岐）である。したがって (1)、(2)、(3) は
Lemma \ref{spaces-more-morphisms-lemma-representable-property-formally-property}
と
Lemma \ref{spaces-more-morphisms-lemma-composition-formally-smooth-etale-unramified}
を組み合わせれば従う。

\medskip\noindent
$b \circ a$ が形式的に滑らかであると仮定し、補題の主張にある図式を考える。
$W = F \times_G T$ とおく。仮定により $W$ は $T$ 上全射エタールな
スキームである。\'Etale Morphisms, Theorem
\ref{etale-theorem-remarkable-equivalence} により、スキーム $W'$ で
$T'$ 上エタールかつ $W = T \times_{T'} W'$ を満たすものが存在する。
アフィン開部分スキーム $U' \subset W'$ で、$t$ が $U' \to T'$ の像に
属するものを選ぶ。$b \circ a$ は形式的に滑らかなので、
% P09-ERR-1068: frozen source line 2622 has "the exist morphisms".
次の図式を可換にする射 $U' \to F$ が存在する。
$$
\xymatrix{
F \ar[d]_{b \circ a} & W \ar[l] & T \times_{T'} U' \ar[d] \ar[l]\\
H & T' \ar[l] & U' \ar[llu] \ar[l]
}
$$
合成 $U' \to F \to G$ を取れば、補題の
% P09-ERR-1069: frozen source line 2630 cites part (5), but this construction proves part (4).
(4) にある写像が得られる。

\medskip\noindent
$f, g : T' \to G$ を、補題の図式に入る二つの点線矢印とする。
$W = F \times_G T$ とおく。仮定により $W$ は $T$ 上全射エタールな
スキームである。\'Etale Morphisms, Theorem
\ref{etale-theorem-remarkable-equivalence} により、スキーム $W'$ で
$T'$ 上エタールかつ $W = T \times_{T'} W'$ を満たすものが存在する。
$a$ は形式的エタールなので、合成
$$
W' \to T' \xrightarrow{f} G
\quad\text{and}\quad
W' \to T' \xrightarrow{g} G
$$
は射 $f', g' : W' \to F$ に持ち上がる
（アフィン開集合上で持ち上げ、一意性によって貼り合わせる）。
ここで $b \circ a : F \to H$ が形式的非分岐ならば $f' = g'$ であり、
したがって $f = g$ となる。実際、$W' \to T'$ はエタール被覆である。これは
% P09-ERR-1070: frozen source line 2647 says part (6); the uniqueness argument proves part (5).
(5) を証明する。

\medskip\noindent
$b \circ a$ が形式的エタールであると仮定する。このとき (4) により
$T'$ 上エタール局所的に図式に入る点線矢印を見つけることができ、
(5) によりこの点線矢印は一意である。したがって局所解を貼り合わせて
(6) を得る。細部は省略する。
\end{proof}

\begin{remark}
\label{spaces-more-morphisms-remark-tempting}
Lemma \ref{spaces-more-morphisms-lemma-etale-on-top} の状況では
「$b$ が形式的に滑らか」$\Leftrightarrow$
「$b \circ a$ が形式的に滑らか」
であると考えたくなる。しかし一般には、おそらくこれは正しくない。
\end{remark}

\begin{lemma}
\label{spaces-more-morphisms-lemma-formally-permanence}
$S$ をスキームとする。
$F, G, H : (\Sch/S)_{fppf}^{opp} \to \textit{Sets}$ とする。
$a : F \to G$, $b : G \to H$ を関手の変換とする。
$b$ は形式的非分岐であると仮定する。
\begin{enumerate}
\item $b \circ a$ が形式的非分岐ならば $a$ は形式的非分岐である。
\item $b \circ a$ が形式的エタールならば $a$ は形式的エタールである。
\item $b \circ a$ が形式的に滑らかならば $a$ は形式的に滑らかである。
\end{enumerate}
\end{lemma}

\begin{proof}
$T \subset T'$ を、二乗零イデアルで定められるアフィンスキームの閉埋め込みとする。
$g' : T' \to G$ と $f : T \to F$ が与えられ、
$g'|_T = a \circ f$ を満たすとする。$b$ は形式的非分岐なので、
$$
\{f' : T' \to F \mid f = f'|_T\text{ and }a \circ f' = g'\}
$$
と
$$
\{f' : T' \to F \mid f = f'|_T\text{ and }b \circ a \circ f' = b \circ g'\}.
$$
の間には一対一対応がある。これから補題は形式的に従う。
\end{proof}








\section{形式的非分岐射}
\label{spaces-more-morphisms-section-formally-unramified}

\noindent
この節では、代数空間の射が形式的非分岐であることの意味を詳しく調べる。

\begin{definition}
\label{spaces-more-morphisms-definition-formally-unramified}
$S$ をスキームとする。代数空間の射 $f : X \to Y$ が $S$ 上で
{\it 形式的非分岐}であるとは、Definition
\ref{spaces-more-morphisms-definition-formally-smooth-etale-unramified} の意味で、
関手の変換として形式的非分岐であることをいう。
\end{definition}

\noindent
Section \ref{spaces-more-morphisms-section-formally-smooth-etale-unramified} で、より一般の
形式的非分岐な関手の変換について証明した結果はここでは再掲しない。
この性質は相対微分加群の消滅によって特徴づけられることが分かる。
Lemma \ref{spaces-more-morphisms-lemma-characterize-formally-unramified} を参照せよ。

\begin{lemma}
\label{spaces-more-morphisms-lemma-formally-unramified}
$S$ をスキームとする。$f : X \to Y$ を $S$ 上の代数空間の射とする。
次は同値である。
\begin{enumerate}
\item $f$ は形式的非分岐である。
\item 任意の図式
$$
\xymatrix{
U \ar[d] \ar[r]_\psi & V \ar[d] \\
X \ar[r]^f & Y
}
$$
で、$U$ と $V$ がスキームであり、縦の矢印がエタールであるものに対して、
スキームの射 $\psi$ は形式的非分岐である（More on Morphisms,
Definition \ref{more-morphisms-definition-formally-unramified} の意味で）。
\item 全射な縦の矢印をもつこのような図式の一つについて、射 $\psi$ は
形式的非分岐である。
\end{enumerate}
\end{lemma}

\begin{proof}
$f$ が形式的非分岐であると仮定する。Lemma
\ref{spaces-more-morphisms-lemma-representable-property-formally-property} により、射
$U \to X$ と $V \to Y$ は形式的非分岐である。したがって Lemma
\ref{spaces-more-morphisms-lemma-composition-formally-smooth-etale-unramified} により、合成
$U \to Y$ は形式的非分岐である。すると Lemma
\ref{spaces-more-morphisms-lemma-formally-permanence} から $U \to V$ は形式的非分岐である。
ゆえに (1) は (2) を含意する。また (2) が (3) を含意することは明らかである。

\medskip\noindent
(3) にある図式が与えられたと仮定する。Lemma
\ref{spaces-more-morphisms-lemma-representable-property-formally-property} により、射
$V \to Y$ は形式的非分岐である。したがって Lemma
\ref{spaces-more-morphisms-lemma-composition-formally-smooth-etale-unramified} により、合成
$U \to Y$ は形式的非分岐である。すると Lemma
\ref{spaces-more-morphisms-lemma-etale-on-top} から $X \to Y$ は形式的非分岐、すなわち (1) が成り立つ。
\end{proof}

\begin{lemma}
\label{spaces-more-morphisms-lemma-formally-unramified-not-affine}
$S$ をスキームとする。
$f : X \to Y$ を $S$ 上の代数空間の形式的非分岐射とする。このとき、
任意の実線部分が可換な図式
$$
\xymatrix{
X \ar[d]_f & T \ar[d]^i \ar[l] \\
Y & T' \ar[l] \ar@{-->}[lu]
}
$$
で、$T \subset T'$ が $S$ 上の代数空間の一次肥厚化であるものに対し、
図式を可換にする点線矢印は高々一つ存在する。言い換えれば、
Definition \ref{spaces-more-morphisms-definition-formally-unramified} における
$T$ がアフィンスキームであるという条件は取り除ける。
\end{lemma}

\begin{proof}
全射エタール射 $U' \to T'$ で、$U'$ がアフィンスキームの非交和である
ものが存在し（Properties of Spaces, Lemma
\ref{spaces-properties-lemma-cover-by-union-affines} を参照）、
射 $T' \to X$ はその $U'$ への制限によって決まるからである。
\end{proof}

\begin{lemma}
\label{spaces-more-morphisms-lemma-composition-formally-unramified}
形式的非分岐射の合成は形式的非分岐である。
\end{lemma}

\begin{proof}
これは形式的である。
\end{proof}

\begin{lemma}
\label{spaces-more-morphisms-lemma-base-change-formally-unramified}
形式的非分岐射の基底変換は形式的非分岐である。
\end{lemma}

\begin{proof}
これは形式的である。
\end{proof}

\begin{lemma}
\label{spaces-more-morphisms-lemma-characterize-formally-unramified}
$S$ をスキームとする。$f : X \to Y$ を $S$ 上の代数空間の射とする。
次は同値である。
\begin{enumerate}
\item $f$ は形式的非分岐である。
\item $\Omega_{X/Y} = 0$ である。
\end{enumerate}
\end{lemma}

\begin{proof}
これは Lemma \ref{spaces-more-morphisms-lemma-formally-unramified}、
More on Morphisms, Lemma
\ref{more-morphisms-lemma-formally-unramified-differentials}、および
Lemma \ref{spaces-more-morphisms-lemma-localize-differentials} を組み合わせたものである。
\end{proof}

\begin{lemma}
\label{spaces-more-morphisms-lemma-unramified-formally-unramified}
$S$ をスキームとする。
$f : X \to Y$ を $S$ 上の代数空間の射とする。
次は同値である。
\begin{enumerate}
\item 射 $f$ は非分岐である。
\item 射 $f$ は局所有限型であり、$\Omega_{X/Y} = 0$ である。
\item 射 $f$ は局所有限型かつ形式的非分岐である。
\end{enumerate}
\end{lemma}

\begin{proof}
図式
$$
\xymatrix{
U \ar[d] \ar[r]_\psi & V \ar[d] \\
X \ar[r]^f & Y
}
$$
で、$U$ と $V$ がスキームであり、縦の矢印がエタールかつ全射であるものを選ぶ。
すると次を得る。
\begin{align*}
f\text{ unramified}
& \Leftrightarrow
\psi\text{ unramified} \\
& \Leftrightarrow
% P09-ERR-1072: frozen source line 2855 omits "of" in "locally finite type".
\psi\text{ locally finite type and }\Omega_{U/V} = 0 \\
& \Leftrightarrow
f\text{ locally finite type and }\Omega_{X/Y} = 0 \\
& \Leftrightarrow
f\text{ locally finite type and formally unramified}
\end{align*}
ここで Morphisms, Lemma \ref{morphisms-lemma-unramified-omega-zero} と
Lemma \ref{spaces-more-morphisms-lemma-characterize-formally-unramified} を用いた。
\end{proof}

\begin{lemma}
\label{spaces-more-morphisms-lemma-universally-injective-unramified}
$S$ をスキームとする。
$f : X \to Y$ を $S$ 上の代数空間の射とする。
次は同値である。
\begin{enumerate}
\item $f$ は非分岐かつ単射である。
\item $f$ は非分岐かつ普遍的単射である。
\item $f$ は局所有限型かつ単射である。
\item $f$ は普遍的単射かつ局所有限型かつ形式的非分岐である。
\end{enumerate}
さらに、この場合 $f$ は表現可能、分離的かつ局所準有限でもある。
\end{lemma}

\begin{proof}
Lemma \ref{spaces-more-morphisms-lemma-unramified-formally-unramified} で、形式的非分岐かつ
局所有限型であることは非分岐であることと同値であると見た。
したがって (4) は (2) と同値である。単射は確かに形式的非分岐なので
(3) は (4) を含意する。(1) が (3) を含意することは明らかである。
最後に (2) が成り立つならば、
$\Delta : X \to X \times_Y X$ は開埋め込みであり
（Morphisms of Spaces, Lemma
\ref{spaces-morphisms-lemma-diagonal-unramified-morphism}）、かつ全射である
（Morphisms of Spaces, Lemma
\ref{spaces-morphisms-lemma-universally-injective}）。したがって同型であり、
すなわち $f$ は単射である。このようにして (2) が (1) を含意することが分かる。
最後に Morphisms of Spaces, Lemmas
\ref{spaces-morphisms-lemma-monomorphism-loc-finite-type-loc-quasi-finite} and
\ref{spaces-morphisms-lemma-locally-quasi-finite-separated-representable}
により、$f$ は表現可能、分離的かつ局所準有限である。
\end{proof}

\begin{lemma}
\label{spaces-more-morphisms-lemma-characterize-closed-immersion}
$S$ をスキームとする。
$f : X \to Y$ を $S$ 上の代数空間の射とする。
次は同値である。
\begin{enumerate}
\item $f$ は閉埋め込みである。
\item $f$ は普遍的閉、非分岐かつ単射である。
\item $f$ は普遍的閉、非分岐かつ普遍的単射である。
\item $f$ は普遍的閉、局所有限型かつ単射である。
\item $f$ は普遍的閉、普遍的単射、局所有限型かつ形式的非分岐である。
\end{enumerate}
\end{lemma}

\begin{proof}
(2) -- (5) の同値性は Lemma
\ref{spaces-more-morphisms-lemma-universally-injective-unramified} から直ちに従う。
さらに (2) -- (5) が成り立つならば $f$ は表現可能である。
同様に (1) が成り立つならば $f$ は表現可能である。
したがって結果はスキームの場合から従う。\'Etale Morphisms, Lemma
\ref{etale-lemma-characterize-closed-immersion} を参照せよ。
\end{proof}






\section{普遍一次肥厚化}
\label{spaces-more-morphisms-section-universal-thickening}

\noindent
$S$ をスキームとする。
$h : Z \to X$ を $S$ 上の代数空間の射とする。
$Z$ の $X$ 上の {\it 普遍一次肥厚化}とは、一次肥厚化
$Z \subset Z'$ で $X$ 上にあるもののうち、任意の一次肥厚化
$T \subset T'$ で $X$ 上にあるものと
実線部分が可換な図式
\begin{equation}
\label{spaces-more-morphisms-equation-universal-first-order-thickening}
\vcenter{
\xymatrix{
& Z \ar[ld] & & T \ar[rd] \ar[ll]^a \\
Z' \ar[rrd] & & & & T' \ar@{..>}[llll]_{a'} \ar[lld]^b \\
 & & X
}
}
\end{equation}
が与えられるたびに、図式を可換にする点線矢印が一意に存在するものをいう。
この状況では $(a, a') : (T \subset T') \to (Z \subset Z')$ は
$X$ 上の肥厚化の射であることに注意せよ。したがって普遍一次肥厚化が
存在すれば、一意な同型を除いて一意である。一般には普遍一次肥厚化は
存在しないが、$h$ が形式的非分岐ならば存在する。これを証明する前に、
スキームの圏における普遍一次肥厚化が代数空間の圏においても
普遍一次肥厚化であることを示そう。

\begin{lemma}
\label{spaces-more-morphisms-lemma-check-universal-first-order-thickening}
$S$ をスキームとする。
$h : Z \to X$ を $S$ 上の代数空間の射とする。
$Z \subset Z'$ を $X$ 上の一次肥厚化とする。
次は同値である。
\begin{enumerate}
\item $Z \subset Z'$ は普遍一次肥厚化である。
\item $T'$ がスキームである任意の図式
(\ref{spaces-more-morphisms-equation-universal-first-order-thickening}) に対し、図式を可換にする
点線矢印が一意に存在する。
\item $T'$ がアフィンスキームである任意の図式
(\ref{spaces-more-morphisms-equation-universal-first-order-thickening}) に対し、図式を可換にする
点線矢印が一意に存在する。
\end{enumerate}
\end{lemma}

\begin{proof}
(1) $\Rightarrow$ (2) $\Rightarrow$ (3) は形式的である。
% P09-ERR-1073: frozen source line 2981 has "Assume (3) a assume".
(3) を仮定し、任意の図式
(\ref{spaces-more-morphisms-equation-universal-first-order-thickening}) が与えられたとする。
表示 $T' = U'/R'$ を選ぶ。Spaces, Definition
\ref{spaces-definition-presentation} を参照せよ。
$U' = \coprod U'_i$ はアフィンスキームの非交和であるとしてよい。したがって
% P09-ERR-1074: frozen source line 2986 writes \times_T' instead of \times_{T'}.
$R' = U' \times_{T'} U' = \coprod_{i, j} U'_i \times_T' U'_j$ である。
各組 $(i, j)$ に対し、アフィン開被覆
$U'_i \times_T' U'_j = \bigcup_k R'_{ijk}$ を選ぶ。
$U_i, R_{ijk}$ を、それぞれ $T$ と $T'$ 上で取ったファイバー積とする。
すると各 $U_i \subset U'_i$ と $R_{ijk} \subset R'_{ijk}$ は
アフィンスキームの一次肥厚化である。
$a_i : U_i \to Z$、それぞれ $a_{ijk} : R_{ijk} \to Z$ を、
$a : T \to Z$ と射 $U_i \to T$、それぞれ $R_{ijk} \to T$ との合成とする。
(3) を $a_i : U_i \to Z$ に適用すると、一意な射
$a'_i : U'_i \to Z'$ を得る。
(3) を $a_{ijk}$ に適用すると、二つの合成
% P09-ERR-1075: frozen source line 2998 uses undefined R'_i and R'_j; U'_i and U'_j appear intended.
$R'_{ijk} \to R'_i \to Z'$ と $R'_{ijk} \to R'_j \to Z'$ は等しい。
したがって $a' = \coprod a'_i : U' = \coprod U'_i \to Z'$ は
商層 $T' = U'/R'$ に降下し、証明が完了する。
\end{proof}

\begin{lemma}
\label{spaces-more-morphisms-lemma-universal-thickening-over-formally-etale}
$S$ をスキームとする。
$Z \to Y \to X$ を $S$ 上の代数空間の射とする。
$Z \subset Z'$ が $Z$ の $Y$ 上の普遍一次肥厚化であり、
$Y \to X$ が形式的エタールならば、$Z \subset Z'$ は
$Z$ の $X$ 上の普遍一次肥厚化である。
\end{lemma}

\begin{proof}
これは形式的である。実際、Lemma
\ref{spaces-more-morphisms-lemma-check-universal-first-order-thickening} により、$T'$ が
アフィンスキームである実線部分が可換な図式
(\ref{spaces-more-morphisms-equation-universal-first-order-thickening}) だけを考えれば十分である。
合成 $T \to Z \to Y$ は $T' \to Y$ へ一意に持ち上がる。
これは $Y \to X$ が形式的エタールと仮定されているからである。
したがって $Z \subset Z'$ が $Y$ 上の普遍一次肥厚化であることから、
求める射 $a' : T' \to Z'$ が得られる。
\end{proof}

\begin{lemma}
\label{spaces-more-morphisms-lemma-etale-morphism-of-universal-thickenings}
$S$ をスキームとする。
$Z \to Y \to X$ を $S$ 上の代数空間の射とする。
$Z \to Y$ はエタールであると仮定する。
\begin{enumerate}
\item $Y \subset Y'$ が $Y$ の $X$ 上の普遍一次肥厚化ならば、
一意なエタール射 $Z' \to Y'$ で $Z = Y \times_{Y'} Z'$ を満たすもの
（Theorem \ref{spaces-more-morphisms-theorem-topological-invariance} を参照）によって、
$Z$ の $X$ 上の普遍一次肥厚化が与えられる。
\item $Z \to Y$ が全射であり、$(Z \subset Z') \to (Y \subset Y')$ が
$X$ 上の一次肥厚化のエタール射で、$Z'$ が $Z$ の $X$ 上の
普遍一次肥厚化ならば、$Y'$ は $Y$ の $X$ 上の普遍一次肥厚化である。
\end{enumerate}
\end{lemma}

\begin{proof}
(1) の証明。Lemma \ref{spaces-more-morphisms-lemma-check-universal-first-order-thickening} により、
$T'$ がアフィンスキームである実線部分が可換な図式
(\ref{spaces-more-morphisms-equation-universal-first-order-thickening}) だけを考えれば十分である。
合成 $T \to Z \to Y$ は $T' \to Y'$ へ一意に持ち上がる。
これは $Y'$ が普遍一次肥厚化だからである。さらに $Z' \to Y'$ がエタールであることから
（Lemma \ref{spaces-more-morphisms-lemma-representable-property-formally-property} を参照）、
$T' \to Y'$ は求める射 $a' : T' \to Z'$ へ持ち上がる。

\medskip\noindent
(2) の証明。$T \subset T'$ を $X$ 上の一次肥厚化とし、
$a : T \to Y$ を射とする。$W = T \times_Y Z$ とおき、射影を
$c : W \to Z$ と書く。一意なエタール射 $W' \to T'$ で
$W = T \times_{T'} W'$ を満たすものを取る。Theorem
\ref{spaces-more-morphisms-theorem-topological-invariance} を参照せよ。
$W' \to T'$ は全射である。これは $Z \to Y$ が全射だからである。
仮定により、一意な射 $c' : W' \to Z'$ を得る。この射は $X$ 上の射で、
$c$ を $W$ 上に延長する。一意性により、$c'$ の
$W' \times_{T'} W'$ への二つの制限は等しい（$c$ の
$W \times_T W$ への二つの制限が等しいからである）。
したがって $c'$ は一意な射 $a' : T' \to Y'$ に降下し、証明が完了する。
\end{proof}

\begin{lemma}
\label{spaces-more-morphisms-lemma-universal-thickening}
$S$ をスキームとする。
$h : Z \to X$ を $S$ 上の代数空間の形式的非分岐射とする。
$Z \subset Z'$ で、$Z$ の $X$ 上の普遍一次肥厚化であるものが存在する。
\end{lemma}

\begin{proof}
任意の可換図式
$$
\xymatrix{
V \ar[d] \ar[r] & U \ar[d] \\
Z \ar[r] & X
}
$$
で、$V$ と $U$ がスキームであり、縦の矢印がエタールであるものを選ぶ。
$V \to U$ はスキームの形式的非分岐射であることに注意せよ。
Lemma \ref{spaces-more-morphisms-lemma-formally-unramified} を参照せよ。
Lemma \ref{spaces-more-morphisms-lemma-check-universal-first-order-thickening} と
More on Morphisms, Lemma
\ref{more-morphisms-lemma-universal-thickening} を組み合わせると、
$V \subset V'$ が存在し、それは $V$ の $U$ 上の普遍一次肥厚化である。
% P09-ERR-1076: frozen source line 3097 cites a nonexistent "part (1)" of this lemma.
Lemma \ref{spaces-more-morphisms-lemma-universal-thickening-over-formally-etale} part (1) により、
$V'$ は $V$ の $X$ 上の普遍一次肥厚化である。

\medskip\noindent
スキーム $U$ と全射エタール射 $U \to X$ を固定する。
上の議論は、任意のエタール射 $V \to Z$ で、$V$ がスキームであり、
$V \to Z \to X$ が $U$ を経由するものに対し、普遍一次肥厚化
$V \subset V'$ が存在し、それが $V$ の $X$ 上にあることを示す
（ただし $V \to X$ の $U$ を
経由する分解の選択には依存しない）。ここで $V$ を、$V \to Z$ が
全射エタールとなるように選べる（Spaces, Lemma
\ref{spaces-lemma-lift-morphism-presentations} を参照）。すると
% P09-ERR-1077: frozen source line 3109 omits "is" before "a scheme".
$R = V \times_Z V$ は $Z$ 上エタールなスキームであり、$R \to X$ も
$U$ を経由する。したがって普遍一次肥厚化
$V \subset V'$ と $R \subset R'$ を $X$ 上に得る。
$V \subset V'$ は普遍一次肥厚化なので、二つの射影 $s, t : R \to V$ は
射 $s', t': R' \to V'$ に持ち上がる。Lemma
\ref{spaces-more-morphisms-lemma-etale-morphism-of-universal-thickenings} により、$R'$ は
$R$ の $X$ 上の普遍一次肥厚化だから、これらの射はエタールである。
% P09-ERR-1078: frozen source line 3119 types the pair (t',s') as R' -> V', not R' -> V' x V'.
このとき $(t', s') : R' \to V'$ はエタール同値関係であり、
$Z' = V'/R'$ とおける。$V' \to Z'$ は全射エタールであり、
% P09-ERR-1079: frozen source line 3121 uses undefined lowercase v'.
$v'$ は $V$ の $X$ 上の普遍一次肥厚化なので、
% P09-ERR-1080: frozen source line 3123 cites a nonexistent part (2) of this lemma.
Lemma \ref{spaces-more-morphisms-lemma-universal-thickening-over-formally-etale} part (2) により、
$Z'$ は $Z$ の $X$ 上の普遍一次肥厚化である。
\end{proof}

\begin{definition}
\label{spaces-more-morphisms-definition-universal-thickening}
$S$ をスキームとする。
$h : Z \to X$ を $S$ 上の代数空間の形式的非分岐射とする。
\begin{enumerate}
\item $Z$ の $X$ 上の {\it 普遍一次肥厚化}とは、Lemma
\ref{spaces-more-morphisms-lemma-universal-thickening} で構成した肥厚化 $Z \subset Z'$ をいう。
\item {\it $Z$ の $X$ 上の余法層}とは、$Z$ の余法層で、その普遍一次肥厚化
$Z'$ におけるものをいう。この肥厚化は $X$ 上にある。
\end{enumerate}
この状況では余法層をしばしば $\mathcal{C}_{Z/X}$ と書く。
\end{definition}

\noindent
したがって層の短完全列
$$
0 \to \mathcal{C}_{Z/X} \to \mathcal{O}_{Z'} \to \mathcal{O}_Z \to 0
$$
が $Z_\etale$ 上にあり、$\mathcal{C}_{Z/X}$ は準連接 $\mathcal{O}_Z$ 加群である。
次の補題は、この定義が $Z \to X$ が埋め込みである場合の定義と
矛盾しないことを示す。

\begin{lemma}
\label{spaces-more-morphisms-lemma-immersion-universal-thickening}
$S$ をスキームとする。
$i : Z \to X$ を $S$ 上の代数空間の埋め込みとする。このとき
\begin{enumerate}
\item $i$ は形式的非分岐である。
\item $Z$ の $X$ 上の普遍一次肥厚化は、Definition
\ref{spaces-more-morphisms-definition-first-order-infinitesimal-neighbourhood} の
$Z$ の $X$ における一次無限小近傍である。
\item Definition \ref{spaces-more-morphisms-definition-conormal-sheaf} の意味での $i$ の余法層は、
Definition \ref{spaces-more-morphisms-definition-universal-thickening} の意味での $i$ の余法層と一致する。
\end{enumerate}
\end{lemma}

\begin{proof}
代数空間の埋め込みは、定義により表現可能な射である。したがって
Morphisms, Lemmas \ref{morphisms-lemma-open-immersion-unramified} and
\ref{morphisms-lemma-closed-immersion-unramified} により、埋め込みは非分岐である
（Spaces, Lemma
\ref{spaces-lemma-representable-transformations-property-implication}
の抽象原理を用いる）。ゆえに Lemma
\ref{spaces-more-morphisms-lemma-unramified-formally-unramified} により形式的非分岐である。
他の主張は Lemmas \ref{spaces-more-morphisms-lemma-first-order-infinitesimal-neighbourhood} and
\ref{spaces-more-morphisms-lemma-infinitesimal-neighbourhood-conormal} と定義を組み合わせれば従う。
\end{proof}

\begin{lemma}
\label{spaces-more-morphisms-lemma-universal-thickening-unramified}
$S$ をスキームとする。
$Z \to X$ を $S$ 上の代数空間の形式的非分岐射とする。
このとき普遍一次肥厚化 $Z'$ は $X$ 上形式的非分岐である。
\end{lemma}

\begin{proof}
$T \subset T'$ を $X$ 上のアフィンスキームの一次肥厚化とする。
可換図式
$$
\xymatrix{
Z' \ar[d] & T \ar[l]^c \ar[d] \\
X & T' \ar[l] \ar[lu]^{a, b}
}
$$
を考える。$T_0 = c^{-1}(Z) \subset T$ および
$T'_a = a^{-1}(Z)$ とおく（スキーム論的に取る）。
$Z'$ は $Z$ の一次肥厚化なので、$T'$ は $T'_a$ の一次肥厚化である。
さらに $c = a|_T$ なので、$T_0 = T \cap T'_a$ である
（スキーム論的に取る）。$T'$ は $T$ の一次肥厚化なので、
$T'_a$ は $T_0$ の一次肥厚化である。ここで $a|_{T'_a}$ と $b|_{T'_a}$ は
$T'_a$ から $Z'$ への $X$ 上の射であり、$T_0$ 上では $Z$ への射として
一致する。したがって $Z'$ の普遍性により
$a|_{T'_a} = b|_{T'_a}$ である。このように $a$ と $b$ は、一次肥厚化
$T'$ からの射である。この肥厚化は $T'_a$ のものであり、両射の
$T'_a$ への制限は $Z$ への射として
一致する。そこで $Z'$ の普遍性をもう一度用いると $a = b$ を得る。
言い換えれば、形式的非分岐射の定義性質が望みどおり $Z' \to X$ に対して
成り立つ。
\end{proof}

\begin{lemma}
\label{spaces-more-morphisms-lemma-universal-thickening-functorial}
% P09-ERR-1082: frozen source line 3219 omits terminal punctuation.
$S$ をスキームとする。
$S$ 上の代数空間の可換図式
$$
\xymatrix{
Z \ar[r]_h \ar[d]_f & X \ar[d]^g \\
W \ar[r]^{h'} & Y
}
$$
を考え、$h$ と $h'$ は形式的非分岐であるとする。
$Z \subset Z'$ を $Z$ の $X$ 上の普遍一次肥厚化とする。
$W \subset W'$ を $W$ の $Y$ 上の普遍一次肥厚化とする。
肥厚化の標準的な射 $(f, f') : (Z, Z') \to (W, W')$ で、
$Y$ 上にあり、次の可換図式に入るものが存在する。
$$
\xymatrix{
& & & Z' \ar[ld] \ar[d]^{f'} \\
Z \ar[rr] \ar[d]_f \ar[rrru] & & X \ar[d] & W' \ar[ld] \\
W \ar[rrru]|!{[rr];[rruu]}\hole \ar[rr] & & Y
}
$$
特に、肥厚化の射 $(f, f')$ は余法層の射
$f^*\mathcal{C}_{W/Y} \to \mathcal{C}_{Z/X}$ を誘導する。
\end{lemma}

\begin{proof}
最初の主張は $W'$ の普遍性から明らかである。余法層上の誘導写像は、
$(Z \subset Z') \to (W \subset W')$ に Lemma
\ref{spaces-more-morphisms-lemma-conormal-functorial} を適用して得られる写像である。
\end{proof}

\begin{lemma}
\label{spaces-more-morphisms-lemma-universal-thickening-fibre-product}
$S$ をスキームとする。代数空間のファイバー積図式
$$
\xymatrix{
Z \ar[r]_h \ar[d]_f & X \ar[d]^g \\
W \ar[r]^{h'} & Y
}
$$
を $S$ 上で考え、$h'$ が形式的非分岐であるとする。このとき $h$ は
形式的非分岐であり、$W \subset W'$ が $W$ の $Y$ 上の普遍一次肥厚化ならば、
$Z = X \times_Y W \subset X \times_Y W'$ は $Z$ の $X$ 上の
普遍一次肥厚化である。特に Lemma
\ref{spaces-more-morphisms-lemma-universal-thickening-functorial} の標準写像
$f^*\mathcal{C}_{W/Y} \to \mathcal{C}_{Z/X}$ は全射である。
\end{lemma}

\begin{proof}
射 $h$ は Lemma \ref{spaces-more-morphisms-lemma-base-change-formally-unramified} により
形式的非分岐である。$X \times_Y W'$ が一次肥厚化であることは明らかである。
$W'$ が普遍性をもつので、これも普遍性をもつことは
（ファイバー積の写像性質により）直ちに確かめられる。これが余法層の写像の
全射性を含意する理由については Lemma
\ref{spaces-more-morphisms-lemma-conormal-functorial-flat} を参照せよ。
\end{proof}

\begin{lemma}
\label{spaces-more-morphisms-lemma-universal-thickening-fibre-product-flat}
$S$ をスキームとする。代数空間のファイバー積図式
$$
\xymatrix{
Z \ar[r]_h \ar[d]_f & X \ar[d]^g \\
W \ar[r]^{h'} & Y
}
$$
を $S$ 上で考え、$h'$ が形式的非分岐かつ $g$ が平坦であるとする。この場合、
対応する普遍一次肥厚化の写像 $Z' \to W'$ は平坦であり、
$f^*\mathcal{C}_{W/Y} \to \mathcal{C}_{Z/X}$ は同型である。
\end{lemma}

\begin{proof}
平坦性は基底変換で保たれる。Morphisms of Spaces, Lemma
\ref{spaces-morphisms-lemma-base-change-flat} を参照せよ。
したがって最初の主張は Lemma
\ref{spaces-more-morphisms-lemma-universal-thickening-fibre-product} における $W'$ の記述から従う。
$X \times_Y W'$ が一次肥厚化であることは明らかである。
$W'$ が普遍性をもつので、これも普遍性をもつことは
（ファイバー積の写像性質により）直ちに確かめられる。これが余法層の写像の
同型性を含意する理由については Lemma
\ref{spaces-more-morphisms-lemma-conormal-functorial-flat} を参照せよ。
\end{proof}

\begin{lemma}
\label{spaces-more-morphisms-lemma-universal-thickening-localize}
普遍一次肥厚化を取る操作はエタール局所化と可換である。より正確には、
形式的非分岐射 $h : Z \to X$ を基礎スキーム $S$ 上の代数空間の射とする。
可換図式
$$
\xymatrix{
V \ar[d] \ar[r] & U \ar[d] \\
Z \ar[r] & X
}
$$
で、縦の矢印がエタールであるものを考える。
$Z'$ を $Z$ の $X$ 上の普遍一次肥厚化とする。このとき $V \to U$ は
形式的非分岐であり、普遍一次肥厚化 $V'$ は $V$ の $U$ 上にあり、$Z'$ 上
エタールである。特に $\mathcal{C}_{Z/X}|_V = \mathcal{C}_{V/U}$ である。
\end{lemma}

\begin{proof}
最初の主張は Lemma \ref{spaces-more-morphisms-lemma-formally-unramified} である。
普遍一次肥厚化の両立性は Lemmas
\ref{spaces-more-morphisms-lemma-universal-thickening-over-formally-etale} and
\ref{spaces-more-morphisms-lemma-etale-morphism-of-universal-thickenings} の帰結である。
\end{proof}

\begin{lemma}
\label{spaces-more-morphisms-lemma-differentials-universally-unramified}
$S$ をスキームとする。$B$ を $S$ 上の代数空間とする。
$h : Z \to X$ を $B$ 上の代数空間の形式的非分岐射とする。
$Z \subset Z'$ を $Z$ の $X$ 上の普遍一次肥厚化とし、その構造射を
$h' : Z' \to X$ とする。標準写像
$$
\text{d}h' : (h')^*\Omega_{X/B} \to \Omega_{Z'/B}
$$
は同型 $h^*\Omega_{X/B} \to \Omega_{Z'/B} \otimes \mathcal{O}_Z$ を誘導する。
\end{lemma}

\begin{proof}
写像 $c_{h'}$ は Lemma \ref{spaces-more-morphisms-lemma-functoriality-differentials} で定義した写像である。
$i : Z \to Z'$ を与えられた閉埋め込みとすれば、$i^*c_{h'}$ は写像
% P09-ERR-1084: frozen source lines 3354--3355 switch the statement's base B to S.
$h^*\Omega_{X/S} \to \Omega_{Z'/S} \otimes \mathcal{O}_Z$ である。
これが同型であることの確認は、エタール局所化によりスキームの場合に帰着する。
Lemma \ref{spaces-more-morphisms-lemma-universal-thickening-localize} および
Lemma \ref{spaces-more-morphisms-lemma-localize-differentials} を参照せよ。この場合の結果は
More on Morphisms, Lemma
\ref{more-morphisms-lemma-differentials-universally-unramified} である。
\end{proof}

\begin{lemma}
\label{spaces-more-morphisms-lemma-universally-unramified-differentials-sequence}
$S$ をスキームとする。$B$ を $S$ 上の代数空間とする。
$h : Z \to X$ を $B$ 上の代数空間の形式的非分岐射とする。
標準的な完全列
$$
\mathcal{C}_{Z/X} \to h^*\Omega_{X/B} \to \Omega_{Z/B} \to 0.
$$
が存在する。最初の矢印は $\text{d}_{Z'/B}$ が誘導する。ここで $Z'$ は
$Z$ の $X$ 上の普遍一次近傍である。
\end{lemma}

\begin{proof}
標準的な完全列
$$
\mathcal{C}_{Z/Z'} \to
\Omega_{Z'/S} \otimes \mathcal{O}_Z \to
\Omega_{Z/S} \to 0.
$$
があることは既知である。ここで基礎が主張と一致しない点も含め、
% P09-ERR-1085: frozen source lines 3381--3383 use S although the statement is relative to B.
Lemma \ref{spaces-more-morphisms-lemma-differentials-relative-immersion} を参照せよ。
したがって Lemma
\ref{spaces-more-morphisms-lemma-differentials-universally-unramified} を適用すれば結果が従う。
\end{proof}

\begin{lemma}
\label{spaces-more-morphisms-lemma-two-unramified-morphisms}
$S$ をスキームとする。代数空間の可換図式
$$
\xymatrix{
Z \ar[r]_i \ar[rd]_j & X \ar[d] \\
& Y
}
$$
を $S$ 上で考え、$i$ と $j$ が形式的非分岐であるとする。
このとき標準的な完全列
$$
\mathcal{C}_{Z/Y} \to
\mathcal{C}_{Z/X} \to
i^*\Omega_{X/Y} \to 0
$$
がある。最初の矢印は Lemma
\ref{spaces-more-morphisms-lemma-universal-thickening-functorial} から、二番目は Lemma
\ref{spaces-more-morphisms-lemma-universally-unramified-differentials-sequence} から得られる。
\end{lemma}

\begin{proof}
写像はすでに定義したので、列の完全性の確認はエタール局所化により
スキームの場合に帰着する。Lemma
\ref{spaces-more-morphisms-lemma-universal-thickening-localize} および
Lemma \ref{spaces-more-morphisms-lemma-localize-differentials} を参照せよ。この場合の結果は
More on Morphisms, Lemma
\ref{more-morphisms-lemma-two-unramified-morphisms} である。
\end{proof}

\begin{lemma}
\label{spaces-more-morphisms-lemma-transitivity-conormal-unramified}
$S$ をスキームとする。
$Z \to Y \to X$ を $S$ 上の代数空間の形式的非分岐射とする。
\begin{enumerate}
\item $Z \subset Z'$ が $Z$ の $X$ 上の普遍一次肥厚化で、
$Y \subset Y'$ が $Y$ の $X$ 上の普遍一次肥厚化ならば、射
$Z' \to Y'$ が存在し、$Y \times_{Y'} Z'$ は $Z$ の $Y$ 上の
普遍一次肥厚化である。
\item 標準的な完全列
$$
i^*\mathcal{C}_{Y/X} \to
\mathcal{C}_{Z/X} \to
\mathcal{C}_{Z/Y} \to 0
$$
があり、その写像は Lemma
\ref{spaces-more-morphisms-lemma-universal-thickening-functorial} から得られる。ここで
$i : Z \to Y$ は最初の射である。
\end{enumerate}
\end{lemma}

\begin{proof}
(1) の写像 $h : Z' \to Y'$ は Lemma
\ref{spaces-more-morphisms-lemma-universal-thickening-functorial} から得られる。
$Y \times_{Y'} Z'$ が $Z$ の $Y$ 上の普遍一次肥厚化であるという主張は、
$Z'$ と $Y'$ の普遍性から明らかである。Lemma
\ref{spaces-more-morphisms-lemma-transitivity-conormal} により完全列
$$
(i')^*\mathcal{C}_{Y \times_{Y'} Z'/Z'} \to
\mathcal{C}_{Z/Z'} \to
\mathcal{C}_{Z/Y \times_{Y'} Z'} \to 0
$$
がある。ここで $i' : Z \to Y \times_{Y'} Z'$ は与えられた射である。
Lemma \ref{spaces-more-morphisms-lemma-conormal-functorial-flat} により全射
$h^*\mathcal{C}_{Y/Y'} \to \mathcal{C}_{Y \times_{Y'} Z'/Z'}$ が存在する。
これを等式
$\mathcal{C}_{Y/Y'} = \mathcal{C}_{Y/X}$、
$\mathcal{C}_{Z/Z'} = \mathcal{C}_{Z/X}$、および
$\mathcal{C}_{Z/Y \times_{Y'} Z'} = \mathcal{C}_{Z/Y}$
と組み合わせれば補題が従う。
\end{proof}













\section{形式的エタール射}
\label{spaces-more-morphisms-section-formally-etale}

\noindent
この節では、代数空間の射が形式的エタールであることの意味を詳しく調べる。

\begin{definition}
\label{spaces-more-morphisms-definition-formally-etale}
$S$ をスキームとする。代数空間の射 $f : X \to Y$ が $S$ 上で
{\it 形式的エタール}であるとは、Definition
\ref{spaces-more-morphisms-definition-formally-smooth-etale-unramified} の意味で、
関手の変換として形式的エタールであることをいう。
\end{definition}

\noindent
Section \ref{spaces-more-morphisms-section-formally-smooth-etale-unramified} で、より一般の
形式的エタールな関手の変換について証明した結果はここでは再掲しない。

\begin{lemma}
\label{spaces-more-morphisms-lemma-formally-etale}
$S$ をスキームとする。$f : X \to Y$ を $S$ 上の代数空間の射とする。
次は同値である。
\begin{enumerate}
\item $f$ は形式的エタールである。
\item 任意の図式
$$
\xymatrix{
U \ar[d] \ar[r]_\psi & V \ar[d] \\
X \ar[r]^f & Y
}
$$
で、$U$ と $V$ がスキームであり、縦の矢印がエタールであるものに対し、
スキームの射 $\psi$ は形式的エタールである（More on Morphisms,
Definition \ref{more-morphisms-definition-formally-etale} の意味で）。
\item 全射な縦の矢印をもつこのような図式の一つについて、射 $\psi$ は
形式的エタールである。
\end{enumerate}
\end{lemma}

\begin{proof}
$f$ が形式的エタールであると仮定する。Lemma
\ref{spaces-more-morphisms-lemma-representable-property-formally-property} により、射
$U \to X$ と $V \to Y$ は形式的エタールである。したがって Lemma
\ref{spaces-more-morphisms-lemma-composition-formally-smooth-etale-unramified} により、合成
$U \to Y$ は形式的エタールである。すると Lemma
\ref{spaces-more-morphisms-lemma-formally-permanence} から $U \to V$ は形式的エタールである。
ゆえに (1) は (2) を含意する。また (2) が (3) を含意することは明らかである。
% P09-ERR-1089: frozen source line 3533 lacks terminal punctuation after "trivially".

\medskip\noindent
(3) にある図式が与えられたと仮定する。Lemma
\ref{spaces-more-morphisms-lemma-representable-property-formally-property} により、射
$V \to Y$ は形式的エタールである。したがって Lemma
\ref{spaces-more-morphisms-lemma-composition-formally-smooth-etale-unramified} により、合成
$U \to Y$ は形式的エタールである。すると Lemma
\ref{spaces-more-morphisms-lemma-etale-on-top} から $X \to Y$ は形式的エタール、すなわち (1) が成り立つ。
\end{proof}

\begin{lemma}
\label{spaces-more-morphisms-lemma-formally-etale-not-affine}
$S$ をスキームとする。
$f : X \to Y$ を $S$ 上の代数空間の形式的エタール射とする。
このとき任意の実線部分が可換な図式
$$
\xymatrix{
X \ar[d]_f & T \ar[d]^i \ar[l]_a \\
Y & T' \ar[l] \ar@{-->}[lu]
}
$$
で、$T \subset T'$ が $Y$ 上の代数空間の一次肥厚化であるものに対し、
図式を可換にする点線矢印がちょうど一つ存在する。言い換えれば、
Definition \ref{spaces-more-morphisms-definition-formally-etale} における $T$ がアフィンである
という条件は取り除ける。
\end{lemma}

\begin{proof}
$U' \to T'$ を全射エタール射とし、$U' = \coprod U'_i$ は
アフィンスキームの非交和であるとする。
$U_i = T \times_{T'} U'_i$ とおく。このとき射
$a'_i : U'_i \to X$ で、その制限 $a'_i|_{U_i}$ が合成
$U_i \to T \to X$ に等しいものを得る。一意性により（Lemma
\ref{spaces-more-morphisms-lemma-formally-unramified-not-affine} を参照）、$a'_i$ と $a'_j$ は
ファイバー積 $U'_i \times_{T'} U'_j$ 上で一致する。したがって
$\coprod a'_i : U' \to X$ は降下して、一意な射 $a' : T' \to X$ を与える。
\end{proof}

\begin{lemma}
\label{spaces-more-morphisms-lemma-composition-formally-etale}
形式的エタール射の合成は形式的エタールである。
\end{lemma}

\begin{proof}
これは形式的である。
\end{proof}

\begin{lemma}
\label{spaces-more-morphisms-lemma-base-change-formally-etale}
形式的エタール射の基底変換は形式的エタールである。
\end{lemma}

\begin{proof}
これは形式的である。
\end{proof}

\begin{lemma}
\label{spaces-more-morphisms-lemma-characterize-formally-etale}
$S$ をスキームとする。
% P09-ERR-1086: frozen source line 3596 lacks terminal punctuation.
$f : X \to Y$ を $S$ 上の代数空間の射とする。
次は同値である。
\begin{enumerate}
\item $f$ は形式的エタールである。
\item $f$ は形式的非分岐であり、$X$ の $Y$ 上の普遍一次肥厚化は $X$ に等しい。
\item $f$ は形式的非分岐であり、$\mathcal{C}_{X/Y} = 0$ である。
\item $\Omega_{X/Y} = 0$ かつ $\mathcal{C}_{X/Y} = 0$ である。
\end{enumerate}
\end{lemma}

\begin{proof}
% P09-ERR-1087: frozen source line 3608 has "assertion only make sense".
実際、最後の主張が意味をなすのは、$\Omega_{X/Y} = 0$ ならば
Lemma \ref{spaces-more-morphisms-lemma-characterize-formally-unramified} と
Definition \ref{spaces-more-morphisms-definition-universal-thickening} によって
$\mathcal{C}_{X/Y}$ が定義されるからにほかならない。
これは (3) と (4) が同値であることも明らかにする。

\medskip\noindent
仮定 (1)、(2)、(3) のいずれからも $f$ が形式的非分岐であることが従う。
したがって $f$ は形式的非分岐であるとしてよい。(1)、(2)、(3) の同値性は、
% P09-ERR-1088: frozen source line 3619 says the universal thickening is over S, but the lemma is relative to Y.
$X'$ の普遍性（これは $X$ の $S$ 上の普遍一次肥厚化である）、および
$X = X' \Leftrightarrow \mathcal{C}_{X/Y} = 0$ から従う。実際、定義により
$\mathcal{C}_{X/Y} = \mathcal{C}_{X/X'}$ は $X$ の $X'$ におけるイデアル層である。
\end{proof}

\begin{lemma}
\label{spaces-more-morphisms-lemma-unramified-flat-formally-etale}
非分岐平坦射は形式的エタールである。
\end{lemma}

\begin{proof}
これはスキームの場合 More on Morphisms, Lemma
\ref{more-morphisms-lemma-unramified-flat-formally-etale} と、エタール局所化
Lemmas \ref{spaces-more-morphisms-lemma-formally-unramified} and \ref{spaces-more-morphisms-lemma-formally-etale}、および
Morphisms of Spaces, Lemma \ref{spaces-morphisms-lemma-flat-local} から従う。
\end{proof}

\begin{lemma}
\label{spaces-more-morphisms-lemma-etale-formally-etale}
$S$ をスキームとする。
$f : X \to Y$ を $S$ 上の代数空間の射とする。
次は同値である。
\begin{enumerate}
\item 射 $f$ はエタールである。
\item 射 $f$ は局所有限表示かつ形式的エタールである。
\end{enumerate}
\end{lemma}

\begin{proof}
これはスキームの場合 More on Morphisms, Lemma
\ref{more-morphisms-lemma-etale-formally-etale} と、エタール局所化
Lemma \ref{spaces-more-morphisms-lemma-formally-etale}、および Morphisms of Spaces, Lemmas
\ref{spaces-morphisms-lemma-finite-presentation-local} and
\ref{spaces-morphisms-lemma-etale-local} から従う。
\end{proof}















\section{写像の無限小変形}
\label{spaces-more-morphisms-section-action-by-derivations}

\noindent
この節では、導分を用いて写像を無限小に動かす方法を説明する。
この節を通じて、肥厚化 $X'$（これは $X$ の肥厚化である）上の層を
$X$ 上の層とみなせることを用いる。Equations
(\ref{spaces-more-morphisms-equation-equivalence-etale-spaces}) and
(\ref{spaces-more-morphisms-equation-fundamental-equivalence}) を参照せよ。

\begin{lemma}
\label{spaces-more-morphisms-lemma-difference-derivation}
$S$ をスキームとする。$B$ を $S$ 上の代数空間とする。
$X \subset X'$ と $Y \subset Y'$ を $B$ 上の代数空間の二つの一次肥厚化とする。
$(a, a'), (b, b') : (X \subset X') \to (Y \subset Y')$ を
$B$ 上の肥厚化の二つの射とする。次を仮定する。
\begin{enumerate}
\item $a = b$ である。
\item 二つの写像 $a^*\mathcal{C}_{Y/Y'} \to \mathcal{C}_{X/X'}$
（Lemma \ref{spaces-more-morphisms-lemma-conormal-functorial}）は等しい。
\end{enumerate}
このとき写像 $(a')^\sharp - (b')^\sharp$ は
$$
\mathcal{O}_{Y'} \to \mathcal{O}_Y \xrightarrow{D}
a_*\mathcal{C}_{X/X'} \to a_*\mathcal{O}_{X'}
$$
と分解する。ここで $D$ は $\mathcal{O}_B$ 導分である。
\end{lemma}

\begin{proof}
$Y$ 上で作業する代わりに $X$ 上で作業する。その利点は引戻し関手
$a^{-1}$ が完全であることである。(1) と (2) を用いると、行が完全な可換図式
$$
\xymatrix{
0 \ar[r] &
\mathcal{C}_{X/X'} \ar[r] &
\mathcal{O}_{X'} \ar[r] &
\mathcal{O}_X \ar[r] & 0 \\
0 \ar[r] &
a^{-1}\mathcal{C}_{Y/Y'} \ar[r] \ar[u] &
a^{-1}\mathcal{O}_{Y'}
\ar[r] \ar@<1ex>[u]^{(a')^\sharp} \ar@<-1ex>[u]_{(b')^\sharp} &
a^{-1}\mathcal{O}_Y \ar[r] \ar[u] & 0
}
$$
を得る。このような状況では、$\mathcal{O}_B$ 代数の写像
$(a')^\sharp$ と $(b')^\sharp$ の差が、$\mathcal{O}_B$ 導分であって
$a^{-1}\mathcal{O}_Y$ から $\mathcal{C}_{X/X'}$ へのものになることは
一般的事実である。関手 $a^{-1}$ と $a_*$ の随伴性により、これは
$\mathcal{O}_B$ 導分であって $\mathcal{O}_Y$ から
$a_*\mathcal{C}_{X/X'}$ へのものと同じである。
細部は省略する。
\end{proof}

\noindent
上の補題の状況では $D$ を
\begin{equation}
\label{spaces-more-morphisms-equation-D}
% P09-ERR-1090: frozen source line 3739 reverses the typed composition; theta follows d_{Y/B}.
D = \text{d}_{Y/B} \circ \theta
\end{equation}
と書けることに注意せよ。ここで $\theta$ は $\mathcal{O}_Y$ 線形写像
$\theta : \Omega_{Y/B} \to a_*\mathcal{C}_{X/X'}$ である。
もちろん、再び随伴により $\theta$ を $\mathcal{O}_X$ 線形写像
$\theta : a^*\Omega_{Y/B} \to \mathcal{C}_{X/X'}$ とみなすこともできる。

\begin{lemma}
\label{spaces-more-morphisms-lemma-action-by-derivations}
$S$ をスキームとする。$B$ を $S$ 上の代数空間とする。
$(a, a') : (X \subset X') \to (Y \subset Y')$ を
$B$ 上の一次肥厚化の射とする。写像
$$
\theta : a^*\Omega_{Y/B} \to \mathcal{C}_{X/X'}
$$
を $\mathcal{O}_X$ 線形とする。このとき対の射
$(b, b') : (X \subset X') \to (Y \subset Y')$ で、Lemma
\ref{spaces-more-morphisms-lemma-difference-derivation} の (1) と (2) を満たし、導分 $D$ と
$\theta$ が Equation (\ref{spaces-more-morphisms-equation-D}) によって関係づけられるものが
一意に存在する。
\end{lemma}

\begin{proof}
写像
$$
\alpha = (a')^\sharp + D : a^{-1}\mathcal{O}_{Y'} \to \mathcal{O}_{X'}
$$
を考える。ここで $D$ は Equation (\ref{spaces-more-morphisms-equation-D}) のものである。
$D$ は $\mathcal{O}_B$ 導分なので、$\alpha$ は $\mathcal{O}_B$ 代数の層の
写像である。構成により $i_X^\sharp \circ \alpha = a^\sharp \circ i_Y^\sharp$
である。ここで $i_X : X \to X'$ と $i_Y : Y \to Y'$ は与えられた
閉埋め込みである。Lemma \ref{spaces-more-morphisms-lemma-first-order-thickening-maps} により、
肥厚化の一意な射
$(a, b') : (X  \subset X') \to (Y \subset Y')$ を $B$ 上に得て、これは
$\alpha = (b')^\sharp$ を満たす。$b = a$ とおけばよい。
\end{proof}

\begin{remark}
\label{spaces-more-morphisms-remark-action-by-derivations}
仮定と記法は Lemma \ref{spaces-more-morphisms-lemma-action-by-derivations} と同じとする。
局所切断 $\theta$ の $a'$ への作用を $\theta \cdot a'$ と書くことがある。
これは、$(a')^\sharp$ と $(\theta \cdot a')^\sharp$ の差が導分 $D$ であり、
その導分が Equation (\ref{spaces-more-morphisms-equation-D}) によって $\theta$ と関係づけられることに
ほかならない。
\end{remark}

\begin{lemma}
\label{spaces-more-morphisms-lemma-sheaf}
$S$ をスキームとする。$B$ を $S$ 上の代数空間とする。
$X \subset X'$ と $Y \subset Y'$ を $B$ 上の一次肥厚化とする。
射 $a : X \to Y$ と写像
$A : a^*\mathcal{C}_{Y/Y'} \to \mathcal{C}_{X/X'}$ が与えられ、後者は
$\mathcal{O}_X$ 加群の写像であると仮定する。対象 $U'$ を
$(X')_{spaces, \etale}$ において取り、
$U = X \times_{X'} U'$ とおき、次を満たす射 $a' : U' \to Y'$ を考える。
\begin{enumerate}
\item $a'$ は $B$ 上の射である。
\item $a'|_U = a|_U$ である。
\item 誘導写像 $a^*\mathcal{C}_{Y/Y'}|_U \to \mathcal{C}_{X/X'}|_U$ は
$A$ の $U$ への制限である。
\end{enumerate}
このとき規則
\begin{equation}
\label{spaces-more-morphisms-equation-sheaf}
U' \mapsto
\{a' : U' \to Y'\text{ such that (1), (2), (3) hold.}\}
\end{equation}
は $(X')_{spaces, \etale}$ 上の集合の層を定める。
\end{lemma}

\begin{proof}
補題の規則を $\mathcal{F}$ と書く。制限写像
$\mathcal{F}(U') \to \mathcal{F}(V')$（ここで
$V' \subset U' \subset X'$）は、$\mathcal{F}$ において、まさに制限写像
$a' \mapsto a'|_{V'}$ である。この定義のもとで $\mathcal{F}$ が層であることは
明らかである。実際、代数空間の射はエタール降下を満たす。Descent on Spaces,
Lemma \ref{spaces-descent-lemma-fpqc-universal-effective-epimorphisms} を参照せよ。
\end{proof}

\begin{lemma}
\label{spaces-more-morphisms-lemma-action-sheaf}
記法と仮定は Lemma \ref{spaces-more-morphisms-lemma-sheaf} と同じとする。
(\ref{spaces-more-morphisms-equation-equivalence-etale-spaces}) によって $X$ 上の層と $X'$ 上の層を
同一視する。層
$$
\SheafHom_{\mathcal{O}_X}(a^*\Omega_{Y/B}, \mathcal{C}_{X/X'})
$$
は層 (\ref{spaces-more-morphisms-equation-sheaf}) に作用する。さらに、任意の対象 $U'$ で、
$(X')_{spaces, \etale}$ の対象であり、層 (\ref{spaces-more-morphisms-equation-sheaf}) がその上で
切断をもつものに対し、この作用は単純推移的である。
\end{lemma}

\begin{proof}
これは Lemmas \ref{spaces-more-morphisms-lemma-difference-derivation},
\ref{spaces-more-morphisms-lemma-action-by-derivations}, and \ref{spaces-more-morphisms-lemma-sheaf} を組み合わせたものである。
\end{proof}

\begin{remark}
\label{spaces-more-morphisms-remark-special-case}
Lemmas \ref{spaces-more-morphisms-lemma-difference-derivation},
\ref{spaces-more-morphisms-lemma-action-by-derivations},
\ref{spaces-more-morphisms-lemma-sheaf}, and
\ref{spaces-more-morphisms-lemma-action-sheaf} の特別な場合として $Y = Y'$ の場合がある。
この場合、写像 $A$ は常に零である。Lemma \ref{spaces-more-morphisms-lemma-sheaf} の層は単に規則
$$
U' \mapsto
\{a' : U' \to Y\text{ over }B\text{ with } a'|_U = a|_U\}
$$
で与えられ、この層に
$\SheafHom_{\mathcal{O}_X}(a^*\Omega_{Y/B}, \mathcal{C}_{X/X'})$ が作用する。
\end{remark}

\begin{remark}
\label{spaces-more-morphisms-remark-another-special-case}
Lemmas \ref{spaces-more-morphisms-lemma-difference-derivation},
\ref{spaces-more-morphisms-lemma-action-by-derivations},
\ref{spaces-more-morphisms-lemma-sheaf}, and
\ref{spaces-more-morphisms-lemma-action-sheaf} の別の特別な場合は、$B$ 自身が肥厚化
$Z \subset Z' = B$ であり、$Y = Z \times_{Z'} Y'$ となる場合である。
図式で表すと
$$
\xymatrix{
(X \subset X') \ar@{..>}[rr]_{(a, ?)} \ar[rd]_{(g, g')} & &
(Y \subset Y') \ar[ld]^{(h, h')} \\
& (Z \subset Z')
}
$$
である。この場合、写像
$A : a^*\mathcal{C}_{Y/Y'} \to \mathcal{C}_{X/X'}$ は $a$ によって決まる。
実際、写像 $h^*\mathcal{C}_{Z/Z'} \to \mathcal{C}_{Y/Y'}$ は全射である
（$Y = Z \times_{Z'} Y'$ と仮定したからである）。したがって引戻し
$g^*\mathcal{C}_{Z/Z'} = a^*h^*\mathcal{C}_{Z/Z'} \to
a^*\mathcal{C}_{Y/Y'}$ は全射であり、合成
$g^*\mathcal{C}_{Z/Z'} \to a^*\mathcal{C}_{Y/Y'} \to \mathcal{C}_{X/X'}$
は $g'$ が誘導する標準写像でなければならない。したがって Lemma
\ref{spaces-more-morphisms-lemma-sheaf} の層は単に規則
$$
U' \mapsto
\{a' : U' \to Y'\text{ over }Z'\text{ with } a'|_U = a|_U\}
$$
で与えられ、この層に
$\SheafHom_{\mathcal{O}_X}(a^*\Omega_{Y/Z}, \mathcal{C}_{X/X'})$ が作用する。
\end{remark}

\begin{lemma}
\label{spaces-more-morphisms-lemma-action-by-derivations-etale-localization}
$S$ をスキームとする。一次肥厚化の可換図式
$$
\vcenter{
\xymatrix{
(T_2 \subset T_2') \ar[d]_{(h, h')} \ar[rr]_{(a_2, a_2')} & &
(X_2 \subset X_2') \ar[d]^{(f, f')} \\
(T_1 \subset T_1') \ar[rr]^{(a_1, a_1')} & &
(X_1 \subset X_1')
}
}
\quad
\begin{matrix}
\text{and a commutative} \\
\text{diagram}
\end{matrix}
\quad
\vcenter{
\xymatrix{
X_2' \ar[r] \ar[d] & B_2 \ar[d] \\
X_1' \ar[r] & B_1
}
}
$$
を $S$ 上の代数空間について考え、$X_2 \to X_1$ と $B_2 \to B_1$ は
エタールであるとする。任意の $\mathcal{O}_{T_1}$ 線形写像
$\theta_1 : a_1^*\Omega_{X_1/B_1} \to \mathcal{C}_{T_1/T'_1}$ に対し、
$\theta_2$ を合成
$$
\xymatrix{
a_2^*\Omega_{X_2/B_2} \ar@{=}[r] &
h^*a_1^*\Omega_{X_1/B_1} \ar[r]^-{h^*\theta_1} &
h^*\mathcal{C}_{T_1/T'_1} \ar[r] &
\mathcal{C}_{T_2/T'_2}
}
$$
とする（等号は証明で説明する）。このとき図式
$$
\xymatrix{
T_2' \ar[rr]_{\theta_2 \cdot a_2'} \ar[d] & & X'_2 \ar[d] \\
T_1' \ar[rr]^{\theta_1 \cdot a_1'} & & X'_1
}
$$
は可換である。ここで作用 $\theta_2 \cdot a_2'$ と
$\theta_1 \cdot a_1'$ は Remark \ref{spaces-more-morphisms-remark-action-by-derivations} のものである。
\end{lemma}

\begin{proof}
% P09-ERR-1091: frozen source lines 3952--3958 replace the statement's B_1,B_2 by undefined S_1,S_2.
等号は、微分層の構成が始域と終域の双方についてエタール局所化と両立するため
得られる同一視 $f^*\Omega_{X_1/S_1} = \Omega_{X_2/S_2}$ から来る。
Lemma \ref{spaces-more-morphisms-lemma-localize-differentials} を参照せよ。実際、これを用いると
$a_2^*\Omega_{X_2/S_2} = a_2^*f^*\Omega_{X_1/S_1} =
h^*a_1^*\Omega_{X_1/S_1}$ を得る。これは
$f \circ a_2 = a_1 \circ h$ だからである。
% P09-ERR-1092: frozen source line 3959 has the malformed phrase "checked on etale locally".
以上を踏まえると、図式の可換性はエタール局所的に確認できる。
したがって $T'_i$, $X'_i$, $B_2$, $B_1$ はスキームであるとしてよく、
この場合、補題は More on Morphisms, Lemma
\ref{more-morphisms-lemma-action-by-derivations-etale-localization} から従う。
別証明：Lemma \ref{spaces-more-morphisms-lemma-first-order-thickening-maps} により、$X_1'$ 上の
環の層のある図式が可換であることを示せば十分である。そこで前述の
More on Morphisms, Lemma
\ref{more-morphisms-lemma-action-by-derivations-etale-localization}
の証明とまったく同様に論じれば、実際に可換であることが分かる。
\end{proof}










\section{代数空間の無限小変形}
\label{spaces-more-morphisms-section-deform}

\noindent
次の簡単な補題は、写像の無限小変形が平坦かどうかを確認するための
便利な道具としてしばしば使われる。

\begin{lemma}
\label{spaces-more-morphisms-lemma-deform}
$S$ をスキームとする。$(f, f') : (X \subset X') \to (Y \subset Y')$ を
$S$ 上の代数空間の一次肥厚化の射とする。$f$ は平坦であると仮定する。
% P09-ERR-1094: frozen source line 3994 lacks punctuation before the enumeration.
次は同値である。
\begin{enumerate}
\item $f'$ は平坦であり、$X = Y \times_{Y'} X'$ である。
\item 標準写像 $f^*\mathcal{C}_{Y/Y'} \to \mathcal{C}_{X/X'}$ は同型である。
\end{enumerate}
\end{lemma}

\begin{proof}
スキーム $V'$ と全射エタール射 $V' \to Y'$ を選ぶ。
スキーム $U'$ と全射エタール射 $U' \to X' \times_{Y'} V'$ を選ぶ。
$U = X \times_{X'} U'$ および $V = Y \times_{Y'} V'$ とおく。
代数空間の平坦射の定義により、誘導写像 $g : U \to V$ はスキームの
平坦射であり、$f'$ が平坦であることと、対応する射
$g' : U' \to V'$ が平坦であることは同値である。また、
% P09-ERR-1093: frozen source line 4010 has the tautology U = V x_{V'} V'; U' is the intended second factor.
$X = Y \times_{Y'} X'$ であることと $U = V \times_{V'} V'$ であることは
同値である。最後に、写像
$f^*\mathcal{C}_{Y/Y'} \to \mathcal{C}_{X/X'}$ が同型であることと、
$g^*\mathcal{C}_{V/V'} \to \mathcal{C}_{U/U'}$ が同型であることは同値である。
したがって補題はスキームの射に対する類似の結果から従う。
More on Morphisms, Lemma \ref{more-morphisms-lemma-deform} を参照せよ。
\end{proof}

\noindent
次の補題は「ファイバーによる平坦性判定法」の「冪零」版である。
Section \ref{spaces-more-morphisms-section-criterion-flat-fibres} を参照せよ。

\begin{lemma}
\label{spaces-more-morphisms-lemma-flatness-morphism-thickenings}
$S$ をスキームとする。可換図式
$$
\xymatrix{
(X \subset X') \ar[rr]_{(f, f')} \ar[rd] & & (Y \subset Y') \ar[ld] \\
& (B \subset B')
}
$$
を $S$ 上の代数空間の肥厚化について考える。次を仮定する。
\begin{enumerate}
\item $X'$ は $B'$ 上平坦である。
\item $f$ は平坦である。
\item $B \subset B'$ は有限次肥厚化である。
\item $X = B \times_{B'} X'$ かつ $Y = B \times_{B'} Y'$ である。
\end{enumerate}
このとき $f'$ は平坦であり、$Y'$ は $B'$ 上で、$f'$ の像に属する
すべての点において平坦である。
\end{lemma}

\begin{proof}
スキーム $U'$ と全射エタール射 $U' \to B'$ を選ぶ。
スキーム $V'$ と全射エタール射
$V' \to U' \times_{B'} Y'$ を選ぶ。
スキーム $W'$ と全射エタール射
$W' \to V' \times_{Y'} X'$ を選ぶ。$U, V, W$ を $U', V', W'$ の
$B \to B'$ による基底変換とする。このとき $f'$ の
平坦性は $W' \to V'$ の平坦性と同値であり、$W \to V$ は平坦であると
与えられている。したがってスキームの肥厚化の図式
$$
\xymatrix{
(W \subset W') \ar[rr] \ar[rd] & & (V \subset V') \ar[ld] \\
& (U \subset U')
}
$$
にスキームの場合の補題を適用できる。More on Morphisms, Lemma
\ref{more-morphisms-lemma-flatness-morphism-thickenings} を参照せよ。
$Y'/B'$ の平坦性に関する、$f'$ の像に属する点での主張も同様に従う。
\end{proof}

\noindent
スキームの射の多くの性質は平坦変形で保たれる。

\begin{lemma}
\label{spaces-more-morphisms-lemma-deform-property}
$S$ をスキームとする。可換図式
$$
\xymatrix{
(X \subset X') \ar[rr]_{(f, f')} \ar[rd] & & (Y \subset Y') \ar[ld] \\
& (B \subset B')
}
$$
を $S$ 上の代数空間の肥厚化について考える。$B \subset B'$ は有限次肥厚化、
$X'$ は $B'$ 上平坦、$X = B \times_{B'} X'$、かつ
$Y = B \times_{B'} Y'$ であると仮定する。このとき次が成り立つ。
\begin{enumerate}
\item $f$ が表現可能であることと $f'$ が表現可能であることは同値である。
\label{spaces-more-morphisms-item-representable}
\item $f$ が平坦であることと $f'$ が平坦であることは同値である。
\label{spaces-more-morphisms-item-flat}
\item $f$ が同型であることと $f'$ が同型であることは同値である。
\label{spaces-more-morphisms-item-isomorphism}
\item $f$ が開埋め込みであることと $f'$ が開埋め込みであることは同値である。
\label{spaces-more-morphisms-item-open-immersion}
\item $f$ が準コンパクトであることと $f'$ が準コンパクトであることは同値である。
\label{spaces-more-morphisms-item-quasi-compact}
\item $f$ が普遍的閉であることと $f'$ が普遍的閉であることは同値である。
\label{spaces-more-morphisms-item-universally-closed}
\item $f$ が（準）分離的であることと $f'$ が（準）分離的であることは同値である。
\label{spaces-more-morphisms-item-separated}
\item $f$ が単射であることと $f'$ が単射であることは同値である。
\label{spaces-more-morphisms-item-monomorphism}
\item $f$ が全射であることと $f'$ が全射であることは同値である。
\label{spaces-more-morphisms-item-surjective}
\item $f$ が普遍的単射であることと $f'$ が普遍的単射であることは同値である。
\label{spaces-more-morphisms-item-universally-injective}
\item $f$ がアフィンであることと $f'$ がアフィンであることは同値である。
\label{spaces-more-morphisms-item-affine}
\item
\label{spaces-more-morphisms-item-finite-type}
$f$ が局所有限型であることと $f'$ が局所有限型であることは同値である。
\item $f$ が局所準有限であることと $f'$ が局所準有限であることは同値である。
\label{spaces-more-morphisms-item-quasi-finite}
\item
\label{spaces-more-morphisms-item-finite-presentation}
$f$ が局所有限表示であることと $f'$ が局所有限表示であることは同値である。
\item
\label{spaces-more-morphisms-item-relative-dimension-d}
$f$ が相対次元 $d$ の局所有限型であることと、$f'$ が相対次元 $d$ の
局所有限型であることは同値である。
\item $f$ が普遍的開であることと $f'$ が普遍的開であることは同値である。
\label{spaces-more-morphisms-item-universally-open}
\item $f$ がシントミックであることと $f'$ がシントミックであることは同値である。
\label{spaces-more-morphisms-item-syntomic}
\item $f$ が滑らかであることと $f'$ が滑らかであることは同値である。
\label{spaces-more-morphisms-item-smooth}
\item $f$ が非分岐であることと $f'$ が非分岐であることは同値である。
\label{spaces-more-morphisms-item-unramified}
\item $f$ がエタールであることと $f'$ がエタールであることは同値である。
\label{spaces-more-morphisms-item-etale}
\item $f$ が固有であることと $f'$ が固有であることは同値である。
\label{spaces-more-morphisms-item-proper}
\item $f$ が整であることと $f'$ が整であることは同値である。
\label{spaces-more-morphisms-item-integral}
\item $f$ が有限であることと $f'$ が有限であることは同値である。
\label{spaces-more-morphisms-item-finite}
\item
\label{spaces-more-morphisms-item-finite-locally-free}
$f$ が（階数 $d$ の）有限局所自由であることと、$f'$ が（階数 $d$ の）
有限局所自由であることは同値である。
% P09-ERR-1097: frozen source line 4138 retains the editorial placeholder "add more here".
\item さらに追加せよ。
\end{enumerate}
\end{lemma}

\begin{proof}
(\ref{spaces-more-morphisms-item-representable}) の場合は Lemma
\ref{spaces-more-morphisms-lemma-thicken-property-morphisms} から従う。

\medskip\noindent
スキーム $U'$ と全射エタール射 $U' \to B'$ を選ぶ。
スキーム $V'$ と全射エタール射 $V' \to U' \times_{B'} Y'$ を選ぶ。
スキーム $W'$ と全射エタール射 $W' \to V' \times_{Y'} X'$ を選ぶ。
$U, V, W$ を $U', V', W'$ の $B \to B'$ による基底変換とする。
スキームの肥厚化の図式
$$
\xymatrix{
(W \subset W') \ar[rr] \ar[rd] & & (V \subset V') \ar[ld] \\
& (U \subset U')
}
$$
を考える。始域と終域の双方でエタール局所的な任意の性質について、結果は
上の図式に適用したスキームの肥厚化の射に対する対応する結果から直ちに従う。
したがって場合
(\ref{spaces-more-morphisms-item-flat}), (\ref{spaces-more-morphisms-item-finite-type}),
(\ref{spaces-more-morphisms-item-quasi-finite}), (\ref{spaces-more-morphisms-item-finite-presentation}),
(\ref{spaces-more-morphisms-item-relative-dimension-d}), (\ref{spaces-more-morphisms-item-syntomic}),
(\ref{spaces-more-morphisms-item-smooth}), (\ref{spaces-more-morphisms-item-unramified}), (\ref{spaces-more-morphisms-item-etale})
は More on Morphisms, Lemma
\ref{more-morphisms-lemma-deform-property} の対応する場合から従う。

\medskip\noindent
$X \to X'$ と $Y \to Y'$ は普遍同相なので、写像 $X \to Y$ と
$X' \to Y'$ の位相に関する任意の問いは同じ答えをもつ。したがって場合
(\ref{spaces-more-morphisms-item-quasi-compact}), (\ref{spaces-more-morphisms-item-universally-closed}),
(\ref{spaces-more-morphisms-item-surjective}), (\ref{spaces-more-morphisms-item-universally-injective}), and
(\ref{spaces-more-morphisms-item-universally-open}) が成り立つ。

\medskip\noindent
残る各場合では、含意
$f\text{ has }P \Rightarrow f'\text{ has }P$
だけを証明する。逆の含意は $P$ が基底変換で安定であることから従う。
Spaces, Lemma
\ref{spaces-lemma-base-change-immersions} および Morphisms of Spaces, Lemmas
\ref{spaces-morphisms-lemma-base-change-separated},
\ref{spaces-morphisms-lemma-base-change-monomorphism},
\ref{spaces-morphisms-lemma-base-change-affine},
\ref{spaces-morphisms-lemma-base-change-proper},
\ref{spaces-morphisms-lemma-base-change-integral}, and
\ref{spaces-morphisms-lemma-base-change-finite-locally-free} を参照せよ。

\medskip\noindent
(\ref{spaces-more-morphisms-item-open-immersion}) の場合。$f$ は開埋め込みであると仮定する。
このとき $f'$ は (\ref{spaces-more-morphisms-item-etale}) によりエタールであり、
(\ref{spaces-more-morphisms-item-universally-injective}) により普遍的単射なので、$f'$ は開埋め込みである。
Morphisms of Spaces, Lemma
\ref{spaces-morphisms-lemma-etale-universally-injective-open} を参照せよ。
この補題を使わず、まず (\ref{spaces-more-morphisms-item-representable}) を用いてスキームの場合に
帰着することもできる。

\medskip\noindent
(\ref{spaces-more-morphisms-item-isomorphism}) の場合。場合
(\ref{spaces-more-morphisms-item-open-immersion}) と (\ref{spaces-more-morphisms-item-surjective}) から従う。

\medskip\noindent
(\ref{spaces-more-morphisms-item-separated}) の場合。Lemma
\ref{spaces-more-morphisms-lemma-thicken-property-morphisms} を参照せよ。

\medskip\noindent
(\ref{spaces-more-morphisms-item-monomorphism}) の場合。$f$ は単射であると仮定する。
対角射 $\Delta_{X'/Y'} : X' \to X' \times_{Y'} X'$ を考える。
$\Delta_{X'/Y'}$ の $B \to B'$ による基底変換は $\Delta_{X/Y}$ であり、
仮定により同型である。(\ref{spaces-more-morphisms-item-isomorphism}) により
$\Delta_{X'/Y'}$ は同型であり、したがって $f'$ は単射である。

\medskip\noindent
(\ref{spaces-more-morphisms-item-affine}) の場合。Lemma
\ref{spaces-more-morphisms-lemma-thicken-property-morphisms} を参照せよ。

\medskip\noindent
(\ref{spaces-more-morphisms-item-proper}) の場合。Lemma
\ref{spaces-more-morphisms-lemma-thicken-property-morphisms-cartesian} を参照せよ。

\medskip\noindent
(\ref{spaces-more-morphisms-item-integral}) の場合。Lemma
\ref{spaces-more-morphisms-lemma-thicken-property-morphisms} を参照せよ。

\medskip\noindent
(\ref{spaces-more-morphisms-item-finite}) の場合。Lemma
\ref{spaces-more-morphisms-lemma-thicken-property-morphisms-cartesian} を参照せよ。

\medskip\noindent
(\ref{spaces-more-morphisms-item-finite-locally-free}) の場合。$f$ は有限局所自由であると仮定する。
(\ref{spaces-more-morphisms-item-finite}) により $f'$ は有限である。
(\ref{spaces-more-morphisms-item-flat}) により $f'$ は平坦である。
(\ref{spaces-more-morphisms-item-finite-presentation}) により $f'$ は局所有限表示である。
したがって Morphisms of Spaces, Lemma
\ref{spaces-morphisms-lemma-finite-flat} により $f'$ は有限局所自由である。
\end{proof}














\noindent
次の補題は「ファイバーによる平坦性判定法」の「局所冪零」版である。
Section \ref{spaces-more-morphisms-section-criterion-flat-fibres} を参照せよ。

\begin{lemma}
\label{spaces-more-morphisms-lemma-flatness-morphism-thickenings-fp-over-ft}
$S$ をスキームとする。可換図式
$$
\xymatrix{
(X \subset X') \ar[rr]_{(f, f')} \ar[rd] & & (Y \subset Y') \ar[ld] \\
& (B \subset B')
}
$$
を $S$ 上の代数空間の肥厚化について考える。次を仮定する。
\begin{enumerate}
\item $Y' \to B'$ は局所有限型である。
\item $X' \to B'$ は平坦かつ局所有限表示である。
\item $f$ は平坦である。
\item $X = B \times_{B'} X'$ かつ $Y = B \times_{B'} Y'$ である。
\end{enumerate}
このとき $f'$ は平坦であり、任意の $y' \in |Y'|$ で $|f'|$ の像に
属するものに対し、射 $Y' \to B'$ は $y'$ で平坦である。
\end{lemma}

\begin{proof}
スキーム $U'$ と全射エタール射 $U' \to B'$ を選ぶ。
スキーム $V'$ と全射エタール射 $V' \to U' \times_{B'} Y'$ を選ぶ。
スキーム $W'$ と全射エタール射 $W' \to V' \times_{Y'} X'$ を選ぶ。
$U, V, W$ を $U', V', W'$ の $B \to B'$ による基底変換とする。
このとき $f'$ の平坦性は $W' \to V'$ の平坦性と同値であり、
$W \to V$ は平坦であると与えられている。したがってスキームの肥厚化の図式
$$
\xymatrix{
(W \subset W') \ar[rr] \ar[rd] & & (V \subset V') \ar[ld] \\
& (U \subset U')
}
$$
にスキームの場合の補題を適用できる。More on Morphisms, Lemma
\ref{more-morphisms-lemma-flatness-morphism-thickenings-fp-over-ft} を参照せよ。
$Y'/B'$ の平坦性に関する、$f'$ の像に属する点での主張も同様に従う。
\end{proof}

\noindent
スキームの射の多くの性質は、上の補題にあるような平坦変形で保たれる。

\begin{lemma}
\label{spaces-more-morphisms-lemma-deform-property-fp-over-ft}
$S$ をスキームとする。可換図式
$$
\xymatrix{
(X \subset X') \ar[rr]_{(f, f')} \ar[rd] & & (Y \subset Y') \ar[ld] \\
& (B \subset B')
}
$$
を $S$ 上の代数空間の肥厚化について考える。
% P09-ERR-1095: frozen source lines 4304--4305 omit "is" in both assumptions.
$Y' \to B'$ は局所有限型、$X' \to B'$ は平坦かつ局所有限表示、
$X = B \times_{B'} X'$、かつ $Y = B \times_{B'} Y'$ であると仮定する。
このとき次が成り立つ。
\begin{enumerate}
\item $f$ が表現可能であることと $f'$ が表現可能であることは同値である。
\label{spaces-more-morphisms-item-representable-fp-over-ft}
\item $f$ が平坦であることと $f'$ が平坦であることは同値である。
\label{spaces-more-morphisms-item-flat-fp-over-ft}
\item $f$ が同型であることと $f'$ が同型であることは同値である。
\label{spaces-more-morphisms-item-isomorphism-fp-over-ft}
\item $f$ が開埋め込みであることと $f'$ が開埋め込みであることは同値である。
\label{spaces-more-morphisms-item-open-immersion-fp-over-ft}
\item $f$ が準コンパクトであることと $f'$ が準コンパクトであることは同値である。
\label{spaces-more-morphisms-item-quasi-compact-fp-over-ft}
\item $f$ が普遍的閉であることと $f'$ が普遍的閉であることは同値である。
\label{spaces-more-morphisms-item-universally-closed-fp-over-ft}
\item $f$ が（準）分離的であることと $f'$ が（準）分離的であることは同値である。
\label{spaces-more-morphisms-item-separated-fp-over-ft}
\item $f$ が単射であることと $f'$ が単射であることは同値である。
\label{spaces-more-morphisms-item-monomorphism-fp-over-ft}
\item $f$ が全射であることと $f'$ が全射であることは同値である。
\label{spaces-more-morphisms-item-surjective-fp-over-ft}
\item $f$ が普遍的単射であることと $f'$ が普遍的単射であることは同値である。
\label{spaces-more-morphisms-item-universally-injective-fp-over-ft}
\item $f$ がアフィンであることと $f'$ がアフィンであることは同値である。
\label{spaces-more-morphisms-item-affine-fp-over-ft}
\item $f$ が局所準有限であることと $f'$ が局所準有限であることは同値である。
\label{spaces-more-morphisms-item-quasi-finite-fp-over-ft}
\item
\label{spaces-more-morphisms-item-relative-dimension-d-fp-over-ft}
$f$ が相対次元 $d$ の局所有限型であることと、$f'$ が相対次元 $d$ の
局所有限型であることは同値である。
\item $f$ が普遍的開であることと $f'$ が普遍的開であることは同値である。
\label{spaces-more-morphisms-item-universally-open-fp-over-ft}
\item $f$ がシントミックであることと $f'$ がシントミックであることは同値である。
\label{spaces-more-morphisms-item-syntomic-fp-over-ft}
\item $f$ が滑らかであることと $f'$ が滑らかであることは同値である。
\label{spaces-more-morphisms-item-smooth-fp-over-ft}
\item $f$ が非分岐であることと $f'$ が非分岐であることは同値である。
\label{spaces-more-morphisms-item-unramified-fp-over-ft}
\item $f$ がエタールであることと $f'$ がエタールであることは同値である。
\label{spaces-more-morphisms-item-etale-fp-over-ft}
\item $f$ が固有であることと $f'$ が固有であることは同値である。
\label{spaces-more-morphisms-item-proper-fp-over-ft}
\item $f$ が有限であることと $f'$ が有限であることは同値である。
\label{spaces-more-morphisms-item-finite-fp-over-ft}
\item
\label{spaces-more-morphisms-item-finite-locally-free-fp-over-ft}
$f$ が（階数 $d$ の）有限局所自由であることと、$f'$ が（階数 $d$ の）
有限局所自由であることは同値である。
% P09-ERR-1098: frozen source line 4354 retains the editorial placeholder "add more here".
\item さらに追加せよ。
\end{enumerate}
\end{lemma}

\begin{proof}
(\ref{spaces-more-morphisms-item-representable-fp-over-ft}) の場合は Lemma
\ref{spaces-more-morphisms-lemma-thicken-property-morphisms} から従う。

\medskip\noindent
スキーム $U'$ と全射エタール射 $U' \to B'$ を選ぶ。
スキーム $V'$ と全射エタール射 $V' \to U' \times_{B'} Y'$ を選ぶ。
スキーム $W'$ と全射エタール射 $W' \to V' \times_{Y'} X'$ を選ぶ。
$U, V, W$ を $U', V', W'$ の $B \to B'$ による基底変換とする。
スキームの肥厚化の図式
$$
\xymatrix{
(W \subset W') \ar[rr] \ar[rd] & & (V \subset V') \ar[ld] \\
& (U \subset U')
}
$$
を考える。始域と終域の双方でエタール局所的な任意の性質について、結果は
上の図式に適用したスキームの肥厚化の射に対する対応する結果から直ちに従う。
したがって場合
(\ref{spaces-more-morphisms-item-flat-fp-over-ft}),
(\ref{spaces-more-morphisms-item-quasi-finite-fp-over-ft}),
(\ref{spaces-more-morphisms-item-relative-dimension-d-fp-over-ft}),
(\ref{spaces-more-morphisms-item-syntomic-fp-over-ft}),
(\ref{spaces-more-morphisms-item-smooth-fp-over-ft}),
(\ref{spaces-more-morphisms-item-unramified-fp-over-ft}),
(\ref{spaces-more-morphisms-item-etale-fp-over-ft})
は More on Morphisms, Lemma
\ref{more-morphisms-lemma-deform-property-fp-over-ft} の対応する場合から従う。

\medskip\noindent
$X \to X'$ と $Y \to Y'$ は普遍同相なので、写像 $X \to Y$ と
$X' \to Y'$ の位相に関する任意の問いは同じ答えをもつ。したがって場合
(\ref{spaces-more-morphisms-item-quasi-compact-fp-over-ft}),
(\ref{spaces-more-morphisms-item-universally-closed-fp-over-ft}),
(\ref{spaces-more-morphisms-item-surjective-fp-over-ft}),
(\ref{spaces-more-morphisms-item-universally-injective-fp-over-ft}), and
(\ref{spaces-more-morphisms-item-universally-open-fp-over-ft}) が成り立つ。

\medskip\noindent
残る各場合では、含意
$f\text{ has }P \Rightarrow f'\text{ has }P$
だけを証明する。逆の含意は $P$ が基底変換で安定であることから従う。
Spaces, Lemma \ref{spaces-lemma-base-change-immersions} および
Morphisms of Spaces, Lemmas
\ref{spaces-morphisms-lemma-base-change-separated},
\ref{spaces-morphisms-lemma-base-change-monomorphism},
\ref{spaces-morphisms-lemma-base-change-affine},
\ref{spaces-morphisms-lemma-base-change-proper},
\ref{spaces-morphisms-lemma-base-change-integral}, and
\ref{spaces-morphisms-lemma-base-change-finite-locally-free} を参照せよ。

\medskip\noindent
(\ref{spaces-more-morphisms-item-open-immersion-fp-over-ft}) の場合。
$f$ は開埋め込みであると仮定する。このとき $f'$ は
(\ref{spaces-more-morphisms-item-etale-fp-over-ft}) によりエタールであり、
(\ref{spaces-more-morphisms-item-universally-injective-fp-over-ft}) により普遍的単射なので、
$f'$ は開埋め込みである。Morphisms of Spaces, Lemma
\ref{spaces-morphisms-lemma-etale-universally-injective-open} を参照せよ。
この補題を使わず、まず (\ref{spaces-more-morphisms-item-representable-fp-over-ft}) を用いて
スキームの場合に帰着することもできる。

\medskip\noindent
(\ref{spaces-more-morphisms-item-isomorphism-fp-over-ft}) の場合。場合
(\ref{spaces-more-morphisms-item-open-immersion-fp-over-ft}) と
(\ref{spaces-more-morphisms-item-surjective-fp-over-ft}) から従う。

\medskip\noindent
(\ref{spaces-more-morphisms-item-separated-fp-over-ft}) の場合。Lemma
\ref{spaces-more-morphisms-lemma-thicken-property-morphisms} を参照せよ。

\medskip\noindent
(\ref{spaces-more-morphisms-item-monomorphism-fp-over-ft}) の場合。$f$ は単射であると仮定する。
対角射 $\Delta_{X'/Y'} : X' \to X' \times_{Y'} X'$ を考える。
$\Delta_{X'/Y'}$ の $B \to B'$ による基底変換は $\Delta_{X/Y}$ であり、
仮定により同型である。(\ref{spaces-more-morphisms-item-isomorphism-fp-over-ft}) により
$\Delta_{X'/Y'}$ は同型であり、したがって $f'$ は単射である。

\medskip\noindent
(\ref{spaces-more-morphisms-item-affine-fp-over-ft}) の場合。Lemma
\ref{spaces-more-morphisms-lemma-thicken-property-morphisms} を参照せよ。

\medskip\noindent
(\ref{spaces-more-morphisms-item-proper-fp-over-ft}) の場合。Lemma
\ref{spaces-more-morphisms-lemma-properties-that-extend-over-thickenings} を参照せよ。

\medskip\noindent
(\ref{spaces-more-morphisms-item-finite-fp-over-ft}) の場合。Lemma
\ref{spaces-more-morphisms-lemma-properties-that-extend-over-thickenings} を参照せよ。

\medskip\noindent
(\ref{spaces-more-morphisms-item-finite-locally-free-fp-over-ft}) の場合。
$f$ は有限局所自由であると仮定する。
(\ref{spaces-more-morphisms-item-finite-fp-over-ft}) により $f'$ は有限である。
(\ref{spaces-more-morphisms-item-flat-fp-over-ft}) により $f'$ は平坦である。
% P09-ERR-1096: frozen source line 4457 omits "of" in "locally of finite presentation".
また Morphisms of Spaces, Lemma
\ref{spaces-morphisms-lemma-finite-presentation-permanence} により、
$f'$ は局所有限表示である。したがって Morphisms of Spaces, Lemma
\ref{spaces-morphisms-lemma-finite-flat} により $f'$ は有限局所自由である。
\end{proof}













\section{形式的に滑らかな射}
\label{spaces-more-morphisms-section-formally-smooth}

\noindent
この節では、代数空間の形式的に滑らかな射 $X \to Y$ の概念を導入する。
このような射は、$T$ がアフィンであるとき、$X$ の $T$ 値点が $T$ の
無限小肥厚化へ持ち上がるという性質によって特徴づけられる。
主結果は、形式的に滑らかかつ局所有限表示な射が滑らかであるというものである。
Lemma \ref{spaces-more-morphisms-lemma-smooth-formally-smooth} を参照せよ。この判定法はしばしば
ヤコビ判定法より使いやすいことが分かる。

\begin{definition}
\label{spaces-more-morphisms-definition-formally-smooth}
$S$ をスキームとする。代数空間の射 $f : X \to Y$ が $S$ 上で
{\it 形式的に滑らか}であるとは、Definition
\ref{spaces-more-morphisms-definition-formally-smooth-etale-unramified} の意味で、
関手の変換として形式的に滑らかであることをいう。
\end{definition}

\noindent
形式的非分岐射および形式的エタール射の場合には、$T'$ がアフィンであるという
条件を取り除けた。Lemmas \ref{spaces-more-morphisms-lemma-formally-unramified-not-affine} and
\ref{spaces-more-morphisms-lemma-formally-etale-not-affine} を参照せよ。形式的に滑らかな射の場合、
これはもはや正しくない。実際、もう少し自然な条件は、$T'$ 上エタール局所的に
点線矢印を補えることだろう。Lemma \ref{spaces-more-morphisms-lemma-smooth-formally-smooth} の証明を
分析すると、これは現在の定義と同値になることが分かる。また、$T'$ 上 fppf
局所的に点線矢印の存在を要求しても十分であるが、その証明は少し難しい。

\medskip\noindent
Section \ref{spaces-more-morphisms-section-formally-smooth-etale-unramified} で、より一般の
形式的に滑らかな関手の変換について証明した結果はここでは再掲しない。

\begin{lemma}
\label{spaces-more-morphisms-lemma-composition-formally-smooth}
形式的に滑らかな射の合成は形式的に滑らかである。
\end{lemma}

\begin{proof}
省略する。
\end{proof}

\begin{lemma}
\label{spaces-more-morphisms-lemma-base-change-formally-smooth}
形式的に滑らかな射の基底変換は形式的に滑らかである。
\end{lemma}

\begin{proof}
省略する。ただし代数版については Algebra, Lemma
\ref{algebra-lemma-base-change-fs} を参照せよ。
\end{proof}

\begin{lemma}
\label{spaces-more-morphisms-lemma-formally-etale-unramified-smooth}
$f : X \to S$ をスキームの射とする。このとき $f$ が形式的エタールであることと、
$f$ が形式的に滑らかかつ形式的非分岐であることは同値である。
\end{lemma}

\begin{proof}
省略する。
\end{proof}

\noindent
次は Lemma \ref{spaces-more-morphisms-lemma-formally-smooth} によって後に置き換えられる補助補題である。

\begin{lemma}
\label{spaces-more-morphisms-lemma-helper-formally-smooth}
$S$ をスキームとする。可換図式
$$
\xymatrix{
U \ar[d] \ar[r]_\psi & V \ar[d] \\
X \ar[r]^f & Y
}
$$
を $S$ 上の代数空間の射について考える。縦の矢印がエタールで、
$f$ が形式的に滑らかならば、
$\psi$ は形式的に滑らかである。
\end{lemma}

\begin{proof}
Lemma \ref{spaces-more-morphisms-lemma-representable-property-formally-property} により、射
$U \to X$ と $V \to Y$ は形式的エタールである。Lemma
\ref{spaces-more-morphisms-lemma-composition-formally-smooth-etale-unramified} により、合成
$U \to Y$ は形式的に滑らかである。Lemma
\ref{spaces-more-morphisms-lemma-formally-permanence} により、$\psi : U \to V$ は形式的に
滑らかである。
\end{proof}

\noindent
次の補題がこの節の主結果である。Limits of Spaces, Proposition
\ref{spaces-limits-proposition-characterize-locally-finite-presentation}
と組み合わせると、代数空間の射 $f : X \to Y$ が滑らかかどうかを、
関手の変換 $X \to Y$ の「単純な」性質で認識できることが従う。

\begin{lemma}[無限小持上げ判定法]
\label{spaces-more-morphisms-lemma-smooth-formally-smooth}
$S$ をスキームとする。
$f : X \to Y$ を $S$ 上の代数空間の射とする。
次は同値である。
\begin{enumerate}
\item 射 $f$ は滑らかである。
\item 射 $f$ は局所有限表示かつ形式的に滑らかである。
\end{enumerate}
\end{lemma}

\begin{proof}
% P09-ERR-1100: frozen source line 4599 switches the lemma's target Y to S.
$f : X \to S$ は局所有限表示かつ形式的に滑らかであると仮定する。
可換図式
$$
\xymatrix{
U \ar[d] \ar[r]_\psi & V \ar[d] \\
X \ar[r]^f & Y
}
$$
で、$U$ と $V$ がスキームであり、縦の矢印がエタールかつ全射であるものを
考える。Lemma \ref{spaces-more-morphisms-lemma-helper-formally-smooth} により
$\psi : U \to V$ は形式的に滑らかである。Morphisms of Spaces, Lemma
\ref{spaces-morphisms-lemma-finite-presentation-local} により、射 $\psi$ は
局所有限表示である。したがってスキームの場合により射 $\psi$ は滑らかである。
More on Morphisms, Lemma
\ref{more-morphisms-lemma-smooth-formally-smooth} を参照せよ。したがって
$f$ は滑らかである。Morphisms of Spaces, Lemma
\ref{spaces-morphisms-lemma-smooth-local} を参照せよ。

\medskip\noindent
逆に、$f : X \to Y$ は滑らかであると仮定する。Definition
\ref{spaces-more-morphisms-definition-formally-smooth} にあるような実線部分が可換な図式
$$
\xymatrix{
X \ar[d]_f & T \ar[d]^i \ar[l]^a \\
Y & T' \ar[l] \ar@{-->}[lu]
}
$$
を考える。点線矢印が存在することを示し、それによって $f$ が形式的に
滑らかであることを証明する。$\mathcal{F}$ を Lemma
\ref{spaces-more-morphisms-lemma-sheaf} の $(T')_{spaces, \etale}$ 上の集合の層とし、ここでは
Remark \ref{spaces-more-morphisms-remark-special-case} で論じた特別な場合を考える。また
$$
\mathcal{H} =
\SheafHom_{\mathcal{O}_T}(a^*\Omega_{X/Y}, \mathcal{C}_{T/T'})
$$
を $\mathcal{O}_T$ 加群の層として $T_{spaces, \etale}$ 上に考え、Lemma
\ref{spaces-more-morphisms-lemma-action-sheaf} の作用
$\mathcal{H} \times \mathcal{F} \to \mathcal{F}$ を備えるものとする。
この作用 $\mathcal{H} \times \mathcal{F} \to \mathcal{F}$ は
$\mathcal{F}$ を擬 $\mathcal{H}$ トーサーにする。Cohomology on Sites,
Definition \ref{sites-cohomology-definition-torsor} を参照せよ。
目標は $\mathcal{F}$ が自明な $\mathcal{H}$ トーサーであることを示すことである。
二つの段階がある。(I) $\mathcal{F}$ がトーサーであることを示すには、
$\mathcal{F}$ がエタール局所的に切断をもつことを示さなければならない。
(II) $\mathcal{F}$ が自明なトーサーであることを示すには
$H^1(T_\etale, \mathcal{H}) = 0$ を示せば十分である。Cohomology on Sites,
Lemma \ref{sites-cohomology-lemma-torsors-h1} を参照せよ。

\medskip\noindent
まず (I) を証明する。そのため可換図式
$$
\xymatrix{
U \ar[d] \ar[r]_\psi & V \ar[d] \\
X \ar[r]^f & Y
}
$$
で、$U$ と $V$ がスキームであり、縦の矢印がエタールかつ全射であるものを
選ぶ。$f$ は滑らかと仮定したので $\psi$ は滑らかであり、したがって
Lemma \ref{spaces-more-morphisms-lemma-representable-property-formally-property} により形式的に
滑らかである。同じ補題により射 $V \to Y$ は形式的エタールである。
したがって Lemma \ref{spaces-more-morphisms-lemma-composition-formally-smooth-etale-unramified}
により合成 $U \to Y$ は形式的に滑らかである。このとき (I) は Lemma
\ref{spaces-more-morphisms-lemma-etale-on-top} part (4) から従う。

\medskip\noindent
最後に (II) を証明する。Lemma
\ref{spaces-more-morphisms-lemma-finite-presentation-differentials} により
% P09-ERR-1101: frozen source lines 4672--4673 use Omega_{X/S}, while the morphism and H are relative to Y.
$\Omega_{X/S}$ は有限表示である。したがって $a^*\Omega_{X/S}$ は有限表示である
（Properties of Spaces, Section
\ref{spaces-properties-section-properties-modules} を参照）。したがって層
$\mathcal{H} =
\SheafHom_{\mathcal{O}_T}(a^*\Omega_{X/Y}, \mathcal{C}_{T/T'})$
は準連接である。Properties of Spaces, Lemma
\ref{spaces-properties-lemma-properties-quasi-coherent} を参照せよ。
ゆえに Descent, Proposition
\ref{descent-proposition-same-cohomology-quasi-coherent} および
Cohomology of Schemes, Lemma
\ref{coherent-lemma-quasi-coherent-affine-cohomology-zero} により、望みどおり
$$
H^1(T_{spaces, \etale}, \mathcal{H}) =
H^1(T_\etale, \mathcal{H}) =
H^1(T, \mathcal{H}) = 0
$$
である。
\end{proof}

\noindent
滑らかな射は強い局所持上げ性質を満たす。Lemma
\ref{spaces-more-morphisms-lemma-smooth-strong-lift} を参照せよ。この補題で $T'$ がアフィンであると
仮定した場合、エタール被覆を取る必要があるかどうかは分からない。
より正確には、代数空間の可換図式
$$
\xymatrix{
X \ar[d] & T \ar[l] \ar[d] \\
Y & T' \ar[l] \ar@{..>}[lu]
}
$$
で、$X \to Y$ が滑らかで $T \to T'$ がアフィンスキームの肥厚化であるものが
与えられたとき、図式を可換にする点線矢印は常に存在するだろうか。
% P09-ERR-1102: frozen source line 4709 has "the does" instead of "then does".
答え、または参考文献を知っている場合は
\href{mailto:stacks.project@gmail.com}{stacks.project@gmail.com} まで連絡されたい。

\begin{lemma}
\label{spaces-more-morphisms-lemma-smooth-strong-lift}
$S$ をスキームとする。可換図式
$$
\xymatrix{
X \ar[d] & T \ar[l] \ar[d] \\
Y & T' \ar[l]
}
$$
を $S$ 上の代数空間について考える。ここで $X \to Y$ は滑らかで、
$T \to T'$ は肥厚化である。このときエタール被覆
$\{T'_i \to T'\}$ が存在し、図式
$$
\xymatrix{
X \ar[d] & T \ar[l] \ar[d] & T \times_{T'} T'_i \ar[l] \ar[d] \\
Y & T' \ar[l] & T'_i \ar[l] \ar@{..>}[llu]
}
$$
を可換にする点線矢印を、すべての $i$ に対して見つけることができる。
\end{lemma}

\begin{proof}
エタール被覆 $\{Y_i \to Y\}$ で、各 $Y_i$ がアフィンであるものを選ぶ。
$T'$ を誘導されるエタール被覆で置き換えると、$Y$ はアフィンであるとしてよい。

\medskip\noindent
$Y$ はアフィンであると仮定する。エタール被覆 $\{X_i \to X\}$ を選ぶ。
これは $T$ のエタール被覆を生じる。この $T$ のエタール被覆は $T'$ の
エタール被覆から来る（Theorem \ref{spaces-more-morphisms-theorem-topological-invariance} および
Section \ref{spaces-more-morphisms-section-thickenings} の議論を参照）。したがって $X$ は
アフィンであるとしてよい。

\medskip\noindent
$X$ と $Y$ はアフィンであると仮定する。$T'$ のエタール被覆をもう一度取り、
$T'$ はアフィンであるとしてよい。この場合、補題は Algebra, Lemma
\ref{algebra-lemma-smooth-strong-lift} から従う。
\end{proof}







\noindent
形式的に滑らかであることが始域上エタール局所的であることを示すため、
もう少し準備をする。まず、形式的に滑らかな射の微分の層が良い性質をもつことを
示す。局所射影的な準連接加群の概念は Properties of Spaces, Section
\ref{spaces-properties-section-locally-projective} で定義されている。

\begin{lemma}
\label{spaces-more-morphisms-lemma-formally-smooth-sheaf-differentials}
$S$ をスキームとする。$f : X \to Y$ を $S$ 上の代数空間の
形式的に滑らかな射とする。このとき $\Omega_{X/Y}$ は $X$ 上局所射影的である。
\end{lemma}

\begin{proof}
図式
$$
\xymatrix{
U \ar[d] \ar[r]_\psi & V \ar[d] \\
X \ar[r]^f & Y
}
$$
で、$U$ と $V$ がアフィン（！）スキームであり、縦の矢印がエタールであるものを
選ぶ。Lemma \ref{spaces-more-morphisms-lemma-helper-formally-smooth} により
$\psi : U \to V$ は形式的に滑らかである。したがって
$\Gamma(V, \mathcal{O}_V) \to \Gamma(U, \mathcal{O}_U)$ は
形式的に滑らかな環準同型である。More on Morphisms, Lemma
\ref{more-morphisms-lemma-affine-formally-smooth} を参照せよ。ゆえに Algebra,
Lemma \ref{algebra-lemma-characterize-formally-smooth-again} により、
$\Gamma(U, \mathcal{O}_U)$ 加群
$\Omega_{\Gamma(U, \mathcal{O}_U)/\Gamma(V, \mathcal{O}_V)}$
は射影的である。したがって $\Omega_{U/V}$ は局所射影的である。
Properties, Section \ref{properties-section-locally-projective} を参照せよ。
$\Omega_{X/Y}|_U = \Omega_{U/V}$ なので、$\Omega_{X/Y}$ も局所射影的である。
（$X$ のエタール被覆を、上のような図式に入るアフィンな $U$ たちによって
見つけられるためである---詳細は省略する。）
\end{proof}

\begin{lemma}
\label{spaces-more-morphisms-lemma-h1-is-zero}
$T$ をアフィンスキームとする。$\mathcal{F}$ と $\mathcal{G}$ を準連接
$\mathcal{O}_T$ 加群で、$T_\etale$ 上にあるものとする。内部 Hom 層
$\mathcal{H} = \SheafHom_{\mathcal{O}_T}(\mathcal{F}, \mathcal{G})$ を
$T_\etale$ 上で考える。$\mathcal{F}$ が局所射影的ならば
$H^1(T_\etale, \mathcal{H}) = 0$ である。
\end{lemma}

\begin{proof}
代数空間上の局所射影的な層の定義（Properties of Spaces, Definition
\ref{spaces-properties-definition-locally-projective} を参照）により、
$\mathcal{F}_{Zar} = \mathcal{F}|_{T_{Zar}}$ はスキーム $T$ 上の
局所射影的な層である。したがって $\mathcal{F}_{Zar}$ は自由な
$\mathcal{O}_{T_{Zar}}$ 加群の直和因子である。そこで
$\mathcal{F} = (\mathcal{F}_{Zar})^a$ であること（Descent, Proposition
\ref{descent-proposition-equivalence-quasi-coherent} を参照）から、
$\mathcal{F}$ は自由な $\mathcal{O}_T$ 加群の直和因子として
$T_\etale$ 上にあると
結論する。したがって $\mathcal{F} = \bigoplus_{i \in I} \mathcal{O}_T$ が
自由加群であると仮定してよい。この場合
$\mathcal{H} = \prod_{i \in I} \mathcal{G}$ は準連接加群の積である。
Cohomology on Sites, Lemma
\ref{sites-cohomology-lemma-cohomology-products} により、アフィンスキーム上の
準連接層のコホモロジーは零なので $H^1 = 0$ と結論する。Descent,
Proposition \ref{descent-proposition-same-cohomology-quasi-coherent} および
Cohomology of Schemes, Lemma
\ref{coherent-lemma-quasi-coherent-affine-cohomology-zero} を参照せよ。
\end{proof}

\begin{lemma}
\label{spaces-more-morphisms-lemma-formally-smooth}
$S$ をスキームとする。$f : X \to Y$ を $S$ 上の代数空間の射とする。
次は同値である。
\begin{enumerate}
\item $f$ は形式的に滑らかである。
\item 任意の図式
$$
\xymatrix{
U \ar[d] \ar[r]_\psi & V \ar[d] \\
X \ar[r]^f & Y
}
$$
で、$U$ と $V$ がスキームであり、縦の矢印がエタールであるものに対して、
スキームの射 $\psi$ は形式的に滑らかである（More on Morphisms,
% P09-ERR-1103: frozen source line 4853 cites the formally-unramified definition for formally smooth.
Definition \ref{more-morphisms-definition-formally-unramified} の意味で）。
\item 全射な縦の矢印をもつこのような図式が一つ存在し、その図式における射
$\psi$ が形式的に滑らかである。
\end{enumerate}
\end{lemma}

\begin{proof}
Lemma \ref{spaces-more-morphisms-lemma-helper-formally-smooth} で、(1) が (2) と (3) を含意することを
見た。(3) を仮定する。$f$ が形式的に滑らかであることの証明は、Lemma
\ref{spaces-more-morphisms-lemma-smooth-formally-smooth} の (1) $\Rightarrow$ (2) の証明と
まったく同様である。

\medskip\noindent
Definition \ref{spaces-more-morphisms-definition-formally-smooth} にあるような実線部分が可換な図式
$$
\xymatrix{
X \ar[d]_f & T \ar[d]^i \ar[l]^a \\
Y & T' \ar[l] \ar@{-->}[lu]
}
$$
を考える。点線矢印が存在することを示し、それによって $f$ が形式的に
滑らかであることを証明する。$\mathcal{F}$ を Lemma \ref{spaces-more-morphisms-lemma-sheaf} の
$(T')_{spaces, \etale}$ 上の集合の層とし、ここでは Remark
\ref{spaces-more-morphisms-remark-special-case} で論じた特別な場合を考える。
$$
\mathcal{H} =
\SheafHom_{\mathcal{O}_T}(a^*\Omega_{X/Y}, \mathcal{C}_{T/T'})
$$
を $\mathcal{O}_T$ 加群の層として $T_{spaces, \etale}$ 上に考え、Lemma
\ref{spaces-more-morphisms-lemma-action-sheaf} の作用
$\mathcal{H} \times \mathcal{F} \to \mathcal{F}$ を備えるものとする。
この作用 $\mathcal{H} \times \mathcal{F} \to \mathcal{F}$ は
$\mathcal{F}$ を擬 $\mathcal{H}$ トーサーにする。Cohomology on Sites,
Definition \ref{sites-cohomology-definition-torsor} を参照せよ。目標は
$\mathcal{F}$ が自明な $\mathcal{H}$ トーサーであることを示すことである。
二つの段階がある。(I) $\mathcal{F}$ がトーサーであることを示すには、
$\mathcal{F}$ がエタール局所的に切断をもつことを示さなければならない。
(II) $\mathcal{F}$ が自明なトーサーであることを示すには
$H^1(T_\etale, \mathcal{H}) = 0$ を示せば十分である。Cohomology on Sites,
Lemma \ref{sites-cohomology-lemma-torsors-h1} を参照せよ。

\medskip\noindent
まず (I) を証明する。そのため、(3) を仮定しているので存在する図式
$$
\xymatrix{
U \ar[d] \ar[r]_\psi & V \ar[d] \\
X \ar[r]^f & Y
}
$$
を考える。ここで $U$ と $V$ はスキーム、縦の矢印はエタールかつ全射であり、
$\psi$ は形式的に滑らかである。Lemma
\ref{spaces-more-morphisms-lemma-representable-property-formally-property} により、射 $V \to Y$ は
形式的エタールである。したがって Lemma
\ref{spaces-more-morphisms-lemma-composition-formally-smooth-etale-unramified} により、合成
$U \to Y$ は形式的に滑らかである。このとき (I) は Lemma
\ref{spaces-more-morphisms-lemma-etale-on-top} part (4) から従う。

\medskip\noindent
最後に (II) を証明する。Lemma
\ref{spaces-more-morphisms-lemma-formally-smooth-sheaf-differentials} により、
% P09-ERR-1104: frozen source line 4919 omits the copula in "Omega_{U/V} locally projective".
$\Omega_{U/V}$ は局所射影的である。したがって $\Omega_{X/Y}$ は
局所射影的である。Descent on Spaces, Lemma
\ref{spaces-descent-lemma-locally-projective-descends} を参照せよ。
ゆえに $a^*\Omega_{X/Y}$ は局所射影的である。Properties of Spaces, Lemma
\ref{spaces-properties-lemma-locally-projective-pullback} を参照せよ。したがって
% P09-ERR-1105: frozen source line 4930 lacks a closing parenthesis before = 0.
$$
H^1(T_\etale, \mathcal{H}) =
H^1(T_\etale,
\SheafHom_{\mathcal{O}_T}(a^*\Omega_{X/Y}, \mathcal{C}_{T/T'}) = 0
$$
である。Lemma \ref{spaces-more-morphisms-lemma-h1-is-zero} により、これで示された。
\end{proof}

\begin{lemma}
\label{spaces-more-morphisms-lemma-descending-property-formally-smooth}
性質 $\mathcal{P}(f) =$「$f$ は形式的に滑らかである」は基底上 fpqc 局所的である。
\end{lemma}

\begin{proof}
$f : X \to Y$ をスキーム $S$ 上の代数空間の射とする。添字集合 $I$ と図式
$$
\xymatrix{
U_i \ar[d] \ar[r]_{\psi_i} & V_i \ar[d] \\
X \ar[r]^f & Y
}
$$
を選び、縦の矢印はエタール、$U_i$ と $V_i$ はアフィンスキームであるとする。
さらに $\coprod U_i \to X$ と $\coprod V_i \to Y$ は全射であると仮定する。
Properties of Spaces, Lemma
\ref{spaces-properties-lemma-cover-by-union-affines} を参照せよ。Lemma
\ref{spaces-more-morphisms-lemma-formally-smooth} により、$f$ が形式的に滑らかであることと、各射
% P09-ERR-1106: frozen source line 4958 has subject-verb disagreement in "each ... are".
$\psi_i$ が形式的に滑らかであることは同値である。したがってアフィンスキームの
射の場合に帰着する。この場合、結果は Algebra, Lemma
\ref{algebra-lemma-descent-formally-smooth} から従う。いくつかの詳細は省略する。
\end{proof}

\begin{lemma}
\label{spaces-more-morphisms-lemma-triangle-differentials-formally-smooth}
$S$ をスキームとする。$f : X \to Y$ と $g : Y \to Z$ を $S$ 上の
代数空間の射とする。$f$ が形式的に滑らかであると仮定する。このとき
$$
0 \to f^*\Omega_{Y/Z} \to \Omega_{X/Z} \to \Omega_{X/Y} \to 0
$$
という Lemma \ref{spaces-more-morphisms-lemma-triangle-differentials} の列は短完全である。
\end{lemma}

\begin{proof}
エタール局所化によりスキームの場合から従う。More on Morphisms, Lemma
\ref{more-morphisms-lemma-triangle-differentials-formally-smooth} と、Lemmas
\ref{spaces-more-morphisms-lemma-formally-smooth} and \ref{spaces-more-morphisms-lemma-localize-differentials} を参照せよ。
\end{proof}

\begin{lemma}
\label{spaces-more-morphisms-lemma-differentials-formally-unramified-formally-smooth}
$S$ をスキームとし、$B$ を $S$ 上の代数空間とする。
$h : Z \to X$ を $B$ 上の代数空間の形式的非分岐射とする。
$Z$ が $B$ 上形式的に滑らかであると仮定する。このとき標準完全列
$$
0 \to \mathcal{C}_{Z/X} \to h^*\Omega_{X/B} \to \Omega_{Z/B} \to 0
$$
は短完全である。Lemma
\ref{spaces-more-morphisms-lemma-universally-unramified-differentials-sequence} を参照せよ。
\end{lemma}

\begin{proof}
$Z \to Z'$ を、$Z$ の $X$ 上の普遍一次肥厚化とする。Lemma
\ref{spaces-more-morphisms-lemma-universally-unramified-differentials-sequence} の証明から、上の列は
$$
\mathcal{C}_{Z/Z'} \to \Omega_{Z'/B} \otimes \mathcal{O}_Z \to
\Omega_{Z/B} \to 0.
$$
と同一視される。
% P09-ERR-1107: frozen source line 5009 switches the assumed formal-smooth base B to S.
$Z \to S$ は形式的に滑らかなので、$Z'$ 上エタール局所的に射
$Z' \to Z$ を $B$ 上で見つけられ、これは包含写像 $Z \to Z'$ の左逆である。
したがって列は
エタール局所的に分裂する。Lemma
\ref{spaces-more-morphisms-lemma-differentials-relative-immersion-section} を参照せよ。
\end{proof}

\begin{lemma}
\label{spaces-more-morphisms-lemma-two-unramified-morphisms-formally-smooth}
$S$ をスキームとする。可換図式
$$
\xymatrix{
Z \ar[r]_i \ar[rd]_j & X \ar[d]^f \\
& Y
}
$$
を $S$ 上の代数空間について考える。ここで $i$ と $j$ は形式的非分岐、
$f$ は形式的に滑らかである。
このとき標準完全列
$$
0 \to
\mathcal{C}_{Z/Y} \to
\mathcal{C}_{Z/X} \to
i^*\Omega_{X/Y} \to 0
$$
は完全かつ局所的に分裂する。Lemma
\ref{spaces-more-morphisms-lemma-two-unramified-morphisms} を参照せよ。
\end{lemma}

\begin{proof}
$Z \to Z'$ を、$Z$ の $X$ 上の普遍一次肥厚化とする。
$Z \to Z''$ を、$Z$ の $Y$ 上の普遍一次肥厚化とする。Lemma
\ref{spaces-more-morphisms-lemma-universally-unramified-differentials-sequence} により、標準射
% P09-ERR-1108: frozen source line 5043 says "here is" instead of "there is".
$Z' \to Z''$ が存在し、可換図式
$$
\xymatrix{
Z \ar[r]_{i'} \ar[rd]_{j'} & Z' \ar[r]_a \ar[d]^k & X \ar[d]^f \\
& Z'' \ar[r]^b & Y
}
$$
を得る。上の列は、形式的非分岐射の余法層に関する定義により、列
$$
\mathcal{C}_{Z/Z''} \to
\mathcal{C}_{Z/Z'} \to
(i')^*\Omega_{Z'/Z''} \to 0
$$
と同一視される。$U'' \to Z''$ を $U''$ がアフィンであるようなエタール射とする。
$U \to Z$ および $U' \to Z'$ を、それぞれ $Z$ および $Z'$ 上エタールな
対応するアフィンスキームとする。$f$ は形式的に滑らかなので、射
$h : U'' \to X$ で、$i$ と $U$ 上で一致し、$f \circ h$ が $b|_{U''}$ に等しい
ものが存在する。$Z'$ は普遍一次肥厚化なので、一意な射
% P09-ERR-1109: frozen source line 5064 has the ill-typed equation g = a circ h; types require a circ g = h.
$g : U'' \to Z'$ で $g = a \circ h$ を満たすものを得る。$Z''$ の普遍性により、
$k \circ g$ は包含写像 $U'' \to Z''$ である。したがって $g$ は $k$ の左逆である。
図式で表すと
$$
\xymatrix{
U \ar[d] \ar[r] & Z' \ar[d]^k \\
U'' \ar[r] \ar[ru]^g & Z''
}
$$
となる。したがって $g$ は写像
$\mathcal{C}_{Z/Z'}|_U \to \mathcal{C}_{Z/Z''}|_U$ を誘導し、これは
$\mathcal{C}_{Z/Z''} \to \mathcal{C}_{Z/Z'}$ の $U$ 上での左逆である。
\end{proof}












\section{ネーター基底上の滑らかさ}
\label{spaces-more-morphisms-section-smooth-Noetherian}

\noindent
この節は More on Morphisms, Section
\ref{more-morphisms-section-smooth-Noetherian} に対応するものである。

\begin{lemma}
\label{spaces-more-morphisms-lemma-lifting-along-artinian-at-point}
$S$ をスキームとする。$f : X \to Y$ を $S$ 上の代数空間の射とし、
$x \in |X|$ とする。$Y$ は局所ネーター的で、$f$ は局所有限型であると仮定する。
次は同値である。
\begin{enumerate}
\item $f$ は $x$ で滑らかである。
\item 任意の実線部分が可換な図式
$$
\xymatrix{
X \ar[d]_f & \Spec(B) \ar[d]^i \ar[l]^-\alpha \\
Y & \Spec(B') \ar[l]_-{\beta} \ar@{-->}[lu]
}
$$
で、$B' \to B$ が局所環の全射、$\Ker(B' \to B)$ の平方が零であり、
$\alpha$ が $\Spec(B)$ の閉点を $x$ に写すものに対して、図式を可換にする
% P09-ERR-1110: frozen source line 5112 lacks punctuation before "there exists".
点線矢印が存在する。
\item (2) と同じ。ただし $B' \to B$ は小拡大全体を動くものとする
（Algebra, Definition \ref{algebra-definition-small-extension} を参照）。
\end{enumerate}
\end{lemma}

\begin{proof}
条件 (1) は、開部分空間 $X' \subset X$ で $X' \to Y$ が滑らかなものが
存在することを意味する。したがって Lemma
\ref{spaces-more-morphisms-lemma-smooth-formally-smooth} により、(1) は条件 (2) と (3) を含意する。
条件 (2) が条件 (3) を含意することは明らかである。(3) を仮定する。
可換図式
$$
\xymatrix{
X \ar[d] & U \ar[l] \ar[d] \\
Y & V \ar[l]
}
$$
で、$U$ と $V$ がアフィン、横の矢印がエタールであり、点 $u \in U$ で
$x$ に写るものが存在する図式を選ぶ。次に、図式
$$
\xymatrix{
X \ar[d] & U \ar[l] \ar[d] & \Spec(B) \ar[d]^i \ar[l]^-\alpha  \\
Y & V \ar[l] & \Spec(B') \ar[l]_-{\beta}
}
$$
を、$u \in U \to V$ に対する (3) のような図式として考える。矢印
$\gamma : \Spec(B') \to X$ を、(3) が $X$ に対して成り立つという仮定から
得られるものとする。$U \to X$ はエタール、したがって形式的エタールなので
（Lemma \ref{spaces-more-morphisms-lemma-etale-formally-etale}）、射 $\gamma$ は $U$ への一意な
持上げで $\alpha$ と両立するものをもつ。さらに $V \to Y$ はエタール、
したがって形式的エタールなので、この持上げは $\beta$ と両立する。
ゆえに (3) は $u \in U \to V$ に対して成り立ち、More on Morphisms, Lemma
\ref{more-morphisms-lemma-lifting-along-artinian-at-point} により
$U \to V$ は $u$ で滑らかであると結論する。これにより $X \to Y$ は
$x$ で滑らかであり、証明が完了する。
\end{proof}

\noindent
小拡大に対する持上げ判定法を、有限型の点を「中心」とする場合だけ確認すれば
よいことが役立つ場合がある。Morphisms of Spaces, Section
\ref{spaces-morphisms-section-points-finite-type} を参照せよ。

\begin{lemma}
\label{spaces-more-morphisms-lemma-lifting-along-artinian}
$S$ をスキームとする。$f : X \to Y$ を $S$ 上の代数空間の射とする。
$Y$ は局所ネーター的で、$f$ は局所有限型であると仮定する。次は同値である。
\begin{enumerate}
\item $f$ は滑らかである。
\item 任意の実線部分が可換な図式
$$
\xymatrix{
X \ar[d]_f & \Spec(B) \ar[d]^i \ar[l]^-\alpha \\
Y & \Spec(B') \ar[l]_-{\beta} \ar@{-->}[lu]
}
$$
で、$B' \to B$ がアルティン局所環の小拡大であり、$\beta$ が有限型（！）で
あるものに対して、図式を可換にする点線矢印が存在する。
\end{enumerate}
\end{lemma}

\begin{proof}
$f$ が滑らかならば、無限小持上げ判定法（Lemma
\ref{spaces-more-morphisms-lemma-smooth-formally-smooth}）により $f$ は形式的に滑らかであり、
(2) が成り立つ。

\medskip\noindent
$f$ が滑らかでないと仮定する。点 $x \in X$ で $f$ が滑らかでないもの全体は、
閉部分集合 $T$ を $|X|$ の中になす。Morphisms of Spaces, Lemma
\ref{spaces-morphisms-lemma-enough-finite-type-points} により、点
$x \in T \subset X$ で $x \in X_{\text{ft-pts}}$ であるものが存在する。
可換図式
$$
\xymatrix{
X \ar[d] & U \ar[l] \ar[d] & u \ar@{|->}[d] \\
Y & V \ar[l] & v
}
$$
で、$U$ と $V$ がアフィン、横の矢印がエタールであり、点 $u \in U$ で
$x$ に写るものが存在する図式を選ぶ。このとき $u$ は $U$ の有限型の点である。
$U \to V$ は点 $u$ で滑らかでないから、More on Morphisms, Lemma
\ref{more-morphisms-lemma-lifting-along-artinian-at-point} により図式
$$
\xymatrix{
X \ar[d] & U \ar[l] \ar[d] & \Spec(B) \ar[d]^i \ar[l]^-\alpha  \\
Y & V \ar[l] & \Spec(B') \ar[l]_-{\beta} \ar@{-->}[lu]
}
$$
で、$B' \to B$ が（アルティン）局所環の小拡大、$B$ の剰余体が
$\kappa(v)$ に等しく、点線矢印が存在しないものがある。$U \to V$ は有限型なので、
$v$ は $V$ の有限型の点である。Morphisms, Lemma
\ref{morphisms-lemma-artinian-finite-type} により、射 $\beta$ は有限型であり、
% P09-ERR-1111: frozen source line 5217 says Spec(B) although beta has domain Spec(B').
したがって合成 $\Spec(B) \to Y$ も有限型である。Lemma
\ref{spaces-more-morphisms-lemma-lifting-along-artinian-at-point} の証明とまったく同じ議論
（$U \to X$ と $V \to Y$ はエタール、したがって形式的エタールであることを
用いる）により、最後に表示した図式の外側の長方形に入る矢印
% P09-ERR-1112: frozen source line 5222 asks for Spec(B) -> X, but that is the given alpha; the dotted lift has domain Spec(B').
$\Spec(B) \to X$ は存在しえない。言い換えれば (2) は成り立たず、証明は完了する。
\end{proof}

\noindent
有用な応用を述べる。

\begin{lemma}
\label{spaces-more-morphisms-lemma-check-smoothness-on-infinitesimal-nbhds}
$S$ をスキームとする。$f : X \to Y$ を $S$ 上の代数空間の射とする。
$f$ は局所有限型で、$Y$ は局所ネーター的であると仮定する。
$Z \subset Y$ を閉部分空間とし、その第 $n$ 無限小近傍を $Z_n \subset Y$ とする。
$X_n = Z_n \times_Y X$ とおく。
\begin{enumerate}
\item $X_n \to Z_n$ がすべての $n$ に対して滑らかならば、$f$ は
$f^{-1}(Z)$ のすべての点で滑らかである。
\item $X_n \to Z_n$ がすべての $n$ に対してエタールならば、$f$ は
$f^{-1}(Z)$ のすべての点でエタールである。
\end{enumerate}
\end{lemma}

\begin{proof}
$X_n \to Z_n$ がすべての $n$ に対して滑らかであると仮定する。
$x \in X$ を $Z$ の点の上にある点とする。Lemma
\ref{spaces-more-morphisms-lemma-lifting-along-artinian-at-point} part (3) にあるような小拡大
$B' \to B$ と射 $\alpha$, $\beta$ が与えられたとする。$B'$ の極大イデアルは
冪零である（$B'$ はアルティン環だから）。したがって、射 $\beta$ は $Z_n$ を
経由し、$\alpha$ は $X_n$ を経由するような適当な $n$ が存在する。ゆえに
$X_n \to Z_n$ の持上げ性質を使えば、図式に必要な点線矢印を得る。
これで (1) が示された。(2) は (1) と、射がエタールであることと相対次元
$0$ の滑らかな射であることが同値であるという事実から従う。
\end{proof}












\section{素朴な余接複体}
\label{spaces-more-morphisms-section-netherlander}

\noindent
この節は Modules on Sites, Section
\ref{sites-modules-section-netherlander} の続きであり、さらにこれは Algebra,
Section \ref{algebra-section-netherlander} の議論の続きである。

\begin{definition}
\label{spaces-more-morphisms-definition-netherlander}
$S$ をスキームとする。$f : X \to Y$ を $S$ 上の代数空間の射とする。
{\it $f$ の素朴な余接複体}とは、環付きトポスの射 $f_{small}$、すなわち
$X$ と $Y$ の小エタールサイトの間の射に対して Modules on Sites, Definition
\ref{sites-modules-definition-cotangent-complex-morphism-ringed-topoi} で定義される
複体である。Properties of Spaces, Lemma
\ref{spaces-properties-lemma-morphism-ringed-topoi} を参照せよ。記法は
$\NL_f$ または $\NL_{X/Y}$ とする。
\end{definition}

\noindent
次の補題群は、この定義が環準同型およびスキームに対する定義と両立し、
$\NL_{X/Y}$ が $D_\QCoh(\mathcal{O}_X)$ の対象であることを示す。

\begin{lemma}
\label{spaces-more-morphisms-lemma-NL-etale-localization}
$S$ をスキームとする。可換図式
$$
\xymatrix{
U \ar[d]_p \ar[r]_g & V \ar[d]^q \\
X \ar[r]^f & Y
}
$$
を $S$ 上の代数空間について考え、$p$ と $q$ はエタールであるとする。
このとき標準的な同一視 $\NL_{X/Y}|_{U_\etale} = \NL_{U/V}$ が
$D(\mathcal{O}_U)$ において存在する。
\end{lemma}

\begin{proof}
素朴な余接複体の構成は引き戻しと可換であり（Modules on Sites, Lemma
\ref{sites-modules-lemma-pullback-NL}）、
$p_{small}^{-1}\mathcal{O}_X = \mathcal{O}_U$ および
$g_{small}^{-1}\mathcal{O}_{V_\etale} =
p_{small}^{-1}f_{small}^{-1}\mathcal{O}_{Y_\etale}$ が成り立つ。後者は
$q_{small}^{-1}\mathcal{O}_{Y_\etale} =
\mathcal{O}_{V_\etale}$ が Properties of Spaces, Lemma
\ref{spaces-properties-lemma-etale-exact-pullback} により成り立つためである。
定義をたどると $\NL_{X/Y}|_{U_\etale} = \NL_{U/V}$ と結論する。
\end{proof}

\begin{lemma}
\label{spaces-more-morphisms-lemma-NL-compare-spaces-schemes}
$S$ をスキームとする。$f : X \to Y$ を $S$ 上の代数空間の射とする。
$X$ と $Y$ はそれぞれスキーム $X_0$ と $Y_0$ によって表現可能であると仮定する。
このとき標準的な同一視 $\NL_{X/Y} = \epsilon^*\NL_{X_0/Y_0}$ が
$D(\mathcal{O}_X)$ において存在する。ここで $\epsilon$ は
Derived Categories of Spaces, Section
\ref{spaces-perfect-section-derived-quasi-coherent-etale} におけるものとし、
$\NL_{X_0/Y_0}$ は More on Morphisms, Definition
\ref{more-morphisms-definition-netherlander} におけるものとする。
\end{lemma}

\begin{proof}
$f_0 : X_0 \to Y_0$ を $f$ に対応するスキームの射とする。標準写像
$\epsilon^{-1}f_0^{-1}\mathcal{O}_{Y_0} \to f_{small}^{-1}\mathcal{O}_Y$
が存在し、これは
$\epsilon^\sharp : \epsilon^{-1}\mathcal{O}_{X_0} \to \mathcal{O}_X$
と両立する。実際、可換図式
$$
\xymatrix{
X_{0, Zar} \ar[d]_{f_0} & X_\etale \ar[l]^\epsilon \ar[d]^f \\
Y_{0, Zar} & Y_\etale \ar[l]_\epsilon
}
$$
がある。Derived Categories of Spaces, Remark
\ref{spaces-perfect-remark-match-total-direct-images} を参照せよ。したがって、
素朴な余接複体の関手性により標準写像
{\small
$$
\epsilon^{-1}\NL_{X_0/Y_0} =
\epsilon^{-1}\NL_{\mathcal{O}_{X_0}/f_0^{-1}\mathcal{O}_{Y_0}} =
\NL_{\epsilon^{-1}\mathcal{O}_{X_0}/\epsilon^{-1}f_0^{-1}\mathcal{O}_{Y_0}}
\to
\NL_{\mathcal{O}_X/f^{-1}_{small}\mathcal{O}_Y} = \NL_{X/Y}
$$
}
を得る。誘導される写像 $\epsilon^*\NL_{X_0/Y_0} \to \NL_{X/Y}$ が
$D(\mathcal{O}_X)$ において同型であることは、幾何学的点における茎で確認して
よい（Properties of Spaces, Theorem
\ref{spaces-properties-theorem-exactness-stalks}）。幾何学的点
$\overline{x} : \Spec(k) \to X_0$ が $x \in X_0$ の上にあり、
% P09-ERR-1113: frozen source line 5366 composes f with a point of X_0 after introducing the corresponding scheme map f_0.
$\overline{y} = f \circ \overline{x}$ が $y \in Y_0$ の上にあるとする。このとき
$$
\NL_{X/Y, \overline{x}} =
\NL_{\mathcal{O}_{X, \overline{x}}/\mathcal{O}_{Y, \overline{y}}}
$$
である。これは $\overline{x}$ における茎を取ることが、
$\overline{x} : \Spec(k) \to X$ による逆像を取ることと同じであり、Modules on
Sites, Lemma \ref{sites-modules-lemma-pullback-NL} を適用できるためである。
一方、
$$
(\epsilon^*\NL_{X_0/Y_0})_{\overline{x}} =
\NL_{X_0/Y_0, x} \otimes_{\mathcal{O}_{X_0, x}}
\mathcal{O}_{X, \overline{x}} =
\NL_{\mathcal{O}_{X_0, x}/\mathcal{O}_{Y_0, y}}
\otimes_{\mathcal{O}_{X_0, x}} \mathcal{O}_{X, \overline{x}}
$$
である。いくつかの詳細は省略する（ヒント：引き戻しの茎が像点における茎で
あることについて Sites, Lemma \ref{sites-lemma-point-morphism-sites} を用い、
加群についての対応する結果として Modules on Sites, Lemma
\ref{sites-modules-lemma-pullback-stalk} も用いる）。
$\mathcal{O}_{X, \overline{x}}$ は $\mathcal{O}_{X_0, x}$ の狭義ヘンゼル化であり、
$\mathcal{O}_{Y, \overline{y}}$ についても同様であることに注意せよ
（Properties of Spaces, Lemma
\ref{spaces-properties-lemma-describe-etale-local-ring}）。したがって結果は
More on Algebra, Lemma \ref{more-algebra-lemma-henselization-NL} から従う。
\end{proof}

\begin{lemma}
\label{spaces-more-morphisms-lemma-netherlander-quasi-coherent}
$S$ をスキームとする。$f : X \to Y$ を $S$ 上の代数空間の射とする。
複体 $\NL_{X/Y}$ のコホモロジー層は準連接であり、次数 $-1$, $0$ 以外では
零で、$\Omega_{X/Y}$ に次数 $0$ で等しい。
\end{lemma}

\begin{proof}
Modules on Sites, Section \ref{sites-modules-section-netherlander} における
素朴な余接複体の構成により、$\NL_{X/Y}$ は次数 $-1$, $0$ に位置する複体であり、
次数 $0$ のコホモロジーは $\Omega_{X/Y}$ である（Section
\ref{spaces-more-morphisms-section-sheaf-differentials} における $\Omega_{X/Y}$ の構成による）。
微分の層は準連接である（Lemma
\ref{spaces-more-morphisms-lemma-module-differentials-quasi-coherent}）。証明を終えるには
$H^{-1}(\NL_{X/Y})$ が準連接であることを示せば十分である。これはエタール
局所的に確認することで従う（Lemma \ref{spaces-more-morphisms-lemma-NL-etale-localization} および
Properties of Spaces, Lemma
\ref{spaces-properties-lemma-characterize-quasi-coherent} により許される）。
まずスキームの場合へ帰着し（Lemma
\ref{spaces-more-morphisms-lemma-NL-compare-spaces-schemes}）、最後にスキームの場合の結果
（More on Morphisms, Lemma
\ref{more-morphisms-lemma-netherlander-quasi-coherent}）を用いる。
\end{proof}

\begin{lemma}
\label{spaces-more-morphisms-lemma-netherlander-fp}
$S$ をスキームとする。$f : X \to Y$ を $S$ 上の代数空間の射とする。
$f$ が局所有限表示ならば、$\NL_{X/Y}$ は $X$ 上エタール局所的に、準同型
$$
\ldots \to 0 \to \mathcal{F}^{-1} \to \mathcal{F}^0 \to 0 \to \ldots
$$
からなる準連接 $\mathcal{O}_X$ 加群の複体と擬同型であり、
$\mathcal{F}^0$ は有限表示、$\mathcal{F}^{-1}$ は有限型である。
\end{lemma}

\begin{proof}
Lemma \ref{spaces-more-morphisms-lemma-NL-etale-localization} により、素朴な余接複体の構成は
エタール局所化と可換である。Lemma
\ref{spaces-more-morphisms-lemma-NL-compare-spaces-schemes} によりスキームの場合へ帰着する。
スキームの場合の結果は More on Morphisms, Lemma
\ref{more-morphisms-lemma-netherlander-fp} である。
\end{proof}

\begin{lemma}
\label{spaces-more-morphisms-lemma-NL-formally-smooth}
$S$ をスキームとする。$f : X \to Y$ を $S$ 上の代数空間の射とする。
次は同値である。
\begin{enumerate}
\item $f$ は形式的に滑らかである。
\item $H^{-1}(\NL_{X/Y}) = 0$ であり、
$H^0(\NL_{X/Y}) = \Omega_{X/Y}$ は局所射影的である。
\end{enumerate}
\end{lemma}

\begin{proof}
Lemma \ref{spaces-more-morphisms-lemma-formally-smooth}, Lemma \ref{spaces-more-morphisms-lemma-NL-etale-localization},
Lemma \ref{spaces-more-morphisms-lemma-NL-compare-spaces-schemes}、およびスキームの場合である
More on Morphisms, Lemma
\ref{more-morphisms-lemma-NL-formally-smooth} から従う。
\end{proof}

\begin{lemma}
\label{spaces-more-morphisms-lemma-NL-formally-etale}
$f : X \to Y$ をスキームの射とする。次は同値である。
\begin{enumerate}
\item $f$ は形式的エタールである。
\item $H^{-1}(\NL_{X/Y}) = H^0(\NL_{X/Y}) = 0$ である。
\end{enumerate}
\end{lemma}

\begin{proof}
(1) を仮定する。形式的エタール射は形式的に滑らかな射である。したがって
$H^{-1}(\NL_{X/Y}) = 0$ である。Lemma
\ref{spaces-more-morphisms-lemma-NL-formally-smooth} を参照せよ。一方、形式的エタール射は
% P09-ERR-1114: frozen source line 5483 says "if formally unramified" instead of "is formally unramified".
形式的非分岐であるから、$\Omega_{X/Y} = 0$ である。Lemma
\ref{spaces-more-morphisms-lemma-characterize-formally-unramified} を参照せよ。逆に (2) が成り立つならば、
$f$ は Lemma \ref{spaces-more-morphisms-lemma-NL-formally-smooth} により形式的に滑らかであり、
Lemma \ref{spaces-more-morphisms-lemma-characterize-formally-unramified} により形式的非分岐である。
したがって Lemmas \ref{spaces-more-morphisms-lemma-formally-etale-unramified-smooth} により形式的
エタールである。
\end{proof}

\begin{lemma}
\label{spaces-more-morphisms-lemma-NL-smooth}
$f : X \to Y$ をスキームの射とする。次は同値である。
\begin{enumerate}
\item $f$ は滑らかである。
\item $f$ は局所有限表示であり、$H^{-1}(\NL_{X/Y}) = 0$ で、
$H^0(\NL_{X/Y}) = \Omega_{X/Y}$ は有限局所自由である。
\end{enumerate}
\end{lemma}

\begin{proof}
Lemma \ref{spaces-more-morphisms-lemma-formally-smooth}, Lemma \ref{spaces-more-morphisms-lemma-NL-etale-localization},
Lemma \ref{spaces-more-morphisms-lemma-NL-compare-spaces-schemes}、およびスキームの場合である
More on Morphisms, Lemma \ref{more-morphisms-lemma-NL-smooth} から従う。
\end{proof}












\section{平坦軌跡の開性}
\label{spaces-more-morphisms-section-open-flat}

\noindent
この節は More on Morphisms, Section
\ref{more-morphisms-section-open-flat} に対応する。代数空間上の準連接加群の
平坦性という概念は Morphisms of Spaces, Section
\ref{spaces-morphisms-section-flat-modules} で定義したことに注意せよ。

\begin{theorem}
\label{spaces-more-morphisms-theorem-openness-flatness}
$S$ をスキームとする。$f : X \to Y$ を $S$ 上の代数空間の射とする。
$\mathcal{F}$ を $X$ 上の準連接層とする。$f$ は局所有限表示であり、
$\mathcal{F}$ は局所有限表示な $\mathcal{O}_X$ 加群であると仮定する。このとき
$$
\{x \in |X| : \mathcal{F}\text{ is flat over }Y\text{ at }x\}
$$
は $|X|$ の開集合である。
\end{theorem}

\begin{proof}
可換図式
$$
\xymatrix{
U \ar[d]_p \ar[r]_\alpha &
V \ar[d]^q \\
X \ar[r]^a & Y
}
$$
で、$U$, $V$ がスキーム、$p$, $q$ が全射かつエタールであるものを選ぶ。
Spaces, Lemma \ref{spaces-lemma-lift-morphism-presentations} を参照せよ。
More on Morphisms, Theorem
\ref{more-morphisms-theorem-openness-flatness} により、集合
$U' = \{u \in |U| : p^*\mathcal{F}\text{ is flat over }V\text{ at }u\}$
は $U$ で開である。Morphisms of Spaces, Definition
\ref{spaces-morphisms-definition-flat-module} により、$U'$ の $|X|$ における像は
定理の集合である。写像 $|U| \to |X|$ は開なので、これで示された。
Properties of Spaces, Lemma
\ref{spaces-properties-lemma-topology-points} を参照せよ。
\end{proof}

\begin{lemma}
\label{spaces-more-morphisms-lemma-flat-locus-base-change}
$S$ をスキームとする。カルテジアン図式
$$
\xymatrix{
X' \ar[r]_{g'} \ar[d]_{f'} & X \ar[d]^f \\
Y' \ar[r]^g & Y
}
$$
を $S$ 上の代数空間について考える。$\mathcal{F}$ を準連接
$\mathcal{O}_X$ 加群とする。
$g$ は平坦、$f$ は局所有限表示、$\mathcal{F}$ は局所有限表示であると仮定する。
このとき
$$
\{x' \in |X'| : (g')^*\mathcal{F}\text{ is flat over }Y'\text{ at }x'\}
$$
は Theorem \ref{spaces-more-morphisms-theorem-openness-flatness} の開部分集合の、連続写像
$|g'| : |X'| \to |X|$ による逆像である。
\end{lemma}

\begin{proof}
Morphisms of Spaces, Lemma
\ref{spaces-morphisms-lemma-base-change-module-flat} から従う。
\end{proof}












\section{ファイバーによる平坦性判定法}
\label{spaces-more-morphisms-section-criterion-flat-fibres}

\noindent
$S$ をスキームとする。$S$ 上の代数空間の可換図式
$$
\xymatrix{
X \ar[rr]_f \ar[dr]_g & & Y \ar[dl]^h \\
& Z
}
$$
と準連接 $\mathcal{O}_X$ 加群 $\mathcal{F}$ を考える。点 $x \in |X|$ が
与えられたとき、$\mathcal{F}$ が $Y$ 上 $x$ で平坦かどうかを問題にする。
$\mathcal{F}$ が $Z$ 上 $x$ で平坦ならば、下の定理はこの問いが、
$\mathcal{F}$ の $X \to Z$ の $g(x)$ 上のファイバーへの制限が
$Y \to Z$ の $g(x)$ 上のファイバー上で平坦かどうかという問いと密接に関係する
ことを述べる。これを意味づけるため、次の準備的な補題を与える。

\begin{lemma}
\label{spaces-more-morphisms-lemma-flat-on-fibres-at-point}
上の状況で次は同値である。
\begin{enumerate}
\item 幾何学的点 $\overline{x}$ を $X$ 上で選び、$x$ の上にあるものとする。
$\overline{y} = f \circ \overline{x}$ および
$\overline{z} = g \circ \overline{x}$ とおく。このとき加群
$\mathcal{F}_{\overline{x}}/
\mathfrak m_{\overline{z}}\mathcal{F}_{\overline{x}}$ は
$\mathcal{O}_{Y, \overline{y}}/
\mathfrak m_{\overline{z}}\mathcal{O}_{Y, \overline{y}}$ 上平坦である。
\item 射 $x : \Spec(K) \to X$ を選び、$x$ の同値類に属するものとする。
$z = g \circ x$, $X_z = \Spec(K) \times_{z, Z} X$,
$Y_z = \Spec(K) \times_{z, Z} Y$ とおき、$\mathcal{F}_z$ を
$\mathcal{F}$ の $X_z$ への引き戻しとする。このとき $\mathcal{F}_z$ は
$x$ で $Y_z$ 上平坦である（Morphisms of Spaces, Definition
\ref{spaces-morphisms-definition-flat-module} の意味で）。
\item 可換図式
$$
\xymatrix{
& & & U \ar[llld]_a \ar[rr] \ar[dr] & & V \ar[llld]_>>>>>>>b \ar[dl] \\
X \ar[rr]_f \ar[dr]_g & & Y \ar[dl]^h &  & W \ar[llld]_c \\
& Z
}
$$
を選ぶ。ここで $U, V, W$ はスキーム、$a, b, c$ はエタールで、点
$u \in U$ は $x$ に写るものとする。$w \in W$ を $u$ の像とする。
$\mathcal{F}_w$ を $\mathcal{F}$ のファイバー $U_w$ への引き戻しとする。
このファイバーは $U \to W$ の $w$ におけるものである。このとき $\mathcal{F}_w$ は
$V_w$ 上 $u$ で平坦である。
\end{enumerate}
\end{lemma}

\begin{proof}
(2) において射 $x : \Spec(K) \to X$ は $K$ 有理点を $X_z$ 上に定めるので、
主張は意味をもつことに注意せよ。さらに (2) の条件は、$\Spec(K) \to X$ の
選び方で $x$ の同値類に属するものの間では変わらない（詳細は省略する。ほかの条件はこの選択に
依存しないので、これは下の議論からも従う）。また (3) のような図式は常に次の
ように選べる。まずスキーム $W$ と全射エタール射 $W \to Z$ を選び、次に
スキーム $V$ と全射エタール射 $V \to W \times_Z Y$ を選び、最後に
スキーム $U$ と全射エタール射 $U \to V \times_Y X$ を選ぶ。これらを選んだ後、
$U \to W$ を合成 $U \to V \to W$ と等しくおき、点 $u \in U$ で $x$ に写るものを
選べる。実際、射 $U \to X$ は全射である。

\medskip\noindent
(3) のような図式と (1) のような幾何学的点
$\overline{x} : \Spec(k) \to X$ の双方が与えられたとする。Properties of Spaces,
Lemma \ref{spaces-properties-lemma-geometric-lift-to-usual} により、幾何学的点
$\overline{u} : \Spec(k) \to U$ を $u$ の上で選び、
$\overline{x} = a \circ \overline{u}$ を満たすようにできる。
$\overline{v} : \Spec(k) \to V$ および
$\overline{w} : \Spec(k) \to W$ を、それぞれ $V$ および $W$ の誘導される
幾何学的点とする。この状況では
$\mathcal{O}_{X, \overline{x}} = \mathcal{O}_{U, u}^{sh}$ であり、$Y$ と $Z$ に
ついても同様である。Properties of Spaces, Lemma
\ref{spaces-properties-lemma-describe-etale-local-ring} を参照せよ。同様に
$$
\mathcal{F}_{\overline{x}} =
(a^*\mathcal{F})_u \otimes_{\mathcal{O}_{U, u}}
\mathcal{O}_{U, u}^{sh}
$$
である。Properties of Spaces, Lemma
\ref{spaces-properties-lemma-stalk-quasi-coherent} を参照せよ。
$\mathcal{F}_w$ の $u$ における茎は
$$
(\mathcal{F}_w)_u = (a^*\mathcal{F})_u/\mathfrak m_w(a^*\mathcal{F})_u
$$
で与えられ、$V_w$ の $v$ における局所環は
$$
\mathcal{O}_{V_w, v} = \mathcal{O}_{V, v}/\mathfrak m_w\mathcal{O}_{V, v}.
$$
で与えられる。また
$\mathfrak m_{\overline{z}} =
\mathfrak m_w \mathcal{O}_{Z, \overline{z}} =
\mathfrak m_w \mathcal{O}_{W, w}^{sh}$
なので、
\begin{align*}
\mathcal{F}_{\overline{x}}/
\mathfrak m_{\overline{z}}\mathcal{F}_{\overline{x}} & =
(a^*\mathcal{F})_u \otimes_{\mathcal{O}_{U, u}}
\mathcal{O}_{X, \overline{x}}/
\mathfrak m_{\overline{z}}\mathcal{O}_{X, \overline{x}} \\
& =
(\mathcal{F}_w)_u \otimes_{\mathcal{O}_{U_w, u}}
\mathcal{O}_{U, u}^{sh}/\mathfrak m_w\mathcal{O}_{U, u}^{sh} \\
& = (\mathcal{F}_w)_u \otimes_{\mathcal{O}_{U_w, u}}
\mathcal{O}_{U_w, \overline{u}}^{sh} \\
& = (\mathcal{F}_w)_{\overline{u}}
\end{align*}
を得る。最後から二番目の等号は Algebra, Lemma
\ref{algebra-lemma-quotient-strict-henselization} により、最後の等号は
Properties of Spaces, Lemma
\ref{spaces-properties-lemma-stalk-quasi-coherent} による。$V$ と $Y$ の構造層に
同じ議論を適用すると
$$
\mathcal{O}_{V_w, \overline{v}}^{sh} =
\mathcal{O}_{V, v}^{sh}/\mathfrak m_w \mathcal{O}_{V, v}^{sh} =
\mathcal{O}_{Y, \overline{y}}/
\mathfrak m_{\overline{z}}\mathcal{O}_{Y, \overline{y}}.
$$
となる。ここで Morphisms of Spaces, Lemma
\ref{spaces-morphisms-lemma-flat-at-point} を用いると、(1) と (3) が同値であると
分かる。

\medskip\noindent
最後に (2) と (3) の同値性を証明する。そのため体拡大 $\tilde K$ を $K$ 上で取り、
射 $\tilde x : \Spec(\tilde K) \to U$ で $u$ の上にあるものを選ぶ
（$u \times_{X, x} \Spec(K)$ は空でないスキームなので可能である）。
% P09-ERR-1116: frozen source line 5748 says "Set ... be the composition" instead of "Let ... be the composition".
$\tilde z : \Spec(\tilde K) \to U \to W$ をその合成とする。可換図式
$$
\xymatrix{
& & & U_w \times_w \tilde z \ar[llld]_a \ar[rr] \ar[dr] & &
V_w \times_w \tilde z \ar[llld]_>>>>>>>b \ar[dl] \\
X_z \ar[rr]_f \ar[dr]_g & & Y_z \ar[dl]^h &  & \tilde z \ar[llld]_c \\
& z
}
$$
を得る。ここで $z = \Spec(K)$ および $w = \Spec(\kappa(w))$ である。
$\mathcal{F}_w$ と $\mathcal{F}_z$ は $U_w \times_w \tilde z$ 上の同じ加群へ
引き戻されることが明らかである。これにより可換図式
$$
\xymatrix{
X_z \ar[d] & U_w \times_w \tilde z \ar[l] \ar[d] \ar[r] & U_w \ar[d] \\
Y_z & V_w \times_w \tilde z \ar[l] \ar[r] & V_w
}
$$
を得る。その二つの正方形はいずれもカルテジアンで、下の横矢印はいずれも
平坦である。左下の横矢印は、射
$Y \times_Z \tilde z \to Y \times_Z z = Y_z$（平坦射の基底変換）、エタール射
$V \times_Z \tilde z \to Y \times_Z \tilde z$、およびエタール射
$V \times_W \tilde z \to V \times_Z \tilde z$ の合成である。したがって
Morphisms of Spaces, Lemma
\ref{spaces-morphisms-lemma-base-change-module-flat} により
$$
\mathcal{F}_z\text{ flat at }x\text{ over }Y_z
\Leftrightarrow
\mathcal{F}|_{U_w \times_w \tilde z}
\text{ flat at }\tilde x\text{ over }V_w \times_w \tilde z
\Leftrightarrow
\mathcal{F}_w\text{ flat at }u\text{ over }V_w
$$
となり、証明が完了する。
\end{proof}

\begin{definition}
\label{spaces-more-morphisms-definition-module-flat-on-fibre}
$S$ をスキームとする。$X \to Y \to Z$ を $S$ 上の代数空間の射とする。
$\mathcal{F}$ を準連接 $\mathcal{O}_X$ 加群とする。$x \in |X|$ を点とし、
その像を $z \in |Z|$ と書く。
\begin{enumerate}
\item Lemma \ref{spaces-more-morphisms-lemma-flat-on-fibres-at-point} の同値な条件が満たされるとき、
{\it $\mathcal{F}$ の $z$ 上のファイバーへの制限は、$x$ で、$Y$ の $z$ 上の
ファイバー上平坦である}という。
\item {\it $X$ の $z$ 上のファイバーは、$x$ で、$Y$ の $z$ 上の
ファイバー上平坦である}とは、Lemma
\ref{spaces-more-morphisms-lemma-flat-on-fibres-at-point} の同値な条件が
$\mathcal{F} = \mathcal{O}_X$ に対して成り立つことをいう。
% P09-ERR-1117: frozen source line 5805 lacks terminal punctuation at the end of item (3).
\item {\it $X$ の $z$ 上のファイバーは $Y$ の $z$ 上のファイバー上平坦である}
とは、すべての $x \in |X|$ で $z$ の上にあるものに対して、$X$ の $z$ 上の
ファイバーが $x$ で $Y$ の $z$ 上のファイバー上平坦であることをいう
\end{enumerate}
\end{definition}

\noindent
この定義を用いると判定法を次のように述べられる。ネーター版は Section
\ref{spaces-more-morphisms-section-flat-over-Noetherian} にある。

\begin{theorem}
\label{spaces-more-morphisms-theorem-criterion-flatness-fibre}
$S$ をスキームとする。$f : X \to Y$ および $Y \to Z$ を $S$ 上の
代数空間の射とする。$\mathcal{F}$ を準連接 $\mathcal{O}_X$ 加群とする。
次を仮定する。
\begin{enumerate}
\item $X$ は $Z$ 上局所有限表示である。
% P09-ERR-1118: frozen source line 5822 omits "is" in the second assumption.
\item $\mathcal{F}$ は有限表示な $\mathcal{O}_X$ 加群である。
\item $Y$ は $Z$ 上局所有限型である。
\end{enumerate}
$x \in |X|$ とし、$y \in |Y|$ および $z \in |Z|$ を $x$ の像とする。
$\mathcal{F}_{\overline{x}} \not = 0$ ならば、次は同値である。
\begin{enumerate}
\item $\mathcal{F}$ は $Z$ 上 $x$ で平坦であり、$\mathcal{F}$ の $z$ 上の
ファイバーへの制限は、$x$ で $Y$ の $z$ 上のファイバー上平坦である。
\item $Y$ は $Z$ 上 $y$ で平坦であり、$\mathcal{F}$ は $Y$ 上 $x$ で平坦である。
\end{enumerate}
さらに、(1) と (2) が成り立つ点 $x$ の集合は
$\text{Supp}(\mathcal{F})$ で開である。
\end{theorem}

\begin{proof}
Lemma \ref{spaces-more-morphisms-lemma-flat-on-fibres-at-point} part (3) のような図式を選ぶ。
定義から、これはスキームの射 $U \to V \to W$、準連接層
$a^*\mathcal{F}$、および点 $u \in U$ に対する対応する定理へ帰着する。
したがって定理はスキームの場合の対応する結果、すなわち More on Morphisms,
Theorem \ref{more-morphisms-theorem-criterion-flatness-fibre} から従う。
\end{proof}

\begin{lemma}
\label{spaces-more-morphisms-lemma-morphism-between-flat}
$S$ をスキームとする。$f : X \to Y$ および $Y \to Z$ を $S$ 上の
代数空間の射とする。次を仮定する。
\begin{enumerate}
\item $X$ は $Z$ 上局所有限表示である。
\item $X$ は $Z$ 上平坦である。
\item すべての $z \in |Z|$ に対して、$X$ の $z$ 上のファイバーは
$Y$ の $z$ 上のファイバー上平坦である。
\item $Y$ は $Z$ 上局所有限型である。
\end{enumerate}
このとき $f$ は平坦である。さらに $f$ が全射ならば、$Y$ は $Z$ 上平坦である。
\end{lemma}

\begin{proof}
Theorem \ref{spaces-more-morphisms-theorem-criterion-flatness-fibre} の特別な場合である。
\end{proof}

\begin{lemma}
\label{spaces-more-morphisms-lemma-base-change-criterion-flatness-fibre}
$S$ をスキームとする。$f : X \to Y$ および $Y \to Z$ を $S$ 上の
代数空間の射とする。$\mathcal{F}$ を準連接 $\mathcal{O}_X$ 加群とする。
次を仮定する。
\begin{enumerate}
\item $X$ は $Z$ 上局所有限表示である。
% P09-ERR-1119: frozen source line 5879 omits "is" in the second assumption.
\item $\mathcal{F}$ は有限表示な $\mathcal{O}_X$ 加群である。
\item $\mathcal{F}$ は $Z$ 上平坦である。
\item $Y$ は $Z$ 上局所有限型である。
\end{enumerate}
このとき集合
$$
A = \{x \in |X| : \mathcal{F} \text{ flat at }x \text{ over }Y\}.
$$
は $|X|$ で開であり、その構成は任意の基底変換と可換である。すなわち、
$Z' \to Z$ が代数空間の射であり、$A'$ が
$X' = X \times_Z Z'$ の点で
$\mathcal{F}' = \mathcal{F} \times_Z Z'$ が
$Y' = Y \times_Z Z'$ 上平坦であるもの全体ならば、$A'$ は $A$ の連続写像
$|X'| \to |X|$ による逆像である。
\end{lemma}

\begin{proof}
一つの証明法は More on Morphisms, Lemma
\ref{more-morphisms-lemma-morphism-between-flat} の証明を代数空間の圏へ
移すことである。ここでは代わりにスキームの場合へ帰着して証明する。すなわち
Lemma \ref{spaces-more-morphisms-lemma-flat-on-fibres-at-point} part (3) のような図式で、
$a$, $b$, $c$ が全射であるものを選ぶ。定義から、これはスキームの射
$U \to V \to W$、準連接層 $a^*\mathcal{F}$、および点 $u \in U$ に対する
対応する定理へ帰着する。注意すべき小さな点は、代数空間の射 $Z' \to Z$ が
与えられたとき、スキーム $W'$ と全射エタール射
$W' \to W \times_Z Z'$ を選ぶことである。次に
$U' = W' \times_W U$ および $V' = W' \times_W V$ とおく。
$a', b', c'$ を $U', V', W'$ から $X', Y', Z'$ への射と書く。この場合、
$A$ と $A'$ はそれぞれ、$U$ と $U'$ の開部分集合で、
$a^*\mathcal{F}$ と $(a')^*\mathcal{F}'$ に付随するものの像である。
これで実際に補題は More on Morphisms,
Lemma \ref{more-morphisms-lemma-morphism-between-flat} へ帰着する。
\end{proof}

\begin{lemma}
\label{spaces-more-morphisms-lemma-base-change-flatness-fibres}
$S$ をスキームとする。$f : X \to Y$ および $Y \to Z$ を $S$ 上の
代数空間の射とする。次を仮定する。
\begin{enumerate}
\item $X$ は $Z$ 上局所有限表示である。
\item $X$ は $Z$ 上平坦である。
\item $Y$ は $Z$ 上局所有限型である。
\end{enumerate}
このとき集合
$$
\{x \in |X| : X\text{ flat at }x \text{ over }Y\}.
$$
は $|X|$ で開であり、その構成は任意の基底変換 $Z' \to Z$ と可換である。
\end{lemma}

\begin{proof}
Lemma \ref{spaces-more-morphisms-lemma-base-change-criterion-flatness-fibre} の特別な場合である。
\end{proof}

\begin{lemma}
\label{spaces-more-morphisms-lemma-flat-and-free-at-point-fibre}
$S$ をスキームとする。$f : X \to Y$ を $S$ 上の代数空間の局所有限表示な射とする。
$\mathcal{F}$ を有限表示な $\mathcal{O}_X$ 加群とする。$x \in |X|$ とし、
その像を $y \in |Y|$ とする。$\mathcal{F}$ が $x$ で $Y$ 上平坦ならば、
次は同値である。
\begin{enumerate}
\item $(\mathcal{F}_{\overline{y}})_{\overline{x}}$ は平坦な
$\mathcal{O}_{X_{\overline{y}}, \overline{x}}$ 加群である。
\item $(\mathcal{F}_{\overline{y}})_{\overline{x}}$ は自由な
$\mathcal{O}_{X_{\overline{y}}, \overline{x}}$ 加群である。
\item $\mathcal{F}_{\overline{y}}$ は $\overline{x}$ の、
$X_{\overline{y}}$ におけるエタール近傍上で有限自由である。
\item $\mathcal{F}$ は $x$ の、$X$ におけるエタール近傍上で有限自由である。
\end{enumerate}
ここで $\overline{x}$ は $X$ の幾何学的点で $x$ の上にあり、
$\overline{y} = f \circ \overline{x}$ である。
\end{lemma}

\begin{proof}
可換図式
$$
\xymatrix{
U \ar[d] \ar[r] & V \ar[d] \\
X \ar[r] & Y
}
$$
で、$U$ と $V$ がスキーム、縦の矢印がエタールであり、点 $u \in U$ で
$x$ に写るものが存在する図式を選ぶ。$v \in V$ を $u$ の像とする。
Lemma \ref{spaces-more-morphisms-lemma-flat-on-fibres-at-point} を
$\text{id} : X \to X$ に $Y$ 上で適用すると、(1) は「$\mathcal{F}|_{U_v}$ が
$U_v$ 上 $u$ で平坦である」という条件へ移される。言い換えると、(1) は
$(\mathcal{F}|_{U_v})_u$ が平坦な $\mathcal{O}_{U_v, u}$ 加群であることと
同値である。スキームの場合（More on Morphisms, Lemma
\ref{more-morphisms-lemma-flat-and-free-at-point-fibre}）により、これは
$\mathcal{F}|_U$ が $u$ のある開近傍上で有限自由であることを含意する。
これで (1) が (4) を含意することが分かる。(4) $\Rightarrow$ (3) および
(2) $\Rightarrow$ (1) は明らかである。(3) $\Rightarrow$ (2) には、
Properties of Spaces, Lemmas
\ref{spaces-properties-lemma-describe-etale-local-ring} and
\ref{spaces-properties-lemma-stalk-quasi-coherent} における局所環と茎の記述を用いる。
\end{proof}

\begin{lemma}
\label{spaces-more-morphisms-lemma-finite-free-open}
$S$ をスキームとする。$f : X \to Y$ を $S$ 上の代数空間の局所有限表示な射とする。
$\mathcal{F}$ を有限表示な $\mathcal{O}_X$ 加群で $Y$ 上平坦なものとする。
このとき集合
$$
\{x \in |X| : \mathcal{F}\text{ free in an \'etale neighbourhood of }x\}
$$
は $|X|$ で開であり、その構成は任意の基底変換 $Y' \to Y$ と可換である。
\end{lemma}

\begin{proof}
開性は明らかである。$Y' \to Y$ を代数空間の射とし、
$X' = Y' \times_Y X$ とおき、$x' \in |X'|$ を $x \in |X|$ の上の点とする。
Lemma \ref{spaces-more-morphisms-lemma-flat-and-free-at-point-fibre} により、$x$ が上の集合に属することと
$(\mathcal{F}_{\overline{y}})_{\overline{x}}$ が平坦な
$\mathcal{O}_{X_{\overline{y}}, \overline{x}}$ 加群であることは同値である。
同様に、$x'$ が引き戻し $\mathcal{F}'$、すなわち $\mathcal{F}$ の $X'$ への
引き戻しに対する対応する集合に属することと
$(\mathcal{F}'_{\overline{y}'})_{\overline{x}'}$ が平坦な
$\mathcal{O}_{X'_{\overline{y}'}, \overline{x}'}$ 加群であることは同値である
（記法は明らかである）。この二つの主張は、Lemma
\ref{spaces-more-morphisms-lemma-flat-on-fibres-at-point} を射
$\text{id} : X \to X$ に $Y$ 上で適用することにより同値である。したがって基底変換に
関する主張が成り立つ。
\end{proof}







\section{ネーター的基底上の平坦性}
\label{spaces-more-morphisms-section-flat-over-Noetherian}

\noindent
ネーターの場合の「ファイバーによる平坦性判定法」を述べる。

\begin{theorem}
\label{spaces-more-morphisms-theorem-criterion-flatness-fibre-Noetherian}
$S$ をスキームとする。
$f : X \to Y$ および $Y \to Z$ を $S$ 上の代数空間の射とする。
$\mathcal{F}$ を準連接 $\mathcal{O}_X$ 加群とする。次を仮定する：
\begin{enumerate}
% P09-ERR-1120: frozen source line 6038 omits "are" in the first assumption.
\item $X$, $Y$, $Z$ は局所ネーター的である。
% P09-ERR-1121: frozen source line 6039 omits "is" in the second assumption.
\item $\mathcal{F}$ は連接 $\mathcal{O}_X$ 加群である。
\end{enumerate}
$x \in |X|$ とし、$y \in |Y|$ および $z \in |Z|$ を $x$ の像とする。
$\mathcal{F}_{\overline{x}} \not = 0$ ならば、次は同値である：
\begin{enumerate}
\item $\mathcal{F}$ は $Z$ 上 $x$ で平坦であり、かつ
$\mathcal{F}$ の $z$ 上のファイバーへの制限は $x$ で、$Y$ の $z$ 上の
ファイバーに対して平坦である。
\item $Y$ は $Z$ 上 $y$ で平坦であり、かつ $\mathcal{F}$ は
$Y$ 上 $x$ で平坦である。
\end{enumerate}
\end{theorem}

\begin{proof}
Lemma \ref{spaces-more-morphisms-lemma-flat-on-fibres-at-point} の (3) にある図を選ぶ。
定義から、これはスキームの射 $U \to V \to W$、準連接層
$a^*\mathcal{F}$、および点 $u \in U$ に対する対応する定理へ帰着する。
したがって、スキームに対する対応する結果である More on Morphisms,
Theorem \ref{more-morphisms-theorem-criterion-flatness-fibre-Noetherian}
から定理が従う。
\end{proof}

\begin{lemma}
\label{spaces-more-morphisms-lemma-morphism-between-flat-Noetherian}
$S$ をスキームとする。
$f : X \to Y$ および $Y \to Z$ を $S$ 上の代数空間の射とする。
次を仮定する：
\begin{enumerate}
\item $X$, $Y$, $Z$ は局所ネーター的である。
\item $X$ は $Z$ 上平坦である。
\item 任意の $z \in |Z|$ に対し、$X$ の $z$ 上のファイバーは
$Y$ の $z$ 上のファイバーに対して平坦である。
\end{enumerate}
このとき $f$ は平坦である。さらに $f$ が全射ならば、$Y$ は $Z$ 上平坦である。
\end{lemma}

\begin{proof}
これは Theorem \ref{spaces-more-morphisms-theorem-criterion-flatness-fibre-Noetherian} の特別な場合である。
\end{proof}

\noindent
滑らかさを調べる場合とまったく同様に、基底がネーター的ならば、アルティン環上で
平坦性を調べれば十分である。その一例を述べる。

\begin{lemma}
\label{spaces-more-morphisms-lemma-flatness-over-Noetherian-ring}
$A$ をネーター環とし、$I \subset A$ をイデアルとする。
$X$ を $S = \Spec(A)$ 上局所有限表示な代数空間とする。
$n \geq 1$ に対して $S_n = \Spec(A/I^n)$ および
$X_n = S_n \times_S X$ とおく。$\mathcal{F}$ を連接 $\mathcal{O}_X$ 加群とする。
% P09-ERR-1122: frozen source line 6094 says the pullback is to X, not X_n.
任意の $n \geq 1$ に対して、引き戻し $\mathcal{F}_n$、すなわち
$\mathcal{F}$ の $X$ への引き戻しが $S_n$ 上平坦ならば、
% P09-ERR-1123: frozen source lines 6095--6096 say flat over X, while the proof says over S.
$\mathcal{F}$ が $X$ 上平坦である（開）軌跡は、$V(I)$ の
$X \to S$ による逆像を含む。
\end{lemma}

\begin{proof}
$\mathcal{F}$ が $S$ 上平坦である軌跡は、Theorem
\ref{spaces-more-morphisms-theorem-openness-flatness} により $|X|$ で開である。
主張は $X$ をエタール被覆の各要素で置き換えても変わらないので、
$X$ がアフィンスキームであると仮定してよい。この場合、結果は
Algebra, Lemma \ref{algebra-lemma-flat-module-powers} から直ちに従う。
いくつかの詳細は省略する。
\end{proof}







\section{正規化再訪}
\label{spaces-more-morphisms-section-normalization}

\noindent
正規化は滑らかな基底変換と可換である。

\begin{lemma}
\label{spaces-more-morphisms-lemma-integral-closure-smooth-pullback}
$S$ をスキームとする。$f : Y \to X$ を $S$ 上の代数空間の滑らかな射とする。
$\mathcal{A}$ を準連接 $\mathcal{O}_X$ 代数の層とする。
$\mathcal{O}_Y$ の $f^*\mathcal{A}$ における整閉包は
$f^*\mathcal{A}'$ に等しい。ここで
$\mathcal{A}' \subset \mathcal{A}$ は $\mathcal{O}_X$ の
$\mathcal{A}$ における整閉包である。
\end{lemma}

\begin{proof}
整閉包の構成（Morphisms of Spaces, Definition
\ref{spaces-morphisms-definition-integral-closure} を参照）により、これは直ちに
$X$ と $Y$ がアフィンである場合へ帰着する。この場合、結果は
Algebra, Lemma \ref{algebra-lemma-integral-closure-commutes-smooth} である。
\end{proof}

\begin{lemma}[正規化は滑らかな基底変換と可換である]
\label{spaces-more-morphisms-lemma-normalization-smooth-localization}
$S$ をスキームとする。
$$
\xymatrix{
Y_2 \ar[r] \ar[d] & Y_1 \ar[d]^f \\
X_2 \ar[r]^\varphi & X_1
}
$$
を $S$ 上の代数空間のファイバー積の図式とする。$f$ は準コンパクトかつ
準分離的であり、$\varphi$ は滑らかであると仮定する。
$Y_i \to X_i' \to X_i$ を $X_i$ の $Y_i$ における正規化とする。
このとき $X_2' \cong X_2 \times_{X_1} X_1'$ である。
\end{lemma}

\begin{proof}
分解 $Y_1 \to X_1' \to X_1$ の $X_2$ への基底変換は
% P09-ERR-1124: frozen source lines 6157--6158 give X_1 as the terminal target after base change.
$Y_2 \to X_2 \times_{X_1} X_1' \to X_1$ という分解であり、
$X_2 \times_{X_1} X_1' \to X_1$ は整である
（Morphisms of Spaces, Lemma \ref{spaces-morphisms-lemma-base-change-integral}）。
したがって Morphisms of Spaces, Lemma
\ref{spaces-morphisms-lemma-characterize-normalization} の普遍性により射
$h : X_2' \to X_2 \times_{X_1} X_1'$ を得る。
$X_2'$ は $\mathcal{O}_{X_2}$ の
$f_{2, *}\mathcal{O}_{Y_2}$ における整閉包の相対スペクトルであることに注意する。
% P09-ERR-1125: frozen source lines 6166--6167 call A' the integral closure of O_{X_2} inside an O_{X_1}-algebra.
$\mathcal{A}' \subset f_{1, *}\mathcal{O}_{Y_1}$ が
$\mathcal{O}_{X_2}$ の整閉包を表すとすると、相対スペクトルの構成は任意の
基底変換と可換なので、$X_2 \times_{X_1} X_1'$ は
$\varphi^*\mathcal{A}'$ の相対スペクトルである。Cohomology of Spaces, Lemma
\ref{spaces-cohomology-lemma-flat-base-change-cohomology} により
$f_{2, *}\mathcal{O}_{Y_2} = \varphi^*f_{1, *}\mathcal{O}_{Y_1}$ である。
したがって結果は Lemma \ref{spaces-more-morphisms-lemma-integral-closure-smooth-pullback} から従う。
\end{proof}









\section{Cohen--Macaulay 射}
\label{spaces-more-morphisms-section-CM}

\noindent
これは More on Morphisms, Section
\ref{more-morphisms-section-CM} の類似である。

\begin{lemma}
\label{spaces-more-morphisms-lemma-CM-local-ring-fibre}
スキームの芽の射に関する性質
\begin{align*}
& \mathcal{P}((X, x) \to (S, s)) = \\
& \text{ファイバーの局所環 }
\mathcal{O}_{X_s, x}
\text{ がネーター的かつ Cohen--Macaulay である}
\end{align*}
は始域と終域の上でエタール局所的である（Descent, Definition
\ref{descent-definition-local-source-target-at-point}）。
\end{lemma}

\begin{proof}
Descent, Definition \ref{descent-definition-local-source-target-at-point}
にある図式が与えられると、ファイバーのエタール射
$U'_{v'} \to U_v$ で $u'$ を $u$ へ写すものを得る。Descent, Lemma
\ref{descent-lemma-etale-on-fiber} を参照せよ。
したがって局所環 $\mathcal{O}_{U'_{v'}, u'}$ と
$\mathcal{O}_{U_v, u}$ の狭義ヘンゼル化は同じである。
More on Algebra, Lemma \ref{more-algebra-lemma-henselization-CM} により結論を得る。
\end{proof}

\begin{definition}
\label{spaces-more-morphisms-definition-CM}
$S$ をスキームとする。
$f : X \to Y$ を $S$ 上の代数空間の射とする。
$f$ のファイバーは局所ネーター的であると仮定する
（Divisors on Spaces, Definition
\ref{spaces-divisors-definition-locally-Noetherian-fibre}）。
\begin{enumerate}
\item $x \in |X|$ とし、$y = f(x)$ とする。$f$ が
\textit{$x$ で Cohen--Macaulay} であるとは、$f$ が $x$ で平坦であり、
Morphisms of Spaces, Lemma
\ref{spaces-morphisms-lemma-local-source-target-at-point} の同値な条件が
Lemma \ref{spaces-more-morphisms-lemma-CM-local-ring-fibre} に記述された性質 $\mathcal{P}$ に対して
成り立つことをいう。
\item $f$ が \textit{Cohen--Macaulay 射}であるとは、$f$ が
$X$ のすべての点で Cohen--Macaulay であることをいう。
\end{enumerate}
\end{definition}

\noindent
次はスキームの場合からの移し替えである。

\begin{lemma}
\label{spaces-more-morphisms-lemma-CM}
$S$ をスキームとする。$f : X \to Y$ を $S$ 上の代数空間の射とする。
$f$ のファイバーは局所ネーター的であると仮定する。次は同値である：
\begin{enumerate}
\item $f$ は Cohen--Macaulay である。
\item $f$ は平坦であり、ある全射エタール射 $V \to Y$ で $V$ が
スキームであるものに対し、$X_V \to V$ のファイバーは
Cohen--Macaulay 代数空間である。
\item $f$ は平坦であり、任意のエタール射 $V \to Y$ で $V$ が
スキームであるものに対し、$X_V \to V$ のファイバーは
Cohen--Macaulay 代数空間である。
\end{enumerate}
$x \in |X|$ を、像 $y \in |Y|$ をもつ点とする。このとき次は同値である：
\begin{enumerate}
\item[(a)] $f$ は $x$ で Cohen--Macaulay である。
\item[(b)] $\mathcal{O}_{Y, \overline{y}} \to \mathcal{O}_{X, \overline{x}}$
は平坦であり、
$\mathcal{O}_{X, \overline{x}}/
\mathfrak m_{\overline{y}}\mathcal{O}_{X, \overline{x}}$ は Cohen--Macaulay である。
\end{enumerate}
\end{lemma}

\begin{proof}
$V \to Y$ を $V$ がスキームであるようなエタール射とし、スキーム $U$ と
全射エタール射 $U \to X \times_Y V$ を選ぶ。可換図式
$$
\xymatrix{
U \ar[d] \ar[r] & V \ar[d] \\
X \ar[r] & Y
}
$$
を考える。$u \in U$ とし、その像をそれぞれ $x \in |X|$、
$y \in |Y|$、$v \in V$ とする。このとき $f$ が $x$ で
Cohen--Macaulay であることと $U \to V$ が $u$ で Cohen--Macaulay
であることは（定義により）同値である。さらに射
$U_v \to X_v = (X_V)_v$ は全射エタールである。したがってスキーム
$U_v$ が Cohen--Macaulay であることと、代数空間 $X_v$ が
Cohen--Macaulay であることは同値である。ゆえに (1), (2), (3) の同値性は、
スキームの射に対する対応する同値性 More on Morphisms, Lemma
\ref{more-morphisms-lemma-CM} から形式的な議論により従う。

\medskip\noindent
(a) と (b) の同値性を証明する。平坦性について対応する同値性は
Morphisms of Spaces, Lemma
\ref{spaces-morphisms-lemma-flat-at-point-etale-local-rings} である。
したがって同値性を証明する際には $f$ が $x$ で平坦であると仮定してよい。
上のような図式と $x, y, u, v$ を考える。このとき
$\mathcal{O}_{Y, \overline{y}} \to \mathcal{O}_{X, \overline{x}}$ は、
局所環の狭義ヘンゼル化の間の写像
$\mathcal{O}_{V, v}^{sh} \to \mathcal{O}_{U, u}^{sh}$ に等しい。
Properties of Spaces, Lemma
\ref{spaces-properties-lemma-describe-etale-local-ring} を参照せよ。
したがって Algebra, Lemma
\ref{algebra-lemma-quotient-strict-henselization} により
$$
\mathcal{O}_{X, \overline{x}}/
\mathfrak m_{\overline{y}}\mathcal{O}_{X, \overline{x}} =
(\mathcal{O}_{U, u}/\mathfrak m_v \mathcal{O}_{U, u})^{sh}
$$
である。ゆえに、ネーター局所環
$\mathcal{O}_{U, u}/\mathfrak m_v \mathcal{O}_{U, u}$ が
Cohen--Macaulay であることと、その狭義ヘンゼル化が
Cohen--Macaulay であることが同値であると示せばよい。これは
More on Algebra, Lemma \ref{more-algebra-lemma-henselization-CM} である。
\end{proof}

\begin{lemma}
\label{spaces-more-morphisms-lemma-composition-CM}
$S$ をスキームとする。$f : X \to Y$ および $g : Y \to Z$ を
$S$ 上の代数空間の射とする。$f$, $g$, $g \circ f$ のファイバーは
局所ネーター的であると仮定する。$x \in |X|$ とし、その像を
$y \in |Y|$ および $z \in |Z|$ とする。
\begin{enumerate}
\item $f$ が $x$ で Cohen--Macaulay であり、$g$ が $f(x)$ で
Cohen--Macaulay ならば、$g \circ f$ は $x$ で Cohen--Macaulay である。
\item $f$ と $g$ が Cohen--Macaulay ならば、$g \circ f$ は
Cohen--Macaulay である。
\item $g \circ f$ が $x$ で Cohen--Macaulay であり、$f$ が $x$ で
平坦ならば、$f$ は $x$ で Cohen--Macaulay であり、$g$ は $f(x)$ で
Cohen--Macaulay である。
% P09-ERR-1126: frozen source line 6323 uses the type-invalid composite f \circ g.
\item $f \circ g$ が Cohen--Macaulay であり、$f$ が平坦ならば、
$f$ は Cohen--Macaulay であり、$g$ は $f$ の像の各点で
Cohen--Macaulay である。
\end{enumerate}
\end{lemma}

\begin{proof}
エタール局所的に作業すれば、これはスキームに対する対応する結果
More on Morphisms, Lemma \ref{more-morphisms-lemma-composition-CM}
から従う。別の方法として Lemma \ref{spaces-more-morphisms-lemma-CM} の (a) と (b) の同値性を
用いてもよい。そこでネーター局所環の局所準同型
$$
\mathcal{O}_{Y, \overline{y}}/
\mathfrak m_{\overline{z}}\mathcal{O}_{Y, \overline{y}}
\longrightarrow
\mathcal{O}_{X, \overline{x}}/
\mathfrak m_{\overline{z}}\mathcal{O}_{X, \overline{x}}
$$
を考える。そのファイバーは
$$
\mathcal{O}_{X, \overline{x}}/
\mathfrak m_{\overline{y}}\mathcal{O}_{X, \overline{x}}
$$
であるから、Algebra, Lemma \ref{algebra-lemma-CM-goes-up} を用いる。
\end{proof}

\begin{lemma}
\label{spaces-more-morphisms-lemma-flat-morphism-from-CM}
$S$ をスキームとする。$f : X \to Y$ を $S$ 上の局所ネーター代数空間の
平坦な射とする。$X$ が Cohen--Macaulay ならば、$f$ は
Cohen--Macaulay であり、$\mathcal{O}_{Y, f(\overline{x})}$ は
すべての $x \in |X|$ に対して Cohen--Macaulay である。
\end{lemma}

\begin{proof}
Lemma \ref{spaces-more-morphisms-lemma-CM} を用いて代数へ移し替えると
（Lemma \ref{spaces-more-morphisms-lemma-composition-CM} の証明と比較せよ）、これは
Algebra, Lemma \ref{algebra-lemma-CM-goes-up} から従う。
\end{proof}

\begin{lemma}
\label{spaces-more-morphisms-lemma-base-change-CM}
$S$ をスキームとする。$f : X \to Y$ を $S$ 上の代数空間の射とする。
$f$ のファイバーは局所ネーター的であると仮定する。
$Y' \to Y$ は局所有限型であるとし、$f' : X' \to Y'$ を $f$ の基底変換とする。
$x' \in |X'|$ を、像 $x \in |X|$ をもつ点とする。
\begin{enumerate}
\item $f$ が $x$ で Cohen--Macaulay ならば、
$f' : X' \to Y'$ は $x'$ で Cohen--Macaulay である。
\item $f$ が $x$ で平坦であり、$f'$ が $x'$ で Cohen--Macaulay ならば、
$f$ は $x$ で Cohen--Macaulay である。
\item $Y' \to Y$ が $f'(x')$ で平坦であり、$f'$ が $x'$ で
Cohen--Macaulay ならば、$f$ は $x$ で Cohen--Macaulay である。
\end{enumerate}
\end{lemma}

\begin{proof}
% P09-ERR-1130: frozen source line 6386 omits "by" in "Denote ... the image".
$y \in |Y|$ および $y' \in |Y'|$ を $x'$ の像と書く。
$V \to Y$ を $V$ がスキームであるような全射エタール射として選ぶ。
$U \to X \times_Y V$ を $U$ がスキームであるような全射エタール射として選ぶ。
% P09-ERR-1127: frozen source line 6390 misspells "surjective" as "surjectiev".
$V' \to Y' \times_Y V$ を $V'$ がスキームであるような全射エタール射として選ぶ。
このとき $U' = U \times_V V'$ は、全射エタール射 $U' \to X'$ を備えた
スキームである。$u' \in U'$ を $x'$ へ写るものとして選び、次いで
$u \in U$ で $u'$ の像を表す。このとき補題は $U \to V$ とその基底変換
$U' \to V'$、および点 $u'$ と $u$ に対する補題から従う
（これは定義から従う）。したがって補題はスキームの場合 More on Morphisms,
Lemma \ref{more-morphisms-lemma-base-change-CM} から従う。
\end{proof}

\begin{lemma}
\label{spaces-more-morphisms-lemma-flat-finite-presentation-CM-open}
$S$ をスキームとする。$f : X \to Y$ を $S$ 上の代数空間の射で、
平坦かつ局所有限表示なものとする。
$$
W = \{x \in |X| : f\text{ が Cohen--Macaulay である点 }x\}
$$
とおく。このとき $W$ は $|X|$ で開であり、$W$ の構成は $f$ の
任意の基底変換と可換である。任意の射 $g : Y' \to Y$ に対し、
基底変換 $f' : X' \to Y'$、すなわち $f$ の基底変換、および射影
$g' : X' \to X$ を考える。
% P09-ERR-1131: frozen source line 6415 redundantly says "is equal to W' =".
このとき対応する集合 $W'$、すなわち射 $f'$ に対するものは
$W' = (g')^{-1}(W)$ に等しい。
\end{lemma}

\begin{proof}
可換図式
$$
\xymatrix{
U \ar[d] \ar[r] & V \ar[d] \\
X \ar[r] & Y
}
$$
で、垂直射がエタールであり、$U$ と $V$ がスキームであるものを選ぶ。
$u \in U$ とし、その像を $x \in |X|$ とする。このとき $f$ が $x$ で
Cohen--Macaulay であることと、$U \to V$ が $u$ で Cohen--Macaulay
であることは（定義により）同値である。したがって射 $U \to V$ の場合へ帰着する。
More on Morphisms, Lemma
\ref{more-morphisms-lemma-flat-finite-presentation-CM-open} を参照せよ。
\end{proof}

\begin{lemma}
\label{spaces-more-morphisms-lemma-lfp-CM-relative-dimension}
$S$ をスキームとする。$f : X \to Y$ を $S$ 上の代数空間の射とする。
$f$ は局所有限表示かつ Cohen--Macaulay であると仮定する。このとき
% P09-ERR-1128: frozen source line 6440 calls algebraic subspaces X_d subschemes.
開かつ閉な部分スキーム $X_d \subset X$ で、
$X = \coprod_{d \geq 0} X_d$ かつ $f|_{X_d} : X_d \to Y$ が
相対次元 $d$ をもつものが存在する。
\end{lemma}

\begin{proof}
可換図式
$$
\xymatrix{
U \ar[d] \ar[r] & V \ar[d] \\
X \ar[r] & Y
}
$$
で、垂直射がエタールであり、$U$ と $V$ がスキームであるものを選ぶ。
このとき $U \to V$ は局所有限表示かつ Cohen--Macaulay である
（定義から直ちに従う）。したがって、開かつ閉な部分スキームへの分解
$U = \coprod_{d \geq 0} U_d$ で、
% P09-ERR-1129: frozen source line 6457 restricts f to U_d although f has domain X.
$f|_{U_d} : U_d \to V$ が相対次元 $d$ をもつものが得られる。
Morphisms, Lemma
\ref{morphisms-lemma-flat-finite-presentation-CM-fibres-relative-dimension}
を参照せよ。$u \in U$ とし、その像を $x \in |X|$ とする。このとき
$f$ が相対次元 $d$ を $x$ でもつことと、$U \to V$ が相対次元 $d$ を
$u$ でもつことは同値である（これは定義から従う）。
このようにして、$U_d$ はある部分集合 $X_d \subset |X|$ の逆像であり、
この部分集合は必然的に開かつ閉であることが分かる。
% P09-ERR-1132: frozen source line 6466 omits "by" in "Denoting X_d ...".
$X_d$ で $X$ の対応する開かつ閉な代数部分空間を表せば、補題が成り立つ。
\end{proof}







\section{Gorenstein 射}
\label{spaces-more-morphisms-section-gorenstein}

\noindent
これは Duality for Schemes, Section
\ref{duality-section-gorenstein-morphisms} の類似である。

\begin{lemma}
\label{spaces-more-morphisms-lemma-gorenstein-local-ring-fibre}
スキームの芽の射に関する性質
\begin{align*}
& \mathcal{P}((X, x) \to (S, s)) = \\
& \text{ファイバーの局所環 }
\mathcal{O}_{X_s, x}
\text{ がネーター的かつ Gorenstein である}
\end{align*}
は始域と終域の上でエタール局所的である（Descent, Definition
\ref{descent-definition-local-source-target-at-point}）。
\end{lemma}

\begin{proof}
Descent, Definition \ref{descent-definition-local-source-target-at-point}
にある図式が与えられると、ファイバーのエタール射
$U'_{v'} \to U_v$ で $u'$ を $u$ へ写すものを得る。Descent, Lemma
\ref{descent-lemma-etale-on-fiber} を参照せよ。
したがって $\mathcal{O}_{U_v, u} \to \mathcal{O}_{U'_{v'}, u'}$ は
エタール環準同型の局所化である。それゆえ、前者がネーター的であることと
後者がネーター的であることは同値である。More on Algebra, Lemma
\ref{more-algebra-lemma-Noetherian-etale-extension} を参照せよ。さらに
$$
\mathcal{O}_{U'_{v'}, u'}/\mathfrak m_u \mathcal{O}_{U'_{v'}, u'}
= \kappa(u')
$$
は Gorenstein 環であるから（Algebra, Lemma
\ref{algebra-lemma-etale-at-prime}）、Dualizing Complexes, Lemma
\ref{dualizing-lemma-flat-under-gorenstein} により
$\mathcal{O}_{U_v, u}$ が Gorenstein であることと
$\mathcal{O}_{U'_{v'}, u'}$ が Gorenstein であることは同値である。
\end{proof}

\begin{definition}
\label{spaces-more-morphisms-definition-gorenstein}
$S$ をスキームとする。
$f : X \to Y$ を $S$ 上の代数空間の射とする。
$f$ のファイバーは局所ネーター的であると仮定する
（Divisors on Spaces, Definition
\ref{spaces-divisors-definition-locally-Noetherian-fibre}）。
\begin{enumerate}
\item $x \in |X|$ とし、$y = f(x)$ とする。$f$ が
\textit{$x$ で Gorenstein} であるとは、$f$ が $x$ で平坦であり、
Morphisms of Spaces, Lemma
\ref{spaces-morphisms-lemma-local-source-target-at-point} の同値な条件が
Lemma \ref{spaces-more-morphisms-lemma-gorenstein-local-ring-fibre} に記述された性質
$\mathcal{P}$ に対して成り立つことをいう。
\item $f$ が \textit{Gorenstein 射}であるとは、$f$ が $X$ の
すべての点で Gorenstein であることをいう。
\end{enumerate}
\end{definition}

\noindent
次はスキームの場合からの移し替えである。

\begin{lemma}
\label{spaces-more-morphisms-lemma-gorenstein}
$S$ をスキームとする。$f : X \to Y$ を $S$ 上の代数空間の射とする。
$f$ のファイバーは局所ネーター的であると仮定する。次は同値である：
\begin{enumerate}
\item $f$ は Gorenstein である。
\item $f$ は平坦であり、ある全射エタール射 $V \to Y$ で $V$ が
スキームであるものに対し、$X_V \to V$ のファイバーは
Gorenstein 代数空間である。
\item $f$ は平坦であり、任意のエタール射 $V \to Y$ で $V$ が
スキームであるものに対し、$X_V \to V$ のファイバーは
Gorenstein 代数空間である。
\end{enumerate}
$x \in |X|$ を、像 $y \in |Y|$ をもつ点とする。このとき次は同値である：
\begin{enumerate}
\item[(a)] $f$ は $x$ で Gorenstein である。
\item[(b)] $\mathcal{O}_{Y, \overline{y}} \to \mathcal{O}_{X, \overline{x}}$
は平坦であり、
$\mathcal{O}_{X, \overline{x}}/
\mathfrak m_{\overline{y}}\mathcal{O}_{X, \overline{x}}$ は Gorenstein である。
\end{enumerate}
\end{lemma}

\begin{proof}
$V \to Y$ を $V$ がスキームであるようなエタール射とし、スキーム $U$ と
全射エタール射 $U \to X \times_Y V$ を選ぶ。可換図式
$$
\xymatrix{
U \ar[d] \ar[r] & V \ar[d] \\
X \ar[r] & Y
}
$$
を考える。$u \in U$ とし、その像をそれぞれ $x \in |X|$、
$y \in |Y|$、$v \in V$ とする。このとき $f$ が $x$ で Gorenstein
であることと、$U \to V$ が $u$ で Gorenstein であることは
（定義により）同値である。さらに射 $U_v \to X_v = (X_V)_v$ は
全射エタールである。したがってスキーム $U_v$ が Gorenstein であることと、
代数空間 $X_v$ が Gorenstein であることは同値である。
ゆえに (1), (2), (3) の同値性は、スキームの射に対する対応する同値性
Duality for Schemes, Lemma \ref{duality-lemma-gorenstein} から
形式的な議論により従う。

\medskip\noindent
(a) と (b) の同値性を証明する。平坦性について対応する同値性は
Morphisms of Spaces, Lemma
\ref{spaces-morphisms-lemma-flat-at-point-etale-local-rings} である。
したがって同値性を証明する際には $f$ が $x$ で平坦であると仮定してよい。
上のような図式と $x, y, u, v$ を考える。このとき
$\mathcal{O}_{Y, \overline{y}} \to \mathcal{O}_{X, \overline{x}}$ は、
局所環の狭義ヘンゼル化の間の写像
$\mathcal{O}_{V, v}^{sh} \to \mathcal{O}_{U, u}^{sh}$ に等しい。
Properties of Spaces, Lemma
\ref{spaces-properties-lemma-describe-etale-local-ring} を参照せよ。
したがって Algebra, Lemma
\ref{algebra-lemma-quotient-strict-henselization} により
$$
\mathcal{O}_{X, \overline{x}}/
\mathfrak m_{\overline{y}}\mathcal{O}_{X, \overline{x}} =
(\mathcal{O}_{U, u}/\mathfrak m_v \mathcal{O}_{U, u})^{sh}
$$
である。ゆえに、ネーター局所環
$\mathcal{O}_{U, u}/\mathfrak m_v \mathcal{O}_{U, u}$ が
Gorenstein であることと、その狭義ヘンゼル化が Gorenstein であることが
同値であると示せばよい。これは Dualizing Complexes, Lemma
\ref{dualizing-lemma-completion-henselization-dualizing} と、自身の上の
双対複体であるネーター局所環としての Gorenstein 局所環の定義から直ちに従う。
\end{proof}

\begin{lemma}
\label{spaces-more-morphisms-lemma-composition-gorenstein}
$S$ をスキームとする。$f : X \to Y$ および $g : Y \to Z$ を
$S$ 上の代数空間の射とする。$f$, $g$, $g \circ f$ のファイバーは
局所ネーター的であると仮定する。$x \in |X|$ とし、その像を
$y \in |Y|$ および $z \in |Z|$ とする。
\begin{enumerate}
\item $f$ が $x$ で Gorenstein であり、$g$ が $f(x)$ で
Gorenstein ならば、$g \circ f$ は $x$ で Gorenstein である。
\item $f$ と $g$ が Gorenstein ならば、$g \circ f$ は Gorenstein である。
\item $g \circ f$ が $x$ で Gorenstein であり、$f$ が $x$ で平坦ならば、
$f$ は $x$ で Gorenstein であり、$g$ は $f(x)$ で Gorenstein である。
% P09-ERR-1133: frozen source line 6624 uses the type-invalid composite f \circ g.
\item $f \circ g$ が Gorenstein であり、$f$ が平坦ならば、$f$ は
Gorenstein であり、$g$ は $f$ の像の各点で Gorenstein である。
\end{enumerate}
\end{lemma}

\begin{proof}
エタール局所的に作業すれば、これはスキームに対する対応する結果
Duality for Schemes, Lemma
\ref{duality-lemma-composition-gorenstein} から従う。別の方法として
Lemma \ref{spaces-more-morphisms-lemma-gorenstein} の (a) と (b) の同値性を用いてもよい。
そこでネーター局所環の局所準同型
$$
\mathcal{O}_{Y, \overline{y}}/
\mathfrak m_{\overline{z}}\mathcal{O}_{Y, \overline{y}}
\longrightarrow
\mathcal{O}_{X, \overline{x}}/
\mathfrak m_{\overline{z}}\mathcal{O}_{X, \overline{x}}
$$
を考える。そのファイバーは
$$
\mathcal{O}_{X, \overline{x}}/
\mathfrak m_{\overline{y}}\mathcal{O}_{X, \overline{x}}
$$
であるから、Dualizing Complexes, Lemma
\ref{dualizing-lemma-flat-under-gorenstein} を用いる。
\end{proof}

\begin{lemma}
\label{spaces-more-morphisms-lemma-flat-morphism-from-gorenstein}
$S$ をスキームとする。$f : X \to Y$ を $S$ 上の局所ネーター代数空間の
平坦な射とする。$X$ が Gorenstein ならば、$f$ は Gorenstein であり、
$\mathcal{O}_{Y, f(\overline{x})}$ はすべての $x \in |X|$ に対して
Gorenstein である。
\end{lemma}

\begin{proof}
Lemma \ref{spaces-more-morphisms-lemma-gorenstein} を用いて代数へ移し替えると
（Lemma \ref{spaces-more-morphisms-lemma-composition-gorenstein} の証明と比較せよ）、これは
Dualizing Complexes, Lemma \ref{dualizing-lemma-flat-under-gorenstein} から従う。
\end{proof}

\begin{lemma}
\label{spaces-more-morphisms-lemma-base-change-gorenstein}
$S$ をスキームとする。$f : X \to Y$ を $S$ 上の代数空間の射とする。
$f$ のファイバーは局所ネーター的であると仮定する。
$Y' \to Y$ は局所有限型であるとする。$f' : X' \to Y'$ を
$f$ の基底変換とする。$x' \in |X'|$ を、像 $x \in |X|$ をもつ点とする。
\begin{enumerate}
\item $f$ が $x$ で Gorenstein ならば、$f' : X' \to Y'$ は
$x'$ で Gorenstein である。
\item $f$ が $x$ で平坦であり、$f'$ が $x'$ で Gorenstein ならば、
$f$ は $x$ で Gorenstein である。
\item $Y' \to Y$ が $f'(x')$ で平坦であり、$f'$ が $x'$ で
Gorenstein ならば、$f$ は $x$ で Gorenstein である。
\end{enumerate}
\end{lemma}

\begin{proof}
% P09-ERR-1134: frozen source line 6688 omits "by" in "Denote ... the image".
$y \in |Y|$ および $y' \in |Y'|$ を $x'$ の像と書く。
$V \to Y$ を $V$ がスキームであるような全射エタール射として選ぶ。
$U \to X \times_Y V$ を $U$ がスキームであるような全射エタール射として選ぶ。
% P09-ERR-1135: frozen source line 6692 misspells "surjective" as "surjectiev".
$V' \to Y' \times_Y V$ を $V'$ がスキームであるような全射エタール射として選ぶ。
このとき $U' = U \times_V V'$ は、全射エタール射 $U' \to X'$ を備えた
スキームである。$u' \in U'$ を $x'$ へ写るものとして選び、次いで
$u \in U$ で $u'$ の像を表す。このとき補題は $U \to V$ とその基底変換
$U' \to V'$、および点 $u'$ と $u$ に対する補題から従う
（これは定義から従う）。したがって補題はスキームの場合
Duality for Schemes, Lemma \ref{duality-lemma-base-change-gorenstein} から従う。
\end{proof}

\begin{lemma}
\label{spaces-more-morphisms-lemma-flat-finite-presentation-gorenstein-open}
$S$ をスキームとする。$f : X \to Y$ を $S$ 上の代数空間の射で、
平坦かつ局所有限表示なものとする。
$$
W = \{x \in |X| : f\text{ が Gorenstein である点 }x\}
$$
とおく。このとき $W$ は $|X|$ で開であり、$W$ の構成は $f$ の
任意の基底変換と可換である。任意の射 $g : Y' \to Y$ に対し、
基底変換 $f' : X' \to Y'$、すなわち $f$ の基底変換、および射影
$g' : X' \to X$ を考える。
% P09-ERR-1136: frozen source line 6717 redundantly says "is equal to W' =".
このとき対応する集合 $W'$、すなわち射 $f'$ に対するものは
$W' = (g')^{-1}(W)$ に等しい。
\end{lemma}

\begin{proof}
% P09-ERR-1137: frozen source lines 6721--6732 omit the chart hypotheses used by the reduction.
可換図式
$$
\xymatrix{
U \ar[d] \ar[r] & V \ar[d] \\
X \ar[r] & Y
}
$$
を選ぶ。$u \in U$ とし、その像を $x \in |X|$ とする。このとき $f$ が
$x$ で Gorenstein であることと、$U \to V$ が $u$ で Gorenstein
であることは（定義により）同値である。したがってスキームの射
$U \to V$ の場合へ帰着する。開性は Duality for Schemes, Lemma
\ref{duality-lemma-flat-finite-presentation-characterize-gorenstein}
で、基底変換との可換性は Duality for Schemes, Lemma
\ref{duality-lemma-flat-lft-base-change-gorenstein} で証明されている。
\end{proof}










\section{Cohen--Macaulay 射の切断}
\label{spaces-more-morphisms-section-slice}

\noindent
$S$ をスキームとし、$X$ を $S$ 上の代数空間とする。
$f_1, \ldots, f_r \in \Gamma(X, \mathcal{O}_X)$ とする。この場合、
% P09-ERR-1138: frozen source line 6754 omits "by" in "denote V(...) the...".
$V(f_1, \ldots, f_r)$ で、$X$ の $f_1, \ldots, f_r$ により切り出される
\textit{閉部分空間}を表す。より正確には、$V(f_1, \ldots, f_r)$ を、
$X$ の閉部分空間であって、$f_1, \ldots, f_r$ により生成される
準連接イデアル層に対応するものとして定義できる。Morphisms of Spaces, Lemma
\ref{spaces-morphisms-lemma-closed-immersion-ideals} を参照せよ。
別の方法として、表示 $X = U/R$ を選び、その閉部分スキーム
$Z \subset U$ で $f_1|U, \ldots, f_r|_U$ により切り出されるものを考えてもよい。
$Z$ が $R$ 不変な閉部分スキームであることは明らかであり
（Groupoids, Definition \ref{groupoids-definition-invariant-open} を参照）、
$V(f_1, \ldots, f_r) = Z/R_Z$ とおいてよい。

\begin{lemma}
\label{spaces-more-morphisms-lemma-slice}
$S$ をスキームとする。カルテシアン図式
$$
\xymatrix{
X \ar[d] & F \ar[l]^p \ar[d] \\
Y & \Spec(k) \ar[l]
}
$$
を考える。ここで $X \to Y$ は $S$ 上の代数空間の射であって平坦かつ
局所有限表示であり、$k$ は $S$ 上の体である。
$f_1, \ldots, f_r \in \Gamma(X, \mathcal{O}_X)$ とし、
$z \in |F|$ は $f_1, \ldots, f_r$ が局所環
$\mathcal{O}_{F, \overline{z}}$ の正則列へ写るような点とする。
このとき $X$ を $p(z)$ を含む開部分空間で置き換えれば、射
$$
V(f_1, \ldots, f_r) \longrightarrow Y
$$
は平坦かつ局所有限表示である。
\end{lemma}

\begin{proof}
$Z = V(f_1, \ldots, f_r)$ とおく。$Z \to X$ が局所有限表示であることは
明らかなので、合成 $Z \to Y$ は局所有限表示である。Morphisms of Spaces,
Lemma \ref{spaces-morphisms-lemma-composition-finite-presentation} を参照せよ。
したがって $Z \to Y$ が $p(z)$ のある近傍で平坦であることを示せば十分である。
$k'/k$ を拡大体とする。このとき
$F' = F \times_{\Spec(k)} \Spec(k')$ は $F$ 上全射かつ平坦であるから、
$z' \in |F'|$ で $z$ へ写るものを見つけられ、局所環準同型
$\mathcal{O}_{F, \overline{z}} \to \mathcal{O}_{F', \overline{z}'}$ は平坦である。
Morphisms of Spaces, Lemma
\ref{spaces-morphisms-lemma-flat-at-point-etale-local-rings} を参照せよ。
したがって $f_1, \ldots, f_r$ の $\mathcal{O}_{F', \overline{z}'}$ における像も
正則列である。Algebra, Lemma \ref{algebra-lemma-flat-increases-depth} を参照せよ。
ゆえに証明中で $k$ を拡大体に置き換えてよい。特に、$z \in |F|$ が
切断 $z : \Spec(k) \to F$、すなわち構造射 $F \to \Spec(k)$ の切断から
来ると仮定してよい。

\medskip\noindent
スキーム $V$ と全射エタール射 $V \to Y$ を選ぶ。スキーム $U$ と
全射エタール射 $U \to X \times_Y V$ を選ぶ。必要ならば $k$ をさらに
拡大して、$\Spec(k) \to F \to X$ が $U$ を経由すると仮定してよい
（$U \to X$ は全射だからである）。そのような分解
$u : \Spec(k) \to U$ をとり、
% P09-ERR-1139: frozen source line 6815 omits "by" in "denote v ... the image".
$v \in V$ で $u$ の像を表す。射
$$
U_v \times_{\Spec(\kappa(v))} \Spec(k) =
U \times_V \Spec(k) \to U \times_Y \Spec(k) \to F
$$
はエタールである（第一の射は $V \to V \times_Y V$ の基底変換であり、
第二の射は $U \to X$ の基底変換である）。さらに構成により、点
$u : \Spec(k) \to U$ は左端の空間の点を与え、それは右端で $z$ へ写る。
% P09-ERR-1140: frozen source line 6823 writes "left most" instead of "leftmost".
したがって元 $f_1, \ldots, f_r$ は次の写像の右側の局所環における
正則列へ写る：
$$
\mathcal{O}_{U_v, u}
\longrightarrow
% P09-ERR-1142: frozen source line 6830 omits the closing parenthesis in Spec(kappa(v)).
\mathcal{O}_{U_v \times_{\Spec(\kappa(v)} \Spec(k), \overline{u}}
=
\mathcal{O}_{U \times_V \Spec(k), \overline{u}}.
$$
ところが表示された矢印は平坦であるから（More on Flatness, Lemma
\ref{flat-lemma-flat-up-down-henselization} と Morphisms of Spaces, Lemma
\ref{spaces-morphisms-lemma-flat-at-point-etale-local-rings} を合わせて用いる）、
Algebra, Lemma \ref{algebra-lemma-flat-increases-depth} により
% P09-ERR-1141: frozen source line 6841 uses singular "maps" with the plural subject.
$f_1, \ldots, f_r$ は $\mathcal{O}_{U_v, u}$ の正則列へ写ることが分かる。
More on Morphisms, Lemma \ref{more-morphisms-lemma-slice-given-elements}
により、スキームの射
$$
V(f_1, \ldots, f_r) \times_X U = V(f_1|_U, \ldots, f_r|_U) \to V
$$
はある開近傍 $U'$、すなわち $u$ の開近傍で平坦である。$X' \subset X$ を
$|U'| \to |X|$ の像に対応する開部分空間とする（Properties of Spaces, Lemmas
\ref{spaces-properties-lemma-topology-points} and
\ref{spaces-properties-lemma-open-subspaces} を参照）。
$V(f_1, \ldots, f_r) \cap X' \to Y$ は平坦であると結論できる
（Morphisms of Spaces, Definition \ref{spaces-morphisms-definition-flat} を参照）。
実際、可換図式
$$
\xymatrix{
V(f_1, \ldots, f_r) \times_X U' \ar[d]_a \ar[r] & V \ar[d]^b \\
V(f_1, \ldots, f_r) \cap X' \ar[r] & Y
}
$$
があり、$a, b$ はエタールで $a$ は全射である。
\end{proof}






\section{被約ファイバー}
\label{spaces-more-morphisms-section-reduced}

\noindent
この節は More on Morphisms, Section
\ref{more-morphisms-section-reduced} の類似である。

\begin{lemma}
\label{spaces-more-morphisms-lemma-geometrically-reduced-fibre}
$S$ をスキームとする。$f : X \to Y$ を $S$ 上の代数空間の射とする。
$y \in |Y|$ とする。次は同値である：
\begin{enumerate}
\item ある射 $\Spec(k) \to Y$ で $y$ の同値類に属するものに対して、
代数空間 $X_k$ は $k$ 上幾何学的被約である。
\item すべての射 $\Spec(k) \to Y$ で $y$ の同値類に属するものに対して、
代数空間 $X_k$ は $k$ 上幾何学的被約である。
\item すべての射 $\Spec(k) \to Y$ で $y$ の同値類に属するものに対して、
代数空間 $X_k$ は被約である。
\end{enumerate}
\end{lemma}

\begin{proof}
これは Spaces over Fields, Lemma
\ref{spaces-over-fields-lemma-geometrically-reduced-upstairs} と、
% P09-ERR-1143: frozen source line 6898 cites the equivalence relation defining |X|, not |Y|.
Properties of Spaces, Section \ref{spaces-properties-section-points}
で与えられる $|X|$ を定義する同値関係の定義から直ちに従う。
\end{proof}

\begin{definition}
\label{spaces-more-morphisms-definition-geometrically-reduced-fibre}
$S$ をスキームとする。$f : X \to Y$ を $S$ 上の代数空間の射とする。
$y \in |Y|$ とする。\textit{$f : X \to Y$ の $y$ におけるファイバーが
幾何学的被約である}とは、Lemma \ref{spaces-more-morphisms-lemma-geometrically-reduced-fibre} の
同値な条件が成り立つことをいう。
\end{definition}

\noindent
必要な標準的補題を述べる。

\begin{lemma}
\label{spaces-more-morphisms-lemma-base-change-fibres-geometrically-reduced}
$S$ をスキームとする。$f : X \to Y$ および $g : Y' \to Y$ を
$S$ 上の代数空間の射とする。
% P09-ERR-1144: frozen source line 6919 omits "by" in "Denote f' ... the base change".
$f' : X' \to Y'$ で、$f$ の $g$ による基底変換を表す。このとき
\begin{align*}
\{y' \in |Y'| :
\text{ファイバー }f' : X' \to Y'\text{ は点 }y'
\text{ で幾何学的被約である}\} \\
= g^{-1}(\{y \in |Y| :
\text{ファイバー }f : X \to Y\text{ は点 }y
\text{ で幾何学的被約である}\}).
\end{align*}
\end{lemma}

\begin{proof}
$y' \in |Y'|$ に対し、射 $\Spec(k) \to Y'$ で
$y'$ の同値類に属するものを選ぶ。このとき $g(y')$ は合成
$\Spec(k) \to Y' \to Y$ により表される。したがって
$X' \times_{Y'} \Spec(k) = X \times_Y \Spec(k)$ であり、結果は定義から従う。
\end{proof}

\begin{lemma}
\label{spaces-more-morphisms-lemma-geometrically-reduced-constructible}
$S$ をスキームとする。$f : X \to Y$ を $S$ 上の代数空間の射で、
準コンパクトかつ局所有限表示なものとする。このとき集合
$$
E = \{y \in |Y| : \text{ }f : X \to Y\text{ の }y
\text{ におけるファイバーが幾何学的被約である}\}
$$
はエタール局所的に構成可能である。
\end{lemma}

\begin{proof}
アフィンスキーム $V$ とエタール射 $V  \to Y$ を選ぶ。主張の意味は、
$E$ の $|V|$ における逆像が構成可能であるということである。
Lemma \ref{spaces-more-morphisms-lemma-base-change-fibres-geometrically-reduced} により $Y$ を $V$ で
置き換えてよい。すなわち $Y$ がアフィンスキームであると仮定してよい。
このとき $X$ は準コンパクトである。アフィンスキーム $U$ と全射エタール射
$U \to X$ を選ぶ。射 $\Spec(k) \to Y$ に対し、ファイバー間の射
$U_k \to X_k$ は全射エタールである。したがって $U_k$ が $k$ 上
幾何学的被約であることと $X_k$ が $k$ 上幾何学的被約であることは同値である。
Spaces over Fields, Lemma
\ref{spaces-over-fields-lemma-geometrically-reduced-etale-local} を参照せよ。
ゆえに集合 $E$、すなわち $X \to Y$ に対するものは、集合 $E$、すなわち
$U \to Y$ に対するものと同じである。
このようにして、補題はスキームの場合 More on Morphisms, Lemma
\ref{more-morphisms-lemma-geometrically-reduced-constructible} から従う。
\end{proof}

\begin{lemma}
\label{spaces-more-morphisms-lemma-proper-flat-over-dvr-reduced-fibre}
$X$ を離散付値環 $R$ 上の代数空間とし、その構造射
$X \to \Spec(R)$ は固有かつ平坦であるとする。特殊ファイバーが被約ならば、
$X$ と一般ファイバー $X_\eta$ はともに被約である。
\end{lemma}

\begin{proof}
エタール射 $U \to X$ で $U$ がアフィンスキームであるものを選ぶ。
このとき $U$ は $R$ 上有限型である。$u \in U$ を特殊ファイバー内の点とする。
局所環 $A = \mathcal{O}_{U, u}$ は $R$ 上本質的有限型であり、したがって
ネーター的である。$\pi \in R$ を一様化元とする。$X$ は $R$ 上平坦なので、
$\pi \in \mathfrak m_A$ は $A$ 上の非零因子であり、$X$ の特殊ファイバーは
被約なので $A/\pi A$ は被約である。
$a \in A$, $a \not = 0$ ならば、ある $n \geq 0$ と元 $a' \in A$ が存在し、
$a = \pi^n a'$ かつ $a' \not \in \pi A$ である。これは Krull の交叉定理
（Algebra, Lemma \ref{algebra-lemma-intersect-powers-ideal-module-zero}）から従う。
$a$ が冪零ならば、$a'$ も冪零である。実際、$\pi$ は非零因子である。
しかし $a'$ は被約環 $A/\pi A$ の非零元へ写るので、これは不可能である。
ゆえに $A$ は被約である。したがって $u$ の、$U$ 内のある開近傍で
被約なものが存在する（小さな詳細は省略する；$U$ がネーター的であることを使う）。
これにより、エタール射 $U \to X$ で $U$ が被約スキームであり、
$X$ の特殊ファイバーのすべての点が像に属するものを見つけられる。
$X$ は $R$ 上固有なので、
$U \to X$ は全射である。したがって $X$ は被約である。
$U \to \Spec(R)$ の一般ファイバーも被約なので（アフィン部分上では局所化を
とることにより計算される）、一般ファイバーについても同じ結論を得る。
\end{proof}

\begin{lemma}
\label{spaces-more-morphisms-lemma-geometrically-reduced-open}
$S$ をスキームとする。$f : X \to Y$ を $S$ 上の代数空間の射とする。
$f$ が平坦、固有、かつ有限表示ならば、集合
$$
E = \{y \in |Y| : \text{ }f : X \to Y\text{ の }y
\text{ におけるファイバーが幾何学的被約である}\}
$$
は $|Y|$ で開である。
\end{lemma}

\begin{proof}
Lemma \ref{spaces-more-morphisms-lemma-base-change-fibres-geometrically-reduced} により $E$ の構成は
基底変換と可換である。$|Y|$ の部分集合が開であることを調べるには、
$Y$ をエタール被覆の各要素で置き換えてよい。したがって $Y$ がアフィンであると
仮定してよい。このとき $Y$ は $\mathbf{Z}$ 上有限型なアフィンスキームの
余フィルター極限である。ゆえに $X \to Y$ は $X_0 \to Y_0$ の基底変換で、
$Y_0$ が有限型 $\mathbf{Z}$ 代数のスペクトルであり、$X_0 \to Y_0$ が
平坦かつ固有であるものと仮定してよい。Limits of Spaces, Lemma
\ref{spaces-limits-lemma-descend-finite-presentation},
\ref{spaces-limits-lemma-descend-flat}, and
\ref{spaces-limits-lemma-eventually-proper} を参照せよ。$E$ の構成は基底変換と
可換なので（上を参照）、基底がネーター的であると仮定してよい。

\medskip\noindent
$Y$ がネーター的であると仮定する。この集合は Lemma
\ref{spaces-more-morphisms-lemma-geometrically-reduced-constructible} により構成可能である。
したがって、この集合が一般化の下で安定であることを示せば十分である
（Topology, Lemma \ref{topology-lemma-characterize-closed-Noetherian}）。
Properties, Lemma \ref{properties-lemma-locally-Noetherian-specialization-dvr}
により、$Y = \Spec(R)$ で $R$ が離散付値環であり、閉ファイバー $X_y$ が
幾何学的被約である場合へ帰着する。示すべきことは、一般ファイバー
$X_\eta$ が幾何学的被約であることである。

\medskip\noindent
そうでなければ、有限拡大 $L$、すなわち $R$ の分数体の有限拡大で
$X_L$ が被約でないものが
存在する。Spaces over Fields, Lemmas
\ref{spaces-over-fields-lemma-perfect-reduced}（標数零）and
\ref{spaces-over-fields-lemma-geometrically-reduced-positive-characteristic}
（正標数）を参照せよ。離散付値環 $R' \subset L$ で、分数体 $L$ をもち
$R$ を支配するものが存在する。Algebra, Lemma
\ref{algebra-lemma-integral-closure-Dedekind} を参照せよ。
$R$ を $R'$ で置き換えると、Lemma
\ref{spaces-more-morphisms-lemma-proper-flat-over-dvr-reduced-fibre} へ帰着する。
\end{proof}







\section{ファイバーの連結成分}
\label{spaces-more-morphisms-section-connected}

\noindent
この節は More on Morphisms, Section
\ref{more-morphisms-section-connected} の類似である。

\begin{lemma}
\label{spaces-more-morphisms-lemma-base-change-fibres-nr-geometrically-connected-components}
$S$ をスキームとする。$f : X \to Y$ を $S$ 上の代数空間の射とする。
$$
n_{X/Y} : |Y| \to \{0, 1, 2, 3, \ldots, \infty\}
$$
を、$y \in Y$ に $X_k$ の連結成分の個数を対応させる関数とする。ここで
$\Spec(k) \to Y$ は $y$ の同値類に属し、$k$ は代数閉体である。
% P09-ERR-1145: frozen source line 7076 writes y in Y instead of y in |Y|.
これは良定義であり、$g : Y' \to Y$ が射ならば
$$
% P09-ERR-1146: frozen source line 7081 composes the point-set function with g rather than |g|.
n_{X'/Y'} = n_{X/Y} \circ g
$$
である。ここで $X' \to Y'$ は $f$ の基底変換である。
\end{lemma}

\begin{proof}
% P09-ERR-1147: frozen source line 7087 writes y' in Y' and y in Y without point bars.
$y' \in Y'$ の像が $y \in Y$ であると仮定する。$\Spec(k') \to Y'$ を
$y'$ の同値類に属し、$k'$ が代数閉体であるものとしてとる。このとき可換図式
$$
\xymatrix{
\Spec(K) \ar[r] \ar[rd] &
\Spec(k') \ar[r] & Y' \ar[d] \\
& \Spec(k) \ar[r] & Y
}
$$
を選べる。ここで $K$ は代数閉体である。スキームの射
$$
\xymatrix{
X'_{k'} & (X'_{k'})_K = (X_k)_K \ar[l] \ar[r] & X_k
}
$$
は連結成分の間に全単射を誘導するので、結果が従う。Spaces over Fields, Lemma
\ref{spaces-over-fields-lemma-separably-closed-field-connected-components}
を参照せよ。これを用いて関数が良定義であることを証明するには $Y' = Y$ とする。
\end{proof}







\section{ファイバーの次元}
\label{spaces-more-morphisms-section-dimension}

\noindent
この節は More on Morphisms, Section
\ref{more-morphisms-section-dimension} の類似である。

\begin{lemma}
\label{spaces-more-morphisms-lemma-dimension-fibre}
$S$ をスキームとする。$f : X \to Y$ を $S$ 上の代数空間の有限型射とする。
$y \in |Y|$ とする。次の量は等しい：
\begin{enumerate}
\item $d = -\infty$（$y$ が $|f|$ の像に属さない場合）とし、そうでなければ、
最小の整数 $d$ であって、$f$ が相対次元 $\leq d$ を、すべての
$x \in |X|$、すなわち $y$ へ写る点でもつもの。
\item 代数空間 $X_k = \Spec(k) \times_Y X$ の次元。ここで任意の射
$\Spec(k) \to Y$ は $y$ を定義する同値類に属する。
\end{enumerate}
\end{lemma}

\begin{proof}
これを解釈するには Morphisms of Spaces, Definition
\ref{spaces-morphisms-definition-dimension-fibre},
Properties of Spaces, Definition
\ref{spaces-properties-definition-dimension},
Properties of Spaces, Definition
\ref{spaces-properties-definition-dimension-at-point} を参照する必要がある。
固定した射 $\Spec(k) \to Y$ に対して (1) と (2) の数が等しいことを示す。
エタール射 $V \to Y$ で $V$ がアフィンスキームであるものと、点
$v \in V$ で $y$ へ写るものを選ぶ。$V \times_Y \Spec(k) \to \Spec(k)$ は全射エタールなので
（Properties of Spaces, Lemma
\ref{spaces-properties-lemma-points-cartesian}）、有限分離拡大 $k'/k$ と
（Morphisms, Lemma \ref{morphisms-lemma-etale-over-field}）、可換図式
$$
\xymatrix{
\Spec(k') \ar[r] \ar[d] & V \ar[d] \\
\Spec(k) \ar[r] & Y
}
$$
を見つけられる。問題の次元は Properties of Spaces, Lemma
\ref{spaces-properties-lemma-dimension-decent-invariant-under-etale} と
Morphisms of Spaces, Lemma
\ref{spaces-morphisms-lemma-dimension-fibre-after-base-change} により変わらないので、
$X \to Y$ を $V \times_Y X \to V$ で置き換え、$X_k$ を
$X_{k'} = \Spec(k') \times_V (V \times_Y X)$ で置き換えてよい。
したがって $Y$ がアフィンスキームであると仮定してよい。
この場合、$k = \kappa(y)$ と仮定してよい。実際、$X_{\kappa(y)}$ と $X_k$ の
次元は、上記の Morphisms of Spaces, Lemma
\ref{spaces-morphisms-lemma-dimension-fibre-after-base-change} により同じである。
また、代数空間 $Z$、すなわち体 $K$ 上のものに対し、$Z$ の相対次元は点
$z \in |Z|$ で、定義により
$\dim_z(Z)$ と同じである。$Y$ はアフィンで $k = \kappa(y)$ と仮定する。
% P09-ERR-1148: frozen source line 7176 omits a conjunction after "X is quasi-compact".
$X$ は準コンパクトであり、アフィンスキーム $U$ と
% P09-ERR-1149: frozen source line 7177 says "an surjective".
全射エタール射 $U \to X$ を選べる。このとき
$\dim(X_k) = \dim(U_k) = \max \dim_u(U_k)$ は定義により (1) の数に等しい。
\end{proof}

\begin{lemma}
\label{spaces-more-morphisms-lemma-base-change-dimension-fibres}
$S$ をスキームとする。$f : X \to Y$ を $S$ 上の代数空間の有限型射とする。
$$
n_{X/Y} : |Y| \to \{-\infty, 0, 1, 2, 3, \ldots\}
$$
を、$y \in |Y|$ に Lemma \ref{spaces-more-morphisms-lemma-dimension-fibre} で論じた整数を
対応させる関数とする。$g : Y' \to Y$ が射ならば
$$
n_{X'/Y'} = n_{X/Y} \circ |g|
$$
である。ここで $X' \to Y'$ は $f$ の基底変換である。
\end{lemma}

\begin{proof}
これは Lemma \ref{spaces-more-morphisms-lemma-dimension-fibre} から直ちに従う。
\end{proof}

\begin{lemma}
\label{spaces-more-morphisms-lemma-dimension-fibres-flat}
$S$ をスキームとする。$f : X \to Y$ を $S$ 上の代数空間の
平坦かつ有限表示な射とする。$n_{X/Y}$ を、Lemma
\ref{spaces-more-morphisms-lemma-base-change-dimension-fibres} で導入した、$Y$ 上の $f$ の
ファイバーの次元を与える関数とする。このとき $n_{X/Y}$ は下半連続である。
\end{lemma}

\begin{proof}
全射エタール射 $V \to Y$ で $V$ がスキームであるものをとる。
% P09-ERR-1150: frozen source line 7215 uses undefined g for the chosen map V -> Y.
合成 $n_{X/Y} \circ |g|$ が下半連続であることを示せれば、$|g|$ は開なので
補題が従う。したがって $Y$ がスキームであると仮定してよい。
% P09-ERR-1151: frozen source line 7218 says V may be assumed affine after reducing to Y.
局所的に作業して $V$ がアフィンスキームであると仮定してよい。
このときアフィンスキーム $U$ と全射エタール射 $U \to X$ を選べる。
すると $n_{X/Y} = n_{U/Y}$ である。したがって $X$ と $Y$ はともに
スキームであると仮定してよい。この場合、補題は More on Morphisms, Lemma
\ref{more-morphisms-lemma-dimension-fibres-flat} から従う。
\end{proof}

\begin{lemma}
\label{spaces-more-morphisms-lemma-dimension-fibres-proper}
$S$ をスキームとする。$f : X \to Y$ を $S$ 上の代数空間の固有射とする。
$n_{X/Y}$ を、Lemma \ref{spaces-more-morphisms-lemma-base-change-dimension-fibres} で導入した、
$Y$ 上の $f$ のファイバーの次元を与える関数とする。このとき
$n_{X/Y}$ は上半連続である。
\end{lemma}

\begin{proof}
$Z_d = \{x \in |X| :
\text{ }f\text{ のファイバーは }x\text{ で次元 }> d\}$
とおく。このとき $Z_d$ は Morphisms of Spaces, Lemma
\ref{spaces-morphisms-lemma-openness-bounded-dimension-fibres} により
$|X|$ の閉部分集合である。$f$ は固有なので $f(Z_d)$ は $|Y|$ で閉である。
$y \in f(Z_d) \Leftrightarrow n_{X/Y}(y) > d$ だから、補題が成り立つ。
\end{proof}

\begin{lemma}
\label{spaces-more-morphisms-lemma-dimension-fibres-proper-flat}
$S$ をスキームとする。$f : X \to Y$ を $S$ 上の代数空間の
固有、平坦、かつ有限表示な射とする。$n_{X/Y}$ を、Lemma
\ref{spaces-more-morphisms-lemma-base-change-dimension-fibres} で導入した、$Y$ 上の $f$ の
ファイバーの次元を与える関数とする。このとき $n_{X/Y}$ は局所定数である。
\end{lemma}

\begin{proof}
これは Lemmas \ref{spaces-more-morphisms-lemma-dimension-fibres-flat} and
\ref{spaces-more-morphisms-lemma-dimension-fibres-proper} の直ちに得られる帰結である。
\end{proof}





\section{鎖状代数空間}
\label{spaces-more-morphisms-section-catenary}

\noindent
この節では Decent Spaces, Section
\ref{decent-spaces-section-catenary} で始めた議論を続ける。
次の補題は、その次の補題の証明で用いる。

\begin{lemma}
\label{spaces-more-morphisms-lemma-construct-glueing}
$S$ をスキームとする。$f : X \to Y$ を $S$ 上の代数空間の整射とする。
$y \in |Y|$ を、閉埋め込み $y : \Spec(k) \to Y$ により表せる点とする。
このとき分解 $X \to X' \to Y$、すなわち $f$ の分解で、次を満たすものが存在する：
\begin{enumerate}
\item $X' \to Y$ は整である。
\item $X \to X'$ は $X' \setminus X'_y$ 上で同型である。
\item $X'_y$ はただ一つの点 $x'$ をもち、$\kappa(x') = k$ である。
\end{enumerate}
さらに $f$ が有限で $Y$ が局所ネーター的ならば、$X' \to Y$ は有限である。
\end{lemma}

\begin{proof}
Morphisms of Spaces, Lemma \ref{spaces-morphisms-lemma-pushforward} により、層
$f_*\mathcal{O}_X$、$(X_y \to Y)_*\mathcal{O}_{X_y}$、および
$y_*\mathcal{O}_{\Spec(k)}$ は準連接 $\mathcal{O}_Y$ 代数の層である。写像
$$
% P09-ERR-1152: frozen source line 7295 uses f_* O_Y instead of f_* O_X.
f_*\mathcal{O}_Y \longrightarrow
(X_y \to Y)_*\mathcal{O}_{X_y} \longleftarrow
y_*\mathcal{O}_{\Spec(k)}
$$
を考える。ファイバー積は準連接 $\mathcal{O}_Y$ 代数の層
$\mathcal{A}'$ であり、$X' \to Y$ を $\mathcal{A}'$ の $Y$ 上の
相対スペクトルとして定義できる。Morphisms, Lemma
\ref{morphisms-lemma-affine-equivalence-algebras} を参照せよ。
この構成は任意の基底変換と可換である。特に開部分空間
$|Y| \setminus \{y\}$ 上で射 $X \to X'$ が同型であり、
$|Y| \setminus \{y\}$ 上で射 $X' \to Y$ が整であることは明らかである。
$y$ の小さな近傍で主張を証明することだけが残る。
アフィンスキーム $V = \Spec(R)$ とエタール射 $\varphi : V \to Y$ で、
$y$ が $\varphi$ の像に属するものを選ぶ。このとき $V_y$ は $V$ の
閉部分スキームで $k$ 上エタールであり、したがって有限個の点からなり、
各剰余体は $k$ 上分離的である（Decent Spaces, Remark
\ref{decent-spaces-remark-recall} を参照）。$V$ を縮小して、ただ一つの閉点
$v = \Spec(l) \to V$ で $y$ へ写るものが存在し、$l/k$ は有限分離的であると
仮定してよい。$V \times_Y X = \Spec(C)$ と書け、$R \to C$ は整環準同型である。
上で述べた基底変換との可換性により
% P09-ERR-1153: frozen source lines 7319 and 7332 use undefined U and Y' in the base change.
$U \times_X Y' = \Spec(C')$ であり、ここで
$$
C' = C \times_{C \otimes_R l} l
$$
である。$R \to l$ は全射なので $C \to C \otimes_R l$ も全射であり、
これは More on Algebra, Situation
\ref{more-algebra-situation-module-over-fibre-product} で研究した種類の
ファイバー積である（$A, A', B, B'$ はそれぞれ
$C \otimes_R l, C, l, C'$ に対応する）。$C'$ は $R$ 部分代数、すなわち
$C$ の部分代数であり、
したがって $R$ 上整であることに注意する。これで (1) が証明された。
最後に More on Algebra, Lemma
\ref{more-algebra-lemma-points-of-fibre-product} により、
$V \times_X Y' = \Spec(C')$ はただ一つの点 $y''$ をもち、それは $v$ の上にあり、
その剰余体は $l$ である（これは明らかな全射 $C' \to l$ に対応する）。
したがって
% P09-ERR-1154: frozen source line 7335 uses X_y where the constructed fibre must be X'_y.
$X_y \times_{\Spec(k)} \Spec(l)$ は剰余体 $l$ をもつただ一つの点をもつ。
$l/k$ は有限分離的なので、これは $X'_y$ が剰余体 $k$ をもつただ一つの点を
もつことを含意する。すなわち (3) が成り立つ。

\medskip\noindent
最後の主張を証明するため、$Y$ が局所ネーター的ならば $R$ はネーター環であり、
$f$ が有限ならば $R \to C$ は有限であることに注意する。このとき $C'$ は
More on Algebra, Lemma \ref{more-algebra-lemma-fibre-product-finite-type}
により有限型 $R$ 代数である。これにより $X' \to Y$ は有限である。
\end{proof}

\begin{lemma}
\label{spaces-more-morphisms-lemma-universally-catenary-dimension-function}
$S$ をスキームとし、$B$ を $S$ 上の代数空間とする。
$\delta : |B| \to \mathbf{Z}$ を関数とする。$B$ は decent、局所ネーター的、
かつ普遍的鎖状であり、$\delta$ は次元関数であると仮定する。
$X$ を $B$ 上の decent な代数空間で、その構造射
$f : X \to B$ が局所有限型であるものとする。このとき
$\delta_X : |X| \to \mathbf{Z}$ を規則
$$
% P09-ERR-1155: frozen source lines 7358 and 7381 run "degreeof" together.
\delta_X(x) = \delta(f(x)) + \text{超越次数 }x/f(x)
$$
により定義する（Morphisms of Spaces, Definition
\ref{spaces-morphisms-definition-dimension-fibre}）。このとき $\delta_X$ は次元関数である。
\end{lemma}

\begin{proof}
問題は $B$ 上局所的である。したがって $B$ が準コンパクトであると仮定してよい。
Decent Spaces, Lemma
\ref{decent-spaces-lemma-locally-Noetherian-decent-quasi-separated}
により $B$ は準分離的である。Limits of Spaces, Proposition
\ref{spaces-limits-proposition-there-is-a-scheme-finite-over} により、
% P09-ERR-1156: frozen source line 7372 targets X, while the finite scheme cover must target B.
有限全射 $\pi : Y \to X$ で $Y$ がスキームであるものを選べる。
主張：$\delta_Y$ は次元関数である。

\medskip\noindent
この主張は補題を含意する。補題にある $X \to B$ に対し、
$Z = Y \times_B X$ とおき、射影を $p : Z \to Y$ および $q : Z \to X$ とする。
このとき
$$
\delta_Z(z) = \delta_Y(p(z)) + \text{超越次数 }z/p(z)
$$
および $\delta_Z(z) = \delta_X(q(z))$ が成り立つ。これは Morphisms of Spaces,
Lemma \ref{spaces-morphisms-lemma-dimension-fibre-at-a-point-additive} と、
有限射に対してこれらの超越次数が零であることから従う。
Decent Spaces, Lemma \ref{decent-spaces-lemma-scheme-with-dimension-function}
と主張により、$\delta_Z$ は次元関数である。次いで Decent Spaces, Lemma
\ref{decent-spaces-lemma-check-dimension-function-finite-cover} により
$\delta_X$ は次元関数であることが分かる。

\medskip\noindent
主張を証明する。特殊化 $y \leadsto y'$、$y \not = y'$、すなわち
ネータースキーム $Y$ の点の特殊化を考える。
$\delta_Y(y) > \delta_Y(y')$ である。実際、$Y$ のファイバー内の点の間には
特殊化がない
（Decent Spaces, Lemma
\ref{decent-spaces-lemma-conditions-on-fibre-and-qf} を参照）。
これを特殊化の鎖に用いると
$$
\delta_Y(y) - \delta_Y(y') \geq
\text{codim}(\overline{\{y'\}}, \overline{\{y\}})
$$
を得る。等号を示すことが課題である。Properties, Lemma
\ref{properties-lemma-locally-Noetherian-closed-point} により特殊化
$y' \leadsto y_0$ を選べる。
$$
\delta_Y(y) - \delta_Y(y_0) =
\text{codim}(\overline{\{y_0\}}, \overline{\{y\}})
$$
を示せば十分である。なぜなら、これは $y \leadsto y'$ と
$y' \leadsto y_0$ の双方について等号を含意するからである。

\medskip\noindent
特殊化の極大鎖
$y = y_c \leadsto y_{c - 1} \leadsto \ldots \leadsto y_0$、すなわち $Y$ におけるもの
を選ぶ。$b = \pi(y)$ および $b_0 = \pi(y_0)$ とおく。
特殊化の極大鎖
$b = b_e \leadsto b_{e - 1} \leadsto \ldots \leadsto b_0$、すなわち $|B|$ におけるもの
を選ぶ。$e = c$ を示さなければならない。
$\pi$ は閉なので（Morphisms of Spaces, Lemma
\ref{spaces-morphisms-lemma-finite-proper}）、特殊化の列
$y = y'_e \leadsto y'_{e - 1} \leadsto \ldots \leadsto y'_0$
で
$b = b_e \leadsto b_{e - 1} \leadsto \ldots \leadsto b_0$
へ写るものを見つけられる。
$y'_e \leadsto y'_{e - 1} \leadsto \ldots \leadsto y'_0$
も極大鎖であることに注意する。$y_0 = y'_0$ ならば、$Y$ は鎖状なので、
望みどおり $e = c$ と結論する。次の段落では、技巧によりこの場合へ帰着し、
同じ仕方で結論する。

\medskip\noindent
$\pi$ は閉なので、$b_0$ は $|B|$ の閉点である。Decent Spaces, Lemma
\ref{decent-spaces-lemma-decent-space-closed-point} により、$b_0$ を閉埋め込み
$b_0 : \Spec(k) \to B$ で表せる。Lemma \ref{spaces-more-morphisms-lemma-construct-glueing}
により分解
$$
% P09-ERR-1157: frozen source lines 7443--7445 factor through X rather than B.
Y \to Y' \to X
$$
で、$\pi' : Y' \to X$ が有限であり、$Y \to Y'$ は $y_0$ と $y'_0$ を
同じ点へ写し、$b_0$ の逆像の外で同型であるような射であるものを見つけられる。
% P09-ERR-1158: frozen source line 7446 uses singular "map" with subject "a morphism".
（もちろん $Y'$ はスキームではないが、以下の議論には差し支えない。）
特殊化の極大鎖
$y_c \leadsto y_{c - 1} \leadsto \ldots \leadsto y_0$ および
$y'_e \leadsto y'_{e - 1} \leadsto \ldots \leadsto y'_0$ は、$Y'$ において
始点と終点が同じ特殊化の極大鎖へ写ることが明らかである。
$B$ は普遍的鎖状なので $|Y'|$ は鎖状であり、先と同様に結論する。
\end{proof}







\section{射のエタール局所化}
\label{spaces-more-morphisms-section-etale-localization}

\noindent
この節は More on Morphisms, Section
\ref{more-morphisms-section-etale-localization} の類似である。

\begin{lemma}
\label{spaces-more-morphisms-lemma-etale-splits-off-quasi-finite-part}
$S$ をスキームとする。$f : X \to Y$ を $S$ 上の代数空間の射とする。
$y \in |Y|$ とし、$x_1, \ldots, x_n \in |X|$ を $y$ へ写る点とする。
次を仮定する：
\begin{enumerate}
\item $f$ は局所有限型である。
\item $f$ は分離的である。
\item $f$ は $x_1, \ldots, x_n$ で準有限である。
\item $f$ は準コンパクトであるか、または $Y$ は decent である。
\end{enumerate}
このとき点付き代数空間のエタール射 $(U, u) \to (Y, y)$ と、
開かつ閉な部分空間への分解
$$
U \times_Y X = W \amalg V
$$
で、射 $V \to U$ が有限であり、$|V| \to |U|$ の $u$ 上のファイバーの
各点がいずれかの $x_i$ へ写り、$|W| \to |U|$ の $u$ 上のファイバーが
いずれの $x_i$ へ写る点も含まないものが存在する。
\end{lemma}

\begin{proof}
$(U, u) \to (Y, y)$ を代数空間のエタール射とし、$w \in |U \times_Y X|$
のうち $u \in |U|$ へ写り、かついずれかの $x_i \in |X|$ へ写るもの全体を考える。
Decent Spaces, Lemma \ref{decent-spaces-lemma-qf-and-qc-finite-fibre}
（$f$ が有限型の場合）または Decent Spaces, Lemma
\ref{decent-spaces-lemma-decent-finite-fibre}（$Y$ が decent の場合）により、
この集合は有限である。したがって $f$ を基底変換 $U \times_Y X \to U$ で
置き換え、$x_1, \ldots, x_n$ をこれらの $w$ の集合で置き換えてよい。
特に $Y$ がアフィンスキームであると仮定し、実際そう仮定する。このとき
$X$ は分離的代数空間である。

\medskip\noindent
アフィンスキーム $Z$ とエタール射 $Z \to X$ で、$x_1, \ldots, x_n$ が
$|Z| \to |X|$ の像に属するものを選ぶ。$|Z| \to |X|$ のファイバーは有限である。
Properties of Spaces, Lemma
\ref{spaces-properties-lemma-finite-fibres-presentation}（またはより一般的な議論
Decent Spaces, Section
\ref{decent-spaces-section-reasonable-decent}）を参照せよ。
$\{z_1, \ldots, z_m\} \subset |Z|$ を
$\{x_1, \ldots, x_n\}$ の逆像とする。More on Morphisms, Lemma
\ref{more-morphisms-lemma-etale-splits-off-quasi-finite-part-technical}
により、エタール射 $(U, u) \to (Y, y)$ で
$U \times_Y Z = Z_1 \amalg Z_2$、$Z_1 \to U$ が有限、かつ
% P09-ERR-1159: frozen source line 7519 labels the fibre (Z_1)_y instead of (Z_1)_u.
$(Z_1)_y = \{z_1, \ldots, z_m\}$ となるものが存在する。
$U$ がアフィンであると仮定してよく、したがって $Z_1$ もアフィンである。

\medskip\noindent
$f$ は分離的なので、像 $V$、すなわち $Z_1 \to X$ の像は開かつ閉である
（Morphisms of Spaces, Lemma
\ref{spaces-morphisms-lemma-universally-closed-permanence}）。
$W = X \setminus V$ とおけば補題にある分解を得る。証明を終えるには
$V \to U$ が有限であることを示さなければならない。
$Z_1 \to V$ は全射エタールなので、$V$ は $Z_1$ をエタール同値関係
$R = Z_1 \times_V Z_1$ で割った商である。Spaces, Lemma
\ref{spaces-lemma-space-presentation} を参照せよ。
$f$ は分離的なので $V \to U$ は分離的であり、$R$ は
$Z_1 \times_U Z_1$ の閉部分である。$Z_1 \to U$ は有限なので、射影
$s, t : R \to Z_1$ は有限である。したがって Groupoids, Proposition
\ref{groupoids-proposition-finite-flat-equivalence} により $V$ は
アフィンスキームである。Morphisms, Lemma
\ref{morphisms-lemma-image-proper-is-proper} により $V \to U$ は固有であり、
Morphisms, Lemma \ref{morphisms-lemma-finite-proper} により $V \to U$ は有限である。
これで証明が完了する。
\end{proof}

\begin{lemma}
\label{spaces-more-morphisms-lemma-etale-splits-off-just-one-quasi-finite-part}
$S$ をスキームとする。$f : X \to Y$ を $S$ 上の代数空間の射とする。
$x \in |X|$ を、像 $y \in |Y|$ をもつ点とする。次を仮定する：
\begin{enumerate}
\item $f$ は局所有限型である。
\item $f$ は分離的である。
\item $f$ は $x$ で準有限である。
\end{enumerate}
このとき点付き代数空間のエタール射 $(U, u) \to (Y, y)$ と、
開かつ閉な部分空間への分解
$$
U \times_Y X = W \amalg V
$$
で、射 $V \to U$ が有限であり、点 $v \in |V|$ で、$x$ へ（$|X|$ において）、
$u$ へ（$|U|$ において）写るものが存在するものがある。
\end{lemma}

\begin{proof}
スキーム $U$、点 $u \in U$、およびエタール射 $U \to Y$ で
$u$ を $y$ へ写すものを選ぶ。点 $x' \in |U \times_Y X|$ で、
$x$ へ（$|X|$ において）、$u$ へ（$|U|$ において）写るものが存在する
（Properties of Spaces, Lemma
\ref{spaces-properties-lemma-points-cartesian}）。証明を終えるには、射
$U \times_Y X \to U$ と点 $x'$ に Lemma
\ref{spaces-more-morphisms-lemma-etale-splits-off-quasi-finite-part} を適用する。
$U$ はスキームであり、したがって $u$ はモノ射
$\Spec(\kappa(u)) \to U$ から来るので、この補題を適用できる。
\end{proof}










\section{ザリスキーの主定理}
\label{spaces-more-morphisms-section-structure-quasi-finite}

\noindent
この節では、前節の結果を用いて代数空間の射に対する
ザリスキーの主定理を証明する。この節は More on Morphisms, Section
\ref{more-morphisms-section-application-etale-neighbourhoods} の類似である。

\begin{lemma}
\label{spaces-more-morphisms-lemma-finite-type-separated}
$S$ をスキームとする。$f : X \to Y$ を $S$ 上の代数空間の有限型かつ
分離的な射とする。$Y'$ を、$Y$ の $X$ における正規化とする。図式は
$$
\xymatrix{
X \ar[rd]_f \ar[rr]_{f'} & & Y' \ar[ld]^\nu \\
& Y &
}
$$
である。このとき開部分空間 $U' \subset Y'$ で、次を満たすものが存在する：
\begin{enumerate}
\item $(f')^{-1}(U') \to U'$ は同型である。
\item $(f')^{-1}(U') \subset X$ は、$f$ が準有限である点全体の集合である。
\end{enumerate}
\end{lemma}

\begin{proof}
Morphisms of Spaces, Lemma
\ref{spaces-morphisms-lemma-locally-finite-type-quasi-finite-part}
により、開部分空間 $U \subset X$ が存在し、その点は $|X|$ のうち
$f$ が準有限である点に対応する。証明すべきことは次である：
\begin{enumerate}
\item[(a)] $|U| \to |Y'|$ の像は $|U'|$ であり、ここで $U'$ は
$Y'$ のある開部分空間である。
\item[(b)] $U = f^{-1}(U')$ である。
\item[(c)] $U \to U'$ は同型である。
\end{enumerate}
$U$ の形成は任意の基底変換と可換であり
（Morphisms of Spaces, Lemma
\ref{spaces-morphisms-lemma-locally-finite-type-quasi-finite-part}）、
正規化 $Y'$ の形成は滑らかな基底変換と可換であり
（Lemma \ref{spaces-more-morphisms-lemma-normalization-smooth-localization}）、
エタール射は開であり、さらに「同型であること」は基底上 fpqc 局所的である
（Descent on Spaces, Lemma
\ref{spaces-descent-lemma-descending-property-isomorphism}）。したがって
(a), (b), (c) を $Y$ 上エタール局所的に証明すれば十分である
（細部の一部は省略する）。ゆえに $Y$ はアフィンスキームであると仮定してよい。
すると $Y'$ も（アフィン）スキームである。

\medskip\noindent
$x \in |U|$ とする。主張：開近傍 $f'(x) \in V \subset Y'$ で、
$(f')^{-1}V \to V$ が同型となるものが存在する。まず、この主張から補題が従うことを示す。
実際、このとき $(f')^{-1}V \cong V$ はスキームであり（$Y'$ の開部分として）、
$Y$ 上局所有限型である（$X$ の開部分空間として）。また $v \in V$ に対し、
剰余体拡大 $\kappa(v)/\kappa(\nu(v))$ は代数的である
（$V \subset Y'$ かつ $Y'$ は $Y$ 上整だからである）。したがって
$V \to Y$ のファイバーは離散的であり（Morphisms, Lemma
\ref{morphisms-lemma-algebraic-residue-field-extension-closed-point-fibre}）、
$(f')^{-1}V \to Y$ は局所準有限である
（Morphisms, Lemma \ref{morphisms-lemma-locally-quasi-finite-fibres}）。
これより $(f')^{-1}V \subset U$ および $V \subset U'$ が従う。
$x$ は任意であったから、(a), (b), (c) が成り立つ。

\medskip\noindent
$y = f(x) \in |Y|$ とする。$(T, t) \to (Y, y)$ を点付きスキームの
エタール射とする。下付き添字 ${}_T$ で $T$ への基底変換を表す。
$z \in X_T$ を、ファイバー $X_t$ に属し $x$ の上にある点とする。
$U_T \subset X_T$ は $f_T$ が準有限である点全体の集合であることに注意する。
Morphisms of Spaces, Lemma
\ref{spaces-morphisms-lemma-locally-finite-type-quasi-finite-part} を参照せよ。
また
$$
X_T \xrightarrow{f'_T} Y'_T \xrightarrow{\nu_T} T
$$
は $T$ の $X_T$ における正規化である。Lemma
\ref{spaces-more-morphisms-lemma-normalization-smooth-localization} を参照せよ。
$z \in U_T \subset X_T \to Y'_T \to T$ に対して主張が成り立つと仮定する。
すなわち、開近傍 $f'_T(z) \in V' \subset Y'_T$ で
$(f'_T)^{-1}V' \to V'$ が同型となるものを取れると仮定する。
射 $Y'_T \to Y'$ はエタールなので、その像 $V \subset Y'$ は、$V'$ の像として
開である。$f'(x) \in V$ であることに注意する。実際 $f'_T(z) \in V'$ である。また
$$
\xymatrix{
(f'_T)^{-1}V' \ar[r] \ar[d] & (f')^{-1}(V) \ar[d] \\
V' \ar[r] & V
}
$$
はファイバー積の図式である（$Y'_T \times_{Y'} X = X_T$ だからである）。
左の垂直射は同型であり、$\{V' \to V\}$ はエタール被覆なので、
Descent on Spaces, Lemma
\ref{spaces-descent-lemma-descending-property-isomorphism}
により右の垂直射も同型である。言い換えれば、
$x \in U \subset X \to Y' \to Y$ に対して主張が成り立つ。

\medskip\noindent
前段落の結果により、$x \in |U|$ に対する主張を証明する際、
$Y$ をエタール近傍 $T$、すなわち $y = f(x)$ の近傍で置き換え、さらに
$x$ を、$x$ の上にある $T \times_Y X$ の任意の点で置き換えてよい。
したがって、開かつ閉な部分空間への分解
$$
X = V \amalg W
$$
で、$V \to Y$ が有限かつ $x \in V$ となるものが存在すると仮定してよい。
Lemma \ref{spaces-more-morphisms-lemma-etale-splits-off-quasi-finite-part} を参照せよ。
$X$ は $V$ と $W$ の $Y$ 上の直和であり、$V \to Y$ は有限なので、
$Y$ の $X$ における正規化は次の射である：
% P09-ERR-1160: frozen source line 7699 sends the normalization factor to S rather than Y.
$$
X = V \amalg W \longrightarrow V \amalg W' \longrightarrow S
$$
ここで $W'$ は $Y$ の $W$ における正規化である。Morphisms of Spaces, Lemmas
\ref{spaces-morphisms-lemma-normalization-in-disjoint-union},
\ref{spaces-morphisms-lemma-finite-integral}, and
\ref{spaces-morphisms-lemma-normalization-in-integral}
を参照せよ。これで主張が従い、証明が完了する。
\end{proof}

\noindent
次の補題は Morphisms of Spaces, Lemma
\ref{spaces-morphisms-lemma-quasi-finite-separated-quasi-affine}
と重複している。同じ補題を二つ置く理由は、証明がいくぶん異なるためである。
以下の証明は、上で与えた代数空間の非表現可能な射に対する
ザリスキーの主定理に基づく。一方、Morphisms of Spaces, Lemma
\ref{spaces-morphisms-lemma-quasi-finite-separated-quasi-affine}
の証明は Morphisms of Spaces, Proposition
\ref{spaces-morphisms-proposition-locally-quasi-finite-separated-over-scheme}
に基づき、スキームの射の場合へ帰着する。

\begin{lemma}
\label{spaces-more-morphisms-lemma-quasi-finite-separated-quasi-affine}
$S$ をスキームとする。$f : X \to Y$ を $S$ 上の代数空間の射とする。
$f$ は準有限かつ分離的であると仮定する。$Y'$ を、$Y$ の $X$ における
正規化とする。図式は
$$
\xymatrix{
X \ar[rd]_f \ar[rr]_{f'} & & Y' \ar[ld]^\nu \\
& Y &
}
$$
である。このとき $f'$ は準コンパクトな開埋入であり、$\nu$ は整である。
特に $f$ は準アフィンである。
\end{lemma}

\begin{proof}
これは Lemma \ref{spaces-more-morphisms-lemma-finite-type-separated} から従う。実際、その補題により
開部分空間 $U' \subset Y'$ で、$(f')^{-1}(U') = X$（！）かつ
$X \to U'$ が同型となるものが存在する！ 言い換えれば、$f'$ は開埋入である。
$f'$ が準コンパクトであることに注意する。実際 $f$ は準コンパクトであり、
$\nu : Y' \to Y$ は分離的である（Morphisms of Spaces, Lemma
\ref{spaces-morphisms-lemma-quasi-compact-permanence}）。
したがって、任意のアフィンスキーム $Z$ と射 $Z \to Y$ に対し、
ファイバー積 $Z \times_Y X$ はアフィンスキーム $Z \times_Y Y'$ の
準コンパクトな開部分スキームである。ゆえに定義により $f$ は準アフィンである。
\end{proof}

\begin{lemma}[ザリスキーの主定理]
\label{spaces-more-morphisms-lemma-quasi-finite-separated-pass-through-finite}
$S$ をスキームとする。$f : X \to Y$ を $S$ 上の代数空間の射とする。
$f$ は準有限かつ分離的であり、$Y$ は準コンパクトかつ準分離的であると仮定する。
このとき分解
$$
\xymatrix{
X \ar[rd]_f \ar[rr]_j & & T \ar[ld]^\pi \\
& Y &
}
$$
で、$j$ が準コンパクトな開埋入、$\pi$ が有限となるものが存在する。
\end{lemma}

\begin{proof}
$X \to Y' \to Y$ を Lemma
\ref{spaces-more-morphisms-lemma-quasi-finite-separated-quasi-affine} の結論にあるものとする。
Limits of Spaces, Lemma
\ref{spaces-limits-lemma-integral-algebra-directed-colimit-finite}
により、
% P09-ERR-1161: frozen source line 7778 types the finite algebras over O_X rather than O_Y.
$\nu_*\mathcal{O}_{Y'} = \colim_{i \in I} \mathcal{A}_i$ を、有限な準連接
$\mathcal{O}_X$-代数 $\mathcal{A}_i \subset \nu_*\mathcal{O}_{Y'}$ の有向余極限として書ける。
$\pi_i : T_i = \underline{\Spec}_Y(\mathcal{A}_i) \to Y$
は各 $i$ に対し有限射である。遷移射 $T_{i'} \to T_i$ はアフィンであり、
$Y' = \lim T_i$ であることに注意する。

\medskip\noindent
Limits of Spaces, Lemma \ref{spaces-limits-lemma-descend-opens}
により、ある $i$ と準コンパクトな開部分 $U_i \subset T_i$ で、その
$Y'$ における逆像が $f'(X)$ に等しいものが存在する。
$i' \geq i$ に対し、$U_{i'}$ を $U_i$ の $T_{i'}$ における逆像とする。
このとき $X \cong f'(X) = \lim_{i' \geq i} U_{i'}$ である。Limits of Spaces, Lemma
\ref{spaces-limits-lemma-directed-inverse-system-has-limit} を参照せよ。
Limits of Spaces, Lemma
\ref{spaces-limits-lemma-finite-type-eventually-closed} により、
$X \to U_{i'}$ は、ある $i' \geq i$ に対し閉埋入である。
（実際、$X \cong U_{i'}$ は十分大きな $i'$ に対して成り立つが、これは不要である。）
したがって $X \to T_{i'}$ は埋入である。Morphisms of Spaces, Lemma
\ref{spaces-morphisms-lemma-factor-the-other-way}
により、これを $X \to T \to T_{i'}$ と分解できる。ここで第一の射は開埋入、
第二の射は閉埋入である。これで証明が完了する。
\end{proof}

\begin{lemma}
\label{spaces-more-morphisms-lemma-quasi-finite-separated-pass-through-finite-addendum}
Lemma \ref{spaces-more-morphisms-lemma-quasi-finite-separated-pass-through-finite} と同じ記法と仮定を用いる。
さらに $f$ は局所有限表示であると仮定する。このとき $T$ が $Y$ 上有限かつ
有限表示となるように分解を選べる。
\end{lemma}

\begin{proof}
Limits of Spaces, Lemma
\ref{spaces-limits-lemma-finite-in-finite-and-finite-presentation}
により、$T = \lim T_i$ と書ける。ここで全ての $T_i$ は $Y$ 上有限かつ有限表示であり、
遷移射 $T_{i'} \to T_i$ は閉埋入である。Limits of Spaces, Lemma
\ref{spaces-limits-lemma-descend-opens} により、ある $i$ と開部分スキーム
% P09-ERR-1162: frozen source line 7821 calls U_i an open subscheme of an algebraic space T_i.
$U_i \subset T_i$ で、その $T$ における逆像が $X$ となるものが存在する。
Limits of Spaces, Lemma
\ref{spaces-limits-lemma-finite-type-eventually-closed}
により、$X \cong U_i$ は十分大きな $i$ に対し成り立つ。
$T$ を $T_i$ で置き換えれば証明が完了する。
\end{proof}












\section{ザリスキーの主定理の応用 I}
\label{spaces-more-morphisms-section-applications-zmt}

\noindent
第一の応用は、有限射をファイバーが有限な固有射として特徴づけることである。

\begin{lemma}
\label{spaces-more-morphisms-lemma-characterize-finite}
$S$ をスキームとする。$f : X \to Y$ を $S$ 上の代数空間の射とする。
次は同値である：
\begin{enumerate}
\item $f$ は有限である。
\item $f$ は固有かつ局所準有限である。
\item $f$ は固有であり、$|X_k|$ は任意の射 $\Spec(k) \to Y$ に対し
離散空間である。ここで $k$ は体である。
\item $f$ は普遍的閉、分離的、局所有限型であり、$|X_k|$ は任意の射
$\Spec(k) \to Y$ に対し離散空間である。ここで $k$ は体である。
\end{enumerate}
\end{lemma}

\begin{proof}
Morphisms of Spaces, Lemmas \ref{spaces-morphisms-lemma-finite-proper},
\ref{spaces-morphisms-lemma-finite-quasi-finite} により
(1) $\Rightarrow$ (2) である。Morphisms of Spaces, Lemma
\ref{spaces-morphisms-lemma-locally-quasi-finite} により
(2) $\Rightarrow$ (3) である。定義により (3) は (4) を含意する。

\medskip\noindent
(4) を仮定する。$f$ は普遍的閉なので準コンパクトである
（Morphisms of Spaces, Lemma
\ref{spaces-morphisms-lemma-universally-closed-quasi-compact}）。
点 $y$ を $|Y|$ から取る。$y$ を射 $\Spec(k) \to Y$ で表す。
$|X_k|$ は準コンパクトな離散空間なので、有限離散空間である。
写像 $|X_k| \to |X|$ は、$|X| \to |Y|$ の $y$ 上のファイバーへ全射である
（Properties of Spaces, Lemma \ref{spaces-properties-lemma-points-cartesian}）。
Morphisms of Spaces, Lemma \ref{spaces-morphisms-lemma-quasi-finite-at-point}
により、$X \to Y$ は $|X| \to |Y|$ の $y$ 上のファイバーの全ての点で
準有限である。初等エタール近傍 $(U, u) \to (Y, y)$ と分解
$X_U = V \amalg W$ を、$|X|$ の $y$ の上にある全ての点に適合するように選ぶ。
Lemma \ref{spaces-more-morphisms-lemma-etale-splits-off-quasi-finite-part} を参照せよ。
$W_u = \emptyset$ であることに注意する。実際、$|X| \to |Y|$ の
$y$ 上のファイバーの全ての点を用いた。$f$ は普遍的閉なので、
$|W|$ の $|U|$ における像は $u$ を含まない閉集合である。
$U$ を縮小した後、$W = \emptyset$ と仮定してよい。言い換えれば、
$X_U = V$ は $U$ 上有限である。$y \in |Y|$ は任意であったから、
エタール射の族 $\{U_i \to Y\}$ で、その像が $Y$ を被覆し、
基底変換 $X_{U_i} \to U_i$ が有限となるものが存在する。
Morphisms of Spaces, Lemma \ref{spaces-morphisms-lemma-integral-local}
により $f$ は有限である。
\end{proof}

\noindent
帰結として、次の有用な結果を得る。

\begin{lemma}
\label{spaces-more-morphisms-lemma-proper-finite-fibre-finite-in-neighbourhood}
$S$ をスキームとする。$f : X \to Y$ を $S$ 上の代数空間の射とする。
$y \in |Y|$ とする。次を仮定する：
\begin{enumerate}
\item $f$ は固有である。
\item $f$ は、全ての $x \in |X|$ で準有限であり、それらの点は $y$ の上にある
（Decent Spaces, Lemma \ref{decent-spaces-lemma-conditions-on-fibre-and-qf}）。
\end{enumerate}
このとき開近傍 $V \subset Y$ で $y$ を含み、
$f|_{f^{-1}(V)} : f^{-1}(V) \to V$ が有限となるものが存在する。
\end{lemma}

\begin{proof}
Morphisms of Spaces, Lemma
\ref{spaces-morphisms-lemma-locally-finite-type-quasi-finite-part}
により、$f$ が準有限である点全体は開集合 $U \subset X$ である。
$Z = X \setminus U$ とおく。このとき $y \not \in f(Z)$ である。
$f$ は固有なので、集合 $f(Z) \subset Y$ は閉である。任意の開近傍
$V \subset Y$ で $y$ を含み、次を満たすものを選ぶ：
% P09-ERR-1163: frozen source line 7925 intersects Z subset X directly with V subset Y.
$Z \cap V = \emptyset$。すると $f^{-1}(V) \to V$ は局所準有限かつ固有である。
したがって Lemma \ref{spaces-more-morphisms-lemma-characterize-finite} により
$f^{-1}(V) \to V$ は有限である。
\end{proof}

\begin{lemma}
\label{spaces-more-morphisms-lemma-flat-proper-family-cannot-collapse-fibre}
\begin{slogan}
固有な族の一つのファイバーを潰すと、近傍のファイバーも潰れざるを得ない。
\end{slogan}
$S$ をスキームとする。代数空間の射からなる可換図式
% P09-ERR-1164: frozen source line 7942 uses singular "morphism" for a diagram containing several morphisms.
$$
\xymatrix{
X \ar[rr]_h \ar[rd]_f & & Y \ar[ld]^g \\
& B
}
$$
を考える。この図式は $S$ 上にある。$b \in B$ とし、$\Spec(k) \to B$ を
$b$ の同値類に属する射とする。
次を仮定する：
\begin{enumerate}
\item $X \to B$ は固有射である。
\item $Y \to B$ は分離的かつ局所有限型である。
\item 次のいずれかが成り立つ：
\begin{enumerate}
\item $|X_k| \to |Y_k|$ の像は有限である。
\item $|f|^{-1}(\{b\})$ の $|Y|$ における像は有限であり、
$B$ は decent である。
\end{enumerate}
\end{enumerate}
このとき開部分空間 $B' \subset B$ で $b$ を含み、
$X_{B'} \to Y_{B'}$ が閉部分空間 $Z \subset Y_{B'}$ を経由して分解し、
この閉部分空間が $B'$ 上有限となるものが存在する。
\end{lemma}

\begin{proof}
$Z \subset Y$ を $h$ のスキーム論的像とする。Morphisms of Spaces, Section
\ref{spaces-morphisms-section-scheme-theoretic-image} を参照せよ。
Morphisms of Spaces, Lemma
\ref{spaces-morphisms-lemma-scheme-theoretic-image-is-proper}
により、$X \to Z$ は全射であり、$Z \to B$ は固有である。したがって
$$
\{x \in |X|\text{ lying over }b\} \to
\{z \in |Z|\text{ lying over }b\}
$$
および $|X_k| \to |Z_k|$ は全射である。(3)(a) または (3)(b) のいずれからも、
% P09-ERR-1165: frozen source lines 7973--7974 omit "at" after "quasi-finite".
$Z \to B$ は $|Z|$ の $b$ の上にある全ての点で準有限であることが分かる。
Decent Spaces, Lemma \ref{decent-spaces-lemma-conditions-on-fibre-and-qf}
を参照せよ。ゆえに Lemma
\ref{spaces-more-morphisms-lemma-proper-finite-fibre-finite-in-neighbourhood} により、
$Z \to B$ は $b$ の開近傍上有限である。
\end{proof}








\section{シュタイン分解}
\label{spaces-more-morphisms-section-stein-factorization}

\noindent
シュタイン分解とは、固有射 $f : X \to S$ が
$f_*\mathcal{O}_X = \mathcal{O}_S$ を満たすならば連結なファイバーをもつ、
という主張である。

\begin{lemma}
\label{spaces-more-morphisms-lemma-stein-universally-closed}
$S$ をスキームとする。$f : X \to Y$ を、$S$ 上の代数空間の普遍的閉かつ
準分離的な射とする。次の分解が存在する：
$$
\xymatrix{
X \ar[rr]_{f'} \ar[rd]_f & & Y' \ar[dl]^\pi \\
& Y &
}
$$
これは次の性質をもつ：
\begin{enumerate}
\item 射 $f'$ は普遍的閉、準コンパクト、準分離的、かつ全射である。
\item 射 $\pi : Y' \to Y$ は整である。
\item $f'_*\mathcal{O}_X = \mathcal{O}_{Y'}$ である。
\item $Y' = \underline{\Spec}_Y(f_*\mathcal{O}_X)$ である。
\item $Y'$ は、$Y$ の $X$ における正規化であり、Morphisms of Spaces, Definition
\ref{spaces-morphisms-definition-normalization-X-in-Y}
で定義されたものである。
\end{enumerate}
分解 $f = \pi \circ f'$ の形成は平坦な基底変換と可換である。
\end{lemma}

\begin{proof}
Morphisms of Spaces, Lemma
\ref{spaces-morphisms-lemma-universally-closed-quasi-compact}
により、射 $f$ は準コンパクトである。$Y'$ を $Y$ の $X$ における正規化として
定義するだけで、(5) と (2) は自動的に成り立つ。Morphisms of Spaces, Lemma
\ref{spaces-morphisms-lemma-normalization-in-universally-closed}
により (4) が成り立つ。射 $f'$ は Morphisms of Spaces, Lemma
\ref{spaces-morphisms-lemma-universally-closed-permanence}
により普遍的閉である。また Morphisms of Spaces, Lemma
\ref{spaces-morphisms-lemma-quasi-compact-permanence}
により準コンパクトであり、Morphisms of Spaces, Lemma
\ref{spaces-morphisms-lemma-compose-after-separated}
により準分離的である。

\medskip\noindent
残る主張を示すには、基底 $Y$ がアフィンであると仮定してよい
（正規化を取ることはエタール局所化と可換だからである）。
$Y = \Spec(R)$ と書く。このとき $Y' = \Spec(A)$ であり、
$A = \Gamma(X, \mathcal{O}_X)$ は整な $R$-代数である。したがって
$f'_*\mathcal{O}_X$ が $\mathcal{O}_{Y'}$ であることは明らかである
（$f'_*\mathcal{O}_X$ は Morphisms of Spaces, Lemma
\ref{spaces-morphisms-lemma-pushforward} により準連接であり、ゆえに
$\widetilde{A}$ に等しい）。これで (3) が証明された。

\medskip\noindent
$f'$ が全射であることを示そう。$f'$ は普遍的閉なので（上を参照）、
$f'$ の像は閉部分集合 $V(I) \subset Y' = \Spec(A)$ である。
$h \in I$ を取る。このとき
% P09-ERR-1166: frozen source line 8052 uses f^sharp although h is a function on Y', not Y.
$h|_X = f^\sharp(h)$ は $X$ の構造層の大域切断で、全ての点で消える。
$X$ は準コンパクトなので、これは $h|_X$ が冪零切断であることを意味する。すなわち
% P09-ERR-1167: frozen source line 8054 omits the restriction underscore in h^n|X.
$h^n|X = 0$ となる $n > 0$ が存在する。しかし
$A = \Gamma(X, \mathcal{O}_X)$ なので $h^n = 0$ である。言い換えれば、
% P09-ERR-1168: frozen source line 8056 says Jacobson radical although the argument proves nilradical containment.
$I$ は $A$ のジェイコブソン根基に含まれる。したがって、望むとおり
$V(I) = Y'$ である。
\end{proof}

\begin{lemma}
\label{spaces-more-morphisms-lemma-stein-universally-closed-residue-fields}
Lemma \ref{spaces-more-morphisms-lemma-stein-universally-closed} において、さらに $f$ は局所有限型で、
$Y$ はアフィンであると仮定する。このとき $y \in Y$ に対するファイバー
$\pi^{-1}(\{y\}) = \{y_1, \ldots, y_n\}$ は有限であり、体拡大
$\kappa(y_i)/\kappa(y)$ は有限である。
\end{lemma}

\begin{proof}
$\pi^{-1}(\{y\})$ の点の間には特殊化が存在しないことを思い出そう。
Algebra, Lemma \ref{algebra-lemma-integral-no-inclusion} を参照せよ。
$f'$ は全射なので、$|X_y| \to \pi^{-1}(\{y\})$ は全射である。
$X_y$ は体上有限型の準分離的代数空間であることに注意する
（準コンパクト性は参照した補題の証明で示された）。したがって
$|X_y|$ はネーター位相空間である（Morphisms of Spaces, Lemma
\ref{spaces-morphisms-lemma-finite-presentation-noetherian}）。
ここで位相的な議論（省略）により $\pi^{-1}(\{y\})$ は有限である。
各 $i$ に対し、有限型の点 $x_i \in |X_y|$ で $y_i$ へ写るものを選べる
（Morphisms of Spaces, Lemma
\ref{spaces-morphisms-lemma-enough-finite-type-points}）。
$\kappa(y_i)/\kappa(y)$ が有限であると結論する。実際、$x_i$ は有限型射
$\Spec(k_i) \to X_y$ で表せる（有限型の点の定義による）。したがって
$\Spec(k_i) \to y = \Spec(\kappa(y))$ は有限型である
（有限型射の合成だからである）。ゆえに $k_i/\kappa(y)$ は有限である
（Morphisms, Lemma \ref{morphisms-lemma-point-finite-type}）。
\end{proof}

\noindent
$f : X \to Y$ を代数空間の射とし、
$\overline{y} : \Spec(k) \to Y$ を幾何学的点とする。このとき
{\it $f$ の $\overline{y}$ 上のファイバー}とは、代数空間
$X_{\overline{y}} = X \times_{Y, \overline{y}} \Spec(k)$ のことであり、これは $k$ 上にある。
$Y$ がスキームで $y \in Y$ が点である場合は、通常どおりファイバーを
$X_y = X \times_Y \Spec(\kappa(y))$ と表す。

\begin{lemma}
\label{spaces-more-morphisms-lemma-characterize-geometrically-connected-fibres}
$S$ をスキームとする。$f : X \to Y$ を $S$ 上の代数空間の射とする。
$\overline{y}$ を $Y$ の幾何学的点とする。このとき $X_{\overline{y}}$ が
連結であることと、任意のエタール近傍
$(V, \overline{v}) \to (Y, \overline{y})$ で $V$ がスキームであるものに対し、基底変換
% P09-ERR-1169: frozen source line 8104 drops the bar from the geometric-point fibre X_{\bar v}.
$X_V \to V$ が連結なファイバー $X_v$ をもつこととは同値である。
\end{lemma}

\begin{proof}
$\overline{y}$ のエタール近傍の圏は余フィルター的であり、スキームからなる
余終集合を含む（Properties of Spaces, Lemma
\ref{spaces-properties-lemma-cofinal-etale}）。したがって $Y$ をこれらの近傍の
一つで置き換え、$Y$ がスキームであると仮定してよい。
$y \in Y$ を $\overline{y}$ に対応する点とする。このとき $X_y$ が
$\kappa(y)$ 上幾何学的に連結であることと、$X_{\overline{y}}$ が連結であること、
および $(X_y)_{k'}$ が連結であることとは同値である。ここで $k'$ は
$\kappa(y)$ の任意の有限分離拡大である。Spaces over Fields, Section
\ref{spaces-over-fields-section-geometrically-connected}、特に Lemma
\ref{spaces-over-fields-lemma-characterize-geometrically-disconnected} を参照せよ。
More on Morphisms, Lemma
\ref{more-morphisms-lemma-realize-prescribed-residue-field-extension-etale}
により、アフィンなエタール近傍 $(V, v) \to (Y, y)$ で、
% P09-ERR-1170: frozen source line 8122 uses undefined s and u instead of the chosen y and v.
$\kappa(s) \subset \kappa(u)$ が、与えられた任意の有限分離拡大
$\kappa(s) \subset k'$ と同一視されるものが存在する。これで補題が従う。
\end{proof}

\begin{theorem}[シュタイン分解；ネーターの場合]
\label{spaces-more-morphisms-theorem-stein-factorization-Noetherian}
$S$ をスキームとする。$f : X \to Y$ を $S$ 上の代数空間の固有射とし、
$Y$ は局所ネーターであるとする。次の分解が存在する：
$$
\xymatrix{
X \ar[rr]_{f'} \ar[rd]_f & & Y' \ar[dl]^\pi \\
& Y &
}
$$
これは次の性質をもつ：
\begin{enumerate}
\item 射 $f'$ は固有であり、幾何学的ファイバーは連結である。
\item 射 $\pi : Y' \to Y$ は有限である。
\item $f'_*\mathcal{O}_X = \mathcal{O}_{Y'}$ である。
\item $Y' = \underline{\Spec}_Y(f_*\mathcal{O}_X)$ である。
\item $Y'$ は $Y$ の $X$ における正規化である。Morphisms, Definition
% P09-ERR-1171: frozen source line 8144 cites the schemes Morphisms definition rather than Morphisms of Spaces.
\ref{morphisms-definition-normalization-X-in-Y} を参照せよ。
\end{enumerate}
\end{theorem}

\begin{proof}
$f = \pi \circ f'$ を Lemma \ref{spaces-more-morphisms-lemma-stein-universally-closed} の分解とする。
Lemma \ref{spaces-more-morphisms-lemma-stein-universally-closed} の結論に加えて、$f'$ は分離的であり
（Morphisms of Spaces, Lemma
\ref{spaces-morphisms-lemma-compose-after-separated}）、有限型でもある
（Morphisms of Spaces, Lemma
\ref{spaces-morphisms-lemma-permanence-finite-type}）。したがって $f'$ は固有である。
Cohomology of Spaces, Lemma
\ref{spaces-cohomology-lemma-proper-pushforward-coherent}
により、$f_*\mathcal{O}_X$ は連接 $\mathcal{O}_Y$-加群である。
したがって $\pi$ は有限であり、すなわち (2) が成り立つ。

\medskip\noindent
これで、最も興味深い主張、すなわち $f'$ の幾何学的ファイバーが連結であること
以外は全て証明された。上の議論から、$Y$ を $Y'$ で置き換えてよいことは明らかである。
すると $Y$ は局所ネーターであり、$f : X \to Y$ は固有で、
$f_*\mathcal{O}_X = \mathcal{O}_Y$ である。$\overline{y}$ を $Y$ の幾何学的点とする。
ここで形式関数定理、より正確には Cohomology of Spaces, Lemma
\ref{spaces-cohomology-lemma-formal-functions-stalk} を適用する。これにより
$$
\mathcal{O}^\wedge_{Y, \overline{y}} =
\lim_n H^0(X_n, \mathcal{O}_{X_n})
$$
を得る。ここで
$X_n = \Spec(\mathcal{O}_{Y, \overline{y}}/\mathfrak m_{\overline{y}}^n) \times_Y X$ である。
$X_1 = X_{\overline{y}} \to X_n$ は（有限次数の）肥厚化であることに注意する。
したがって $X_n$ の台位相空間は $X_{\overline{y}}$ のものに等しい。
ゆえに、$X_{\overline{y}} = T_1 \amalg T_2$ が空でない開かつ閉な部分空間の
直和であれば、同様に $X_n = T_{1, n} \amalg T_{2, n}$ となる。これは全ての
$n$ に対して成り立つ。すると
$H^0(X_n, \mathcal{O}_{X_n})$ は非自明な冪等元 $e_{1, n}$ を含む。
これは恒等値 $1$ を $T_{1, n}$ 上でとり、恒等値 $0$ を $T_{2, n}$ 上でとる関数である。
$e_{1, n + 1}$ の制限が $e_{1, n}$ となることは、$X_n$ 上で明らかである。
したがって $e_1 = \lim e_{1, n}$ は極限の非自明な冪等元である。
これは $\mathcal{O}^\wedge_{Y, \overline{y}}$ が局所環であることに矛盾する。
よって仮定は誤りであり、望むとおり $X_{\overline{y}}$ は連結である。
\end{proof}

\begin{theorem}[シュタイン分解；一般の場合]
\label{spaces-more-morphisms-theorem-stein-factorization-general}
$S$ をスキームとする。$f : X \to Y$ を $S$ 上の代数空間の固有射とする。
次の分解が存在する：
$$
\xymatrix{
X \ar[rr]_{f'} \ar[rd]_f & & Y' \ar[dl]^\pi \\
& Y &
}
$$
これは次の性質をもつ：
\begin{enumerate}
\item 射 $f'$ は固有であり、幾何学的ファイバーは連結である。
\item 射 $\pi : Y' \to Y$ は整である。
\item $f'_*\mathcal{O}_X = \mathcal{O}_{Y'}$ である。
\item $Y' = \underline{\Spec}_Y(f_*\mathcal{O}_X)$ である。
\item $Y'$ は $Y$ の $X$ における正規化である
（Morphisms of Spaces, Definition
\ref{spaces-morphisms-definition-normalization-X-in-Y}）。
\end{enumerate}
\end{theorem}

\begin{proof}
Lemma \ref{spaces-more-morphisms-lemma-stein-universally-closed} を適用して射 $f' : X \to Y'$ を得る。
Lemma \ref{spaces-more-morphisms-lemma-stein-universally-closed} の結論に加えて、$f'$ は分離的であり
（Morphisms of Spaces, Lemma
\ref{spaces-morphisms-lemma-compose-after-separated}）、有限型でもある
（Morphisms of Spaces, Lemma
\ref{spaces-morphisms-lemma-permanence-finite-type}）。したがって $f'$ は固有である。
この時点で、$f'$ が連結な幾何学的ファイバーをもつという主張を除き、
全ての主張を証明した。

\medskip\noindent
議論から、$Y$ を $Y'$ で置き換えてよいことは明らかである。すると
$f : X \to Y$ は固有で、$f_*\mathcal{O}_X = \mathcal{O}_Y$ である。
これらの条件は平坦な基底変換で保たれることに注意する
（Morphisms of Spaces, Lemma \ref{spaces-morphisms-lemma-base-change-proper}
および Cohomology of Spaces, Lemma
\ref{spaces-cohomology-lemma-flat-base-change-cohomology}）。
$\overline{y}$ を $Y$ の幾何学的点とする。Lemma
\ref{spaces-more-morphisms-lemma-characterize-geometrically-connected-fibres}
と今述べたことにより、$Y$ がスキーム、$y \in Y$ が点であり、
% P09-ERR-1172: frozen source line 8243 calls the morphism f an algebraic space over Y.
$f : X \to Y$ が $Y$ 上の固有な代数空間で、
$f_*\mathcal{O}_X = \mathcal{O}_Y$ である場合に帰着する。この場合、
ファイバー $X_y$ が連結であることを示せばよい。$Y$ を $y$ のアフィン近傍で
置き換えることにより、$Y = \Spec(R)$ がアフィンであると仮定してよい。
このとき $f_*\mathcal{O}_X = \mathcal{O}_Y$ は、環準同型
$R \to \Gamma(X, \mathcal{O}_X)$ が全単射であることを意味する。

\medskip\noindent
Limits of Spaces, Lemma
\ref{spaces-limits-lemma-proper-limit-of-proper-finite-presentation-noetherian}
により、$(X \to Y) = \lim (X_i \to Y_i)$ と書ける。ここで
$X_i \to Y_i$ は固有かつ有限表示で、$Y_i$ はネーターである。十分大きな
$i$ に対し $Y_i$ はアフィンである（Limits of Spaces, Lemma
\ref{spaces-limits-lemma-limit-is-affine}）。$Y_i = \Spec(R_i)$ と書き、
$R'_i = \Gamma(X_i, \mathcal{O}_{X_i})$ とおく。環準同型
$R_i \to R_i' \to R$ があることに注意する。第一の準同型があるのは
$X_i$ が $R_i$ 上の代数空間だからであり、第二の準同型があるのは
$X \to X_i$ があり、$R = \Gamma(X, \mathcal{O}_X)$ だからである。また
$R = \colim R'_i$ であることに注意する。Limits of Spaces, Lemma
\ref{spaces-limits-lemma-descend-section} を参照せよ。このとき
$$
\xymatrix{
X \ar[d]  \ar[r] & X_i \ar[d] \\
Y \ar[r] & Y'_i \ar[r] & Y_i
}
$$
は可換であり、$Y'_i = \Spec(R'_i)$ である。$y'_i \in Y'_i$ を $y$ の像とする。
$X_y = \lim X_{i, y'_i}$ である。実際、$X = \lim X_i$、$Y = \lim Y'_i$、かつ
$\kappa(y) = \colim \kappa(y'_i)$ である。
いま $X_y = U \amalg V$ とし、$U$ と $V$ は開かつ閉であるとする。
このとき $U, V$ は開集合 $U_i, V_i$ の逆像であり、これらは
$X_{i, y'_i}$ にある
（Limits of Spaces, Lemma \ref{spaces-limits-lemma-descend-opens}）。
Theorem \ref{spaces-more-morphisms-theorem-stein-factorization-Noetherian} により
$X_i \to Y'_i$ のファイバーは連結なので、$U$ または $V$ の一方は空である。
これで証明が完了する。
\end{proof}

\noindent
次は一つの応用である。

\begin{lemma}
\label{spaces-more-morphisms-lemma-geometrically-connected-fibres-towards-normal}
$S$ をスキームとする。$f : X \to Y$ を $S$ 上の代数空間の射とする。
次を仮定する：
\begin{enumerate}
\item $f$ は固有である。
\item $Y$ は整であり（Spaces over Fields, Definition
\ref{spaces-over-fields-definition-integral-algebraic-space}）、その生成点を
$\xi$ とする。
\item $Y$ は正規である。
\item $X$ は被約である。
\item $|X|$ の各既約成分の全ての生成点は $\xi$ へ写る。
\item $H^0(X_\xi, \mathcal{O}) = \kappa(\xi)$ である。
\end{enumerate}
このとき $f_*\mathcal{O}_X = \mathcal{O}_Y$ であり、$f$ は幾何学的に連結な
ファイバーをもつ。
\end{lemma}

\begin{proof}
Theorem \ref{spaces-more-morphisms-theorem-stein-factorization-general} を適用し、分解
$X \to Y' \to Y$ を得る。$Y' = Y$ を示せば十分である。
$Y' \times_Y V \to V$ が同型であることを示せば十分である。ここで
$V \to Y$ はエタール射で、$V$ はアフィン整スキームである。
Spaces over Fields, Lemma
\ref{spaces-over-fields-lemma-normal-integral-cover-by-affines} を参照せよ。
$Y'$ の形成はエタールな基底変換と可換である。Morphisms of Spaces, Lemma
\ref{spaces-morphisms-lemma-properties-normalization} を参照せよ。
$X \times_Y V$ の生成点は $X$ の生成点の上にある
（Decent Spaces, Lemma \ref{decent-spaces-lemma-decent-generic-points}）。
したがって仮定 (5) により $V$ の生成点へ写る。さらに条件 (6) は
$V \to Y$ による基底変換で保たれる。例えば平坦基底変換
（Cohomology of Spaces, Lemma
\ref{spaces-cohomology-lemma-flat-base-change-cohomology}）を用いればよい。
ゆえに $Y$ が正規整アフィンスキームである場合に補題を証明すれば十分である。

\medskip\noindent
$Y$ が正規整アフィンスキームであると仮定する。Morphisms, Lemma
\ref{morphisms-lemma-finite-birational-over-normal} を適用して
$Y' \to Y$ が同型であることを示す。実際、$Y'$ は被約である。これは $X$ が被約だからである
（Morphisms of Spaces, Lemma
\ref{spaces-morphisms-lemma-normalization-in-reduced}）。上で引用した定理により、
射 $Y' \to Y$ は整である。$Y$ は decent であり、$X \to Y$ は分離的なので、
$X$ も decent である。これを見るには Decent Spaces, Lemmas
\ref{decent-spaces-lemma-properties-trivial-implications} および
\ref{decent-spaces-lemma-property-over-property} を用いる。仮定 (5)、
Morphisms of Spaces, Lemma \ref{spaces-morphisms-lemma-normalization-generic}、
および Decent Spaces, Lemma \ref{decent-spaces-lemma-decent-generic-points}
により、$|Y'|$ の各既約成分の全ての生成点は $\xi$ へ写る。
一方、$Y'$ は $f_*\mathcal{O}_X$ の相対スペクトルなので、スキーム論的ファイバー
$Y'_\xi$ は $H^0(X_\xi, \mathcal{O})$ のスペクトルであり、仮定により
$\kappa(\xi)$ に等しい。したがって $Y'$ は整スキームで、その関数体は
$Y$ の関数体に等しい。これで証明が完了する。
\end{proof}

\noindent
次はもう一つの応用である。

\begin{lemma}
\label{spaces-more-morphisms-lemma-proper-flat-nr-geom-conn-comps-lower-semicontinuous}
$S$ をスキームとする。$X \to Y$ を $S$ 上の代数空間の射とする。
% P09-ERR-1173: frozen source lines 8352--8355 use f without naming X -> Y as f.
$f$ が固有、平坦、かつ有限表示ならば、関数
$n_{X/Y} : |Y| \to \mathbf{Z}$、すなわち $f$ のファイバーの
幾何学的連結成分の個数を数えるもの
（Lemma \ref{spaces-more-morphisms-lemma-base-change-fibres-nr-geometrically-connected-components}）
は下半連続である。
\end{lemma}

\begin{proof}
問題は $Y$ 上エタール局所的なので、$Y$ はアフィンスキームであると仮定してよい。
$y \in Y$ とし、
% P09-ERR-1174: frozen source lines 8362 and 8366 switch from n_{X/Y} to n_{X/S}.
$n = n_{X/S}(y)$ とおく。$n < \infty$ であることに注意する。実際、
$X \to Y$ の $y$ における幾何学的ファイバーは体上の固有代数空間であり、
したがってネーターであり、ゆえに連結成分の個数は有限である。
開近傍 $V$ で $y$ を含み、$n_{X/S}|_V \geq n$ となるものを見つけなければならない。
$X \to Y' \to Y$ を Theorem
\ref{spaces-more-morphisms-theorem-stein-factorization-general} にあるシュタイン分解とする。
Lemma \ref{spaces-more-morphisms-lemma-stein-universally-closed-residue-fields} により、点
$y'_1, \ldots, y'_m \in Y'$ は $y$ の上にあり、その個数は有限である。また拡大
$\kappa(y'_i)/\kappa(y)$ は有限である。More on Morphisms, Lemma
\ref{more-morphisms-lemma-etale-makes-integral-split}
によれば、$Y$ を $y$ のエタール近傍で置き換えた後、
$Y' = V_1 \amalg \ldots \amalg V_m$ がスキームとなり、
$y'_i \in V_i$ かつ $\kappa(y'_i)/\kappa(y)$ が純非分離となると仮定してよい。
すると代数空間 $X_{y_i'}$ は $\kappa(y)$ 上幾何学的に連結であり、
したがって $m = n$ である。代数空間
$X_i = (f')^{-1}(V_i)$、$i = 1, \ldots, n$ は $Y$ 上平坦かつ有限表示である。
したがって $X_i \to Y$ の像は開である（Morphisms of Spaces, Lemma
\ref{spaces-morphisms-lemma-fppf-open}）。ゆえに $y$ の近傍において
$n_{X/Y}$ は少なくとも $n$ である。
\end{proof}

\begin{lemma}
\label{spaces-more-morphisms-lemma-proper-flat-geom-red}
$S$ をスキームとする。$f : X \to Y$ を $S$ 上の代数空間の射とする。
次を仮定する：
\begin{enumerate}
\item $f$ は固有、平坦、かつ有限表示である。
\item $f$ の幾何学的ファイバーは被約である。
\end{enumerate}
% P09-ERR-1175: frozen source lines 8394 and 8418 use n_{X/S} instead of n_{X/Y}.
このとき、関数 $n_{X/S} : |Y| \to \mathbf{Z}$、すなわち $f$ のファイバーの
幾何学的連結成分の個数を数えるもの
（Lemma \ref{spaces-more-morphisms-lemma-base-change-fibres-nr-geometrically-connected-components}）
は局所定数である。
\end{lemma}

\begin{proof}
Lemma \ref{spaces-more-morphisms-lemma-proper-flat-nr-geom-conn-comps-lower-semicontinuous}
により、関数 $n_{X/Y}$ は下半連続である。
% P09-ERR-1176: frozen source line 8403 misspells semicon(t)inuous as semincontinuous.
したがって、上半連続であることを示せば十分である。そのためには
$Y$ 上エタール局所的に議論してよいので、$Y$ がアフィンスキームであると
仮定してよい。$y \in Y$ に対し、$\kappa(y)$-代数
$$
A = H^0(X_y, \mathcal{O}_{X_y})
$$
を考える。Spaces over Fields, Lemma
\ref{spaces-over-fields-lemma-proper-geometrically-reduced-global-sections}
および $X_y$ が幾何学的に被約であることにより、$A$ は
$\kappa(y)$ の有限分離拡大の有限直積である。したがって
$A \otimes_{\kappa(y)} \kappa(\overline{y})$ は
$\beta_0(y) = \dim_{\kappa(y)} A$ 個の $\kappa(\overline{y})$ のコピーの直積である。
ゆえに $X_{\overline{y}}$ は $\beta_0(y)$ 個の連結成分をもつ。
言い換えれば、$n_{X/S} = \beta_0$ である。これは $Y$ 上の関数としての等式である。
Derived Categories of Spaces, Lemma
\ref{spaces-perfect-lemma-jump-loci-geometric} により、$n_{X/Y}$ は上半連続である。
これで証明が完了する。
\end{proof}

\begin{lemma}
\label{spaces-more-morphisms-lemma-stein-factorization-etale}
$S$ をスキームとする。$f : X \to Y$ を $S$ 上の代数空間の固有射とする。
$X \to Y' \to Y$ を $f$ のシュタイン分解とする
（Theorem \ref{spaces-more-morphisms-theorem-stein-factorization-general}）。
$f$ が有限表示、平坦、かつ幾何学的に被約なファイバーをもつならば
（Definition \ref{spaces-more-morphisms-definition-geometrically-reduced-fibre}）、
$Y' \to Y$ は有限エタールである。
\end{lemma}

\begin{proof}
シュタイン分解の形成は平坦基底変換と可換である。Lemma
\ref{spaces-more-morphisms-lemma-stein-universally-closed} を参照せよ。したがって
$Y$ 上エタール局所的に議論してよく、$Y$ がアフィンスキームであると仮定してよい。
すると $Y'$ はアフィンスキームで、$Y' \to Y$ は整である。

\medskip\noindent
$y \in Y$ とする。$n$ を幾何学的ファイバー $X_{\overline{y}}$ の連結成分の個数とする。
$n < \infty$ であることに注意する。実際、$X \to Y$ の $y$ における
幾何学的ファイバーは体上の固有代数空間なのでネーターであり、
したがって連結成分の個数は有限である。Lemma
\ref{spaces-more-morphisms-lemma-stein-universally-closed-residue-fields} により、点
$y'_1, \ldots, y'_m \in Y'$ は有限個で $y$ の上にあり、各 $i$ に対して
% P09-ERR-1177: frozen source line 8452 omits a conjunction before the finite-extension clause.
有限型の点 $x_i \in |X_y|$ で $y'_i$ へ写るものを選べ、拡大
$\kappa(y'_i)/\kappa(y)$ は有限である。したがって More on Morphisms, Lemma
\ref{more-morphisms-lemma-etale-makes-integral-split}
により、$Y$ を $y$ のエタール近傍で置き換えた後、
$Y' = V_1 \amalg \ldots \amalg V_m$ がスキームとなり、
$y'_i \in V_i$ かつ $\kappa(y'_i)/\kappa(y)$ が純非分離となると仮定してよい。
この場合、代数空間 $X_{y_i'}$ は $\kappa(y)$ 上幾何学的に連結なので、
$m = n$ である。代数空間 $X_i = (f')^{-1}(V_i)$、$i = 1, \ldots, n$ は
$Y$ 上固有、平坦、有限表示で、幾何学的に被約なファイバーをもつ。
各射 $X_i \to Y$ に対して補題を証明すれば十分である。これにより
$X_{\overline{y}}$ が連結である場合へ帰着される。

\medskip\noindent
$X_{\overline{y}}$ が連結であると仮定する。Lemma
\ref{spaces-more-morphisms-lemma-proper-flat-geom-red} により、$X \to Y$ は $y$ の近傍で
幾何学的に連結なファイバーをもつ。したがって $X \to Y$ のファイバーが
幾何学的に連結であると仮定してよい。このとき Derived Categories of Spaces, Lemma
\ref{spaces-perfect-lemma-proper-flat-geom-red-connected} により
$f_*\mathcal{O}_X = \mathcal{O}_Y$ である。これで証明が完了する。
\end{proof}

\noindent
次の補題の証明がスキームに対するシュタイン分解を用いるため、
この節に置かれている。

\begin{lemma}
\label{spaces-more-morphisms-lemma-split-off-proper-part-henselian}
$(A, I)$ をヘンゼル対とする。$X$ を $A$ 上分離的かつ有限型の代数空間とする。
$X_0 = X \times_{\Spec(A)} \Spec(A/I)$ とおく。
$Y \subset X_0$ を開かつ閉な部分空間で、$Y \to \Spec(A/I)$ が固有となるものとする。
このとき、開かつ閉な部分空間 $W \subset X$ で、$A$ 上固有かつ
$W \times_{\Spec(A)} \Spec(A/I) = Y$ となるものが存在する。
\end{lemma}

\begin{proof}
基底変換を $T \mapsto T_0$ と表す。これは $\Spec(A/I) \to \Spec(A)$ による。
チャウの補題の弱い版（Cohomology of Spaces, Lemma
\ref{spaces-cohomology-lemma-weak-chow} の形）により、全射固有射
$\varphi : X' \to X$ で、$X'$ が $\mathbf{P}^n_A$ への埋入をもつものが存在する。
$Y' = \varphi^{-1}(Y)$ とおく。これは $X'_0$ の開かつ閉な部分スキームである。
More on Morphisms, Lemma
\ref{more-morphisms-lemma-split-off-proper-part-henselian}
により、補題は $(X', Y')$ に対して成り立つ。$W' \subset X'$ を $A$ 上固有な
開かつ閉な部分スキームで、$Y' = W'_0$ となるものとする。
Morphisms of Spaces, Lemma
\ref{spaces-morphisms-lemma-universally-closed-permanence} により
$Q_1 = \varphi(|W'|) \subset |X|$ および
$Q_2 = \varphi(|X' \setminus W'|) \subset |X|$ は閉部分集合である。
また Morphisms of Spaces, Lemma
\ref{spaces-morphisms-lemma-image-proper-is-proper} により、$Q_1$ 上の任意の
閉部分空間構造は $A$ 上固有である。$Q_1 \cap Q_2$ の $\Spec(A)$ における像は
閉である。$(A, I)$ はヘンゼル的なので、$Q_1 \cap Q_2$ が空でなければ、
$Q_1 \cap Q_2$ は $\Spec(A/I)$ の上に点をもつ。これは
$W'_0 = Y' = \varphi^{-1}(Y)$ なので不可能である。したがって $Q_1$ は
$|X|$ の開かつ閉な部分集合である。$W \subset X$ を対応する開かつ閉な
部分空間とする。このとき $W$ は $A$ 上固有で、$W_0 = Y$ である。
\end{proof}










\section{開部分からの性質の延長}
\label{spaces-more-morphisms-section-extending-properties}

\noindent
この節では、次の形の結果をいくつか集める：$f : X \to Y$ が代数空間の
平坦射であり、$f$ が $Y$ の稠密開部分上である性質を満たすならば、$f$ は
$Y$ 全体上で同じ性質を満たす。

\begin{lemma}
\label{spaces-more-morphisms-lemma-flat-finite-type-finitely-presented-over-dense-open}
$S$ をスキームとする。$f : X \to Y$ を $S$ 上の代数空間の射とする。
$\mathcal{F}$ を準連接 $\mathcal{O}_X$-加群とする。
$V \subset Y$ を開部分空間とする。次を仮定する：
\begin{enumerate}
\item $f$ は局所有限表示である。
\item $\mathcal{F}$ は有限型であり、$Y$ 上平坦である。
\item $V \to Y$ は準コンパクトかつスキーム論的に稠密である。
\item $\mathcal{F}|_{f^{-1}V}$ は有限表示である。
\end{enumerate}
このとき $\mathcal{F}$ は有限表示である。
\end{lemma}

\begin{proof}
$\mathcal{F}$ の、$X$ 上全射エタールなスキームへの引き戻しが有限表示であることを
証明すれば十分である。したがって $X$ はスキームであると仮定してよい。
同様に、$Y$ を $Y$ 上全射エタールなスキームで置き換えてよい
% P09-ERR-1178: frozen source line 8559 contains an extra "and" in "etale and over Y".
（このスキームにおける $V$ の逆像はスキーム論的に稠密である。Morphisms of Spaces,
Section \ref{spaces-morphisms-section-scheme-theoretic-closure} を参照せよ）。
こうしてスキームの場合へ帰着される。これは More on Flatness, Lemma
\ref{flat-lemma-flat-finite-type-finitely-presented-over-dense-open} である。
\end{proof}

\begin{lemma}
\label{spaces-more-morphisms-lemma-flat-finite-type-finitely-presented-over-dense-open-X}
$S$ をスキームとする。$f : X \to Y$ を $S$ 上の代数空間の射とする。
$V \subset Y$ を開部分空間とする。次を仮定する：
\begin{enumerate}
\item $f$ は局所有限型かつ平坦である。
\item $V \to Y$ は準コンパクトかつスキーム論的に稠密である。
\item $f|_{f^{-1}V} : f^{-1}V \to V$ は局所有限表示である。
\end{enumerate}
% P09-ERR-1179: frozen source line 8579 says "is of locally of finite presentation".
このとき $f$ は局所有限表示である。
\end{lemma}

\begin{proof}
証明は Lemma \ref{spaces-more-morphisms-lemma-flat-finite-type-finitely-presented-over-dense-open}
の証明と同一である。ただし More on Flatness, Lemma
\ref{flat-lemma-flat-finite-type-finitely-presented-over-dense-open-X} を用いる。
\end{proof}

\begin{lemma}
\label{spaces-more-morphisms-lemma-flat-fp-dimension-over-dense-open}
$S$ をスキームとする。$f : X \to Y$ を $S$ 上の代数空間の平坦かつ
局所有限型な射とする。$V \subset Y$ を開部分空間で、$|V| \subset |Y|$ が
稠密かつ $X_V \to V$ の相対次元が $\leq d$ となるものとする。
さらに次のいずれかを仮定する：
\begin{enumerate}
\item $f$ は局所有限表示である。
\item $V \to Y$ は準コンパクトである。
\end{enumerate}
このとき $f : X \to Y$ の相対次元は $\leq d$ である。
\end{lemma}

\begin{proof}
$Y$ をその被約化で置き換えてよいので、$Y$ は被約であると仮定してよい。
すると $V$ は $Y$ においてスキーム論的に稠密である。Morphisms of Spaces, Lemma
\ref{spaces-morphisms-lemma-quasi-compact-immersion} を参照せよ。
定義により、相対次元が $\leq d$ であるという性質はエタール被覆上で検査できる。
% P09-ERR-1180: frozen source line 8610 says "Sections" before a single section reference.
Morphisms of Spaces, Sections \ref{spaces-morphisms-section-relative-dimension} を参照せよ。
したがって、$f$ の相対次元が $\leq d$ であることを証明すれば十分である。
その証明では $X$ を $X$ 上全射エタールなスキームで置き換える。
同様に、$Y$ を $Y$ 上全射エタールなスキームで置き換えてよい。
% P09-ERR-1181: frozen source line 8614 again contains an extra "and" in "etale and over Y".
このスキームにおける $V$ の逆像はスキーム論的に稠密である。Morphisms of Spaces,
Section \ref{spaces-morphisms-section-scheme-theoretic-closure} を参照せよ。
スキームのスキーム論的に稠密な開部分は特に稠密なので、スキームの場合へ帰着される。
これは More on Flatness, Lemma
\ref{flat-lemma-flat-finite-presentation-dimension-over-dense-open} である。
\end{proof}

\begin{lemma}
\label{spaces-more-morphisms-lemma-proper-flat-finite-over-dense-open}
$S$ をスキームとする。$f : X \to Y$ を $S$ 上の代数空間の平坦かつ固有な射とする。
$V \to Y$ を開部分空間で、$|V| \subset |Y|$ が稠密かつ
$X_V \to V$ が有限となるものとする。さらに $f$ が局所有限表示であるか、
$V \to Y$ が準コンパクトであるならば、$f$ は有限である。
\end{lemma}

\begin{proof}
Lemma \ref{spaces-more-morphisms-lemma-flat-fp-dimension-over-dense-open} により、$f$ のファイバーの
次元は零である。Morphisms of Spaces, Lemma
\ref{spaces-morphisms-lemma-locally-quasi-finite-rel-dimension-0}
により、これは $f$ が局所準有限であることを含意する。Morphisms of Spaces, Lemma
\ref{spaces-morphisms-lemma-locally-quasi-finite-separated-representable}
により、これは $f$ が表現可能であることを含意する。$f$ が有限であるかどうかは
$Y$ 上エタール局所的に検査できるので、$Y$ はスキームであると仮定してよい。
$f$ は表現可能なので、スキームの場合へ帰着される。これは More on Flatness, Lemma
\ref{flat-lemma-proper-flat-finite-over-dense-open} である。
\end{proof}

\begin{lemma}
\label{spaces-more-morphisms-lemma-zariski}
$S$ をスキームとする。$f : X \to Y$ を $S$ 上の代数空間の射とする。
$V \subset Y$ を開部分空間とする。次を仮定する：
\begin{enumerate}
\item $f$ は分離的、局所有限型、かつ平坦である。
\item $f^{-1}(V) \to V$ は同型である。
\item $V \to Y$ は準コンパクトかつスキーム論的に稠密である。
\end{enumerate}
このとき $f$ は開埋入である。
\end{lemma}

\begin{proof}
Lemma \ref{spaces-more-morphisms-lemma-flat-finite-type-finitely-presented-over-dense-open-X}
を適用すると、$f$ は局所有限表示である。Lemma
\ref{spaces-more-morphisms-lemma-flat-fp-dimension-over-dense-open} を適用すると、$f$ の相対次元は
$\leq 0$ である。Morphisms of Spaces, Lemma
\ref{spaces-morphisms-lemma-locally-quasi-finite-rel-dimension-0}
により、これは $f$ が局所準有限であることを含意する。Morphisms of Spaces, Lemma
\ref{spaces-morphisms-lemma-locally-quasi-finite-separated-representable}
により、これは $f$ が表現可能であることを含意する。Descent on Spaces, Lemma
\ref{spaces-descent-lemma-descending-property-open-immersion}
により、$f$ が開埋入であるかどうかは $Y$ 上エタール局所的に検査できる。
したがって $Y$ はスキームであると仮定してよい。$f$ は表現可能なので、
スキームの場合へ帰着される。これは More on Flatness, Lemma
\ref{flat-lemma-zariski} である。
\end{proof}







\section{ブローアップと平坦性}
\label{spaces-more-morphisms-section-blowup-flat}

\noindent
More on Flatness, Section \ref{flat-section-blowup-flat} の仕事を繰り返す代わりに、
More on Flatness, Lemma \ref{flat-lemma-push-ideal} の類似を証明する。
これは適切なブローアップを見つける問題が、しばしば基底上エタール局所的であることを示す。

\begin{lemma}
\label{spaces-more-morphisms-lemma-push-ideal}
$S$ をスキームとする。$X$ を $S$ 上準コンパクトかつ準分離的な代数空間とする。
$\varphi : W \to X$ を準コンパクト、分離的、エタールな射とする。
$U \subset X$ を準コンパクトな開部分空間とする。
$\mathcal{I} \subset \mathcal{O}_W$ を有限型準連接イデアル層で、
$V(\mathcal{I}) \cap \varphi^{-1}(U) = \emptyset$ を満たすものとする。
このとき有限型準連接イデアル層 $\mathcal{J} \subset \mathcal{O}_X$ で、
次を満たすものが存在する：
\begin{enumerate}
\item $V(\mathcal{J}) \cap U = \emptyset$。
\item $\varphi^{-1}(\mathcal{J})\mathcal{O}_W = \mathcal{I} \mathcal{I}'$。
ここで $\mathcal{I}' \subset \mathcal{O}_W$ はある有限型準連接イデアルである。
\end{enumerate}
\end{lemma}

\begin{proof}
分解 $W \to Y \to X$ を選ぶ。ここで $j : W \to Y$ は準コンパクトな開埋入、
$\pi : Y \to X$ は有限表示の有限射である
（Lemma \ref{spaces-more-morphisms-lemma-quasi-finite-separated-pass-through-finite-addendum}）。
$V = j(W) \cup \pi^{-1}(U) \subset Y$ とおく。$\mathcal{I}$（$W \cong j(W)$ 上）と
$\mathcal{O}_{\pi^{-1}(U)}$ は貼り合わされ、
有限型準連接イデアル層 $\mathcal{I}_1 \subset \mathcal{O}_V$ をなすことに注意する。
Limits of Spaces, Lemma \ref{spaces-limits-lemma-extend} により、
有限型準連接イデアル層 $\mathcal{I}_2 \subset \mathcal{O}_Y$ で、
$\mathcal{I}_2|_V = \mathcal{I}_1$ となるものが存在する。言い換えれば、
$\mathcal{I}_2 \subset \mathcal{O}_Y$ は有限型準連接イデアル層であり、
$V(\mathcal{I}_2)$ は $\pi^{-1}(U)$ と交わらず、
$j^{-1}\mathcal{I}_2 = \mathcal{I}$ である。対応する閉埋入を
% P09-ERR-1182: frozen source line 8731 says "Denote i ... the" rather than "Denote by i ...".
$i : Z \to Y$ と表す。これは有限表示である（Morphisms of Spaces, Lemma
\ref{spaces-morphisms-lemma-closed-immersion-finite-presentation}）。
特に合成 $\tau = \pi \circ i : Z \to X$ は有限かつ有限表示である
（Morphisms of Spaces, Lemmas
\ref{spaces-morphisms-lemma-composition-finite-presentation} and
\ref{spaces-morphisms-lemma-composition-integral}）。

\medskip\noindent
$\mathcal{F} = \tau_*\mathcal{O}_Z$ とおき、これを準連接
$\mathcal{O}_X$-加群とみなす。Descent on Spaces, Lemma
\ref{spaces-descent-lemma-finite-finitely-presented-module}
により、$\mathcal{F}$ は有限表示 $\mathcal{O}_X$-加群である。
$\mathcal{J} = \text{Fit}_0(\mathcal{F})$ とおく。
% P09-ERR-1183: frozen source line 8747 contains an unresolved editorial placeholder for a reference.
（環付きトポス上の Fitting 加群への参照をここに挿入する。）これは $X$ 上の
有限型準連接イデアル層である（$\mathcal{F}$ は有限表示である。More on Algebra,
Lemma \ref{more-algebra-lemma-fitting-ideal-basics} を参照せよ）。
補題の (1) は成り立つ。実際、$|\tau|(|Z|) \cap |U| = \emptyset$ である
（$\mathcal{I}_2$ の選び方による）。また零加群の第 $0$ Fitting イデアルは
構造層に等しい。(2) を証明するには、Fitting イデアルを取ることは基底変換と
可換なので
$\varphi^{-1}(\mathcal{J})\mathcal{O}_W = \text{Fit}_0(\varphi^*\mathcal{F})$
であることに注意する。一方、$\varphi : W \to X$ は分離的かつエタールなので、
$(1, j) : W \to W \times_X Y$ は開かつ閉な埋入である。したがって
$W \times_Y Z = V(\mathcal{I}) \amalg Z'$ であり、ある有限かつ有限表示な
代数空間の射 $\tau' : Z' \to W$ が存在する。ゆえに
\begin{align*}
\text{Fit}_0(\varphi^*\mathcal{F}) & =
\text{Fit}_0((W \times_Y Z \to W)_*\mathcal{O}_{W \times_Y Z}) \\
& =
\text{Fit}_0(\mathcal{O}_W/\mathcal{I}) \cdot
\text{Fit}_0(\tau'_*\mathcal{O}_{Z'}) \\
& =
\mathcal{I} \cdot \text{Fit}_0(\tau'_*\mathcal{O}_{Z'})
\end{align*}
である。第二の等号には More on Algebra, Lemma
\ref{more-algebra-lemma-fitting-ideal-basics} を環付きトポス上の層へ移したものを用いた。
$\mathcal{I}' = \text{Fit}_0(\tau'_*\mathcal{O}_{Z'})$ とおけば、補題の証明が完了する。
\end{proof}

\begin{theorem}
\label{spaces-more-morphisms-theorem-flatten-module}
$S$ をスキームとする。$B$ を $S$ 上準コンパクトかつ準分離的な代数空間とする。
$X$ を $B$ 上の代数空間とする。$\mathcal{F}$ を $X$ 上の準連接加群とする。
$U \subset B$ を準コンパクトな開部分空間とする。次を仮定する：
\begin{enumerate}
\item $X$ は準コンパクトである。
\item $X$ は $B$ 上局所有限表示である。
\item $\mathcal{F}$ は有限型加群である。
\item $\mathcal{F}_U$ は有限表示である。
\item $\mathcal{F}_U$ は $U$ 上平坦である。
\end{enumerate}
このとき $U$-許容ブローアップ $B' \to B$ で、狭義変換 $\mathcal{F}'$、
すなわち $\mathcal{F}$ の狭義変換が有限表示の $\mathcal{O}_{X \times_B B'}$-加群であり、
$B'$ 上平坦となるものが存在する。
\end{theorem}

\begin{proof}
アフィンスキーム $V$ と全射エタール射 $V \to X$ を選ぶ。狭義変換は
エタール局所化と可換なので（Divisors on Spaces, Lemma
\ref{spaces-divisors-lemma-strict-transform-local}）、$X$ を $V$ で置き換えた結果を
証明すれば十分である。したがって、
% P09-ERR-1184: frozen source line 8802 calls this theorem a lemma.
$X \to B$ は（補題の仮定に加えて）表現可能であると仮定してよい。

\medskip\noindent
$X \to B$ が表現可能であると仮定する。アフィンスキーム $W$ と全射エタール射
$\varphi : W \to B$ を選ぶ。$X \times_B W$ はスキームであることに注意する。
スキームの場合（More on Flatness, Theorem \ref{flat-theorem-flatten-module}）により、
有限型準連接イデアル層 $\mathcal{I} \subset \mathcal{O}_W$ で、次を満たすものを取れる：
(a) $|V(\mathcal{I})| \cap |\varphi^{-1}(U)| = \emptyset$、および (b)
$\mathcal{F}|_{X \times_B W}$ の、ブローアップ $W' \to W$、すなわち
$\mathcal{I}$ に沿うものに関する狭義変換は $W'$ 上平坦となり、有限表示加群である。
Lemma \ref{spaces-more-morphisms-lemma-push-ideal} にある有限型イデアル層
$\mathcal{J} \subset \mathcal{O}_B$ を選ぶ。$B' \to B$ を $\mathcal{J}$ の
ブローアップとする。このブローアップが機能すると主張する。実際、$B' \to B$ が
$U$-許容であることは明らかである（イデアル $\mathcal{J}$ の選び方による）。
さらに基底変換 $B' \times_B W \to W$ は、$W$ の
$\varphi^{-1}\mathcal{J} = \mathcal{I}\mathcal{I}'$ に沿うブローアップである
（ブローアップと平坦基底変換との両立性については Divisors on Spaces, Lemma
\ref{spaces-divisors-lemma-flat-base-change-blowing-up} を参照せよ）。
したがって分解
$$
W \times_B B' \to W' \to W
$$
が存在し、第一の射もブローアップである。Divisors on Spaces, Lemma
\ref{spaces-divisors-lemma-blowing-up-two-ideals} を参照せよ。
$\mathcal{F}'$（$B' \times_B X$ 上にある）の
$W \times_B B' \times_B X$ への制限は、
$\mathcal{F}|_{X \times_B W}$ の狭義変換である
（Divisors on Spaces, Lemma \ref{spaces-divisors-lemma-strict-transform-local}）。
したがってこれは、上に表示した二つのブローアップによる
$\mathcal{F}|_{X \times_B W}$ の二度反復した狭義変換である
（Divisors on Spaces, Lemma
\ref{spaces-divisors-lemma-strict-transform-composition-blowups}）。
第一のブローアップの後、構成により我々の層は既に基底上平坦かつ有限表示である。
ゆえに第二の狭義変換の後も同じことが成り立つ（これは Divisors on Spaces, Lemma
\ref{spaces-divisors-lemma-strict-transform-flat} により引き戻しだからである）。
したがって $\mathcal{F}'$ の $B' \times_B X$ のエタール被覆への制限は
所望の性質をもち、定理が証明された。
\end{proof}








\section{応用}
\label{spaces-more-morphisms-section-applications-flattening-by-blowing-up}

\noindent
この節では、ブローアップによる平坦化の結果を応用する。

\begin{lemma}
\label{spaces-more-morphisms-lemma-flat-after-blowing-up}
$S$ をスキームとする。$f : X \to B$ を $S$ 上の代数空間の射とする。
$U \subset B$ を開部分空間とする。次を仮定する：
\begin{enumerate}
\item $B$ は準コンパクトかつ準分離的である。
\item $U$ は準コンパクトである。
\item $f : X \to B$ は有限型かつ準分離的である。
\item $f^{-1}(U) \to U$ は平坦かつ局所有限表示である。
\end{enumerate}
このとき $U$-許容ブローアップ $B' \to B$ で、狭義変換 $X'$（$X$ の
狭義変換）が $B'$ 上平坦かつ有限表示となるものが存在する。
\end{lemma}

\begin{proof}
$B' \to B$ を $U$-許容ブローアップとする。$X$ の狭義変換は $B'$ 上
準コンパクトかつ準分離的であることに注意する。これは $X$ が $B$ 上
準コンパクトかつ準分離的だからである。したがって、狭義変換が平坦かつ
局所有限表示となる $U$-許容
ブローアップを見つけることだけを考えればよい。$X$ は $B$ 上局所有限表示では
ないため、Theorem \ref{spaces-more-morphisms-theorem-flatten-module} を直接適用することはできない。

\medskip\noindent
アフィンスキーム $V$ と全射エタール射 $V \to X$ を選ぶ。
（これは $X$ が、準コンパクト空間 $B$ 上有限型の空間として準コンパクトだから可能である。）
このとき射 $V \to B$ に対して結果を示せば十分である
（狭義変換はエタール局所化と可換する。Divisors on Spaces, Lemma
\ref{spaces-divisors-lemma-strict-transform-local} を参照せよ）。
したがって $X \to B$ は有限型であるだけでなく分離的でもあると仮定してよい。
この場合、閉埋入 $i : X \to Y$ で、$Y \to B$ が分離的かつ有限表示となるものを
見つけられる。Limits of Spaces, Proposition
\ref{spaces-limits-proposition-separated-closed-in-finite-presentation} を参照せよ。

\medskip\noindent
$\mathcal{F} = i_*\mathcal{O}_X$ に、$Y/B$ 上で Theorem
\ref{spaces-more-morphisms-theorem-flatten-module} を適用する。すると $U$-許容ブローアップ
% P09-ERR-1185: frozen source line 8899 omits "the" before "strict transform".
$B' \to B$ で、$\mathcal{F}$ の狭義変換が $B'$ 上平坦かつ有限表示となるものを得る。
$X'$ を $X$ の、ブローアップ $B' \to B$ による狭義変換とする。
$i' : X' \to Y \times_B B'$ を誘導された射とする。狭義変換を取ることは
アフィン射に沿う押し出しと可換なので（Divisors on Spaces, Lemma
\ref{spaces-divisors-lemma-strict-transform-affine}）、$i'_*\mathcal{O}_{X'}$ は
$B'$ 上平坦であり、$\mathcal{O}_{Y \times_B B'}$-加群として有限表示である。
したがって $X' \to B'$ は平坦かつ局所有限表示である。先の注意により補題が従う。
\end{proof}

\begin{lemma}
\label{spaces-more-morphisms-lemma-finite-after-blowing-up}
$S$ をスキームとする。$f : X \to B$ を $S$ 上の代数空間の射とする。
$U \subset B$ を開部分空間とする。次を仮定する：
\begin{enumerate}
\item $B$ は準コンパクトかつ準分離的である。
\item $U$ は準コンパクトである。
\item $f : X \to B$ は固有である。
\item $f^{-1}(U) \to U$ は有限局所自由である。
\end{enumerate}
このとき $U$-許容ブローアップ $B' \to B$ で、狭義変換 $X'$（$X$ の
狭義変換）が $B'$ 上有限局所自由となるものが存在する。
\end{lemma}

\begin{proof}
Lemma \ref{spaces-more-morphisms-lemma-flat-after-blowing-up} により、$X \to B$ は平坦かつ有限表示であると
仮定してよい。必要なら $B$ を $U$-許容ブローアップで置き換えることにより、
$U \subset B$ はスキーム論的に稠密であると仮定してよい。このとき Lemma
\ref{spaces-more-morphisms-lemma-proper-flat-finite-over-dense-open} により $f$ は有限である。
したがって Morphisms of Spaces, Lemma
\ref{spaces-morphisms-lemma-finite-flat} により $f$ は有限局所自由である。
\end{proof}

\begin{lemma}
\label{spaces-more-morphisms-lemma-isomorphism-after-blowing-up}
$S$ をスキームとする。$f : X \to B$ を $S$ 上の代数空間の射とする。
$U \subset B$ を開部分空間とする。次を仮定する：
\begin{enumerate}
\item $B$ は準コンパクトかつ準分離的である。
\item $U$ は準コンパクトである。
\item $f : X \to B$ は固有である。
\item $f^{-1}(U) \to U$ は同型である。
\end{enumerate}
このとき $U$-許容ブローアップ $B' \to B$ で、狭義変換 $X'$（$X$ の
狭義変換）が $B'$ へ同型に写るものが存在する。
\end{lemma}

\begin{proof}
Lemma \ref{spaces-more-morphisms-lemma-flat-after-blowing-up} により、$X \to B$ は平坦かつ有限表示であると
仮定してよい。必要なら $B$ を $U$-許容ブローアップで置き換えることにより、
$U \subset B$ はスキーム論的に稠密であると仮定してよい。このとき Lemma
\ref{spaces-more-morphisms-lemma-proper-flat-finite-over-dense-open} により $f$ は有限であり、Lemma
\ref{spaces-more-morphisms-lemma-zariski} により開埋入である。したがって $f$ は、像が閉で稠密開集合
$U$ を含む開埋入であり、それゆえ $f$ は同型である。
\end{proof}

\begin{lemma}
\label{spaces-more-morphisms-lemma-zariski-after-blowup}
$S$ をスキームとする。$f : X \to B$ を $S$ 上の代数空間の射とする。
$U \subset B$ を開部分空間とする。次を仮定する：
\begin{enumerate}
% P09-ERR-1186: frozen source line 8966 omits "is" after B.
\item $B$ は準コンパクトかつ準分離的である。
\item $U$ は準コンパクトである。
% P09-ERR-1187: frozen source line 8968 omits the list-separating comma.
\item $f$ は有限型である。
\item $f^{-1}(U) \to U$ は同型である。
\end{enumerate}
このとき $U$-許容ブローアップ $B' \to B$ で、$U$ が $B'$ において
スキーム論的に稠密であり、かつ狭義変換 $X'$（$X$ の狭義変換）が $B'$ の
開部分空間へ同型に写るものが存在する。
\end{lemma}

\begin{proof}
この補題は Lemma \ref{spaces-more-morphisms-lemma-isomorphism-after-blowing-up} の一般化である。
$U$-許容ブローアップの合成は $U$-許容なので（Divisors on Spaces, Lemma
\ref{spaces-divisors-lemma-composition-admissible-blowups}）、段階的に進められる。
有限型準連接イデアル層 $\mathcal{I} \subset \mathcal{O}_B$ で
$|B| \setminus |U| = |V(\mathcal{I})|$ を満たすものを選ぶ。$B$ を
$B$ の $\mathcal{I}$ に沿うブローアップで置き換え、$X$ を $X$ の狭義変換で
置き換える。この置換の後、$B \setminus U$ は有効 Cartier 因子 $D$ の台である
（Divisors on Spaces, Lemma
\ref{spaces-divisors-lemma-blowing-up-gives-effective-Cartier-divisor}）。
特に $U$ は $B$ においてスキーム論的に稠密である
（Divisors on Spaces, Lemma
\ref{spaces-divisors-lemma-complement-effective-Cartier-divisor}）。
次に、もう一度 $U$-許容ブローアップを行い、$X \to B$ が平坦かつ有限表示となる
状況にする。Lemma \ref{spaces-more-morphisms-lemma-flat-after-blowing-up} を参照せよ。$U$ は依然として
$B$ においてスキーム論的に稠密であることに注意する。したがって Lemma
\ref{spaces-more-morphisms-lemma-zariski} により $X \to B$ は開埋入である。
\end{proof}

\noindent
次の補題は、モディフィケーションがブローアップによって支配されることを述べる。

\begin{lemma}
\label{spaces-more-morphisms-lemma-dominate-modification-by-blowup}
$S$ をスキームとする。$f : X \to B$ を $S$ 上の代数空間の射とする。
$U \subset B$ を開部分空間とする。次を仮定する：
\begin{enumerate}
\item $B$ は準コンパクトかつ準分離的である。
\item $U$ は準コンパクトである。
\item $f : X \to B$ は固有である。
% P09-ERR-1188: frozen source line 9011 has "us" instead of "is".
\item $f^{-1}(U) \to U$ は同型である。
\end{enumerate}
このとき $U$-許容ブローアップ $B' \to B$ で $X$ を支配するもの、すなわち
ブローアップ射の分解 $B' \to X \to B$ が存在するものがある。
\end{lemma}

\begin{proof}
Lemma \ref{spaces-more-morphisms-lemma-isomorphism-after-blowing-up} により、$U$-許容ブローアップ
$B' \to B$ で、狭義変換 $X'$ が $B'$ へ同型に写るものを見つけられる。
そこで $B' = X' \to X$ を分解として用いればよい。
\end{proof}

\begin{lemma}
\label{spaces-more-morphisms-lemma-get-section-after-blowup}
$S$ をスキームとする。$X$, $Y$ を $S$ 上の代数空間とする。
$U \subset W \subset Y$ を開部分空間とする。$f : X \to W$ および
$s : U \to X$ を $f \circ s = \text{id}_U$ を満たす射とする。次を仮定する：
\begin{enumerate}
\item $f$ は固有である。
\item $Y$ は準コンパクトかつ準分離的である。
\item $U$ と $W$ は準コンパクトである。
\end{enumerate}
このとき $U$-許容ブローアップ $b : Y' \to Y$ と、射
$s' : b^{-1}(W) \to X$ で、$s$ を延長し
$f \circ s' = b|_{b^{-1}(W)}$ を満たすものが存在する。
\end{lemma}

\begin{proof}
$X$ を $s$ のスキーム論的像で置き換えてよく、そうする。このとき $X \to W$ は
$U$ 上同型である。Morphisms of Spaces, Lemma
\ref{spaces-morphisms-lemma-scheme-theoretic-image-of-partial-section} を参照せよ。
Lemma \ref{spaces-more-morphisms-lemma-dominate-modification-by-blowup} により、$U$-許容ブローアップ
$W' \to W$ と延長 $W' \to X$（$s$ の延長）が存在する。最後に Divisors on Spaces,
Lemma \ref{spaces-divisors-lemma-extend-admissible-blowups} を適用し、
$W' \to W$ を $U$-許容な $Y$ のブローアップへ延長すれば証明が完了する。
\end{proof}















\section{Chow の補題}
\label{spaces-more-morphisms-section-chow}

\noindent
この節では Chow の補題（Lemma \ref{spaces-more-morphisms-lemma-chow-noetherian-separated}）を証明する。
容易に証明でき、多くの応用に十分な Chow の補題の弱い形については、
Cohomology of Spaces, Section \ref{spaces-cohomology-section-weak-chow} を
参照することを勧める。

\medskip\noindent
代数空間の射の射影性はまだ定義していないので、補題の主張（例えば Lemma
\ref{spaces-more-morphisms-lemma-blowup-to-find-embedding} を参照せよ）では射影的な射ではなく、
表現可能な固有射を用いる。

\begin{lemma}
\label{spaces-more-morphisms-lemma-find-common-blowups}
$S$ をスキームとする。$Y$ を $S$ 上準コンパクトかつ準分離的な代数空間とする。
$U \to X_1$ と $U \to X_2$ を $Y$ 上の代数空間の開埋入とし、
% P09-ERR-1189: frozen source line 9087 omits "are" after the three subjects.
$U$, $X_1$, $X_2$ は $Y$ 上有限型かつ分離的であると仮定する。このとき
次の可換図式が存在する：
$$
\xymatrix{
X_1' \ar[d] \ar[r] & X & X_2' \ar[l] \ar[d] \\
X_1 & U \ar[l] \ar[lu] \ar[u] \ar[ru] \ar[r] & X_2
}
$$
これは $Y$ 上の代数空間の図式であり、$X_i' \to X_i$ は $U$-許容
ブローアップ、$X_i' \to X$ は開埋入、$X$ は $Y$ 上分離的かつ有限型である。
\end{lemma}

\begin{proof}
証明を通じて、すべての代数空間は $Y$ 上分離的かつ有限型であるとする。
特に、これらの代数空間は準コンパクトかつ準分離的であり、それらの間の射は
準コンパクトかつ分離的である。Morphisms of Spaces, Sections
\ref{spaces-morphisms-section-separation-axioms} and
\ref{spaces-morphisms-section-quasi-compact} を参照せよ。
このような空間の埋入 $U \to W$ が $Y$ 上にあるとき、スキーム論的像
$Z$（$U$ の $W$ における像）は $W$ の閉部分空間であり、$U \to Z$ は開埋入、
$U \subset Z$ はスキーム論的に稠密、かつ $|U| \subset |Z|$ は稠密であることを
用いる。Morphisms of Spaces, Lemma
\ref{spaces-morphisms-lemma-quasi-compact-immersion} を参照せよ。

\medskip\noindent
$X_{12} \subset X_1 \times_Y X_2$ を $U \to X_1 \times_Y X_2$ の
スキーム論的像とする。Morphisms of Spaces, Lemma
\ref{spaces-morphisms-lemma-scheme-theoretic-image-of-partial-section} により、射影
$p_i : X_{12} \to X_i$ は同型 $p_i^{-1}(U) \to U$ を誘導する。
$U$-許容ブローアップ $X_i^i \to X_i$ で、狭義変換 $X_{12}^i$
（$X_{12}$ の狭義変換）が $X_i^i$ の開部分空間と同型になるものを選ぶ。Lemma
\ref{spaces-more-morphisms-lemma-zariski-after-blowup} を参照せよ。$\mathcal{I}_i \subset
\mathcal{O}_{X_i}$ を対応する有限型準連接イデアル層とする。
$X_{12}^i \to X_{12}$ は
$p_i^{-1}\mathcal{I}_i \mathcal{O}_{X_{12}}$ に沿うブローアップであることを
思い出そう。Divisors on Spaces, Lemma
\ref{spaces-divisors-lemma-strict-transform} を参照せよ。
$X_{12}'$ を $X_{12}$ の
$p_1^{-1}\mathcal{I}_1 p_2^{-1}\mathcal{I}_2 \mathcal{O}_{X_{12}}$ に沿う
ブローアップとする。これが何を意味するかについては Divisors on Spaces, Lemma
\ref{spaces-divisors-lemma-blowing-up-two-ideals} を参照せよ。次の可換図式を得る：
$$
\xymatrix{
X_{12}' \ar[d] \ar[r] & X_{12}^2 \ar[d] \\
X_{12}^1 \ar[r] & X_{12}
}
$$
ここですべての射は $U$-許容ブローアップである。$X_{12}^i \subset X_i^i$ は
開部分空間なので、$U$-許容ブローアップ $X_i' \to X_i^i$ で
$X_{12}' \to X_{12}^i$ に制限されるものを選べる。Divisors on Spaces, Lemma
\ref{spaces-divisors-lemma-extend-admissible-blowups} を参照せよ。
このとき $X_{12}' \subset X_i'$ は開部分空間であり、図式
$$
\xymatrix{
X_{12}' \ar[d] \ar[r] & X_i' \ar[d] \\
X_{12}^i \ar[r] & X_i^i
}
$$
は可換で、縦の射はブローアップ、横の射は開埋入である。
$X'_{12} \to X_1' \times_Y X_2'$ は埋入かつ固有であることに注意する
（$X'_{12} \to X_{12}$ が固有、$X_{12} \to X_1 \times_Y X_2$ が閉、かつ
$X_1' \times_Y X_2' \to X_1 \times_Y X_2$ が分離的であることを用い、
Morphisms of Spaces, Lemma
\ref{spaces-morphisms-lemma-universally-closed-permanence} を適用せよ）。
したがって $X'_{12} \to  X_1' \times_Y X_2'$ は閉埋入である。
$X$ を、$X_1'$ と $X_2'$ を共通の開部分空間 $X_{12}'$ に沿って貼り合わせて
定義すれば、$X \to Y$ は有限型かつ分離的である\footnote{対角射
$X \to X \times_Y X$ の閉性は、四つの開部分 $X'_i \times_Y X'_j$
（$X \times_Y X$ の開部分）上で確認でき、そこでは明らかだからである。}。
$U$-許容ブローアップの合成は $U$-許容ブローアップなので
（Divisors on Spaces, Lemma
\ref{spaces-divisors-lemma-composition-admissible-blowups}）、補題が証明された。
\end{proof}

\begin{lemma}
\label{spaces-more-morphisms-lemma-blowup-to-find-embedding}
$S$ をスキームとする。$f : X \to Y$ を $S$ 上の代数空間の射とする。
$U \subset X$ を開部分空間とする。次を仮定する：
\begin{enumerate}
\item $U$ は準コンパクトである。
\item $Y$ は準コンパクトかつ準分離的である。
\item 埋入 $U \to \mathbf{P}^n_Y$ が $Y$ 上に存在する。
\item $f$ は有限型かつ分離的である。
\end{enumerate}
このとき次の可換図式が存在する：
$$
\xymatrix{
& U \ar[ld] \ar[d] \ar[rd] \ar[rrd] \\
X \ar[rd] & X' \ar[l] \ar[d] \ar[r] & Z' \ar[ld] \ar[r] & Z \ar[ld] \\
& Y & \mathbf{P}^n_Y \ar[l]
}
$$
ここで、始点が $U$ である射は開埋入、$X' \to X$ は $U$-許容
ブローアップ、$X' \to Z'$ は開埋入、$Z' \to Y$ は代数空間の固有かつ
表現可能な射である。より正確には、$Z' \to Z$ は $U$-許容ブローアップであり、
$Z \to \mathbf{P}^n_Y$ は閉埋入である。
\end{lemma}

\begin{proof}
$Z \subset \mathbf{P}^n_Y$ を埋入 $U \to \mathbf{P}^n_Y$ のスキーム論的像とする。
$U \to \mathbf{P}^n_Y$ は準コンパクトなので、$U \subset Z$ は
（スキーム論的に）稠密な開部分空間である
（Morphisms of Spaces, Lemma
\ref{spaces-morphisms-lemma-quasi-compact-immersion}）。
Lemma \ref{spaces-more-morphisms-lemma-find-common-blowups} を適用して次の図式を得る：
$$
\xymatrix{
X' \ar[d] \ar[r] & \overline{X}' & Z' \ar[l] \ar[d] \\
X & U \ar[l] \ar[lu] \ar[u] \ar[ru] \ar[r] & Z
}
$$
その性質は同補題の主張に列挙されたとおりである。$X' \to X$ と $Z' \to Z$ は
$U$-許容ブローアップなので、$U$ は $X'$ と $Z'$ の双方において
スキーム論的に稠密な開部分である
（Divisors on Spaces, Lemmas
\ref{spaces-divisors-lemma-blowing-up-gives-effective-Cartier-divisor} and
\ref{spaces-divisors-lemma-complement-effective-Cartier-divisor} を参照せよ）。
$Z' \to Z \to Y$ は固有なので、$Z' \subset \overline{X}'$ は閉部分空間である
（Morphisms of Spaces, Lemma
\ref{spaces-morphisms-lemma-universally-closed-permanence} を参照せよ）。
したがって $X' \subset Z'$ が（スキーム論的に）成り立ち、それゆえ $X'$ は
$Z'$ の開部分空間である（細部は省略する）。これで補題が証明された。
\end{proof}

\begin{lemma}
\label{spaces-more-morphisms-lemma-chow-noetherian}
$S$ をスキームとする。$f : X \to Y$ を $S$ 上の代数空間の射とする。
% P09-ERR-1190: frozen source line 9227 omits "is" after f.
$f$ は分離的かつ有限型であり、$Y$ は Noetherian であると仮定する。
このとき稠密開部分空間 $U \subset X$ と次の可換図式が存在する：
$$
\xymatrix{
& U \ar[ld] \ar[d] \ar[rd] \ar[rrd] \\
X \ar[rd] & X' \ar[l] \ar[d] \ar[r] & Z' \ar[ld] \ar[r] & Z \ar[ld] \\
& Y & \mathbf{P}^n_Y \ar[l]
}
$$
ここで、始点が $U$ である射は開埋入、$X' \to X$ は $U$-許容
ブローアップ、$X' \to Z'$ は開埋入、$Z' \to Y$ は代数空間の固有かつ
表現可能な射である。より正確には、$Z' \to Z$ は $U$-許容ブローアップであり、
$Z \to \mathbf{P}^n_Y$ は閉埋入である。
\end{lemma}

\begin{proof}
Limits of Spaces, Lemma
\ref{spaces-limits-lemma-embedding-into-affine-over-qs} により、稠密開部分空間
$U \subset X$ と埋入 $U \to \mathbf{A}^n_Y$ が $Y$ 上に存在する。
これを開埋入 $\mathbf{A}^n_Y \to \mathbf{P}^n_Y$ と合成すると、Lemma
\ref{spaces-more-morphisms-lemma-blowup-to-find-embedding} の状況を得るので、結果が従う。
\end{proof}

\begin{remark}
\label{spaces-more-morphisms-remark-chow-Noetherian}
Lemmas \ref{spaces-more-morphisms-lemma-blowup-to-find-embedding} and \ref{spaces-more-morphisms-lemma-chow-noetherian} において、
射 $g : Z' \to Y$ は射影的射の合成である。おそらく（代数空間に対する
Morphisms, Lemma \ref{morphisms-lemma-ample-composition} の類似により）、
$g$-豊富な可逆層が $Z'$ 上に存在する。これが必要になれば、ここで主張し証明する。
\end{remark}

\noindent
次の結果は \cite[IV Theorem 3.1]{Kn} である。埋入
$X' \to \mathbf{P}^n_Y$ は準コンパクトなので、$X' \to Z' \to
\mathbf{P}^n_Y$ と分解でき、第一の射は開埋入、第二の射は閉埋入であることに
注意する（Morphisms of Spaces, Lemma
\ref{spaces-morphisms-lemma-quasi-compact-immersion}）。

\begin{lemma}[Chow の補題]
\label{spaces-more-morphisms-lemma-chow-noetherian-separated}
\begin{reference}
\cite[IV Theorem 3.1]{Kn}
\end{reference}
$S$ をスキームとする。$f : X \to Y$ を $S$ 上の代数空間の射とする。
% P09-ERR-1191: frozen source line 9280 omits "is" after f.
$f$ は分離的かつ有限型であり、$Y$ は分離的かつ Noetherian であると仮定する。
このとき次の可換図式が存在する：
$$
\xymatrix{
X \ar[rd] & X' \ar[l] \ar[d] \ar[r] & \mathbf{P}^n_Y \ar[ld] \\
& Y
}
$$
ここで $X' \to X$ は、ある $U$ に対する許容ブローアップであり、
$U \subset X$ は稠密開集合である。また射 $X' \to \mathbf{P}^n_Y$ は埋入である。
\end{lemma}

\begin{proof}
証明のこの第一段落では、補題を $Y$ が $\Spec(\mathbf{Z})$ 上有限型である場合へ
帰着する。基礎スキーム $S$ を $\Spec(\mathbf{Z})$ で置き換えてよく、そうする。
$Y = \lim Y_i$ を $\Spec(\mathbf{Z})$ 上有限型の分離的代数空間の有向極限として
書ける。Limits of Spaces, Proposition
\ref{spaces-limits-proposition-approximate} and Lemma
\ref{spaces-limits-lemma-descend-separated} を参照せよ。十分大きいすべての $i$ に
対し、分離的有限型射 $X_i \to Y_i$ で $X = Y \times_{Y_i} X_i$ を満たすものを
見つけられる。Limits of Spaces, Lemmas
\ref{spaces-limits-lemma-descend-finite-presentation} and
\ref{spaces-limits-lemma-descend-separated-morphism} を参照せよ。
$\eta_1, \ldots, \eta_n$ を $|X|$ の既約成分の生成点とする
（$X$ は Noetherian 代数空間 $Y$ 上有限型かつ分離的な代数空間として
Noetherian であり、したがって $|X|$ は Noetherian 位相空間である）。
Limits of Spaces, Lemma \ref{spaces-limits-lemma-topology-limit} により、
$\eta_1, \ldots, \eta_n$ の $|X_i|$ における像は十分大きい $i$ に対して
相異なる。$X_i$ を射 $X \to X_i$（これは準コンパクトであり、実際アフィンである）
のスキーム論的像で置き換えてよい。この置換の後、$\eta_1, \ldots, \eta_n$ の
$|X_i|$ における像は $|X_i|$ の既約成分の生成点である。Morphisms of Spaces,
Lemma \ref{spaces-morphisms-lemma-quasi-compact-scheme-theoretic-image} を参照せよ。
以上を踏まえ、次の図式を見つけられると仮定する：
$$
\xymatrix{
X_i \ar[rd] & X_i' \ar[l] \ar[d] \ar[r] & \mathbf{P}^n_{Y_i} \ar[ld] \\
% P09-ERR-1192: frozen source line 9323 has Y although the diagram is over Y_i.
& Y
}
$$
ここで $X_i' \to X_i$ は、ある $U_i$ に対する許容ブローアップであり、
$U_i \subset X_i$ は稠密開集合である。また射 $X_i' \to \mathbf{P}^n_{Y_i}$ は埋入である。
このとき狭義変換 $X' \to X$（$X$ の、$X_i' \to X_i$ に関する狭義変換）は
$U$-許容ブローアップである。ここで $U \subset X$ は $U_i$ の $X$ における逆像である。
添字 $i$ を注意深く選んだことにより、$\eta_1, \ldots, \eta_n \in |U|$ が従い、
$U \subset X$ は稠密である。さらに $X' \to \mathbf{P}^n_Y$ は埋入である。
実際、$X'$ は
$X_i' \times_{X_i} X = X_i' \times_{Y_i} Y$ において閉であり、後者は
$\mathbf{P}^n_Y$ への埋入を備える。したがって次の段落の状況へ帰着した。

\medskip\noindent
$Y$ は $\Spec(\mathbf{Z})$ 上分離的かつ有限型であると仮定する。このとき
$X \to \Spec(\mathbf{Z})$ も分離的かつ有限型である。Lemma
\ref{spaces-more-morphisms-lemma-chow-noetherian} を $X \to \Spec(\mathbf{Z})$ に適用し、
稠密開部分空間 $U \subset X$ と次の可換図式を得る：
$$
\xymatrix{
& U \ar[ld] \ar[d] \ar[rd] \ar[rrd] \\
X \ar[rd] & X' \ar[l] \ar[d] \ar[r] & Z' \ar[ld] \ar[r] & Z \ar[ld] \\
& \Spec(\mathbf{Z}) & \mathbf{P}^n_\mathbf{Z} \ar[l]
}
$$
これは補題に列挙されたすべての性質をもつ。$Z$ は豊富な可逆層、すなわち
$\mathcal{O}_{\mathbf{P}^n}(1)|_Z$ をもつことに注意する。したがって
Morphisms, Lemma
\ref{morphisms-lemma-projective-over-quasi-projective-is-H-projective} により、
$Z' \to Z$ は H-射影的な射である。ゆえに Morphisms, Lemma
\ref{morphisms-lemma-H-projective-composition} により
$Z' \to \Spec(\mathbf{Z})$ は H-射影的である。したがって閉埋入
$Z' \to \mathbf{P}^m_{\Spec(\mathbf{Z})}$ が、ある $m \geq 0$ に対して存在する。
それゆえ対角射
$$
X' \to Y \times \mathbf{P}^m_\mathbf{Z} = \mathbf{P}^m_Y
$$
は埋入であり（$\mathbf{P}^m_\mathbf{Z}$ への射影との合成が埋入だからである）、
結論を得る。
\end{proof}






\section{Chow の補題の変種}
\label{spaces-more-morphisms-section-chows-lemma}

\noindent
この節では、非 Noetherian 代数空間の間の射を扱う Chow の補題のいくつかの
変種を証明する。Noetherian 版は Lemma \ref{spaces-more-morphisms-lemma-chow-noetherian} および
Lemma \ref{spaces-more-morphisms-lemma-chow-noetherian-separated} である。

\begin{lemma}
\label{spaces-more-morphisms-lemma-chow-finite-type}
$S$ をスキームとする。$Y$ を $S$ 上準コンパクトかつ準分離的な代数空間とする。
$f : X \to Y$ を有限型の分離的射とする。このとき次の可換図式が存在する：
$$
\xymatrix{
X \ar[rd] & X' \ar[l] \ar[d] \ar[r] & \overline{X}' \ar[ld] \\
& Y
}
$$
ここで $X' \to X$ は固有全射、$X' \to \overline{X}'$ は開埋入、
% P09-ERR-1193: frozen source line 9396 omits the article before "proper".
$\overline{X}' \to Y$ は代数空間の固有かつ表現可能な射である。
\end{lemma}

\begin{proof}
Limits of Spaces, Proposition
\ref{spaces-limits-proposition-separated-closed-in-finite-presentation} により、
閉埋入 $X \to X_1$ で、$X_1$ が $Y$ 上分離的かつ有限表示となるものを
見つけられる。$X_1 \to Y$ に対して主張を証明すれば $X$ に対して結果が
従うことは明らかである。したがって $X$ は $Y$ 上有限表示であると仮定してよい。

\medskip\noindent
基礎スキーム $S$ を $\Spec(\mathbf{Z})$ で置き換えてよく、そうする。
$Y = \lim_i Y_i$ を $\Spec(\mathbf{Z})$ 上有限型の準分離的代数空間の有向極限として
書く。Limits of Spaces, Proposition
\ref{spaces-limits-proposition-approximate} を参照せよ。Limits of Spaces, Lemma
\ref{spaces-limits-lemma-descend-finite-presentation} により、添字 $i \in I$ と
有限表示のスキーム $X_i \to Y_i$ で $X = Y \times_{Y_i} X_i$ を満たすものを
見つけられる。Limits of Spaces, Lemma
\ref{spaces-limits-lemma-descend-separated-morphism} により、$X_i \to Y_i$ は
分離的であると仮定してよい。$X_i$ に対して $Y_i$ 上で主張を証明すれば、
$X$ に対する主張が成り立つことは明らかである。$X_i \to Y_i$ の場合は
Lemma \ref{spaces-more-morphisms-lemma-chow-noetherian} によって扱われる。
\end{proof}

\begin{lemma}
\label{spaces-more-morphisms-lemma-chow-finite-type-separated}
$S$ をスキームとする。$f : X \to Y$ を $S$ 上の代数空間の射とする。
% P09-ERR-1194: frozen source line 9434 omits "is" after f.
$f$ は分離的かつ有限型であり、$Y$ は分離的かつ準コンパクトであると仮定する。
このとき次の可換図式が存在する：
$$
\xymatrix{
X \ar[rd] & X' \ar[l] \ar[d] \ar[r] & \mathbf{P}^n_Y \ar[ld] \\
& Y
}
$$
% P09-ERR-1195: frozen source line 9442 omits the article before "proper".
ここで $X' \to X$ は固有全射であり、射 $X' \to \mathbf{P}^n_Y$ は埋入である。
\end{lemma}

\begin{proof}
Limits of Spaces, Proposition
\ref{spaces-limits-proposition-separated-closed-in-finite-presentation} により、
閉埋入 $X \to X_1$ で、$X_1$ が $Y$ 上分離的かつ有限表示となるものを
見つけられる。$X_1 \to Y$ に対して主張を証明すれば $X$ に対して結果が
従うことは明らかである。したがって $X$ は $Y$ 上有限表示であると仮定してよい。

\medskip\noindent
基礎スキーム $S$ を $\Spec(\mathbf{Z})$ で置き換えてよく、そうする。
$Y = \lim_i Y_i$ を $\Spec(\mathbf{Z})$ 上有限型の準分離的代数空間の有向極限として
書く。Limits of Spaces, Proposition
\ref{spaces-limits-proposition-approximate} を参照せよ。Limits of Spaces, Lemma
\ref{spaces-limits-lemma-descend-separated} により、$Y_i$ はすべての $i$ に対して
分離的であると仮定してよい。Limits of Spaces, Lemma
\ref{spaces-limits-lemma-descend-finite-presentation} により、添字 $i \in I$ と
有限表示のスキーム $X_i \to Y_i$ で $X = Y \times_{Y_i} X_i$ を満たすものを
見つけられる。Limits of Spaces, Lemma
\ref{spaces-limits-lemma-descend-separated-morphism} により、$X_i \to Y_i$ は
分離的であると仮定してよい。$X_i$ に対して $Y_i$ 上で主張を証明すれば、
$X$ に対する主張が成り立つことは明らかである。$X_i \to Y_i$ の場合は
Lemma \ref{spaces-more-morphisms-lemma-chow-noetherian-separated} によって扱われる。
\end{proof}






\section{Grothendieck の存在定理}
\label{spaces-more-morphisms-section-existence-proper}

\noindent
この節では代数空間に対する Grothendieck の存在定理を論じる。「形式代数空間」の
理論を展開する代わりに、「Noetherian 進形式空間上の連接加群」という概念を
置き換える言語を一時的に少し整備する。

\medskip\noindent
$S$ をスキームとする。$X$ を $S$ 上の Noetherian 代数空間とする。
$\mathcal{I} \subset \mathcal{O}_X$ を準連接イデアル層とする。以下では、
逆系 $(\mathcal{F}_n)$（連接 $\mathcal{O}_X$-加群の逆系）で次を満たすものを考える：
\begin{enumerate}
\item $\mathcal{F}_n$ は $\mathcal{I}^n$ によって零化される。
\item 遷移射は同型
$\mathcal{F}_{n + 1}/\mathcal{I}^n\mathcal{F}_{n + 1} \to \mathcal{F}_n$
を誘導する。
\end{enumerate}
このような逆系の射 $\alpha : (\mathcal{F}_n) \to (\mathcal{G}_n)$ とは、単に
射の両立する系 $\alpha_n : \mathcal{F}_n \to \mathcal{G}_n$ のことである。
これらの逆系の圏を $\textit{Coh}(X, \mathcal{I})$ と表すことにする。
% P09-ERR-1196: frozen source line 9508 has the ungrammatical phrase "will parallel to".
これらの系について、スキームの場合の対応する結果と平行する理論を展開する。
Cohomology of Schemes, Sections \ref{coherent-section-existence},
\ref{coherent-section-existence-proper},
\ref{coherent-section-existence-proper-support}, and
\ref{coherent-section-algebraization} を参照せよ。

\medskip\noindent
関手性。$f : X \to Y$ をスキーム $S$ 上の Noetherian 代数空間の射とし、
$\mathcal{J} \subset \mathcal{O}_Y$ を準連接イデアル層とする。
$\mathcal{I} = f^{-1}\mathcal{J}\mathcal{O}_X$ とおく。この状況では関手
$$
f^* : \textit{Coh}(Y, \mathcal{J}) \longrightarrow \textit{Coh}(X, \mathcal{I})
$$
があり、$(\mathcal{G}_n)$ を $(f^*\mathcal{G}_n)$ へ送る。Cohomology of Schemes,
Lemma \ref{coherent-lemma-inverse-systems-pullback} と比較せよ。$f$ がエタールなら、
これは単に系を $X$ へ制限することだとみなせる。Properties of Spaces,
Equation \ref{spaces-properties-equation-restrict-modules} を参照せよ。

\medskip\noindent
エタール降下。$S$ をスキームとする。$U_0 \to X$ を Noetherian 代数空間の
全射エタール射とする。$U_1 = U_0 \times_X U_0$ および
$U_2 = U_0 \times_X U_0 \times_X U_0$ とおく。
$\mathcal{I} \subset \mathcal{O}_{X}$ を準連接イデアル層とする。
$\mathcal{I}_i = \mathcal{I}|_{U_i}$ とおく。この状況では圏の図式
$$
\xymatrix{
\textit{Coh}(X, \mathcal{I}) \ar[r] &
\textit{Coh}(U_0, \mathcal{I}_0) \ar@<0.5ex>[r] \ar@<-0.5ex>[r] &
\textit{Coh}(U_1, \mathcal{I}_1) \ar@<1ex>[r] \ar[r] \ar@<-1ex>[r] &
\textit{Coh}(U_2, \mathcal{I}_2)
}
$$
を得る。% P09-ERR-1197: frozen source line 9545 says "an the first arrow".
そして第一の射は $\textit{Coh}(X, \mathcal{I})$ を図式右側のホモトピー極限として
表示する。より正確には、{it 降下データ}、すなわち組
$((\mathcal{G}_n), \varphi)$ が与えられ、$(\mathcal{G}_n)$ が
$\textit{Coh}(U_0, \mathcal{I}_0)$ の対象であり、
$\varphi : \text{pr}_0^*(\mathcal{G}_n) \to \text{pr}_1^*(\mathcal{G}_n)$ が
$\textit{Coh}(U_1, \mathcal{I}_1)$ における同型で、
$\textit{Coh}(U_2, \mathcal{I}_2)$ におけるコサイクル条件を満たすとする。
このとき一意な対象 $(\mathcal{F}_n)$（$\textit{Coh}(X, \mathcal{I})$ の対象）で、
それに付随する標準的降下データが $((\mathcal{G}_n), \varphi)$ と同型であるものが
存在する。Descent on Spaces, Definition
\ref{spaces-descent-definition-descent-datum-effective-quasi-coherent} と比較せよ。
この主張の証明は、降下データ $(\mathcal{G}_n, \varphi_n)$ に、変化する $n$ ごとに
Descent on Spaces, Proposition
\ref{spaces-descent-proposition-fpqc-descent-quasi-coherent} を適用すれば直ちに従う。

\begin{lemma}
\label{spaces-more-morphisms-lemma-inverse-systems-abelian}
$S$ をスキームとする。$X$ を $S$ 上の Noetherian 代数空間とし、
$\mathcal{I} \subset \mathcal{O}_X$ を準連接イデアル層とする。
\begin{enumerate}
\item 圏 $\textit{Coh}(X, \mathcal{I})$ は Abel 圏である。
\item $\textit{Coh}(X, \mathcal{I})$ における完全性はエタール局所的に確認できる。
\item Noetherian 代数空間の任意の平坦射 $f : X' \to X$ に対し、関手
$f^* : \textit{Coh}(X, \mathcal{I}) \to
\textit{Coh}(X', f^{-1}\mathcal{I}\mathcal{O}_{X'})$ は完全である。
\end{enumerate}
\end{lemma}

\begin{proof}
(1) の証明。アフィンスキーム $U_0$ と全射エタール射 $U_0 \to X$ を選ぶ。
上のエタール降下の議論と同様に $U_1 = U_0 \times_X U_0$ および
$U_2 = U_0 \times_X U_0 \times_X U_0$ とおく。圏
$\textit{Coh}(U_i, \mathcal{I}_i)$ は Abel 圏であり
（Cohomology of Schemes, Lemma
\ref{coherent-lemma-inverse-systems-abelian}）、引き戻し関手
$\textit{Coh}(U_0, \mathcal{I}_0) \to \textit{Coh}(U_1, \mathcal{I}_1)$ および
$\textit{Coh}(U_1, \mathcal{I}_1) \to \textit{Coh}(U_2, \mathcal{I}_2)$ は完全関手である
（Cohomology of Schemes, Lemma
\ref{coherent-lemma-inverse-systems-pullback}）。したがって、
$\textit{Coh}(X, \mathcal{I})$ を降下データの圏として記述したことから補題が
形式的に従う。いくつかの細部は省略する。Groupoids, Lemma
\ref{groupoids-lemma-abelian} の証明と比較せよ。

\medskip\noindent
(2) は前段落の議論から直ちに従う。(3) の状況では可換図式
$$
\xymatrix{
U' \ar[d] \ar[r] & U \ar[d] \\
X' \ar[r] & X
}
$$
を選ぶ。ここで $U'$ と $U$ はアフィンスキームであり、縦の射は全射エタールである。
このとき $U' \to U$ は Noetherian スキームの平坦射である
（Morphisms of Spaces, Lemma \ref{spaces-morphisms-lemma-flat-local}）。したがって
引き戻し関手
$\textit{Coh}(U, \mathcal{I}\mathcal{O}_U) \to
\textit{Coh}(U', \mathcal{I}\mathcal{O}_{U'})$ は Cohomology of Schemes, Lemma
\ref{coherent-lemma-inverse-systems-pullback} により完全である。
% P09-ERR-1198: frozen source line 9610 writes Coh(X,O_X) where the defined category uses I.
$\textit{Coh}(X, \mathcal{O}_X)$ における完全性は $U$ 上で確認でき、
$X', U'$ についても同様なので、主張が従う。
\end{proof}

\begin{lemma}
\label{spaces-more-morphisms-lemma-inverse-systems-surjective}
$S$ をスキームとする。$X$ を $S$ 上の Noetherian 代数空間とし、
$\mathcal{I} \subset \mathcal{O}_X$ を準連接イデアル層とする。写像
$(\mathcal{F}_n) \to (\mathcal{G}_n)$ が $\textit{Coh}(X, \mathcal{I})$ において
全射であることと、$\mathcal{F}_1 \to \mathcal{G}_1$ が全射であることは同値である。
\end{lemma}

\begin{proof}
Lemma \ref{spaces-more-morphisms-lemma-inverse-systems-abelian} により、$X$ のアフィンエタール被覆上で
確認できる。したがってスキームの場合へ帰着し、これは Cohomology of Schemes,
Lemma \ref{coherent-lemma-inverse-systems-surjective} である。
\end{proof}

\noindent
$S$ をスキームとする。$X$ を $S$ 上の Noetherian 代数空間とし、
$\mathcal{I} \subset \mathcal{O}_X$ を準連接イデアル層とする。関手
\begin{equation}
\label{spaces-more-morphisms-equation-completion-functor}
\textit{Coh}(\mathcal{O}_X) \longrightarrow \textit{Coh}(X, \mathcal{I}), \quad
\mathcal{F} \longmapsto \mathcal{F}^\wedge
\end{equation}
があり、連接 $\mathcal{O}_X$-加群 $\mathcal{F}$ に対象
$\mathcal{F}^\wedge = (\mathcal{F}/\mathcal{I}^n\mathcal{F})$
（$\textit{Coh}(X, \mathcal{I})$ の対象）を対応させる。

\begin{lemma}
\label{spaces-more-morphisms-lemma-exact}
関手 (\ref{spaces-more-morphisms-equation-completion-functor}) は完全である。
\end{lemma}

\begin{proof}
Lemma \ref{spaces-more-morphisms-lemma-inverse-systems-abelian} により、これを $X$ 上エタール局所的に
確認すれば十分である。したがってスキームの場合へ帰着し、これは
Cohomology of Schemes, Lemma \ref{coherent-lemma-exact} である。
\end{proof}

\begin{lemma}
\label{spaces-more-morphisms-lemma-completion-internal-hom}
$S$ をスキームとする。$X$ を $S$ 上の Noetherian 代数空間とし、
$\mathcal{I} \subset \mathcal{O}_X$ を準連接イデアル層とする。
$\mathcal{F}$, $\mathcal{G}$ を連接 $\mathcal{O}_X$-加群とする。
$\mathcal{H} = \SheafHom_{\mathcal{O}_X}(\mathcal{F}, \mathcal{G})$ とおく。このとき
$$
\lim H^0(X, \mathcal{H}/\mathcal{I}^n\mathcal{H}) =
\Mor_{\textit{Coh}(X, \mathcal{I})}
(\mathcal{F}^\wedge, \mathcal{G}^\wedge).
$$
\end{lemma}

\begin{proof}
$\mathcal{H}$ は $X_\etale$ 上の層であり、かつ
$\textit{Coh}(X, \mathcal{I})$ の対象にはエタール降下があるので、これを
エタール局所的に証明すれば十分である。したがってスキームの場合へ帰着し、これは
Cohomology of Schemes, Lemma
\ref{coherent-lemma-completion-internal-hom} である。
\end{proof}

\noindent
この節の残りを通じて焦点を当てる設定を導入する。

\begin{situation}
\label{spaces-more-morphisms-situation-existence}
ここで $A$ はイデアル $I$ に関して完備な Noetherian 環である。また
$f : X \to \Spec(A)$ は代数空間の有限型分離的射であり、
$\mathcal{I} = I\mathcal{O}_X$ である。
\end{situation}

\noindent
この状況では
$$
\textit{Coh}_{\text{support proper over } A}(\mathcal{O}_X)
$$
% P09-ERR-1199: frozen source lines 9689-9693 use "denote ... be" instead of "denote by ... the".
で $\textit{Coh}(\mathcal{O}_X)$ の充満部分圏を表す。その対象は、連接
$\mathcal{O}_X$-加群で台が $\Spec(A)$ 上固有なもの、同値な言い方では
スキーム論的台が $\Spec(A)$ 上固有なものからなる。Derived Categories of Spaces,
Lemma \ref{spaces-perfect-lemma-module-support-proper-over-base} を参照せよ。
同様に
$$
\textit{Coh}_{\text{support proper over } A}(X, \mathcal{I})
$$
を $\textit{Coh}(X, \mathcal{I})$ の充満部分圏で、対象 $(\mathcal{F}_n)$ のうち
$\mathcal{F}_1$ の台が $\Spec(A)$ 上固有なものからなるものとする。商加群の台は
もとの加群の台に含まれるので、(\ref{spaces-more-morphisms-equation-completion-functor}) は関手
\begin{equation}
\label{spaces-more-morphisms-equation-completion-functor-proper-over-A}
\textit{Coh}_{\text{support proper over }A}(\mathcal{O}_X)
\longrightarrow
\textit{Coh}_{\text{support proper over }A}(X, \mathcal{I})
\end{equation}
を誘導する。最初の結果は、この関手が忠実充満であるというものである。

\begin{lemma}
\label{spaces-more-morphisms-lemma-fully-faithful}
Situation \ref{spaces-more-morphisms-situation-existence} において、$\mathcal{F}$, $\mathcal{G}$ を連接
$\mathcal{O}_X$-加群とする。$\mathcal{F}$ と $\mathcal{G}$ の台の共通部分が
$\Spec(A)$ 上固有であると仮定する。このとき (\ref{spaces-more-morphisms-equation-completion-functor})
から来る写像
$$
\Mor_{\textit{Coh}(\mathcal{O}_X)}(\mathcal{F}, \mathcal{G})
\longrightarrow
\Mor_{\textit{Coh}(X, \mathcal{I})}
(\mathcal{F}^\wedge, \mathcal{G}^\wedge)
$$
は全単射である。特に (\ref{spaces-more-morphisms-equation-completion-functor-proper-over-A}) は忠実充満である。
\end{lemma}

\begin{proof}
% P09-ERR-1200: frozen source lines 9735 and 9744 reverse F and G relative to the lemma's Hom map.
$\mathcal{H} = \SheafHom_{\mathcal{O}_X}(\mathcal{G}, \mathcal{F})$ とおく。
これは連接 $\mathcal{O}_X$-加群である。% P09-ERR-1201: frozen source line 9737 says "restriction of schemes" instead of "restriction to schemes".
実際、$X$ 上エタールなスキームへの制限は Modules, Lemma
\ref{modules-lemma-internal-hom-locally-kernel-direct-sum} により連接である。
Lemma \ref{spaces-more-morphisms-lemma-completion-internal-hom} により写像
$$
\lim_n H^0(X, \mathcal{H}/\mathcal{I}^n\mathcal{H})
\to
\Mor_{\textit{Coh}(X, \mathcal{I})}
(\mathcal{G}^\wedge, \mathcal{F}^\wedge)
$$
は全単射である。$i : Z \to X$ を $\mathcal{H}$ のスキーム論的台とする。
$Z$ は閉部分空間で、$|Z|$ は $\mathcal{F}$ と $\mathcal{G}$ の台の共通部分に
含まれることは明らかである。したがって仮定により $Z \to \Spec(A)$ は固有である
（Derived Categories of Spaces, Section
\ref{spaces-perfect-section-proper-over-base} を参照せよ）。
$\mathcal{H} = i_*\mathcal{H}'$ と書く。ここで $\mathcal{O}_Z$-加群
$\mathcal{H}'$ は連接である。
次が成り立つ：
$i_*(\mathcal{H}'/I^n\mathcal{H}') = \mathcal{H}/I^n\mathcal{H}$。
したがって
\begin{align*}
\lim_n H^0(X, \mathcal{H}/\mathcal{I}^n\mathcal{H})
& =
\lim_n H^0(Z, \mathcal{H}'/\mathcal{I}^n\mathcal{H}') \\
& =
H^0(Z, \mathcal{H}') \\
& =
H^0(X, \mathcal{H}) \\
& = 
\Mor_{\textit{Coh}(\mathcal{O}_X)}(\mathcal{F}, \mathcal{G})
\end{align*}
を得る。第二の等号には形式関数定理
（Cohomology of Spaces, Lemma
\ref{spaces-cohomology-lemma-spell-out-theorem-formal-functions}）を用いた。
これで補題が証明された。
\end{proof}

\begin{remark}
\label{spaces-more-morphisms-remark-inverse-systems-kernel-cokernel-annihilated-by}
$S$ をスキームとする。$X$ を $S$ 上の Noetherian 代数空間とし、
$\mathcal{I}, \mathcal{K} \subset \mathcal{O}_X$ を準連接イデアル層とする。
$\alpha : (\mathcal{F}_n) \to (\mathcal{G}_n)$ を
$\textit{Coh}(X, \mathcal{I})$ の射とする。アフィンスキーム
$U = \Spec(A)$ と全射エタール射 $U \to X$ が与えられたとき、
% P09-ERR-1202: frozen source lines 9780-9782 omit "by" in both Denote constructions.
$I, K \subset A$ で制限 $\mathcal{I}|_U, \mathcal{K}|_U$ に対応するイデアルを
表す。また $\alpha_U : M \to N$ で、有限 $A^\wedge$-加群の射であり、
Cohomology of Schemes, Lemma
\ref{coherent-lemma-inverse-systems-affine} を介して $\alpha|_U$ に対応するものを表す。
次が同値であると主張する：
\begin{enumerate}
\item ある整数 $t \geq 1$ が存在し、$\Ker(\alpha_n)$ と $\Coker(\alpha_n)$ が
$\mathcal{K}^t$ により、すべての $n \geq 1$ に対して零化される。
\item 上のような任意の（またはいずれかの）アフィン開集合
$\Spec(A) = U \subset X$ に対して、加群 $\Ker(\alpha_U)$ と
$\Coker(\alpha_U)$ が $K^t$ により、ある整数 $t \geq 1$ に対して零化される。
\end{enumerate}
これらの同値な条件が成り立つとき、$\alpha$ を「核と余核が
$\mathcal{K}$ のある冪で零化される写像」と呼ぶ。同値性については
Cohomology of Schemes, Remark
\ref{coherent-remark-inverse-systems-kernel-cokernel-annihilated-by} を参照せよ。
\end{remark}

\begin{lemma}
\label{spaces-more-morphisms-lemma-existence-easy}
$S$ をスキームとする。$X$ を $S$ 上の Noetherian 代数空間とし、
$\mathcal{I} \subset \mathcal{O}_X$ を準連接イデアル層とする。
$\mathcal{G}$ を連接 $\mathcal{O}_X$-加群、$(\mathcal{F}_n)$ を
$\textit{Coh}(X, \mathcal{I})$ の対象とし、
$\alpha : (\mathcal{F}_n) \to \mathcal{G}^\wedge$ を、核と余核が
$\mathcal{I}$ のある冪で零化される写像とする。このとき次の三つ組
$(\mathcal{F}, a, \beta)$ が、一意な同型を除いて一意に存在する：
\begin{enumerate}
\item $\mathcal{F}$ は連接 $\mathcal{O}_X$-加群である。
\item $a : \mathcal{F} \to \mathcal{G}$ は $\mathcal{O}_X$-加群の写像で、
その核と余核は $\mathcal{I}$ のある冪で零化される。
\item $\beta : (\mathcal{F}_n) \to \mathcal{F}^\wedge$ は同型である。
\item $\alpha = a^\wedge \circ \beta$ である。
\end{enumerate}
\end{lemma}

\begin{proof}
$\textit{Coh}(X, \mathcal{I})$ と $\textit{Coh}(\mathcal{O}_X)$ の対象の一意性と
エタール降下により、$(\mathcal{F}, a, \beta)$ を $X$ 上エタール局所的に
構成すれば十分である。したがってスキームの場合へ帰着し、これは
Cohomology of Schemes, Lemma \ref{coherent-lemma-existence-easy} である。
\end{proof}

\begin{lemma}
\label{spaces-more-morphisms-lemma-existence-tricky}
Situation \ref{spaces-more-morphisms-situation-existence} において、
$\mathcal{K} \subset \mathcal{O}_X$ を準連接イデアル層とする。
$X_e \subset X$ を $\mathcal{K}^e$ で切り出される閉部分空間とし、
$\mathcal{I}_e = \mathcal{I}\mathcal{O}_{X_e}$ とおく。$(\mathcal{F}_n)$ を
$\textit{Coh}_{\text{support proper over } A}(X, \mathcal{I})$ の対象とする。
次を仮定する：
\begin{enumerate}
\item 関手
$\textit{Coh}_{\text{support proper over } A}(\mathcal{O}_{X_e})
\to \textit{Coh}_{\text{support proper over } A}(X_e, \mathcal{I}_e)$
は、すべての $e \geq 1$ に対して圏同値である。
\item 対象 $\mathcal{H}$（
$\textit{Coh}_{\text{support proper over } A}(\mathcal{O}_X)$ の対象）と写像
$\alpha : (\mathcal{F}_n) \to \mathcal{H}^\wedge$ が存在し、その核と余核は
$\mathcal{K}$ のある冪で零化される。
\end{enumerate}
このとき $(\mathcal{F}_n)$ は
(\ref{spaces-more-morphisms-equation-completion-functor-proper-over-A}) の本質像に属する。
\end{lemma}

\begin{proof}
この証明では、閉埋入 $i : Z \to X$ に対し関手 $i_*$ が、$Z$ 上の連接加群の圏と、
$X$ 上の連接加群のうち $Z$ のイデアル層で零化されるものの圏との圏同値を与えることを、
今後断りなく用いる。Cohomology of Spaces, Lemma
\ref{spaces-cohomology-lemma-i-star-equivalence} を参照せよ。特に
$$
\textit{Coh}_{\text{support proper over } A}(\mathcal{O}_{X_e})
\subset
\textit{Coh}_{\text{support proper over } A}(\mathcal{O}_X)
$$
を、$\mathcal{K}^e$ で零化される加群からなる充満部分圏とみなし、また
$$
\textit{Coh}_{\text{support proper over } A}(X_e, \mathcal{I}_e)
\subset
\textit{Coh}_{\text{support proper over } A}(X, \mathcal{I})
$$
を、$\mathcal{K}^e$ で零化される対象からなる充満部分圏とみなす。
% P09-ERR-1203: frozen source line 9861 says "subcategory of consisting of".
さらに (1) により、これら二圏は完備化関手
(\ref{spaces-more-morphisms-equation-completion-functor-proper-over-A}) のもとで圏同値である。

\medskip\noindent
この圏同値を適用すると、連接 $\mathcal{O}_X$-加群 $\mathcal{G}_e$ で
$\mathcal{K}^e$ により零化され、系
$(\mathcal{F}_n/\mathcal{K}^e\mathcal{F}_n)$（
$\textit{Coh}_{\text{support proper over } A}(X, \mathcal{I})$ の系）に
対応するものを得る。写像
$\mathcal{F}_n/\mathcal{K}^{e + 1}\mathcal{F}_n \to
\mathcal{F}_n/\mathcal{K}^e\mathcal{F}_n$ は標準写像
$\mathcal{G}_{e + 1} \to \mathcal{G}_e$ に対応し、これは同型
$\mathcal{G}_{e + 1}/\mathcal{K}^e\mathcal{G}_{e + 1} \to \mathcal{G}_e$ を
誘導する。したがって対象 $(\mathcal{G}_e)$（圏
$\textit{Coh}_{\text{support proper over } A}(X, \mathcal{K})$ の対象）を得る。
写像 $\alpha$ は写像の系
$$
\mathcal{F}_n/\mathcal{K}^e\mathcal{F}_n
\longrightarrow
\mathcal{H}/(\mathcal{I}^n + \mathcal{K}^e)\mathcal{H}
$$
を誘導し、したがって（再び圏同値により）写像
$\mathcal{G}_e \to \mathcal{H}/\mathcal{K}^e\mathcal{H}$ を誘導する。
仮定 (2) により存在する整数 $t \geq 1$ で、$\mathcal{K}^t$ がすべての写像
$\mathcal{F}_n \to \mathcal{H}/\mathcal{I}^n\mathcal{H}$ の核と余核を
零化するものを取る。このとき $\mathcal{K}^{2t}$ が写像
$\mathcal{F}_n/\mathcal{K}^e\mathcal{F}_n \to
\mathcal{H}/(\mathcal{I}^n + \mathcal{K}^e)\mathcal{H}$ の核と余核を零化する
（細部は省略する。Cohomology of Schemes, Remark
\ref{coherent-remark-inverse-systems-kernel-cokernel-annihilated-by} を参照せよ）。
そこから $\mathcal{K}^{4t}$ は写像
$$
\mathcal{G}_e
\longrightarrow
\mathcal{H}/\mathcal{K}^e\mathcal{H},
$$
の核と余核を零化すると結論する
（細部は省略する。Cohomology of Schemes, Remark
\ref{coherent-remark-inverse-systems-kernel-cokernel-annihilated-by} を参照せよ）。
Lemma \ref{spaces-more-morphisms-lemma-existence-easy} を適用して、連接 $\mathcal{O}_X$-加群
$\mathcal{F}$、写像 $a : \mathcal{F} \to \mathcal{H}$、および同型
$\beta : (\mathcal{G}_e) \to (\mathcal{F}/\mathcal{K}^e\mathcal{F})$
（$\textit{Coh}(X, \mathcal{K})$ における同型）を得る。
逆向きに辿ると、与えられた $n$ に対して三つ組
$(\mathcal{F}/\mathcal{I}^n\mathcal{F}, a \bmod \mathcal{I}^n, \beta
\bmod \mathcal{I}^n)$ は、射 $\alpha_n \bmod \mathcal{K}^e :
(\mathcal{F}_n/\mathcal{K}^e\mathcal{F}_n) \to
(\mathcal{H}/(\mathcal{I}^n + \mathcal{K}^e)\mathcal{H})$
（$\textit{Coh}(X, \mathcal{K})$ の射）に対する
補題のような三つ組である。したがって Lemma \ref{spaces-more-morphisms-lemma-existence-easy} の一意性から、
視野にあるすべての射と両立する標準同型
$\mathcal{F}/\mathcal{I}^n\mathcal{F} \to \mathcal{F}_n$ を得る。

\medskip\noindent
補題の証明を完了するには、$\mathcal{F}$ の台が $A$ 上固有であることを示す必要がある。
構成により $a : \mathcal{F} \to \mathcal{H}$ の核は $\mathcal{K}$ のある冪で
零化される。したがってこの核の台は $\mathcal{G}_1$ の台に含まれる。
$\mathcal{G}_1$ は
$\textit{Coh}_{\text{support proper over } A}(\mathcal{O}_{X_1})$ の対象なので、
これは $A$ 上固有である。$\mathcal{H}$ の台が $A$ 上固有であるという事実と
合わせると、Derived Categories of Spaces, Lemma
\ref{spaces-perfect-lemma-union-closed-proper-over-base} により
$\mathcal{F}$ の台は $A$ 上固有であると結論する。
\end{proof}

\begin{lemma}
\label{spaces-more-morphisms-lemma-inverse-systems-push-pull}
$S$ をスキームとする。$f : X \to Y$ を $S$ 上の Noetherian 代数空間の
表現可能な固有射とする。$\mathcal{J}, \mathcal{K} \subset \mathcal{O}_Y$ を
準連接イデアル層とする。$f$ は $V = Y \setminus V(\mathcal{K})$ 上同型であると
仮定する。$\mathcal{I} = f^{-1}\mathcal{J} \mathcal{O}_X$ とおく。
$(\mathcal{G}_n)$ を $\textit{Coh}(Y, \mathcal{J})$ の対象、$\mathcal{F}$ を
連接 $\mathcal{O}_X$-加群とし、
$\beta : (f^*\mathcal{G}_n)  \to \mathcal{F}^\wedge$ を
$\textit{Coh}(X, \mathcal{I})$ における同型とする。このとき写像
$$
\alpha :
(\mathcal{G}_n)
\longrightarrow
(f_*\mathcal{F})^\wedge
$$
で、$\textit{Coh}(Y, \mathcal{J})$ に属し、その核と余核が
$\mathcal{K}$ のある冪で零化されるものが存在する。
\end{lemma}

\begin{proof}
$f$ は固有射なので、$f_*\mathcal{F}$ は連接 $\mathcal{O}_Y$-加群である
（Cohomology of Spaces, Lemma
\ref{spaces-cohomology-lemma-proper-pushforward-coherent}）。
したがって補題の主張には意味がある。合成
$$
\gamma_n : \mathcal{G}_n \to
f_*f^*\mathcal{G}_n \to
f_*(\mathcal{F}/\mathcal{I}^n\mathcal{F}).
$$
を考える。ここで第一の写像は随伴写像であり、第二は $f_*\beta_n$ である。
補題のような一意な $\alpha$ が存在し、合成
$$
\mathcal{G}_n \xrightarrow{\alpha_n}
f_*\mathcal{F}/\mathcal{J}^nf_*\mathcal{F} \to
f_*(\mathcal{F}/\mathcal{I}^n\mathcal{F})
$$
が $\gamma_n$ に、すべての $n$ に対して等しくなると主張する。一意性と
$\textit{Coh}(Y, \mathcal{J})$ に対するエタール降下により、これを $Y$ 上
エタール局所的に証明すれば十分である。したがって $Y$ は Noetherian 環の
スペクトルであると仮定してよい。$f$ は表現可能なので、$X$ もスキームである。
ゆえにスキームの場合へ帰着する。Cohomology of Schemes, Lemma
\ref{coherent-lemma-inverse-systems-push-pull} の証明を参照せよ。
\end{proof}

\begin{theorem}[Grothendieck の存在定理]
\label{spaces-more-morphisms-theorem-grothendieck-existence}
Situation \ref{spaces-more-morphisms-situation-existence} において、関手
(\ref{spaces-more-morphisms-equation-completion-functor-proper-over-A}) は圏同値である。
\end{theorem}

\begin{proof}
定理の証明では、Cohomology of Spaces, Lemma
\ref{spaces-cohomology-lemma-i-star-equivalence} の圏同値を今後断りなく用いる。
Lemma \ref{spaces-more-morphisms-lemma-fully-faithful} により関手は忠実充満である。したがって、関手が
本質的全射であることを証明すればよい。

\medskip\noindent
集合 $\Xi$ を、準連接イデアル層 $\mathcal{K} \subset \mathcal{O}_X$ で、
すべての対象 $(\mathcal{F}_n)$（
$\textit{Coh}_{\text{support proper over }A}(X, \mathcal{I})$ の対象）のうち
$\mathcal{K}$ によって零化されるものに対して主張が成り立つもの全体とする。
$(0)$ が $\Xi$ に属することを示したい。そうでなければ、$X$ は Noetherian なので、
極大準連接イデアル層 $\mathcal{K}$ で $\Xi$ に属さないものが存在する。
Cohomology of Spaces, Lemma \ref{spaces-cohomology-lemma-acc-coherent} を参照せよ。
$X$ を $X$ の閉部分スキーム（$\mathcal{K}$ に対応するもの）で置き換えると、すべての非零
$\mathcal{K}$ は $\Xi$ に属すると仮定してよい。$(\mathcal{F}_n)$ を
$\textit{Coh}_{\text{support proper over }A}(X, \mathcal{I})$ の対象とする。
この対象が本質像に属することを示し、それにより定理の証明を完了する。

\medskip\noindent
Chow の補題（Lemma \ref{spaces-more-morphisms-lemma-chow-noetherian-separated}）を適用して、
固有全射 $f : Y \to X$ で、稠密開集合 $U \subset X$ 上同型となり、$Y$ が
$A$ 上 H-準射影的となるものを見つける。$Y$ はスキームで、$f$ は表現可能である
ことに注意する。開埋入 $j : Y \to Y'$ で、$Y'$ が $A$ 上射影的となるものを選ぶ。
Morphisms, Lemma \ref{morphisms-lemma-H-quasi-projective-open-H-projective} を参照せよ。
$T_n$ を $\mathcal{F}_n$ のスキーム論的台とする。$|T_n| = |T_1|$ なので、
$T_n$ は $A$ 上、すべての $n$ に対して固有であることに注意する
（Morphisms of Spaces, Lemma
\ref{spaces-morphisms-lemma-image-proper-is-proper}）。このとき
$f^*\mathcal{F}_n$ は閉部分スキーム $f^{-1}T_n$ 上に台をもち、後者は $A$ 上固有である
（Morphisms of Spaces, Lemma \ref{spaces-morphisms-lemma-composition-proper} と
$f$ の固有性による）。特に合成 $f^{-1}T_n \to Y \to Y'$ は閉である
（Morphisms, Lemma \ref{morphisms-lemma-image-proper-scheme-closed}）。
$T'_n \subset Y'$ を対応する閉部分スキームとする。これは開部分スキーム $Y$ に
含まれ、$f^{-1}T_n$ に $Y$ の閉部分スキームとして等しい。
$\mathcal{F}_n'$ を連接 $\mathcal{O}_{Y'}$-加群で、$f^*\mathcal{F}_n$ を
$Y'$ 上の連接加群として閉埋入 $f^{-1}T_n = T'_n \subset Y'$ を介して
みなしたものに対応するものとする。このとき $(\mathcal{F}_n')$ は
$\textit{Coh}(Y', I\mathcal{O}_{Y'})$ の対象である。Grothendieck の存在定理の
射影的な場合（Cohomology of Schemes, Lemma
\ref{coherent-lemma-existence-projective}）により、連接 $\mathcal{O}_{Y'}$-加群
$\mathcal{F}'$ と同型 $(\mathcal{F}')^\wedge \cong (\mathcal{F}'_n)$
（$\textit{Coh}(Y', I\mathcal{O}_{Y'})$ における同型）が存在する。
$Z' \subset Y'$ を $\mathcal{F}'$ のスキーム論的台とする。
$\mathcal{F}'/I\mathcal{F}' = \mathcal{F}'_1$ なので、集合論的に
$Z' \cap V(I\mathcal{O}_{Y'}) = T'_1$ である。構造射
$p' : Y' \to \Spec(A)$ は固有なので、
$p'(Z' \cap (Y' \setminus Y))$ は $\Spec(A)$ において閉である。空でなければ、
$V(I)$ の点を含む。実際、$I$ は $A$ の Jacobson 根基に含まれる
（Algebra, Lemma \ref{algebra-lemma-radical-completion}）。しかし上で
$Z' \cap (p')^{-1}V(I) = T'_1 \subset Y$ を見たので、$Z' \subset Y$ と結論する。
したがって $\mathcal{F}'|_Y$ は $Y$ の閉部分スキーム上に台をもち、その閉部分スキームは
$A$ 上固有である。

\medskip\noindent
$\mathcal{K}$ を、被約な補集合 $X \setminus U$ を切り出す準連接イデアル層とする。
Cohomology of Spaces, Lemma
\ref{spaces-cohomology-lemma-proper-pushforward-coherent} により、
$\mathcal{O}_X$-加群 $\mathcal{H} = f_*\mathcal{F}'$ は連接であり、Lemma
\ref{spaces-more-morphisms-lemma-inverse-systems-push-pull} により、射
$\alpha : (\mathcal{F}_n) \to \mathcal{H}^\wedge$
（圏 $\textit{Coh}_{\text{support proper over } A}(X, \mathcal{I})$ の射）で、その核と余核が
$\mathcal{K}$ のある冪で零化されるものが存在する。$Z_0 \subset X$ を
$\mathcal{H}$ のスキーム論的台とする。$|Z_0| \subset f(|Z'|)$ であることは明らかである。
したがって $Z_0 \to \Spec(A)$ は固有である
（Morphisms of Spaces, Lemma
\ref{spaces-morphisms-lemma-image-proper-is-proper}）。ゆえに $\mathcal{H}$ は
$\textit{Coh}_{\text{support proper over } A}(\mathcal{O}_X)$ の対象である。
イデアル層 $\mathcal{K}^e$ の各々が $\Xi$ の元なので、Lemma
\ref{spaces-more-morphisms-lemma-existence-tricky} の仮定が満たされ、結論を得る。
\end{proof}

\begin{remark}[Grothendieck の存在定理の展開]
\label{spaces-more-morphisms-remark-reformulate-existence-theorem}
$A$ をイデアル $I$ に関して完備な Noetherian 環とする。
$S = \Spec(A)$ および $S_n = \Spec(A/I^n)$ と書く。
$X \to S$ を分離的かつ有限型な代数空間の射とする。$n \geq 1$ に対し
$X_n = X \times_S S_n$ とおく。図式は次のとおりである：
$$
\xymatrix{
X_1 \ar[r]_{i_1} \ar[d] & X_2 \ar[r]_{i_2} \ar[d] & X_3 \ar[r] \ar[d] &
\ldots & X \ar[d] \\
S_1 \ar[r] & S_2 \ar[r] & S_3 \ar[r] & \ldots & S
}
$$
この状況では系 $(\mathcal{F}_n, \varphi_n)$ で次を満たすものを考える：
\begin{enumerate}
\item $\mathcal{F}_n$ は連接 $\mathcal{O}_{X_n}$-加群である。
\item $\varphi_n : i_n^*\mathcal{F}_{n + 1} \to \mathcal{F}_n$ は同型である。
\item $\text{Supp}(\mathcal{F}_1)$ は $S_1$ 上固有である。
\end{enumerate}
Theorem \ref{spaces-more-morphisms-theorem-grothendieck-existence} は、完備化関手
$$
\begin{matrix}
\text{連接 }\mathcal{O}_X\text{-加群 }\mathcal{F} \\
\text{台が固有となる基礎 }A
\end{matrix}
\quad
\longrightarrow
\quad
\begin{matrix}
\text{上のような系 }(\mathcal{F}_n) \\
\text{上記のとおり}
\end{matrix}
$$
が圏同値であることを述べる。特に $X$ が $A$ 上固有な場合には、台に関する
条件を省ける。
\end{remark}







\section{Grothendieck の代数化定理}
\label{spaces-more-morphisms-section-algebraization}

\noindent
この節は Cohomology of Schemes, Section
\ref{coherent-section-algebraization} の類似である。ただし、豊富な可逆層を備えた
固有代数空間の変形の代数化に関する結果は欠けている。豊富な可逆層を備えた
固有代数空間は既にスキームだからである。相対次元 $1$ の固有代数空間については
代数化の結果がある。最初の結果は、Grothendieck の存在定理を閉部分スキームと
有限射の言葉へ翻訳したものである。

\begin{lemma}
\label{spaces-more-morphisms-lemma-algebraize-formal-closed-subscheme}
$A$ をイデアル $I$ に関して完備な Noetherian 環とする。
$S = \Spec(A)$ および $S_n = \Spec(A/I^n)$ と書く。
$X \to S$ を分離的かつ有限型な代数空間の射とする。$n \geq 1$ に対し
$X_n = X \times_S S_n$ とおく。Cartesian な正方形からなる代数空間の可換図式
$$
\xymatrix{
Z_1 \ar[r] \ar[d] & Z_2 \ar[r] \ar[d] & Z_3 \ar[r] \ar[d] & \ldots \\
X_1 \ar[r]^{i_1} & X_2 \ar[r]^{i_2} & X_3 \ar[r] & \ldots
}
$$
が与えられたとする。次を仮定する：
\begin{enumerate}
\item $Z_1 \to X_1$ は閉埋入である。
\item $Z_1 \to S_1$ は固有である。
\end{enumerate}
このとき代数空間の閉埋入 $Z \to X$ で、$Z_n = Z \times_S S_n$ が
すべての $n \geq 1$ に対して成り立つものが存在する。さらに $Z$ は $S$ 上固有である。
\end{lemma}

\begin{proof}
縦の射を $j_n : Z_n \to X_n$ と書こう。主張中の正方形は Cartesian なので、
$j_n$ の $X_1$ への基底変換は $j_1$ である。したがって Limits of Spaces, Lemma
\ref{spaces-limits-lemma-check-closed-infinitesimally} により、$j_n$ は閉埋入である。
$\mathcal{F}_n = j_{n, *}\mathcal{O}_{Z_n}$ とおく。このとき $j_n^\sharp$ は全射
$\mathcal{O}_{X_n} \to \mathcal{F}_n$ である。再び正方形が Cartesian であることから、
$\mathcal{F}_{n + 1}$ の $X_n$ への引き戻しは $\mathcal{F}_n$ である。したがって
Remark \ref{spaces-more-morphisms-remark-reformulate-existence-theorem} の形の Grothendieck の存在定理により、
写像 $\mathcal{O}_X \to \mathcal{F}$（連接 $\mathcal{O}_X$-加群の写像）で、その $X_n$ への
制限が $\mathcal{O}_{X_n} \to \mathcal{F}_n$ を復元するものが存在する。さらに
$\mathcal{F}$ の台は $S$ 上固有である。完備化関手は完全なので
（Lemma \ref{spaces-more-morphisms-lemma-exact}）、$\mathcal{O}_X \to \mathcal{F}$ は全射である。
したがって $\mathcal{F} = \mathcal{O}_X/\mathcal{J}$ である。ここで
$\mathcal{J}$ はある準連接イデアル層である。$Z = V(\mathcal{J})$ とおけば
証明が完了する。
\end{proof}





























\begin{lemma}
\label{spaces-more-morphisms-lemma-algebraize-formal-algebraic-space-finite-over-proper}
$A$ をイデアル $I$ に関して完備な Noetherian 環とする。
$S = \Spec(A)$ および $S_n = \Spec(A/I^n)$ と書く。
$X \to S$ を分離的かつ有限型な代数空間の射とする。$n \geq 1$ に対し
$X_n = X \times_S S_n$ とおく。Cartesian な正方形からなる代数空間の可換図式
$$
\xymatrix{
Y_1 \ar[r] \ar[d] & Y_2 \ar[r] \ar[d] & Y_3 \ar[r] \ar[d] & \ldots \\
X_1 \ar[r]^{i_1} & X_2 \ar[r]^{i_2} & X_3 \ar[r] & \ldots
}
$$
が与えられたとする。次を仮定する：
\begin{enumerate}
\item $Y_1 \to X_1$ は有限射である。
\item $Y_1 \to S_1$ は固有である。
\end{enumerate}
このとき代数空間の有限射 $Y \to X$ で、$Y_n = Y \times_S S_n$ が
すべての $n \geq 1$ に対して成り立つものが存在する。さらに $Y$ は $S$ 上固有である。
\end{lemma}

\begin{proof}
縦の射を $f_n : Y_n \to X_n$ と書こう。主張中の正方形は Cartesian なので、
$f_n$ の $X_1$ への基底変換は $f_1$ である。したがって Lemma
\ref{spaces-more-morphisms-lemma-thicken-property-morphisms-cartesian} により $f_n$ は有限射である。
$\mathcal{F}_n = f_{n, *}\mathcal{O}_{Y_n}$ とおく。正方形が Cartesian であることから、
$\mathcal{F}_{n + 1}$ の $X_n$ への引き戻しは $\mathcal{F}_n$ である。したがって
Remark \ref{spaces-more-morphisms-remark-reformulate-existence-theorem} の形の Grothendieck の存在定理により、
連接 $\mathcal{O}_X$-加群 $\mathcal{F}$ で、その $X_n$ への制限が
$\mathcal{F}_n$ を復元するものが存在する。さらに $\mathcal{F}$ の台は $S$ 上固有である。
完備化関手は忠実充満なので（Theorem \ref{spaces-more-morphisms-theorem-grothendieck-existence}）、乗法写像
$\mathcal{F}_n \otimes_{\mathcal{O}_{X_n}} \mathcal{F}_n \to \mathcal{F}_n$ は
貼り合わさって $\mathcal{F}$ 上の代数構造を与える。
$Y = \underline{\Spec}_X(\mathcal{F})$ とおけば証明が完了する。
\end{proof}

\begin{lemma}
\label{spaces-more-morphisms-lemma-algebraize-morphism}
$A$ をイデアル $I$ に関して完備な Noetherian 環とする。
$S = \Spec(A)$ および $S_n = \Spec(A/I^n)$ と書く。$X$, $Y$ を $S$ 上の
代数空間とする。$n \geq 1$ に対し $X_n = X \times_S S_n$ および
$Y_n = Y \times_S S_n$ とおく。両立する可換図式の系
$$
\xymatrix{
& & X_{n + 1} \ar[rd] \ar[rr]_{g_{n + 1}} & & Y_{n + 1} \ar[ld] \\
X_n \ar[rru] \ar[rd] \ar[rr]_{g_n} & & Y_n \ar[rru] \ar[ld] & S_{n + 1} \\
& S_n \ar[rru]
}
$$
が与えられたとする。次を仮定する：
\begin{enumerate}
\item $X \to S$ は固有である。
% P09-ERR-1204: frozen source line 10255 says "separated of finite type".
\item $Y \to S$ は分離的かつ有限型である。
\end{enumerate}
このとき代数空間の一意な射 $g : X \to Y$（$S$ 上の射）で、$g_n$ が $g$ の
$S_n$ への基底変換となるものが存在する。
\end{lemma}

\begin{proof}
射 $(1, g_n) : X_n \to X_n \times_S Y_n$ は閉埋入である。実際、
$Y_n \to S_n$ は分離的である
（Morphisms of Spaces, Lemma \ref{spaces-morphisms-lemma-section-immersion}）。
したがって Lemma \ref{spaces-more-morphisms-lemma-algebraize-formal-closed-subscheme} により、
閉部分空間 $Z \subset X \times_S Y$ で $S$ 上固有であり、その $S_n$ への基底変換が
$X_n \subset X_n \times_S Y_n$ を復元するものが存在する。第一射影
$p : Z \to X$ は固有射である（$Z$ は $S$ 上固有である。Morphisms of Spaces,
Lemma \ref{spaces-morphisms-lemma-universally-closed-permanence} を参照せよ）。
その $S_n$ への基底変換は、すべての $n$ に対して同型である。特に
$p : Z \to X$ は $Z$ のある開部分空間上準有限であり、この開部分空間は
% P09-ERR-1205: frozen source line 10274 refers to undefined Z_0 although indices begin at 1.
$Z_0$ のすべての点を含む。例えば Morphisms of Spaces, Lemma
\ref{spaces-morphisms-lemma-locally-finite-type-quasi-finite-part} による。
$Z$ は $S$ 上固有なので、この開近傍は $Z$ 全体である。したがって Zariski の
主定理により $p : Z \to X$ は有限であると結論する
（例えば Lemma \ref{spaces-more-morphisms-lemma-quasi-finite-separated-pass-through-finite} を適用し、
$Z$ の $X$ 上の固有性を用いて埋入が閉埋入であることを見よ）。Theorem
\ref{spaces-more-morphisms-theorem-grothendieck-existence} の圏同値を適用すると、
$p_*\mathcal{O}_Z = \mathcal{O}_X$ である。実際、これは $I^n$ を法として
すべての $n$ に対して成り立つ。したがって $p$ は同型であり、射 $g$ を
合成 $X \cong Z \to Y$ として得る。
一意性の証明は省略する。
\end{proof}

\begin{remark}
\label{spaces-more-morphisms-remark-weaken-separation-axioms-question}
Grothendieck の代数化定理（Lemma \ref{spaces-more-morphisms-lemma-algebraize-morphism} の形）において、
標的に課す分離公理を弱めても済むか、という問いを立てることができる。
より正確に述べよう。$A$, $I$, $S$, $S_n$, $X$, $Y$, $X_n$, $Y_n$,
および $g_n$ を Lemma \ref{spaces-more-morphisms-lemma-algebraize-morphism} の主張におけるものとし、
次を仮定する：
\begin{enumerate}
\item $X \to S$ は固有である。
\item $Y \to S$ は局所有限型である。
\end{enumerate}
代数空間の射 $g : X \to Y$ で $S$ 上の射であり、$g_n$ が $g$ の $S_n$ への
基底変換となるものは存在するだろうか？一般の場合の答えは分かっていない。答えをご存じなら
\href{mailto:stacks.project@gmail.com}{stacks.project@gmail.com} までメールされたい。
$Y \to S$ が分離的ならば、この結果は補題から従う（$X$ が $S$ 上有限型である
場合へ直ちに帰着できる。これは $g_1$ の像を含む準コンパクトな開部分空間を
選ぶことによる）。
$Y \to S$ が準分離的であることだけを仮定しても、この結果は成り立つ。
まず、前と同様に $Y$ は準コンパクトかつ準分離的であるとしてよい。
このとき \cite{Bhatt-Algebraize} または \cite{Hall-Rydh-coherent} を用いて
$(g_n)$ を代数化できる。実際、第一の文献を適用するには次を用いる：
% P09-ERR-1206: the frozen source reverses X/Y in both pullback arguments and conclusions, and also changes the base A to R.
$$
D_{perf}(X) \to \lim D_{perf}(X_n)
\xrightarrow{\lim Lg_n^*}
\lim D_{perf}(Y_n) = D_{perf}(Y)
$$
ここで最後の段階では、固有代数空間 $Y$ の $R$ 上での導来圏に対する
Grothendieck の存在定理を用いる
（Flatness on Spaces, Remark
\ref{spaces-flat-remark-correct-generality} と比較せよ）。
引用した論文は、この矢印が所望の射 $Y \to X$ を定めることを示している。
第二の文献を適用するには、連接加群について同じ議論を用いる：
$$
\textit{Coh}(\mathcal{O}_X) \to \lim \textit{Coh}(\mathcal{O}_{X_n})
\xrightarrow{\lim g_n^*}
\lim \textit{Coh}(\mathcal{O}_{Y_n}) = \textit{Coh}(\mathcal{O}_Y)
$$
ここで最後の等号は Grothendieck の存在定理
（Theorem \ref{spaces-more-morphisms-theorem-grothendieck-existence}）の帰結である。
第二の文献によれば、この関手は射 $Y \to X$（$R$ 上の射）に対応する。
この一般化が必要になったときには、ここで結果を正確に述べ、注意深く証明することにする。
\end{remark}

\begin{lemma}
\label{spaces-more-morphisms-lemma-formal-algebraic-space-proper-reldim-1}
$(A, \mathfrak m, \kappa)$ を完備局所 Noetherian 環とする。
$S = \Spec(A)$ および $S_n = \Spec(A/\mathfrak m^n)$ とおく。
次の可換図式を考える：
$$
\xymatrix{
X_1 \ar[r]_{i_1} \ar[d] & X_2 \ar[r]_{i_2} \ar[d] & X_3 \ar[r] \ar[d] &
\ldots \\
S_1 \ar[r] & S_2 \ar[r] & S_3 \ar[r] & \ldots
}
$$
% P09-ERR-1207: the frozen statement supplies no properness hypothesis, but the proof immediately assumes X_1 is proper over kappa.
この図式は代数空間からなり、各正方形は Cartesian であるとする。
$\dim(X_1) \leq 1$ ならば、スキームの射影的射 $X \to S$ と、同型
$X_n \cong X \times_S S_n$ が存在し、これらの同型は $i_n$ と両立する。
\end{lemma}

\begin{proof}
Spaces over Fields, Lemma
\ref{spaces-over-fields-lemma-codim-1-point-in-schematic-locus} により、代数空間
$X_1$ はスキームである。したがって $X_1$ は次元 $\leq 1$ の
$\kappa$ 上の固有スキームである。Varieties, Lemma
\ref{varieties-lemma-dim-1-proper-projective} により、$X_1$ は $\kappa$ 上
H-射影的である。$\mathcal{L}_1$ を $X_1$ 上の豊富な可逆層とする。

\medskip\noindent
$\mathcal{L}_1$ が可逆層の両立する系 $\{\mathcal{L}_n\}$ へ持ち上がり、
それらが $\{X_n\}$ 上にあることを示そう。Lemma
\ref{spaces-more-morphisms-lemma-thickening-scheme} により $X_n$ もスキームであることに注意せよ。
$X_1 \to X_n$ は基礎位相空間の同相を誘導することを思い出す。
以下の証明では、$X_n$ 上の層と $X_1$ 上の層を区別しない。
持ち上げ $\mathcal{L}_n$ が $X_n$ に与えられたとする。完全列
$$
1 \to
(1 + \mathfrak m^n\mathcal{O}_{X_{n + 1}})^* \to
\mathcal{O}_{X_{n + 1}}^* \to \mathcal{O}_{X_n}^* \to 1
$$
を考える。これは $X_{n + 1}$ 上の層の完全列である。
$\mathcal{L}_n$ の $H^1(X_n, \mathcal{O}_{X_n}^*)$ における類
（Cohomology, Lemma \ref{cohomology-lemma-h1-invertible} を参照）が
$H^1(X_{n + 1}, \mathcal{O}_{X_{n + 1}}^*)$ の元へ持ち上がるための必要十分条件は、
次に属する障害類
$$
H^2(X_{n + 1}, (1 + \mathfrak m^n\mathcal{O}_{X_{n + 1}})^*)
$$
が零であることである。
$X_1$ は次元 $\leq 1$ の Noetherian スキームなので、このコホモロジー群は消える
（Cohomology, Proposition
\ref{cohomology-proposition-vanishing-Noetherian}）。

\medskip\noindent
Grothendieck の代数化定理
（Cohomology of Schemes, Theorem \ref{coherent-theorem-algebraization}）により、
スキームの射影的射 $X \to S = \Spec(A)$ と、両立する同型の系
$X_n = S_n \times_S X$ を得る。
\end{proof}

\begin{lemma}
\label{spaces-more-morphisms-lemma-projective-over-complete}
$(A, \mathfrak m, \kappa)$ を完備 Noetherian 局所環とする。
$X$ を $\Spec(A)$ 上の代数空間とする。
$X \to \Spec(A)$ が固有であり $\dim(X_\kappa) \leq 1$ ならば、
$X$ は $A$ 上射影的なスキームである。
\end{lemma}

\begin{proof}
$X_n = X \times_{\Spec(A)} \Spec(A/\mathfrak m^n)$ とおく。
Lemma \ref{spaces-more-morphisms-lemma-formal-algebraic-space-proper-reldim-1} により、射影的射
$Y \to \Spec(A)$ と両立する同型
$$
Y \times_{\Spec(A)} \Spec(A/\mathfrak m^n) \cong
X \times_{\Spec(A)} \Spec(A/\mathfrak m^n)
$$
が存在する。Lemma \ref{spaces-more-morphisms-lemma-algebraize-morphism} により $X \cong Y$ であり、
証明が完了する。
\end{proof}










\section{正則埋入}
\label{spaces-more-morphisms-section-regular-immersions}

\noindent
本節は、代数空間の射について Divisors, Section
\ref{divisors-section-regular-immersions} に対応する節である。読者には、まずその節で
スキームの正則埋入について読んでおくことを勧める。

\medskip\noindent
Divisors, Section \ref{divisors-section-regular-immersions} では、スキームの射に対して
四種類の正則埋入を定義した。このうち、標的上エタール位相について局所的であるのは
（分かっている限り）三種類だけであり、通常の正則埋入は、例によってそうではない。
このため、代数空間の射については正則埋入を実際に定義することができない。
% P09-ERR-1211: frozen lines 10440--10441 omit "the" and mismatch singular/plural in "a Koszul-regular immersions".
（この種の煩わしさが、Grothendieck とその学派に、もとの正則埋入という概念を
Koszul 正則埋入で置き換えるよう促した。\cite[Exposee VII, Definition 1.4]{SGA6}
を参照せよ。）しかし、Koszul 正則、$H_1$-正則、および準正則埋入は定義できる。
もう一つ注意すべきことがある。Koszul 正則埋入は任意の基底変換で保たれないため、
これを定義するのに Morphisms of Spaces, Section
\ref{spaces-morphisms-section-representable} の方策を用いることはできない。
同様に、Koszul 正則埋入は始域上エタール局所的ではないため、Morphisms of Spaces,
Lemma \ref{spaces-morphisms-lemma-local-source-target} を用いて定義することもできない。
そこで、この補題を次の補題で置き換える。

\begin{lemma}
\label{spaces-more-morphisms-lemma-representable-etale-local-target}
$\mathcal{P}$ を、標的上エタール局所的であるスキームの射の性質とする。
$S$ をスキームとする。$f : X \to Y$ を $S$ 上の代数空間の表現可能な射とする。
可換図式
$$
\xymatrix{
X \times_Y V \ar[d] \ar[r] & V \ar[d] \\
X \ar[r]^f & Y
}
$$
で、$V$ がスキームであり $V \to Y$ がエタールであるものを考える。
次の条件は同値である：
\begin{enumerate}
\item 上のような任意の図式について、射影 $X \times_Y V \to V$ は性質
$\mathcal{P}$ をもつ。
\item 上のような図式のうち $V \to Y$ が全射であるものが一つ存在して、射影
$X \times_Y V \to V$ が性質 $\mathcal{P}$ をもつ。
\end{enumerate}
$X$ と $Y$ が表現可能ならば、これらはさらに、$f$ が（スキームの射として）
性質 $\mathcal{P}$ をもつことと同値である。
\end{lemma}

\begin{proof}
(1) と (2) の同値性を証明しよう。(1) $\Rightarrow$ (2) は直ちに従う。
補題における二つの図式
$$
\xymatrix{
X \times_Y V \ar[d] \ar[r] & V \ar[d] \\
X \ar[r]^f & Y
}
\quad\quad
\xymatrix{
X \times_Y V' \ar[d] \ar[r] & V' \ar[d] \\
X \ar[r]^f & Y
}
$$
が与えられたとする。$V \to Y$ は全射であり、$X \times_Y V \to V$ は性質
$\mathcal{P}$ をもつと仮定する。(2) が (1) を含意することを示すには、
$X \times_Y V' \to V'$ が $\mathcal{P}$ をもつことを証明しなければならない。
そのため、次の図式を考える：
$$
\xymatrix{
X \times_Y V \ar[d] &
(X \times_Y V) \times_X (X \times_Y V') \ar[l] \ar[d] \ar[r] &
X \times_Y V' \ar[d] \\
V &
V \times_Y V' \ar[l] \ar[r] &
V'
}
$$
% P09-ERR-1208: frozen line 10507 says local on the source, contradicting the lemma's target-local hypothesis.
$\mathcal{P}$ は始域上エタール局所的であるという仮定により、$\mathcal{P}$ は
エタール基底変換で保たれる。Descent, Lemma
\ref{descent-lemma-pullback-property-local-target} を参照せよ。したがって左の垂直射が
$\mathcal{P}$ をもてば、中央の垂直射もそうである。
% P09-ERR-1209: frozen line 10511 uses undefined U and U' where the diagram has V and V'.
$U \times_X U' \to U'$ は全射かつエタールであり（したがって $U'$ のエタール被覆を
定める）、また $\mathcal{P}$ は標的上エタール位相について局所的であると仮定したので、
% P09-ERR-1210: frozen line 10514 says the left vertical arrow, but descent is meant to prove the right vertical arrow has P.
左の垂直射が $\mathcal{P}$ をもつことが従う。

\medskip\noindent
$X$ と $Y$ が表現可能ならば、エタール被覆として
$\text{id}_Y : Y \to Y$ を取ることができ、補題の最後の主張が成り立つことが分かる。
\end{proof}

\noindent
「Koszul 正則（resp.\ $H_1$-正則、resp.\ 準正則）埋入であること」は、標的上
fpqc 局所的なスキームの射の性質である。Descent, Lemma
\ref{descent-lemma-descending-property-regular-immersion} を参照せよ。
したがって、次の定義は意味をもつ。

\begin{definition}
\label{spaces-more-morphisms-definition-regular-immersion}
$S$ をスキームとする。$i : X \to Y$ を $S$ 上の代数空間の射とする。
\begin{enumerate}
\item $i$ が {it Koszul 正則埋入}であるとは、$i$ が表現可能であり、
Lemma \ref{spaces-more-morphisms-lemma-representable-etale-local-target} の同値な条件が
$\mathcal{P}(f) =$「$f$ は Koszul 正則埋入である」という性質について成り立つことをいう。
\item $i$ が {it $H_1$-正則埋入}であるとは、$i$ が表現可能であり、
Lemma \ref{spaces-more-morphisms-lemma-representable-etale-local-target} の同値な条件が
$\mathcal{P}(f) =$「$f$ は $H_1$-正則埋入である」という性質について成り立つことをいう。
\item $i$ が {it 準正則埋入}であるとは、$i$ が表現可能であり、
Lemma \ref{spaces-more-morphisms-lemma-representable-etale-local-target} の同値な条件が
$\mathcal{P}(f) =$「$f$ は準正則埋入である」という性質について成り立つことをいう。
\end{enumerate}
\end{definition}

\begin{lemma}
\label{spaces-more-morphisms-lemma-regular-quasi-regular-immersion}
$S$ をスキームとする。$i : Z \to X$ を $S$ 上の代数空間の埋入とする。
次の含意が成り立つ：
$i$ が Koszul 正則 $\Rightarrow$
$i$ が $H_1$-正則 $\Rightarrow$
$i$ が準正則。
\end{lemma}

\begin{proof}
定義により、この補題は直ちに Divisors, Lemma
\ref{divisors-lemma-regular-quasi-regular-immersion} に帰着する。
\end{proof}

\begin{lemma}
\label{spaces-more-morphisms-lemma-regular-immersion-noetherian}
$S$ をスキームとする。$i : Z \to X$ を $S$ 上の代数空間の埋入とする。
$X$ は局所 Noetherian であると仮定する。このとき
$i$ が Koszul 正則 $\Leftrightarrow$
$i$ が $H_1$-正則 $\Leftrightarrow$
$i$ が準正則。
\end{lemma}

\begin{proof}
Definition \ref{spaces-more-morphisms-definition-regular-immersion}（および Properties of Spaces,
Section \ref{spaces-properties-section-types-properties} における局所 Noetherian
代数空間の定義）により、これは直ちにスキームの場合、すなわち Divisors, Lemma
\ref{divisors-lemma-regular-immersion-noetherian} に移される。
\end{proof}

\begin{lemma}
\label{spaces-more-morphisms-lemma-flat-base-change-regular-immersion}
\begin{slogan}
正則埋入は平坦基底変換で安定である。
\end{slogan}
$S$ をスキームとする。$i : Z \to X$ を Koszul 正則、$H_1$-正則、または
準正則埋入とし、これは $S$ 上の代数空間の射であるとする。$X' \to X$ を $S$ 上の代数空間の
平坦射とする。このとき基底変換 $i' : Z \times_X X' \to X'$ は Koszul 正則、
$H_1$-正則、または準正則埋入である。
\end{lemma}

\begin{proof}
Definition \ref{spaces-more-morphisms-definition-regular-immersion}（および Morphisms of Spaces,
Section \ref{spaces-morphisms-section-flat} における代数空間の平坦射の定義）により、
この補題はスキームの場合に帰着する。Divisors, Lemma
\ref{divisors-lemma-flat-base-change-regular-immersion} を参照せよ。
\end{proof}

\begin{lemma}
\label{spaces-more-morphisms-lemma-quasi-regular-immersion}
$S$ をスキームとする。$i : Z \to X$ を $S$ 上の代数空間の埋入とする。
$i$ が準正則埋入であるための必要十分条件は、次の条件が満たされることである：
\begin{enumerate}
\item $i$ は局所有限表示である。
\item 余法層 $\mathcal{C}_{Z/X}$ は有限局所自由である。
\item 写像 (\ref{spaces-more-morphisms-equation-conormal-algebra-quotient}) は同型である。
\end{enumerate}
\end{lemma}

\begin{proof}
エタール局所化によりスキームの場合
（Divisors, Lemma \ref{divisors-lemma-quasi-regular-immersion}）から従う
（Definition \ref{spaces-more-morphisms-definition-regular-immersion} および Lemma
\ref{spaces-more-morphisms-lemma-etale-conormal-algebra} を用いよ）。
\end{proof}

\begin{lemma}
\label{spaces-more-morphisms-lemma-transitivity-conormal-quasi-regular}
$S$ をスキームとする。$Z \to Y \to X$ を $S$ 上の代数空間の埋入とする。
$Z \to Y$ が $H_1$-正則であると仮定する。このとき Lemma
\ref{spaces-more-morphisms-lemma-transitivity-conormal} の標準列
$$
0 \to i^*\mathcal{C}_{Y/X} \to
\mathcal{C}_{Z/X} \to
\mathcal{C}_{Z/Y} \to 0
$$
は完全であり、（エタール）局所的に分裂する。
\end{lemma}

\begin{proof}
$\mathcal{C}_{Z/Y}$ は有限局所自由なので（Lemma
\ref{spaces-more-morphisms-lemma-quasi-regular-immersion} および Lemma
\ref{spaces-more-morphisms-lemma-regular-quasi-regular-immersion} を参照）、列が完全であることを証明すれば
十分である。一般にこの列はすでに右完全なので、最初の写像が単射であることを
示せば十分である。$X$ 上でエタール局所化した後、これはスキームの場合に帰着する。
Divisors, Lemma
\ref{divisors-lemma-transitivity-conormal-quasi-regular} を参照せよ。
\end{proof}

\noindent
準正則埋入の合成は準正則とは限らない。Algebra, Remark
\ref{algebra-remark-join-quasi-regular-sequences} を参照せよ。
他の種類の正則埋入は合成で保たれる。

\begin{lemma}
\label{spaces-more-morphisms-lemma-composition-regular-immersion}
$S$ をスキームとする。$i : Z \to Y$ および $j : Y \to X$ を $S$ 上の
代数空間の埋入とする。
\begin{enumerate}
\item $i$ と $j$ が Koszul 正則埋入ならば、$j \circ i$ もそうである。
\item $i$ と $j$ が $H_1$-正則埋入ならば、$j \circ i$ もそうである。
\item $i$ が $H_1$-正則埋入で $j$ が準正則埋入ならば、$j \circ i$ は
準正則埋入である。
\end{enumerate}
\end{lemma}

\begin{proof}
スキームの場合から直ちに従う。Divisors, Lemma
\ref{divisors-lemma-composition-regular-immersion} を参照せよ。
\end{proof}

\begin{lemma}
\label{spaces-more-morphisms-lemma-permanence-regular-immersion}
$S$ をスキームとする。$i : Z \to Y$ および $j : Y \to X$ を $S$ 上の
代数空間の埋入とする。$j$ は局所有限表示であり、Lemma
\ref{spaces-more-morphisms-lemma-transitivity-conormal} の列
$$
0 \to i^*\mathcal{C}_{Y/X} \to
\mathcal{C}_{Z/X} \to
\mathcal{C}_{Z/Y} \to 0
$$
が完全かつ局所的に分裂すると仮定する。
\begin{enumerate}
\item $j \circ i$ が準正則埋入ならば、$i$ も準正則埋入である。
\item $j \circ i$ が $H_1$-正則埋入ならば、$i$ もそうである。
\item $j$ と $j \circ i$ がともに Koszul 正則埋入ならば、$i$ も
Koszul 正則埋入である。
\end{enumerate}
\end{lemma}

\begin{proof}
スキームの場合から直ちに従う。Divisors, Lemma
\ref{divisors-lemma-permanence-regular-immersion} を参照せよ。
\end{proof}

\begin{lemma}
\label{spaces-more-morphisms-lemma-extra-permanence-regular-immersion-noetherian}
$S$ をスキームとする。$i : Z \to Y$ および $j : Y \to X$ を $S$ 上の
代数空間の埋入とする。$X$ は局所 Noetherian であると仮定する。
次の条件は同値である：
\begin{enumerate}
\item $i$ と $j$ は Koszul 正則埋入である。
\item $i$ と $j \circ i$ は Koszul 正則埋入である。
\item $j \circ i$ は Koszul 正則埋入であり、余法列
$$
0 \to i^*\mathcal{C}_{Y/X} \to
\mathcal{C}_{Z/X} \to
\mathcal{C}_{Z/Y} \to 0
$$
は完全かつ局所的に分裂する。
\end{enumerate}
\end{lemma}

\begin{proof}
スキームの場合から直ちに従う。Divisors, Lemma
\ref{divisors-lemma-extra-permanence-regular-immersion-noetherian} を参照せよ。
\end{proof}














\section{相対的擬連接性}
\label{spaces-more-morphisms-section-relative-pseudo-coherence}

\noindent
本節は More on Morphisms, Section
\ref{more-morphisms-section-relative-pseudo-coherence} に対応する節である。
ただし、代数空間についてこの題材を扱うにあたり、コホモロジー層が準連接である
導来圏の対象だけを用いることにした。理由は二つある：(1) これにより叙述が大幅に
簡潔になること、(2) 現時点では、より一般の概念を使う場面がないことである。

\begin{remark}
\label{spaces-more-morphisms-remark-match-relative-pseudo-coherence}
$S$ をスキームとする。$f : X \to Y$ を、局所有限型な $S$ 上の表現可能な
代数空間の射とする。$f_0 : X_0 \to Y_0$ を $f$ を表現するスキームの射とする
（不格好だが一時的な記法である）。このとき $f_0$ は局所有限型である。
$E$ が $D_\QCoh(\mathcal{O}_X)$ の対象ならば、$E$ は一意な対象
$E_0$（$D_\QCoh(\mathcal{O}_{X_0})$ のもの）の引き戻しである。Derived
Categories of Spaces, Lemma
\ref{spaces-perfect-lemma-derived-quasi-coherent-small-etale-site} を参照せよ。
この状況では、「$E$ は $m$-擬連接（$Y$ に関して）である」という句を、More on
Morphisms, Section \ref{more-morphisms-section-relative-pseudo-coherence} で定義された
意味で「$E_0$ は $m$-擬連接（$Y_0$ に関して）である」という意味に取る。
\end{remark}

\begin{lemma}
\label{spaces-more-morphisms-lemma-qcoh-relative-pseudo-coherence-characterize}
$S$ をスキームとする。$f : X \to Y$ を、局所有限型な $S$ 上の代数空間の
射とする。$m \in \mathbf{Z}$ とする。$E \in D_\QCoh(\mathcal{O}_X)$ とする。
Remark \ref{spaces-more-morphisms-remark-match-relative-pseudo-coherence} で説明した記法のもとで、
次の条件は同値である：
\begin{enumerate}
\item 任意の可換図式
$$
\xymatrix{
U \ar[d] \ar[r] & V \ar[d] \\
X \ar[r] & Y
}
$$
で、$U$, $V$ がスキームであり垂直射がエタールであるものについて、複体
$E|_U$ は $m$-擬連接（$V$ に関して）である。
\item (1) のような可換図式のうち $U \to X$ が全射であるものが一つ存在して、
複体 $E|_U$ が $m$-擬連接（$V$ に関して）である。
\item (1) のような任意の可換図式で $U$ と $V$ がアフィンであるものについて、
$R\Gamma(U, E)$ は $\mathcal{O}_X(U)$-加群の複体であり、
$m$-擬連接（$\mathcal{O}_Y(V)$ に関して）である。
\end{enumerate}
\end{lemma}

\begin{proof}
(1) が (3) を含意することは More on Morphisms, Lemma
\ref{more-morphisms-lemma-qcoh-relative-pseudo-coherence-characterize} から従う。

\medskip\noindent
(3) を仮定する。$U \to X$ が全射である (1) のような任意の可換図式を取る。
アフィン開被覆 $V = \bigcup V_j$ と、アフィン開被覆
$(U \to V)^{-1}(V_j) = \bigcup U_{ij}$ を選ぶ。(3) と More on Morphisms,
Lemma \ref{more-morphisms-lemma-qcoh-relative-pseudo-coherence-characterize} により、
$E|_U$ は $m$-擬連接（$V$ に関して）である。したがって (3) は (2) を含意する。

\medskip\noindent
(2) を仮定する。可換図式
$$
\xymatrix{
U \ar[d] \ar[r] & V \ar[d] \\
X \ar[r] & Y
}
$$
で、$U$, $V$ がスキーム、垂直射がエタール、射 $U \to X$ が全射であり、
$E|_U$ が $m$-擬連接（$V$ に関して）であるものを選ぶ。次に、エタールな
垂直射をもつ第二の可換図式
$$
\xymatrix{
U' \ar[d] \ar[r] & V' \ar[d] \\
X \ar[r] & Y
}
$$
で、$U', V'$ がスキームであるものが与えられたとする。
$E|_{U'}$ が $m$-擬連接（$V'$ に関して）であることを示したい。
射 $U'' = U \times_X U' \to U'$ は全射かつエタールであり、$U'' \to V'$ は
$V'' = V' \times_Y V$ を経由する。後者は $V'$ 上エタールである。
したがって $E|_{U''}$ が $m$-擬連接（$V''$ に関して）であることを示せば十分である。
More on Morphisms, Lemmas
\ref{more-morphisms-lemma-relative-pseudo-coherent-descends-fppf} および
\ref{more-morphisms-lemma-relative-pseudo-coherent-post-compose} を参照せよ。
後者の補題をもう一度用いると、$E|_{U''}$ が $m$-擬連接（$V$ に関して）であることを
示せば十分である。これは More on Morphisms, Lemma
\ref{more-morphisms-lemma-pull-relative-pseudo-coherent} と、スキームのエタール射が
擬連接であるという事実から従う。More on Morphisms, Lemma
\ref{more-morphisms-lemma-flat-finite-presentation-pseudo-coherent} を参照せよ。
\end{proof}

\begin{definition}
\label{spaces-more-morphisms-definition-relative-pseudo-coherence}
$S$ をスキームとする。$f : X \to Y$ を、局所有限型な $S$ 上の代数空間の
射とする。$E$ を $D_\QCoh(\mathcal{O}_X)$ の対象とする。$\mathcal{F}$ を
準連接 $\mathcal{O}_X$-加群とする。$m \in \mathbf{Z}$ を固定する。
\begin{enumerate}
\item $E$ が {it $m$-擬連接（$Y$ に関して）}であるとは、Lemma
\ref{spaces-more-morphisms-lemma-qcoh-relative-pseudo-coherence-characterize} の同値な条件が満たされることをいう。
\item $E$ が {it 擬連接（$Y$ に関して）}であるとは、$E$ が
$m$-擬連接（$Y$ に関して）であることが、すべての $m \in \mathbf{Z}$ について
成り立つことをいう。
\item $\mathcal{F}$ が {it $m$-擬連接（$Y$ に関して）}であるとは、
$\mathcal{F}$ を $D_\QCoh(\mathcal{O}_X)$ の対象とみなしたものが
$m$-擬連接（$Y$ に関して）であることをいう。
\item $\mathcal{F}$ が {it 擬連接（$Y$ に関して）}であるとは、
$\mathcal{F}$ を $D_\QCoh(\mathcal{O}_X)$ の対象とみなしたものが
擬連接（$Y$ に関して）であることをいう。
\end{enumerate}
\end{definition}

\noindent
基底に関して擬連接な複体の性質の大半は、スキームの場合の対応する性質から
直ちに従う。必要に応じて、関連する補題をここに加えることにする。

\begin{lemma}
\label{spaces-more-morphisms-lemma-relative-pseudo-coherent-is-moot}
$S$ をスキームとする。$f : X \to Y$ を $S$ 上の代数空間の射とする。
% P09-ERR-1212: frozen source line 10862 omits "be" in "Let E in".
$E$ を $D_\QCoh(\mathcal{O}_X)$ の対象とする。$f$ が平坦かつ局所有限表示ならば、
次の条件は同値である：
\begin{enumerate}
\item $E$ は擬連接（$Y$ に関して）である。
\item $E$ は $X$ 上擬連接である。
\end{enumerate}
\end{lemma}

\begin{proof}
エタール局所化と定義により、$X$ と $Y$ はスキームであるとしてよい。
スキームの場合には More on Morphisms, Lemma
\ref{more-morphisms-lemma-check-relative-pseudo-coherence-on-charts} から従う。
\end{proof}


















\section{擬連接射}
\label{spaces-more-morphisms-section-pseudo-coherent}

\noindent
本節は、スキームの射について More on Morphisms, Section
\ref{more-morphisms-section-pseudo-coherent} に対応する節である。読者には、まずその節で
スキームの擬連接射について読んでおくことを勧める。

\medskip\noindent
スキームの射の「擬連接」という性質は、始域と標的上エタール局所的である。
これを見るには More on Morphisms, Lemmas
\ref{more-morphisms-lemma-descending-property-pseudo-coherent} および
\ref{more-morphisms-lemma-pseudo-coherent-fppf-local-source}、ならびに Descent,
Lemma \ref{descent-lemma-etale-local-source-target} を用いよ。Morphisms of Spaces,
Lemma \ref{spaces-morphisms-lemma-local-source-target} により、代数空間の擬連接射という
概念を次のように定義できる。また、問題の代数空間が表現可能ならば、この定義は
More on Morphisms, Section \ref{more-morphisms-section-pseudo-coherent} ですでに定義された
概念と一致する。

\begin{definition}
\label{spaces-more-morphisms-definition-pseudo-coherent}
$S$ をスキームとする。$f : X \to Y$ を $S$ 上の代数空間の射とする。
\begin{enumerate}
\item $f$ が {it 擬連接}であるとは、Morphisms of Spaces, Lemma
\ref{spaces-morphisms-lemma-local-source-target} の同値な条件が
$\mathcal{P} =$「擬連接」という性質について成り立つことをいう。
\item $x \in |X|$ とする。$f$ が {it $x$ において擬連接}であるとは、
$X' \subset X$ が $x$ の開近傍で、$f|_{X'} : X' \to Y$ が擬連接となるものが
存在することをいう。
\end{enumerate}
\end{definition}

\noindent
擬連接射の基底変換は、一般には擬連接とは限らないことに注意せよ。

\begin{lemma}
\label{spaces-more-morphisms-lemma-flat-base-change-pseudo-coherent}
擬連接射の平坦基底変換は擬連接である。
\end{lemma}

\begin{proof}
省略する。ヒント：この補題のスキーム版、すなわち More on Morphisms, Lemma
\ref{more-morphisms-lemma-flat-base-change-pseudo-coherent} を用いよ。
\end{proof}

\begin{lemma}
\label{spaces-more-morphisms-lemma-composition-pseudo-coherent}
擬連接射の合成は擬連接である。
\end{lemma}

\begin{proof}
省略する。ヒント：この補題のスキーム版、すなわち More on Morphisms, Lemma
\ref{more-morphisms-lemma-composition-pseudo-coherent} を用いよ。
\end{proof}

\begin{lemma}
\label{spaces-more-morphisms-lemma-pseudo-coherent-finite-presentation}
擬連接射は局所有限表示である。
\end{lemma}

\begin{proof}
定義から直ちに従う。
\end{proof}

\begin{lemma}
\label{spaces-more-morphisms-lemma-flat-finite-presentation-pseudo-coherent}
局所有限表示な平坦射は擬連接である。
\end{lemma}

\begin{proof}
省略する。ヒント：この補題のスキーム版、すなわち More on Morphisms, Lemma
\ref{more-morphisms-lemma-flat-finite-presentation-pseudo-coherent} を用いよ。
\end{proof}

\begin{lemma}
\label{spaces-more-morphisms-lemma-permanence-pseudo-coherent}
$f : X \to Y$ を、基底代数空間 $B$ 上擬連接な代数空間の射とする。
このとき $f$ は擬連接である。
\end{lemma}

\begin{proof}
省略する。ヒント：この補題のスキーム版、すなわち More on Morphisms, Lemma
\ref{more-morphisms-lemma-permanence-pseudo-coherent} を用いよ。
\end{proof}

\begin{lemma}
\label{spaces-more-morphisms-lemma-Noetherian-pseudo-coherent}
$S$ をスキームとする。$f : X \to Y$ を $S$ 上の代数空間の射とする。
$Y$ が局所 Noetherian ならば、$f$ が擬連接であることと $f$ が局所有限型であることは
同値である。
\end{lemma}

\begin{proof}
省略する。ヒント：この補題のスキーム版、すなわち More on Morphisms, Lemma
\ref{more-morphisms-lemma-Noetherian-pseudo-coherent} を用いよ。
\end{proof}








\section{完全射}
\label{spaces-more-morphisms-section-perfect}

\noindent
本節は、スキームの射について More on Morphisms, Section
\ref{more-morphisms-section-perfect} に対応する節である。読者には、まずその節で
スキームの完全射について読んでおくことを勧める。

\medskip\noindent
スキームの射の「完全」という性質は、始域と標的上エタール局所的である。
これを見るには More on Morphisms, Lemmas
\ref{more-morphisms-lemma-descending-property-perfect} および
\ref{more-morphisms-lemma-perfect-fppf-local-source}、ならびに Descent, Lemma
\ref{descent-lemma-etale-local-source-target} を用いよ。Morphisms of Spaces, Lemma
\ref{spaces-morphisms-lemma-local-source-target} により、代数空間の完全射という概念を
次のように定義できる。また、問題の代数空間が表現可能ならば、この定義は More on
Morphisms, Section \ref{more-morphisms-section-perfect} ですでに定義された概念と一致する。

\begin{definition}
\label{spaces-more-morphisms-definition-perfect}
$S$ をスキームとする。$f : X \to Y$ を $S$ 上の代数空間の射とする。
\begin{enumerate}
\item $f$ が {it 完全}であるとは、Morphisms of Spaces, Lemma
\ref{spaces-morphisms-lemma-local-source-target} の同値な条件が
$\mathcal{P} =$「完全」という性質について成り立つことをいう。
\item $x \in |X|$ とする。$f$ が {it $x$ において完全}であるとは、
$X' \subset X$ が $x$ の開近傍で、$f|_{X'} : X' \to Y$ が完全となるものが
存在することをいう。
\end{enumerate}
\end{definition}

\noindent
完全射は擬連接であり、したがって局所有限表示であることに注意せよ。
完全射の基底変換は、一般には完全とは限らない。

\begin{lemma}
\label{spaces-more-morphisms-lemma-flat-base-change-perfect}
完全射の平坦基底変換は完全である。
\end{lemma}

\begin{proof}
省略する。ヒント：この補題のスキーム版、すなわち More on Morphisms, Lemma
\ref{more-morphisms-lemma-flat-base-change-perfect} を用いよ。
\end{proof}

\begin{lemma}
\label{spaces-more-morphisms-lemma-composition-perfect}
完全射の合成は完全である。
\end{lemma}

\begin{proof}
省略する。ヒント：この補題のスキーム版、すなわち More on Morphisms, Lemma
\ref{more-morphisms-lemma-composition-perfect} を用いよ。
\end{proof}

\begin{lemma}
\label{spaces-more-morphisms-lemma-flat-finite-presentation-perfect}
$S$ をスキームとする。$f : X \to Y$ を $S$ 上の代数空間の射とする。
次の条件は同値である：
\begin{enumerate}
\item $f$ は平坦かつ完全である。
\item $f$ は平坦かつ局所有限表示である。
\end{enumerate}
\end{lemma}

\begin{proof}
省略する。ヒント：この補題のスキーム版、すなわち More on Morphisms, Lemma
\ref{more-morphisms-lemma-flat-finite-presentation-perfect} を用いよ。
\end{proof}

\begin{lemma}
\label{spaces-more-morphisms-lemma-perfect-proper-perfect-direct-image}
$S$ をスキームとする。$Y$ を $S$ 上の Noetherian 代数空間とする。
$f : X \to Y$ を代数空間の完全かつ固有な射とする。
$E \in D(\mathcal{O}_X)$ を完全とする。このとき $Rf_*E$ は
$D(\mathcal{O}_Y)$ の完全対象である。
\end{lemma}

\begin{proof}
Derived Categories of Spaces, Lemma
\ref{spaces-perfect-lemma-perfect-direct-image} が適用できると主張する。
条件 (1) と (2) は直ちに成り立つ。条件 (3) は $X$ 上局所的である。
したがって $X$ と $Y$ はアフィンであり、$E$ は $\mathcal{O}_X$-加群の
厳密完全複体で表されるとしてよい。ゆえに、エタールサイト上で
$\mathcal{O}_X$ が $f^{-1}\mathcal{O}_Y$-加群の層として有限 Tor 次元をもつことを
示せば十分である。Derived Categories of Spaces, Lemma
\ref{spaces-perfect-lemma-tor-dimension-rel} により、これを Zariski サイト上で
確かめれば十分である。スキームの有限型射について、これは完全であることと同値である。
More on Morphisms, Lemma \ref{more-morphisms-lemma-check-perfect-stalks} を参照せよ。
\end{proof}












\section{局所完全交叉射}
\label{spaces-more-morphisms-section-lci}

\noindent
本節は、スキームの射について More on Morphisms, Section
\ref{more-morphisms-section-lci} に対応する節である。読者には、まずその節で
スキームの局所完全交叉射について読んでおくことを勧める。

\medskip\noindent
スキームの射の「局所完全交叉射である」という性質は、始域と標的上エタール
局所的である。これを見るには More on Morphisms, Lemmas
\ref{more-morphisms-lemma-descending-property-lci} および
\ref{more-morphisms-lemma-lci-syntomic-local-source}、ならびに Descent, Lemma
\ref{descent-lemma-etale-local-source-target} を用いよ。Morphisms of Spaces, Lemma
\ref{spaces-morphisms-lemma-local-source-target} により、代数空間の局所完全交叉射という
概念を次のように定義できる。また、問題の代数空間が表現可能ならば、この定義は
More on Morphisms, Section \ref{more-morphisms-section-lci} ですでに定義された概念と
一致する。

\begin{definition}
\label{spaces-more-morphisms-definition-lci}
$S$ をスキームとする。$f : X \to Y$ を $S$ 上の代数空間の射とする。
\begin{enumerate}
\item $f$ が {it Koszul 射}である、または $f$ が {it 局所完全交叉射}であるとは、
Morphisms of Spaces, Lemma \ref{spaces-morphisms-lemma-local-source-target} の同値な条件が
$\mathcal{P}(f) =$「$f$ は局所完全交叉射である」という性質について成り立つことをいう。
\item $x \in |X|$ とする。$f$ が {it $x$ において Koszul} であるとは、
$X' \subset X$ が $x$ の開近傍で、$f|_{X'} : X' \to Y$ が局所完全交叉射となるものが
存在することをいう。
\end{enumerate}
\end{definition}

\noindent
ある意味で、局所完全交叉射の決定的な性質は次の補題の結論である。

\begin{lemma}
\label{spaces-more-morphisms-lemma-lci}
$S$ をスキームとする。$f : X \to Y$ を $S$ 上の代数空間の局所完全交叉射とする。
$P$ を $Y$ 上滑らかな代数空間とする。$U \to X$ を代数空間のエタール射とし、
% P09-ERR-1213: frozen source line 11178 omits "be" in "let i ... an immersion".
$i : U \to P$ を $Y$ 上の代数空間の埋入とする。図式で表せば、
$$
\xymatrix{
X \ar[rd] & U \ar[l] \ar[d] \ar[r]_i & P \ar[ld] \\
& Y
}
$$
である。このとき $i$ は代数空間の Koszul 正則埋入である。
\end{lemma}

\begin{proof}
スキーム $V$ と全射エタール射 $V \to Y$ を選ぶ。スキーム $W$ と全射エタール射
$W \to P \times_Y V$ を選ぶ。$U' = U \times_P W$ とおく。これは $U$ 上
エタールなスキームである。Definition \ref{spaces-more-morphisms-definition-regular-immersion} により、
$U' \to W$ がスキームの Koszul 正則埋入であることを示さなければならない。
上の Definition \ref{spaces-more-morphisms-definition-lci} により、スキームの射 $U' \to V$ は
局所完全交叉射である。したがって結果は More on Morphisms, Lemma
\ref{more-morphisms-lemma-lci} から従う。
\end{proof}

\noindent
すべての Koszul 射に共通する性質を、ここにまとめておくのがよいだろう。

\begin{lemma}
\label{spaces-more-morphisms-lemma-lci-properties}
$S$ をスキームとする。$f : X \to Y$ を $S$ 上の代数空間の局所完全交叉射とする。
このとき、
\begin{enumerate}
\item $f$ は局所有限表示である。
\item $f$ は擬連接である。
\item $f$ は完全である。
\end{enumerate}
\end{lemma}

\begin{proof}
省略する。ヒント：この補題のスキーム版、すなわち More on Morphisms, Lemma
\ref{more-morphisms-lemma-lci-properties} を用いよ。
\end{proof}

\noindent
Koszul 射の基底変換は、一般には Koszul とは限らないことに注意せよ。

\begin{lemma}
\label{spaces-more-morphisms-lemma-flat-base-change-lci}
局所完全交叉射の平坦基底変換は局所完全交叉射である。
\end{lemma}

\begin{proof}
省略する。ヒント：この補題のスキーム版、すなわち More on Morphisms, Lemma
\ref{more-morphisms-lemma-flat-base-change-lci} を用いよ。
\end{proof}

\begin{lemma}
\label{spaces-more-morphisms-lemma-composition-lci}
局所完全交叉射の合成は局所完全交叉射である。
\end{lemma}

\begin{proof}
省略する。ヒント：この補題のスキーム版、すなわち More on Morphisms, Lemma
\ref{more-morphisms-lemma-composition-lci} を用いよ。
\end{proof}

\begin{lemma}
\label{spaces-more-morphisms-lemma-flat-lci}
\begin{slogan}
（代数空間について）シントミックは、平坦かつ lci であることと同値である。
\end{slogan}
$S$ をスキームとする。$f : X \to Y$ を $S$ 上の代数空間の射とする。
次の条件は同値である：
\begin{enumerate}
\item $f$ は平坦かつ局所完全交叉射である。
\item $f$ はシントミックである。
\end{enumerate}
\end{lemma}

\begin{proof}
省略する。ヒント：この補題のスキーム版、すなわち More on Morphisms, Lemma
\ref{more-morphisms-lemma-flat-lci} を用いよ。
\end{proof}

\begin{lemma}
\label{spaces-more-morphisms-lemma-regular-immersion-lci}
$S$ をスキームとする。$S$ 上の代数空間の Koszul 正則埋入は局所完全交叉射である。
\end{lemma}

\begin{proof}
$i : X \to Y$ を $S$ 上の代数空間の Koszul 正則埋入とする。
定義により、全射エタール射 $V \to Y$ で $V$ がスキームであり、
$X \times_Y V$ がスキームで、基底変換 $X \times_Y V \to V$ がスキームの
Koszul 正則埋入となるものが存在する。More on Morphisms, Lemma
\ref{more-morphisms-lemma-regular-immersion-lci} により、$X \times_Y V \to V$ は
局所完全交叉射である。Definition \ref{spaces-more-morphisms-definition-lci} から、$i$ は代数空間の
局所完全交叉射であると結論する。
\end{proof}

\begin{lemma}
\label{spaces-more-morphisms-lemma-lci-permanence}
$S$ をスキームとする。可換図式
$$
\xymatrix{
X \ar[rr]_f \ar[rd] & & Y \ar[ld] \\
& Z
}
$$
を考える。これは $S$ 上の代数空間の射からなるとする。$Y \to Z$ は滑らかであり、
$X \to Z$ は局所完全交叉射であると仮定する。このとき $f : X \to Y$ は
局所完全交叉射である。
\end{lemma}

\begin{proof}
スキーム $W$ と全射エタール射 $W \to Z$ を選ぶ。スキーム $V$ と全射エタール射
$V \to W \times_Z Y$ を選ぶ。スキーム $U$ と全射エタール射
$U \to V \times_Y X$ を選ぶ。このとき $U \to W$ はスキームの局所完全交叉射であり、
$V \to W$ はスキームの滑らかな射である。スキームに対する結果
（More on Morphisms, Lemma \ref{more-morphisms-lemma-lci-permanence}）により、
$U \to V$ は局所完全交叉射であると結論する。定義により、これは $f$ が
局所完全交叉射であることを意味する。
\end{proof}

\begin{lemma}
\label{spaces-more-morphisms-lemma-descending-property-lci}
性質 $\mathcal{P}(f) =$「$f$ は局所完全交叉射である」は、基底上 fpqc 局所的である。
\end{lemma}

\begin{proof}
$S$ をスキームとする。$f : X \to Y$ を $S$ 上の代数空間の射とする。
$\{Y_i \to Y\}$ を fpqc 被覆とする（Topologies on Spaces, Definition
\ref{spaces-topologies-definition-fpqc-covering}）。$f_i : X_i \to Y_i$ を、
$f$ の $Y_i \to Y$ による基底変換とする。$f$ が局所完全交叉射ならば、Lemma
\ref{spaces-more-morphisms-lemma-flat-base-change-lci} により各 $f_i$ は局所完全交叉射である。

\medskip\noindent
逆に、各 $f_i$ が局所完全交叉射であると仮定する。被覆を細分で置き換えてよい
（これも、平坦基底変換が局所完全交叉射であるという性質を保つためである）。
したがって、$Y_i$ は各 $i$ についてスキームであるとしてよい。Topologies on Spaces,
Lemma \ref{spaces-topologies-lemma-refine-fpqc-schemes} を参照せよ。スキーム $V$ と
全射エタール射 $V \to Y$ を選ぶ。スキーム $U$ と全射エタール射
$U \to V \times_Y X$ を選ぶ。$U \to V$ がスキームの局所完全交叉射であることを
示さなければならない。Topologies on Spaces, Lemma
\ref{spaces-topologies-lemma-recognize-fpqc-covering} により、
$\{Y_i \times_Y V \to V\}$ はスキームの fpqc 被覆である。
スキームの場合（More on Morphisms, Lemma
\ref{more-morphisms-lemma-descending-property-lci}）により、
% P09-ERR-1214: frozen display line 11350 targets V although this base change should target V times_Y Y_i, as the following composition does.
$$
U \times_Y Y_i = U \times_V (V \times_Y Y_i) \longrightarrow V
$$
が局所完全交叉射であることを示せば十分である。これは $U \to V$ の
$V \times_Y Y_i \to V$ による基底変換である。この射は合成
$$
U \times_Y Y_i \longrightarrow
(V \times_Y X) \times_Y Y_i = V \times_Y X_i \longrightarrow
V \times_Y Y_i
$$
と書ける。第一の矢印はスキームのエタール射（$U \to V \times_Y X$ の基底変換）であり、
第二の矢印は $f_i$ の平坦基底変換としてスキームの局所完全交叉射である。
局所完全交叉射であることは始域上 syntomic 局所的であり、エタール射は syntomic
なので、結果が従う（More on Morphisms, Lemma
\ref{more-morphisms-lemma-lci-syntomic-local-source} および Morphisms, Lemma
\ref{morphisms-lemma-etale-syntomic}）。
\end{proof}

\begin{lemma}
\label{spaces-more-morphisms-lemma-lci-syntomic-local-source}
性質 $\mathcal{P}(f) =$「$f$ は局所完全交叉射である」は、始域上 syntomic 局所的である。
\end{lemma}

\begin{proof}
これは Descent on Spaces, Lemma
\ref{spaces-descent-lemma-transfer-from-schemes} および More on Morphisms, Lemma
\ref{more-morphisms-lemma-lci-syntomic-local-source} から従う。
\end{proof}

\begin{lemma}
\label{spaces-more-morphisms-lemma-base-change-lci-fibres}
$S$ をスキームとする。可換図式
$$
\xymatrix{
X \ar[rr]_f \ar[rd]_p & & Y \ar[ld]^q \\
& Z
}
$$
を考える。これは $S$ 上の代数空間からなるとする。$p$ と $q$ はともに平坦かつ
局所有限表示であると仮定する。このとき開部分空間 $U(f) \subset X$ で、
$|U(f)| \subset |X|$ が $f$ が Koszul である点の集合となるものが存在する。
さらに、代数空間の任意の射 $Z' \to Z$ について、$f' : X' \to Y'$ が
$f$ の $Z' \to Z$ による基底変換ならば、$U(f')$ は $U(f)$ の逆像である。
ここで逆像は射影 $X' \to X$ によるものとする。
\end{lemma}

\begin{proof}
この補題は More on Morphisms, Lemma
\ref{more-morphisms-lemma-base-change-lci-fibres} に対応するものであり、実際にそこから
導く。Definition \ref{spaces-more-morphisms-definition-lci} により、集合
$\{x \in |X| : f \text{ is Koszul at }x\}$ は $|X|$ で開である。したがって
Properties of Spaces, Lemma \ref{spaces-properties-lemma-open-subspaces} により、これは
開部分空間 $U(f)$（$X$ のもの）に対応する。ゆえに最後の主張だけを証明すればよい。

\medskip\noindent
スキーム $W$ と全射エタール射 $W \to Z$ を選ぶ。スキーム $V$ と全射エタール射
$V \to W \times_Z Y$ を選ぶ。スキーム $U$ と全射エタール射
$U \to V \times_Y X$ を選ぶ。最後に、スキーム $W'$ と全射エタール射
$W' \to W \times_Z Z'$ を選ぶ。$V' = W' \times_W V$ および
$U' = W' \times_W U$ とおく。すると全射エタール射 $V' \to Y'$ および
$U' \to X'$ を得る。
% P09-ERR-1215: frozen lines 11419--11420 omit "that" after "use without further mention".
代数空間のエタール射は付随する位相空間の開写像を誘導する、ということを以下では
断りなく用いる（Properties of Spaces, Lemma
\ref{spaces-properties-lemma-etale-open} を参照せよ）。定義により、$U(f)$ は
$|X|$ における集合 $T$ の像である。ここでこの集合は、$U$ のうちスキームの射
$U \to V$ が Koszul である点の集合である。同様に、$U(f')$ は $|X'|$ における
集合 $T'$ の像である。ここでこの集合は、$U'$ のうちスキームの射 $U' \to V'$ が
Koszul である点の集合である。構成により図式
$$
\xymatrix{
U' \ar[r] \ar[d] & U \ar[d] \\
V' \ar[r] & V
}
$$
は（スキームの圏で）Cartesian である。したがって前述の More on Morphisms,
Lemma \ref{more-morphisms-lemma-base-change-lci-fibres} により、$T'$ は $T$ の逆像である。
$|U'| \to |X'|$ は全射なので、補題が従う。
\end{proof}

\begin{lemma}
\label{spaces-more-morphisms-lemma-unramified-lci}
$S$ をスキームとする。$f : X \to Y$ を $S$ 上の代数空間の局所完全交叉射とする。
$f$ が非分岐であることと $f$ が形式的非分岐であることは同値であり、この場合、
余法層 $\mathcal{C}_{X/Y}$ は $X$ 上有限局所自由である。
\end{lemma}

\begin{proof}
エタール局所化により、スキームの射に対する対応する結果から従う。More on Morphisms,
Lemma \ref{more-morphisms-lemma-unramified-lci} および Lemma
\ref{spaces-more-morphisms-lemma-universal-thickening-localize} を参照せよ。
（この補題の状況では、射 $V \to U$ は定義により非分岐かつ局所完全交叉射であることに
注意せよ。）
\end{proof}

\begin{lemma}
\label{spaces-more-morphisms-lemma-transitivity-conormal-lci}
$S$ をスキームとする。$Z \to Y \to X$ を $S$ 上の代数空間の形式的非分岐射とする。
$Z \to Y$ は局所完全交叉射であると仮定する。Lemma
\ref{spaces-more-morphisms-lemma-transitivity-conormal} の完全列
$$
0 \to i^*\mathcal{C}_{Y/X} \to
\mathcal{C}_{Z/X} \to
\mathcal{C}_{Z/Y} \to 0
$$
は短完全である。
\end{lemma}

\begin{proof}
スキーム $U$ と全射エタール射 $U \to X$ を選ぶ。スキーム $V$ と全射エタール射
$V \to U \times_X Y$ を選ぶ。スキーム $W$ と全射エタール射
$W \to V \times_Y Z$ を選ぶ。Lemma \ref{spaces-more-morphisms-lemma-universal-thickening-localize} により、
射 $W \to V$ と $V \to U$ は形式的非分岐である。さらに、列
$i^*\mathcal{C}_{Y/X} \to \mathcal{C}_{Z/X} \to \mathcal{C}_{Z/Y} \to 0$ は、
$i^*\mathcal{C}_{V/U} \to \mathcal{C}_{W/U} \to \mathcal{C}_{W/V} \to 0$ に
制限される。これは $W \to V \to U$ に対する対応する列である。
したがってスキームに対する結果（More on Morphisms, Lemma
\ref{more-morphisms-lemma-transitivity-conormal-lci}）から結論が従う。実際、定義により
射 $W \to V$ は局所完全交叉射である。
\end{proof}










\section{射が同型となるのはいつか？}
\label{spaces-more-morphisms-section-when-isomorphism}

\noindent
より一般に、「射が性質 $\mathcal{P}$ をもつのはいつか？」と問うことができる。
より正確には次の問いである。代数空間の可換図式
$$
\xymatrix{
X \ar[rr]_f \ar[rd]_p & & Y \ar[ld]^q \\
& Z
}
$$
が与えられたとする。次の二つの性質をもつ代数空間のモノ射 $W \to Z$ は
存在するだろうか？
\begin{enumerate}
\item 基底変換 $f_W : X_W \to Y_W$ は性質 $\mathcal{P}$ をもつ。
\item 代数空間の任意の射 $Z' \to Z$ が $W$ を経由することと、基底変換
$f_{Z'} : X_{Z'} \to Y_{Z'}$ が性質 $\mathcal{P}$ をもつことは同値である。
\end{enumerate}
多くの場合、$W \to Z$ が存在するならば、それは埋入、開埋入、または閉埋入である。

\medskip\noindent
この問いの答えは、射 $f$, $p$, および $q$ の補助的な性質に依存しうる。
一例は $\mathcal{P}(f) =$「$f$ は平坦である」という性質である。スキームの射について、
$Y = S$ の場合を ``More on Flatness'' の章で詳細に論じた。More on Flatness,
Section \ref{flat-section-flattening-functors} から始まる。

\begin{lemma}
\label{spaces-more-morphisms-lemma-where-unramified}
代数空間の可換図式
$$
\xymatrix{
X \ar[rr]_f \ar[rd]_p & & Y \ar[ld]^q \\
& Z
}
$$
を考える。$p$ は局所有限型かつ閉であると仮定する。このとき開部分空間
$W \subset Z$ で、射 $Z' \to Z$ が $W$ を経由することと、基底変換
$f_{Z'} : X_{Z'} \to Y_{Z'}$ が非分岐であることが同値となるものが存在する。
\end{lemma}

\begin{proof}
Morphisms of Spaces, Lemma \ref{spaces-morphisms-lemma-where-unramified} により、
開部分空間 $U(f) \subset X$ が存在し、その点集合は $f$ が非分岐である点からなる。
さらに、$U(f)$ の形成は任意の基底変換と可換する。開部分空間 $W \subset Z$
（Properties of Spaces, Lemma
\ref{spaces-properties-lemma-open-subspaces} を参照）を、その点集合が
$$
|W| = |Z| \setminus |p|\left(|X| \setminus |U(f)|\right)
$$
となるものとする。すなわち、$z \in |Z|$ が $W$ の点であることと、$f$ が
$X$ のうち $z$ の上にあるすべての点で非分岐であることは同値である。
$p$ は閉であると仮定したので、これは開であることに注意せよ。$U(f)$ の形成は
任意の基底変換と可換するので、Properties of Spaces, Lemma
\ref{spaces-properties-lemma-factor-through-open-subspace} を用いれば、$W$ が所望の
普遍性をもつことが直ちに分かる。
\end{proof}

\begin{lemma}
\label{spaces-more-morphisms-lemma-where-unramified-universally-injective}
代数空間の可換図式
$$
\xymatrix{
X \ar[rr]_f \ar[rd]_p & & Y \ar[ld]^q \\
& Z
}
$$
を考える。次を仮定する：
\begin{enumerate}
\item $p$ は局所有限型である。
\item $p$ は閉である。
\item $p_2 : X \times_Y X \to Z$ は閉である。
\end{enumerate}
このとき開部分空間 $W \subset Z$ で、射 $Z' \to Z$ が $W$ を経由することと、
基底変換 $f_{Z'} : X_{Z'} \to Y_{Z'}$ が非分岐かつ普遍単射であることが同値となる
ものが存在する。
\end{lemma}

\begin{proof}
Lemma \ref{spaces-more-morphisms-lemma-where-unramified} で得た開部分空間で $Z$ を置き換えると、
$f$ はすでに非分岐であるとしてよい。この置き換えは仮定 (2) または (3) を
損なわないことに注意せよ。Morphisms of Spaces, Lemma
\ref{spaces-morphisms-lemma-diagonal-unramified-morphism} により、
$\Delta_{X/Y} : X \to X \times_Y X$ は開埋入である。これは任意の基底変換後にも
成り立つ。したがって Morphisms of Spaces, Lemma
\ref{spaces-morphisms-lemma-universally-injective} により、$f_{Z'}$ が普遍単射であることと、
対角射の基底変換 $X_{Z'} \to (X \times_Y X)_{Z'}$ が同型であることは同値である。
開部分空間 $W \subset Z$（Properties of Spaces, Lemma
\ref{spaces-properties-lemma-open-subspaces} を参照）を、その点集合が
$$
|W| = |Z| \setminus
|p_2|\left(|X \times_Y X| \setminus \Im(|\Delta_{X/Y}|)\right)
$$
となるものとする。すなわち、$z \in |Z|$ が $W$ の点であることと、写像
$|X \times_Y X| \to |Z|$ の $z$ 上のファイバーが
$|X| \to |X \times_Y X|$ の像に含まれることは同値である。以上の議論から、
制限 $p^{-1}(W) \to q^{-1}(W)$（$f$ のもの）が非分岐かつ普遍単射であることは明らかである。

\medskip\noindent
逆に、$f_{Z'}$ が非分岐かつ普遍単射であると仮定する。$Z' \to Z$ が $W$ を
経由することを示すには、$|Z'| \to |Z|$ の像が $|W|$ に含まれることを
示せば十分である。Properties of Spaces, Lemma
\ref{spaces-properties-lemma-factor-through-open-subspace} を参照せよ。
したがって、$Z'$ が体のスペクトルである場合を証明すれば十分である。
% P09-ERR-1216: frozen line 11621 omits "by" in the notation declaration "Denote z ... the image".
$z \in |Z|$ で $|Z'| \to |Z|$ の像を表す。以上の議論により、
$$
|X_{Z'}| \longrightarrow |(X \times_Y X)_{Z'}|
$$
は全射である。Properties of Spaces, Lemma
\ref{spaces-properties-lemma-points-cartesian} により、可換図式
$$
\xymatrix{
|X_{Z'}| \ar[d] \ar[r] &
|(X \times_Y X)_{Z'}| \ar[d] \\
|p|^{-1}(\{z\}) \ar[r] &
|p_2|^{-1}(\{z\})
}
$$
の垂直射は全射である。したがって所望のとおり $z \in |W|$ である。
\end{proof}

\begin{lemma}
\label{spaces-more-morphisms-lemma-where-closed-immersion}
代数空間の可換図式
$$
\xymatrix{
X \ar[rr]_f \ar[rd]_p & & Y \ar[ld]^q \\
& Z
}
$$
を考える。次を仮定する：
\begin{enumerate}
\item $p$ は局所有限型である。
\item $p$ は普遍閉である。
\item $q : Y \to Z$ は分離的である。
\end{enumerate}
このとき開部分空間 $W \subset Z$ で、射 $Z' \to Z$ が $W$ を経由することと、
基底変換 $f_{Z'} : X_{Z'} \to Y_{Z'}$ が閉埋入であることが同値となるものが存在する。
\end{lemma}

\begin{proof}
閉埋入を普遍閉、非分岐、かつ普遍単射な射として特徴付ける結果を用いる。
Lemma \ref{spaces-more-morphisms-lemma-characterize-closed-immersion} を参照せよ。まず、$p$ は普遍閉で
$q$ は分離的なので、$f$ は普遍閉である。Morphisms of Spaces, Lemma
\ref{spaces-morphisms-lemma-universally-closed-permanence} を参照せよ。
したがって $f$ の任意の基底変換は普遍閉である。Morphisms of Spaces, Lemma
\ref{spaces-morphisms-lemma-base-change-universally-closed} を参照せよ。
ゆえに補題の証明を終えるには、Lemma
\ref{spaces-more-morphisms-lemma-where-unramified-universally-injective} の仮定が満たされることを
証明すれば十分である。射影 $\text{pr}_0 : X \times_Y X \to X$ は $f$ の
基底変換として普遍閉である。Morphisms of Spaces, Lemma
\ref{spaces-morphisms-lemma-base-change-universally-closed} を参照せよ。
したがって $X \times_Y X \to Z$ は普遍閉射の合成として普遍閉である
（Morphisms of Spaces, Lemma
\ref{spaces-morphisms-lemma-composition-universally-closed} を参照）。
これで補題の証明が完了する。
\end{proof}

\begin{lemma}
\label{spaces-more-morphisms-lemma-where-flat}
代数空間の可換図式
$$
\xymatrix{
X \ar[rr]_f \ar[rd]_p & & Y \ar[ld]^q \\
& Z
}
$$
を考える。次を仮定する：
\begin{enumerate}
\item $p$ は局所有限表示である。
\item $p$ は平坦である。
\item $p$ は閉である。
\item $q$ は局所有限型である。
\end{enumerate}
このとき開部分空間 $W \subset Z$ で、射 $Z' \to Z$ が $W$ を経由することと、
基底変換 $f_{Z'} : X_{Z'} \to Y_{Z'}$ が平坦であることが同値となるものが存在する。
\end{lemma}

\begin{proof}
Lemma \ref{spaces-more-morphisms-lemma-base-change-flatness-fibres} により、集合
$$
A = \{x \in |X| : X\text{ flat at }x \text{ over }Y\}.
$$
は $|X|$ で開であり、その形成は任意の基底変換と可換する。開部分空間
$W \subset Z$（Properties of Spaces, Lemma
\ref{spaces-properties-lemma-open-subspaces} を参照）を、その点集合が
$$
|W| = |Z| \setminus |p|\left(|X| \setminus A\right)
$$
となるものとする。すなわち、$z \in |Z|$ が $W$ の点であることと、写像
$|X| \to |Z|$ の $z$ 上のファイバー全体が $A$ に含まれることは同値である。
$p$ は閉なので、これは開である。$A$ の形成は任意の基底変換と可換するため、
$W$ が所望の性質をもつことが従う。
\end{proof}

\begin{lemma}
\label{spaces-more-morphisms-lemma-where-surjective-flat}
代数空間の可換図式
$$
\xymatrix{
X \ar[rr]_f \ar[rd]_p & & Y \ar[ld]^q \\
& Z
}
$$
を考える。次を仮定する：
\begin{enumerate}
\item $p$ は局所有限表示である。
\item $p$ は平坦である。
\item $p$ は閉である。
\item $q$ は局所有限型である。
\item $q$ は閉である。
\end{enumerate}
このとき開部分空間 $W \subset Z$ で、射 $Z' \to Z$ が $W$ を経由することと、
基底変換 $f_{Z'} : X_{Z'} \to Y_{Z'}$ が全射かつ平坦であることが同値となるものが
存在する。
\end{lemma}

\begin{proof}
Lemma \ref{spaces-more-morphisms-lemma-where-flat} により、$f$ は平坦であるとしてよい。
Morphisms of Spaces, Lemma
\ref{spaces-morphisms-lemma-finite-presentation-permanence} により、$f$ は
局所有限表示であることに注意せよ。したがって Morphisms of Spaces, Lemma
\ref{spaces-morphisms-lemma-fppf-open} により $f$ は開である。開部分空間
$W \subset Z$（Properties of Spaces, Lemma
\ref{spaces-properties-lemma-open-subspaces} を参照）を、その点集合が
$$
|W| = |Z| \setminus |q|\left(|Y| \setminus |f|(|X|)\right).
$$
となるものとする。言い換えると、$z \in |Z|$ に対して $z \in |W|$ であることと、
写像 $|Y| \to |Z|$ の $z$ 上のファイバー全体が $|X| \to |Y|$ の像に
含まれることは同値である。$q$ は閉なので、この集合は $|Z|$ で開である。
構成により射 $X_W \to Y_W$ は全射である。最後に、$X_{Z'} \to Y_{Z'}$ が
全射であると仮定する。$Z' \to Z$ が $W$ を経由することを示すには、
$|Z'| \to |Z|$ の像が $|W|$ に含まれることを示せば十分である。
Properties of Spaces, Lemma
\ref{spaces-properties-lemma-factor-through-open-subspace} を参照せよ。
したがって $Z'$ が体のスペクトルである場合を証明すれば十分である。
% P09-ERR-1217: frozen line 11776 repeats the missing-"by" notation grammar: "Denote z ... the image".
$z \in |Z|$ で $|Z'| \to |Z|$ の像を表す。Properties of Spaces, Lemma
\ref{spaces-properties-lemma-points-cartesian} により、可換図式
$$
\xymatrix{
|X_{Z'}| \ar[d] \ar[r] &
|Y_{Z'}| \ar[d] \\
|p|^{-1}(\{z\}) \ar[r] &
|q|^{-1}(\{z\})
}
$$
の垂直射は全射である。したがって所望のとおり $z \in |W|$ である。
\end{proof}

\begin{lemma}
\label{spaces-more-morphisms-lemma-where-isomorphism}
代数空間の可換図式
$$
\xymatrix{
X \ar[rr]_f \ar[rd]_p & & Y \ar[ld]^q \\
& Z
}
$$
を考える。次を仮定する：
\begin{enumerate}
\item $p$ は局所有限表示である。
\item $p$ は平坦である。
\item $p$ は普遍閉である。
\item $q$ は局所有限型である。
\item $q$ は閉である。
\item $q$ は分離的である。
\end{enumerate}
このとき開部分空間 $W \subset Z$ で、射 $Z' \to Z$ が $W$ を経由することと、
基底変換 $f_{Z'} : X_{Z'} \to Y_{Z'}$ が同型であることが同値となるものが存在する。
\end{lemma}

\begin{proof}
Lemma \ref{spaces-more-morphisms-lemma-where-surjective-flat} により、開部分空間 $W_1 \subset Z$ で、
$f_{Z'}$ が全射かつ平坦であることと $Z' \to Z$ が $W_1$ を経由することが
同値となるものが存在する。Lemma \ref{spaces-more-morphisms-lemma-where-closed-immersion} により、
開部分空間 $W_2 \subset Z$ で、$f_{Z'}$ が閉埋入であることと $Z' \to Z$ が
$W_2$ を経由することが同値となるものが存在する。$W = W_1 \cap W_2$ が
所望のものだと主張する。確かに、$f_{Z'}$ が同型ならば $Z' \to Z$ は $W$ を
経由する。したがって $f_W$ が同型であることを示せば十分である。構成により
$f_W$ は全射、平坦、かつ閉埋入である。特に $f_W$ は表現可能である。
スキームの全射かつ平坦な閉埋入は同型なので（Morphisms, Lemma
\ref{morphisms-lemma-characterize-flat-closed-immersions} を参照）、結論を得る。
（実際には $f_W$ は局所有限表示であり、したがって開でもあるので、望むならこの補題の
使用を避けられることに注意せよ。）
\end{proof}

\begin{lemma}
\label{spaces-more-morphisms-lemma-where-lci}
代数空間の可換図式
$$
\xymatrix{
X \ar[rr]_f \ar[rd]_p & & Y \ar[ld]^q \\
& Z
}
$$
を考える。次を仮定する：
\begin{enumerate}
\item $p$ は平坦かつ局所有限表示である。
\item $p$ は閉である。
\item $q$ は平坦かつ局所有限表示である。
\end{enumerate}
このとき開部分空間 $W \subset Z$ で、射 $Z' \to Z$ が $W$ を経由することと、
基底変換 $f_{Z'} : X_{Z'} \to Y_{Z'}$ が局所完全交叉射であることが同値となる
ものが存在する。
\end{lemma}

\begin{proof}
Lemma \ref{spaces-more-morphisms-lemma-base-change-lci-fibres} により、開部分空間 $U(f) \subset X$ が存在し、
その点集合は $f$ が Koszul である点からなる。さらに、$U(f)$ の形成は任意の
基底変換と可換する。開部分空間 $W \subset Z$（Properties of Spaces, Lemma
\ref{spaces-properties-lemma-open-subspaces} を参照）を、その点集合が
$$
|W| = |Z| \setminus |p|\left(|X| \setminus |U(f)|\right)
$$
となるものとする。すなわち、$z \in |Z|$ が $W$ の点であることと、$f$ が
$X$ のうち $z$ の上にあるすべての点で Koszul であることは同値である。
$p$ は閉であると仮定したので、これは開であることに注意せよ。$U(f)$ の形成は
任意の基底変換と可換するので、Properties of Spaces, Lemma
\ref{spaces-properties-lemma-factor-through-open-subspace} を用いれば、$W$ が所望の
普遍性をもつことが直ちに分かる。
\end{proof}









\section{微分と余法層の完全列}
\label{spaces-more-morphisms-section-exact}

\noindent
本節では、余法層と微分層の完全列に関するいくつかの結果をまとめる。
ある意味では、これらはすべて、代数空間の合成可能な射に付随する余接複体の
三角の実現である。

\medskip\noindent
% P09-ERR-1218: frozen lines 11894--11895 use plural agreement in "each of the maps are".
以下の列の各写像は、Lemma \ref{spaces-more-morphisms-lemma-functoriality-differentials} または Lemma
\ref{spaces-more-morphisms-lemma-universal-thickening-functorial} で構成されたものである。
$S$ をスキームとする。$g : Z \to Y$ および $f : Y \to X$ を $S$ 上の
代数空間の射とする。
\begin{enumerate}
\item 標準完全列
$$
g^*\Omega_{Y/X} \to \Omega_{Z/X} \to \Omega_{Z/Y} \to 0,
$$
が存在する。Lemma \ref{spaces-more-morphisms-lemma-triangle-differentials} を参照せよ。
$g : Z \to Y$ が形式的に滑らかならば、この列は短完全列である。Lemma
\ref{spaces-more-morphisms-lemma-triangle-differentials-formally-smooth} を参照せよ。
\item $g$ が形式的非分岐ならば、標準完全列
$$
\mathcal{C}_{Z/Y} \to g^*\Omega_{Y/X} \to \Omega_{Z/X} \to 0,
$$
が存在する。Lemma \ref{spaces-more-morphisms-lemma-universally-unramified-differentials-sequence} を参照せよ。
$f \circ g : Z \to X$ が形式的に滑らかならば、この列は短完全列である。Lemma
\ref{spaces-more-morphisms-lemma-differentials-formally-unramified-formally-smooth} を参照せよ。
\item $g$ と $f \circ g$ が形式的非分岐ならば、標準完全列
$$
\mathcal{C}_{Z/X} \to \mathcal{C}_{Z/Y} \to g^*\Omega_{Y/X} \to 0,
$$
が存在する。Lemma \ref{spaces-more-morphisms-lemma-two-unramified-morphisms} を参照せよ。
$f : Y \to X$ が形式的に滑らかならば、この列は短完全列である。Lemma
\ref{spaces-more-morphisms-lemma-two-unramified-morphisms-formally-smooth} を参照せよ。
\item $g$ と $f$ が形式的非分岐ならば、標準完全列
$$
g^*\mathcal{C}_{Y/X} \to \mathcal{C}_{Z/X} \to \mathcal{C}_{Z/Y} \to 0.
$$
が存在する。Lemma \ref{spaces-more-morphisms-lemma-transitivity-conormal-unramified} を参照せよ。
$g : Z \to Y$ が局所完全交叉射ならば、この列は短完全列である。Lemma
\ref{spaces-more-morphisms-lemma-transitivity-conormal-lci} を参照せよ。
\end{enumerate}











\section{擬連接複体の特徴付け、II}
\label{spaces-more-morphisms-section-characterize-pseudo-coherent}

\noindent
本節では、コホモロジーによる擬連接複体の特徴付けを論じる。代数空間上の
擬連接複体に関する先の内容は Derived Categories of Spaces, Section
\ref{spaces-perfect-section-spell-out} および Derived Categories of Spaces, Section
\ref{spaces-perfect-section-pseudo-coherent-hocolim} にある。スキームについて本節に
対応するのは More on Morphisms, Section
\ref{more-morphisms-section-characterize-pseudo-coherent} である。基本的な道具は、
Chow の補題の導来版を用いて射影空間の場合へ帰着することである。Lemma
\ref{spaces-more-morphisms-lemma-derived-chow} を参照せよ。

\begin{lemma}
\label{spaces-more-morphisms-lemma-case-of-tor-independence}
$S$ をスキームとする。代数空間の可換図式
$$
\xymatrix{
Z' \ar[d] \ar[r] & Y' \ar[d] \\
X' \ar[r] & B'
}
$$
を考える。これは $S$ 上の図式とする。$B \to B'$ を射とする。$X$ と $Y$ で、
$X'$ と $Y'$ の $B$ への基底変換を表す。$Y' \to B'$ と $Z' \to X'$ は
平坦であると仮定する。このとき $X \times_B Y$ と $Z'$ は
$X' \times_{B'} Y'$ 上 Tor 独立である。
\end{lemma}

\begin{proof}
Derived Categories of Spaces, Lemma
\ref{spaces-perfect-lemma-tor-independent} により、Tor 独立性は
$X \times_B Y$ と $Z'$ 上でエタール局所的に確かめてよい。これにより
\footnote{議論をもう少し詳しく述べる。全射エタール射 $W' \to B'$ で
$W'$ がスキームであるものを選ぶ。全射エタール射
$W \to B \times_{B'} W'$ で $W$ がスキームであるものを選ぶ。全射エタール射
$U' \to X' \times_{B'} W'$ で $U'$ がスキームであるものを選ぶ。
全射エタール射 $V' \to Y' \times_{B'} W'$ で $V'$ がスキームであるものを選ぶ。
$U' \times_{W'} V' \to X' \times_{B'} Y'$ は全射エタールであることに注意せよ。
全射エタール射
$T' \to Z' \times_{X' \times_{B'} Y'} U' \times_{W'} V'$ で $T'$ がスキームで
あるものを選ぶ。
% P09-ERR-1219: frozen line 11998 omits "by" in "Denote U and V the base changes".
$U$ と $V$ で、$U'$ と $V'$ の $W$ への基底変換を表す。このとき補題の主張、
すなわち代数空間として $X \times_B Y$ と $Z'$ が $X' \times_{B'} Y'$ 上
Tor 独立であることは、スキームとして $U \times_W V$ と $T'$ が
$U' \times_{W'} V'$ 上 Tor 独立であることと同値である。したがって、角が
$T', U', V', W'$ である正方形と $W \to W'$ による基底変換について補題を
証明すれば十分である。$Y' \to B'$ と $Z' \to X'$ の平坦性は、
$V' \to W'$ と $T' \to U'$ の平坦性を含意する。}
補題はスキームの場合へ帰着する。これは More on Morphisms, Lemma
\ref{more-morphisms-lemma-case-of-tor-independence} である。
\end{proof}

\begin{lemma}[導来 Chow の補題]
\label{spaces-more-morphisms-lemma-derived-chow}
$A$ を環とする。$X$ を $A$ 上分離的かつ有限表示な代数空間とする。
$x \in |X|$ とする。このとき、$n \geq 0$、閉部分空間
$Z \subset X \times_A \mathbf{P}^n_A$、点 $z \in |Z|$、開部分空間
$V \subset \mathbf{P}^n_A$、および対象 $E$
（$D(\mathcal{O}_{X \times_A \mathbf{P}^n_A})$ のもの）で、次を満たすものが存在する：
\begin{enumerate}
\item $Z \to X \times_A \mathbf{P}^n_A$ は有限表示である。
\item $c : Z \to \mathbf{P}^n_A$ は $V$ 上閉埋入である。
$W = c^{-1}(V)$ とおく。
\item $b : Z \to X$ の $W$ への制限はエタールであり、$z \in W$ かつ
$b(z) = x$ である。
\item $E|_{X \times_A V} \cong
(b, c)_*\mathcal{O}_Z|_{X \times_A V}$ である。
\item $E$ は擬連接であり、$Z$ 上に台をもつ。
\end{enumerate}
\end{lemma}

\begin{proof}
有限型な $\mathbf{Z}$-部分代数 $A' \subset A$ と、代数空間 $X'$（$A'$ 上
分離的かつ有限表示なもの）で、その $A$ への基底変換が $X$ となるものを見つけられる。
Limits of Spaces, Lemmas
\ref{spaces-limits-lemma-descend-finite-presentation} および
\ref{spaces-limits-lemma-descend-separated-morphism} を参照せよ。
$x' \in |X'|$ を $x$ の像とする。$(X'/A', x')$ について補題を証明できれば、
$(X/A, x)$ について補題が従う。実際、$n', Z', z', V', E'$ が
$(X'/A', x')$ に対する解を与えるならば、$n = n'$ とし、
$Z \subset X \times \mathbf{P}^n$ を $Z'$ の逆像とし、$z \in Z$ を $x$ に写る
一意な点とし、$V \subset \mathbf{P}^n_A$ を $V'$ の逆像とし、$E$ を $E'$ の
導来引き戻しとする。Cohomology on Sites, Lemma
\ref{sites-cohomology-lemma-pseudo-coherent-pullback} により $E$ は擬連接であることに
注意せよ。
% P09-ERR-1220: frozen lines 12053 and 12075 say condition (5), but the intervening argument verifies the restriction identity in condition (4).
(5) を確かめることだけが残る。そのため $W = c^{-1}(V)$ および
$W' = (c')^{-1}(V')$ とおき、Cartesian な正方形
$$
\xymatrix{
W \ar[d]_{(b, c)} \ar[r] & W' \ar[d]^{(b', c')} \\
X \times_A V \ar[r] & X' \times_{A'} V'
}
$$
を考える。Lemma \ref{spaces-more-morphisms-lemma-case-of-tor-independence} により、$X \times_A V$ と
$W'$ は $X' \times_{A'} V'$ 上 Tor 独立である。したがって
$(b', c')_*\mathcal{O}_{W'}$ の $X \times_A V$ への導来引き戻しは
$(b, c)_*\mathcal{O}_W$ である。Derived Categories of Spaces, Lemma
\ref{spaces-perfect-lemma-compare-base-change} を参照せよ。ここではさらに、
$R(b', c')_*\mathcal{O}_{Z'} = (b', c')_*\mathcal{O}_{Z'}$ であることを用いている。
これは $(b', c')$ が閉埋入だからであり、$(b, c)_*\mathcal{O}_Z$ についても同様である。
% P09-ERR-1221: frozen lines 12072--12074 use undefined U' and U; the surrounding square and lemma data use X' and X.
$E'|_{U' \times_{A'} V'} =
(b', c')_*\mathcal{O}_{W'}$ なので、
$E|_{U \times_A V} = (b, c)_*\mathcal{O}_W$ を得て、(5) が成り立つ。
これにより、次の段落で述べる状況へ帰着された。

\medskip\noindent
$A$ は $\mathbf{Z}$ 上有限型であると仮定する。エタール射 $U \to X$ で
$U$ がアフィンスキームであるものと、点 $u \in U$（$x$ に写るもの）を選ぶ。
このとき $U$ は $A$ 上有限型である。閉埋入 $U \to \mathbf{A}^n_A$ を選び、
% P09-ERR-1222: frozen lines 12083--12085 omit \"by\" in the notation declaration for j.
$j : U \to \mathbf{P}^n_A$ で、開埋入
$\mathbf{A}^n_A \to \mathbf{P}^n_A$ と合成して得られる埋入を表す。
$Z$ を
$$
(\text{id}_U, j) : U \longrightarrow X \times_A \mathbf{P}^n_A
$$
のスキーム論的閉包とする。$z \in Z$ を $u$ の像とする。
$Y \subset \mathbf{P}^n_A$ を $j$ のスキーム論的閉包とする。このとき
$Z \subset X \times_A Y$ が
$(\text{id}_U, j) : U \to X \times_A Y$ のスキーム論的閉包であることは明らかである。
$X$ は分離的なので、射 $X \times_A Y \to Y$ も分離的である。したがって
Morphisms of Spaces, Lemma
\ref{spaces-morphisms-lemma-scheme-theoretic-image-of-partial-section} により、
$Z \to Y$ は開部分スキーム $j(U) \subset Y$ 上で同型である。
$V \subset \mathbf{P}^n_A$ を $V \cap Y = j(U)$ となる開部分空間として選ぶ。
すると (2) が成り立ち、$W = (\text{id}_U, j)(U)$ であり、したがって (3) が
成り立つ。(1) は $A$ が Noetherian なので成り立つ。

\medskip\noindent
$A$ は Noetherian なので、$X$ と $X \times_A \mathbf{P}^n_A$ は Noetherian
代数空間である。したがってこの場合 $E = (b, c)_*\mathcal{O}_Z$ と取れる。
(4) は明らかであり、(5) については Derived Categories of Spaces, Lemma
\ref{spaces-perfect-lemma-identify-pseudo-coherent-noetherian} を参照せよ。
これで証明が完了する。
\end{proof}

\begin{lemma}
\label{spaces-more-morphisms-lemma-compute-Fourier-Mukai-for-derived-chow}
Lemma \ref{spaces-more-morphisms-lemma-derived-chow} と同様に、$X/A$、$x \in |X|$、および
$n, Z, z, V, E$ が与えられているとする。
任意の $K \in D_\QCoh(\mathcal{O}_X)$ に対して
$$
Rq_*(Lp^*K \otimes^\mathbf{L} E)|_V = R(W \to V)_*K|_W
$$
が成り立つ。ここで $p : X \times_A \mathbf{P}^n_A \to X$ および
$q : X \times_A \mathbf{P}^n_A \to \mathbf{P}^n_A$ は射影であり、
射 $W \to V$ は有限表示な閉埋入 $c|_W : W \to V$ である。
\end{lemma}

\begin{proof}
$W = c^{-1}(V)$ であり、$c$ は $V$ 上閉埋入なので、$c|_W$ は閉埋入である。
$W$ と $V$ は $A$ 上有限表示なので、これは有限表示である。Morphisms of Spaces,
Lemma \ref{spaces-morphisms-lemma-finite-presentation-permanence} を参照せよ。
まず
$$
Rq_*(Lp^*K \otimes^\mathbf{L} E)|_V =
Rq'_*\left((Lp^*K \otimes^\mathbf{L} E)|_{X \times_A V}\right)
$$
を得る。ここで $q' : X \times_A V \to V$ は射影である。これは全直像の構成が
局所化と可換だからである。
% P09-ERR-1223: frozen line 12142 omits "by" in the notation declaration for i.
$i = (b, c)|_W : W \to X \times_A V$ で、与えられた閉埋入を表す。このとき
$$
Rq'_*\left((Lp^*K \otimes^\mathbf{L} E)|_{X \times_A V}\right) =
Rq'_*(Lp^*K|_{X \times_A V} \otimes^\mathbf{L} i_*\mathcal{O}_W)
$$
が性質 (5) により成り立つ。$i$ は閉埋入なので
$i_*\mathcal{O}_W = Ri_*\mathcal{O}_W$ である。Derived Categories of Spaces,
Lemma \ref{spaces-perfect-lemma-cohomology-base-change} を用いると、これは
$$
Rq'_* Ri_* Li^* Lp^*K|_{X \times_A V} =
R(q' \circ i)_* Lb^*K|_W =
R(W \to V)_* K|_W
$$
と書き換えられ、これが求める等式である。（$W$ への制限と、$W \to X$ による
導来引き戻しは、$W$ が $X$ 上エタールなので同じであることに注意せよ。）
\end{proof}

\begin{lemma}
\label{spaces-more-morphisms-lemma-characterize-pseudo-coherent}
$A$ を環とする。$X$ を $A$ 上分離的かつ有限表示な代数空間とする。
$K \in D_\QCoh(\mathcal{O}_X)$ とする。
$R\Gamma(X, E \otimes^\mathbf{L} K)$ が $D(A)$ において擬連接であることが、
任意の擬連接な $E$（$D(\mathcal{O}_X)$ の対象）について成り立つならば、
$K$ は $A$ に関して擬連接である
（Definition \ref{spaces-more-morphisms-definition-relative-pseudo-coherence}）。
\end{lemma}

\begin{proof}
$K \in D_\QCoh(\mathcal{O}_X)$ とし、
$R\Gamma(X, E \otimes^\mathbf{L} K)$ が $D(A)$ において擬連接であることが、
任意の擬連接な $E$（$D(\mathcal{O}_X)$ の対象）について成り立つと仮定する。
$x \in |X|$ とする。$K$ が $A$ に関して擬連接であることを、$x$ のエタール近傍で
示そう。相対擬連接性の定義により、これで補題が従う。

\medskip\noindent
Lemma \ref{spaces-more-morphisms-lemma-derived-chow} と同様に $n, Z, z, V, E$ を選ぶ。
$p : X \times \mathbf{P}^n \to X$ および
$q : X \times \mathbf{P}^n \to \mathbf{P}^n_A$ で射影を表す。
このとき、任意の $i \in \mathbf{Z}$ に対して
\begin{align*}
& R\Gamma(\mathbf{P}^n_A,
Rq_*(Lp^*K \otimes^\mathbf{L} E)
\otimes^\mathbf{L}
\mathcal{O}_{\mathbf{P}^n_A}(i)) \\
& =
R\Gamma(X \times \mathbf{P}^n,
Lp^*K \otimes^\mathbf{L} E \otimes^\mathbf{L}
Lq^*\mathcal{O}_{\mathbf{P}^n_A}(i)) \\
& =
R\Gamma(X,
K \otimes^\mathbf{L} Rp_*(E \otimes^\mathbf{L}
Lq^*\mathcal{O}_{\mathbf{P}^n_A}(i)))
\end{align*}
が Derived Categories of Spaces, Lemma
\ref{spaces-perfect-lemma-cohomology-base-change} により成り立つ。
Derived Categories of Spaces, Lemma
\ref{spaces-perfect-lemma-flat-proper-pseudo-coherent-direct-image-general} により、
複体 $Rp_*(E \otimes^\mathbf{L} Lq^*\mathcal{O}_{\mathbf{P}^n_A}(i))$ は
$X$ 上擬連接である。したがって仮定より、表示式は $D(A)$ の擬連接対象である。
Derived Categories of Schemes, Lemma
\ref{perfect-lemma-pseudo-coherent-on-projective-space} により、
$Rq_*(Lp^*K \otimes^\mathbf{L} E)$ は $\mathbf{P}^n_A$ 上擬連接であると結論する。
Lemma \ref{spaces-more-morphisms-lemma-compute-Fourier-Mukai-for-derived-chow} により
% P09-ERR-1224: frozen line 12219 restricts the q-pushforward to X x_A V although it lives on P^n_A and the preceding lemma restricts it to V.
$$
Rq_*(Lp^*K \otimes^\mathbf{L} E)|_{X \times_A V} =
R(W \to V)_*K|_W
$$
$W \to V$ は $\mathbf{P}^n_A$ の開部分スキームへの閉埋入なので、これは
$K|_W$ が $A$ に関して擬連接であることを意味する。例えば More on Morphisms,
Lemma \ref{more-morphisms-lemma-check-relative-pseudo-coherence-on-charts} を参照せよ。
\end{proof}

\begin{lemma}
\label{spaces-more-morphisms-lemma-characterize-pseudo-coh-improved}
$A$ を環とする。$X$ を $A$ 上分離的かつ有限表示な代数空間とする。
$K \in D_\QCoh(\mathcal{O}_X)$ とする。
$R \Gamma (X, E \otimes ^{\mathbf{L}} K)$ が $D(A)$ において擬連接であることが、
任意の完全な $E \in D(\mathcal{O}_X)$ について成り立つならば、$K$ は
$A$ に関して擬連接である。
\end{lemma}

\begin{proof}
Lemma \ref{spaces-more-morphisms-lemma-characterize-pseudo-coherent} に照らせば、
$R \Gamma (X, E \otimes ^{\mathbf{L}} K)$ が $D(A)$ において擬連接であることを、
任意の擬連接な $E \in D(\mathcal{O}_X)$ について示せば十分である。
Derived Categories of Spaces, Proposition
\ref{spaces-perfect-proposition-detecting-bounded-above} により
$K \in D^-_\QCoh (\mathcal{O}_X)$ が従う。すると Derived Categories of Spaces,
Lemma \ref{spaces-perfect-lemma-perfect-enough} により結果が従う。
\end{proof}













\section{相対完全対象}
\label{spaces-more-morphisms-section-relatively-perfect}

\noindent
本節では \cite{lieblich-complexes} に由来する概念を導入する。この概念はスキームの
射について Derived Categories of Schemes, Section
\ref{perfect-section-relatively-perfect} で論じられている。

\begin{definition}
\label{spaces-more-morphisms-definition-relatively-perfect}
$S$ をスキームとする。$f : X \to Y$ を $S$ 上の代数空間の射で、平坦かつ
局所有限表示なものとする。$E$（$D(\mathcal{O}_X)$ の対象）が
{\it $Y$ に関して完全} または {\it $Y$-完全} であるとは、$E$ が擬連接であり
（Cohomology on Sites, Definition
\ref{sites-cohomology-definition-pseudo-coherent}）、かつ $E$ が局所的に
$D(f^{-1}\mathcal{O}_Y)$ の対象として有限 Tor 次元をもつことをいう
（Cohomology on Sites, Definition
\ref{sites-cohomology-definition-tor-amplitude}）。
\end{definition}

\noindent
議論については Derived Categories of Schemes, Remark
\ref{perfect-remark-discuss-rel-perfect} を参照されたい。ここでは、Lemma
\ref{spaces-more-morphisms-lemma-relative-pseudo-coherent-is-moot} により、$E$ が擬連接であることは
$E$ が $Y$ に関して擬連接であることと同じである、とだけ述べておく。さらに、
$E$ の擬連接性は $E \in D_\QCoh(\mathcal{O}_X)$ を含意する。Derived Categories of
Spaces, Lemma \ref{spaces-perfect-lemma-pseudo-coherent} を参照せよ。

\begin{example}
\label{spaces-more-morphisms-example-relatively-perfect-field}
$k$ を体とする。$X$ を $k$ 上有限表示な代数空間とする（特に $X$ は準コンパクトで
ある）。このとき $E$（$D(\mathcal{O}_X)$ の対象）が $k$-完全であるための必要十分条件は、
定義により有界かつ擬連接であることであり、これは Derived Categories of Spaces,
Lemma \ref{spaces-perfect-lemma-identify-pseudo-coherent-noetherian} により
$D^b_{\textit{Coh}}(X)$ に属することと同値である。したがって相対完全であることは
「ファイバー上完全」であることを意味するものでは{\bf ない}。
\end{example}

\noindent
対応する代数的概念は More on Algebra, Section
\ref{more-algebra-section-relatively-perfect} で研究されている。代数空間に対する概念と
代数的概念は次のように結び付けられる。

\begin{lemma}
\label{spaces-more-morphisms-lemma-affine-locally-rel-perfect}
$S$ をスキームとする。$f : X \to Y$ を $S$ 上の代数空間の射で、平坦かつ
局所有限表示なものとする。$E \in D_\QCoh(\mathcal{O}_X)$ とする。次は同値である：
\begin{enumerate}
\item $E$ は $Y$-完全である。
\item 任意の可換図式
$$
\xymatrix{
U \ar[d] \ar[r]_g & V \ar[d] \\
X \ar[r]^f & Y
}
$$
で、$U$, $V$ がスキームで垂直射がエタールであるものに対して、複体 $E|_U$ は
Derived Categories of Schemes, Definition
\ref{perfect-definition-relatively-perfect} の意味で $V$-完全である。
\item (2) のような可換図式で $U \to X$ が全射であるものが存在し、その図式に対して
複体 $E|_U$ は Derived Categories of Schemes, Definition
\ref{perfect-definition-relatively-perfect} の意味で $V$-完全である。
\item (2) のような任意の可換図式で $U$ と $V$ がアフィンであるものに対して、
複体 $R\Gamma(U, E)$ は $\mathcal{O}_Y(V)$-完全である。
\end{enumerate}
\end{lemma}

\begin{proof}
補題の (2), (3), (4) を意味あるものにするため、$E|_U$（
$D_\QCoh(\mathcal{O}_U)$ の対象）は $E_0$（$D_\QCoh(\mathcal{O}_{U_0})$ の対象）
に対応することに注意する。
ここで $U_0$ は $U$ の台となるスキームを表す。Derived Categories of Spaces,
Lemma \ref{spaces-perfect-lemma-derived-quasi-coherent-small-etale-site} を参照せよ。
さらにこの場合、$E_0$ が擬連接であるための必要十分条件は $E|_U$ が擬連接である
ことである。Derived Categories of Spaces, Lemma
\ref{spaces-perfect-lemma-descend-pseudo-coherent} を参照せよ。また、
$E|_U$ が $f^{-1}\mathcal{O}_Y|_U = g^{-1}\mathcal{O}_V$ 上局所有限 Tor 次元を
もつための必要十分条件は、$E_0$ が $g_0^{-1}\mathcal{O}_{V_0}$ 上局所有限 Tor 次元を
もつことである。Derived Categories of Spaces, Lemma
\ref{spaces-perfect-lemma-tor-dimension-rel} を参照せよ。ここで
$g_0 : U_0 \to V_0$ は $g : U \to V$ を表現するスキームの射である
（記法については Derived Categories of Spaces, Remark
\ref{spaces-perfect-remark-match-total-direct-images} を参照）。最後に、「擬連接である」
ことはエタール局所的であり、もちろん「局所有限 Tor 次元をもつ」こともエタール局所的
であることに注意せよ。したがって $Y$-完全性はエタール局所的に確かめれば十分であり、
上の議論から (1) は (2) を含意し、(3) は (1) を含意する。さらに Derived Categories
of Schemes, Lemma \ref{perfect-lemma-affine-locally-rel-perfect} により (4) は
(2) および (3) と同値なので、証明は完了する。
\end{proof}

\begin{lemma}
\label{spaces-more-morphisms-lemma-triangulated}
$S$ をスキームとする。$f : X \to Y$ を $S$ 上の代数空間の射で、平坦かつ
局所有限表示なものとする。$D(\mathcal{O}_X)$ の $Y$-完全対象からなる充満部分圏は、
飽和した\footnote{Derived Categories, Definition
\ref{derived-definition-saturated}.} 三角部分圏である。
\end{lemma}

\begin{proof}
これは Cohomology on Sites, Lemmas
\ref{sites-cohomology-lemma-cone-pseudo-coherent},
\ref{sites-cohomology-lemma-summands-pseudo-coherent},
\ref{sites-cohomology-lemma-cone-tor-amplitude}, および
\ref{sites-cohomology-lemma-summands-tor-amplitude} から従う。
\end{proof}

\begin{lemma}
\label{spaces-more-morphisms-lemma-perfect-relatively-perfect}
$S$ をスキームとする。$f : X \to Y$ を $S$ 上の代数空間の射で、平坦かつ
局所有限表示なものとする。$D(\mathcal{O}_X)$ の完全対象は $Y$-完全である。
$K, M \in D(\mathcal{O}_X)$ とするとき、$K \otimes_{\mathcal{O}_X}^\mathbf{L} M$ が
$Y$-完全であるのは、$K$ が完全で $M$ が $Y$-完全な場合である。
\end{lemma}

\begin{proof}
Lemma \ref{spaces-more-morphisms-lemma-affine-locally-rel-perfect} を用いてスキームの場合へ帰着し、
Derived Categories of Schemes, Lemma
\ref{perfect-lemma-perfect-relatively-perfect} を適用する。
\end{proof}

\begin{lemma}
\label{spaces-more-morphisms-lemma-base-change-relatively-perfect}
$S$ をスキームとする。$f : X \to Y$ を $S$ 上の代数空間の射で、平坦かつ
局所有限表示なものとする。$g : Y' \to Y$ を $S$ 上の代数空間の射とする。
$X' = Y' \times_Y X$ とおき、
% P09-ERR-1225: frozen line 12407 omits "by" in the notation declaration for g'.
$g' : X' \to X$ で射影を表す。$K \in D(\mathcal{O}_X)$ が $Y$-完全ならば、
$L(g')^*K$ は $Y'$-完全である。
\end{lemma}

\begin{proof}
Lemma \ref{spaces-more-morphisms-lemma-affine-locally-rel-perfect} を用いてスキームの場合へ帰着し、
Derived Categories of Schemes, Lemma
\ref{perfect-lemma-base-change-relatively-perfect} を適用する。
\end{proof}

\begin{situation}
\label{spaces-more-morphisms-situation-relative-descent}
$S$ をスキームとする。$Y = \lim_{i \in I} Y_i$ を $S$ 上の代数空間の有向系の
極限とし、その遷移射 $g_{i'i} : Y_{i'} \to Y_i$ はアフィンであるとする。
$Y_i$ は準コンパクトかつ準分離的であると仮定する（すべての $i \in I$ について）。
$g_i : Y \to Y_i$ で射影を表す。元 $0 \in I$ と有限表示な平坦射
$X_0 \to Y_0$ を固定する。$X_i = Y_i \times_{Y_0} X_0$ および
$X = Y \times_{Y_0} X_0$ とおき、遷移射 $f_{i'i} : X_{i'} \to X_i$ と射影
$f_i : X \to X_i$ を表す。
\end{situation}

\begin{lemma}
\label{spaces-more-morphisms-lemma-relative-descend-homomorphisms}
% P09-ERR-1226: frozen line 12436 is the sentence fragment "In Situation ...".
Situation \ref{spaces-more-morphisms-situation-relative-descent} のもとで考える。
$K_0$ と $L_0$ を $D(\mathcal{O}_{X_0})$ の対象とする。
$K_i = Lf_{i0}^*K_0$ および $L_i = Lf_{i0}^*L_0$（$i \geq 0$）とおき、
$K = Lf_0^*K_0$ および $L = Lf_0^*L_0$ とおく。このとき写像
$$
\colim_{i \geq 0} \Hom_{D(\mathcal{O}_{X_i})}(K_i, L_i)
\longrightarrow
\Hom_{D(\mathcal{O}_X)}(K, L)
$$
は、$K_0$ が擬連接であり、$L_0 \in D_\QCoh(\mathcal{O}_{X_0})$ が
$D((X_0 \to Y_0)^{-1}\mathcal{O}_{Y_0})$ の対象として（局所的に）有限 Tor 次元を
もつならば同型である。
% P09-ERR-1227: frozen lemma ending at line 12449 lacks terminal punctuation.
\end{lemma}

\begin{proof}
$U_0$ を $(X_0)_{spaces, \etale}$ の任意の準コンパクトかつ準分離的な対象とし、
条件 $P$ を、写像
$$
\colim_{i \geq 0} \Hom_{D(\mathcal{O}_{U_i})}(K_i|_{U_i}, L_i|_{U_i})
\longrightarrow
\Hom_{D(\mathcal{O}_U)}(K|_U, L|_U)
$$
が同型であることとする。ここで $U = X \times_{X_0} U_0$ および
$U_i = X_i \times_{X_0} U_0$ である。$P$ が各 $U_0$ について成り立つことを示す。

\medskip\noindent
$(U_0 \subset W_0, V_0 \to W_0)$ が $(X_0)_{spaces, \etale}$ の基本識別正方形で、
$P$ が $U_0, V_0, U_0 \times_{W_0} V_0$ について成り立つと仮定する。
このとき導来圏の hom に対する Mayer--Vietoris により $P$ は $W_0$ についても
成り立つ。Derived Categories of Spaces, Lemma
\ref{spaces-perfect-lemma-mayer-vietoris-hom} を参照せよ。

\medskip\noindent
まず $U_0 = W_0 \times_{Y_0} X_0$ であり、$W_0$ が
$(Y_0)_{spaces, \etale}$ の準コンパクトかつ準分離的な対象である場合を考える。
これらの $W_0$ と前段落に Derived Categories of Spaces, Lemma
\ref{spaces-perfect-lemma-induction-principle} の帰納原理を適用すると、$P$ を
$U_0 = W_0 \times_{Y_0} X_0$ について証明すれば十分であり、ここで $W_0$ は
アフィンとしてよい。言い換えれば、$Y_0$ がアフィンの場合へ帰着した。次に帰納原理をもう一度、
$X_0$ のすべての準コンパクトかつ準分離的な開部分に適用し、$X_0$ もアフィンである
場合へ帰着する。

\medskip\noindent
$X_0$ と $Y_0$ がアフィンならば、Derived Categories of Schemes, Lemma
\ref{perfect-lemma-relative-descend-homomorphisms} で証明されたスキームの場合に戻る。
主張をスキームとスキーム上の加群の導来圏だけを含む主張へ移すには、
Derived Categories of Spaces, Lemmas
\ref{spaces-perfect-lemma-pseudo-coherent},
\ref{spaces-perfect-lemma-derived-quasi-coherent-small-etale-site},
\ref{spaces-perfect-lemma-descend-pseudo-coherent}, および
\ref{spaces-perfect-lemma-tor-dimension-rel} を用いればよい。
\end{proof}

\begin{lemma}
\label{spaces-more-morphisms-lemma-descend-relatively-perfect}
Situation \ref{spaces-more-morphisms-situation-relative-descent} において、$Y$-完全対象（
$D(\mathcal{O}_X)$ のもの）の圏は、$Y_i$-完全対象（$D(\mathcal{O}_{X_i})$ のもの）
の圏の余極限である。
\end{lemma}

\begin{proof}
$U_0$ を $(X_0)_{spaces, \etale}$ の任意の準コンパクトかつ準分離的な対象とし、
条件 $P$ を、関手
$$
\colim_{i \geq 0} D_{Y_i\text{-perfect}}(\mathcal{O}_{U_i})
\longrightarrow
D_{Y\text{-perfect}}(\mathcal{O}_U)
$$
が圏同値であることとする。ここで $U = X \times_{X_0} U_0$ および
$U_i = X_i \times_{X_0} U_0$ である。Lemma
\ref{spaces-more-morphisms-lemma-relative-descend-homomorphisms} により、この関手が充満忠実であることは
すでに分かっている。したがって本質的全射性を証明すれば十分である。

\medskip\noindent
$(U_0 \subset W_0, V_0 \to W_0)$ が $(X_0)_{spaces, \etale}$ の基本識別正方形で、
$P$ が $U_0, V_0, U_0 \times_{W_0} V_0$ について成り立つと仮定する。
$P$ が $W_0$ について成り立つと主張する。記法
$U_i = X_i \times_{X_0} U_0$、$U = X \times_{X_0} U_0$ を用い、$V_0$ と
$W_0$ についても同様の記法を用いる。また、$f_i$ を、すべての射
$U \to U_i$、$V \to V_i$、$U \times_W V \to U_i \times_{W_i} V_i$、および
$W \to W_i$ を表す記号として濫用する。
% P09-ERR-1228: frozen lines 12529, 12532, and 12533 use "an Y-perfect" / "an Y_i-perfect" instead of "a".
$E$ を $Y$-完全対象（$D(\mathcal{O}_W)$ の対象）とする。目標は、$E$ がこの関手の
本質像に属することを示すことである。仮定により、$i \geq 0$、$Y_i$-完全対象
$E_{U, i}$（$U_i$ 上のもの）、$Y_i$-完全対象 $E_{V, i}$（$V_i$ 上のもの）、および同型
$Lf_i^*E_{U, i} \to E|_U$ と $Lf_i^*E_{V, i} \to E|_V$ を見つけられる。
$$
a : E_{U, i} \to (Rf_{i, *}E)|_{U_i}
\quad\text{and}\quad
b : E_{V, i} \to (Rf_{i, *}E)|_{V_i}
$$
を、同型 $Lf_i^*E_{U, i} \to E|_U$ および
$Lf_i^*E_{V, i} \to E|_V$ に随伴する写像とする。充満忠実性により、$i$ を大きく
した後、同型
$c : E_{U, i}|_{U_i \times_{W_i} V_i} \to E_{V, i}|_{U_i \times_{W_i} V_i}$
で、その引き戻しが同一視
$$
Lf_i^*E_{U, i}|_{U \times_W V} \to E|_{U \times_W V} \to
Lf_i^*E_{V, i}|_{U \times_W V}.
$$
となるものを見つけられる。Derived Categories of Spaces, Lemma
\ref{spaces-perfect-lemma-glue} を適用して、対象 $E_i$（$W_i$ 上のもの）と写像
$d : E_i \to Rf_{i, *}E$ で、写像 $a$ および $b$ に制限されるもの（それぞれ
$U_i$ および $V_i$ 上）を得る。このとき $E_i$ が $Y_i$-完全であること（相対完全性は
エタール局所的な性質だから）と、$d$ が同型 $Lf_i^*E_i \to E$ に随伴することは
明らかである。

\medskip\noindent
Lemma \ref{spaces-more-morphisms-lemma-relative-descend-homomorphisms} の証明で用いたものと全く同じ議論を、
帰納原理（Derived Categories of Spaces, Lemma
\ref{spaces-perfect-lemma-induction-principle}）とともに用い、$X_0$ と $Y_0$ がともに
アフィンである場合へ帰着する。まず $(Y_0)_{spaces, \etale}$ の準コンパクトかつ
準分離的な対象について作業して $Y_0$ がアフィンの場合へ帰着し、次に
% P09-ERR-1229: frozen line 12568 uses singular "object" after "quasi-compact and quasi-separated" in the corresponding plural construction.
$(X_0)_{spaces, \etale}$ の準コンパクトかつ準分離的な対象について作業して
$X_0$ がアフィンの場合へ帰着する。アフィンの場合、結果は Derived Categories of
Schemes, Lemma \ref{perfect-lemma-descend-relatively-perfect} であるスキームの場合から従う。
スキームの場合への移行は Lemma \ref{spaces-more-morphisms-lemma-affine-locally-rel-perfect} により行われる。
\end{proof}

\begin{lemma}
\label{spaces-more-morphisms-lemma-derived-pushforward-rel-perfect}
$S$ をスキームとする。$f : X \to Y$ を $S$ 上の代数空間の射で、平坦、固有、かつ
有限表示なものとする。$E \in D(\mathcal{O}_X)$ を $Y$-完全とする。このとき
$Rf_*E$ は $D(\mathcal{O}_Y)$ の完全対象であり、その構成は任意の基底変換と可換である。
\end{lemma}

\begin{proof}
基底変換に関する主張は Derived Categories of Spaces, Lemma
\ref{spaces-perfect-lemma-base-change-tensor} である（$\mathcal{G}^\bullet$ を
$\mathcal{O}_X$（次数 $0$）とする）。したがって $Rf_*E$ が完全対象であることを示せば
十分である。極限の議論により、$Y$ が Noetherian アフィンである場合へ帰着する。

\medskip\noindent
問題は $Y$ 上エタール局所的なので、$Y$ はアフィンと仮定してよい。
$Y = \Spec(R)$ と書く。$R = \colim R_i$ を Noetherian 環 $R_i$ のフィルター付き
余極限として書く。Limits of Spaces, Lemma
\ref{spaces-limits-lemma-descend-finite-presentation} により、ある $i$ と代数空間
$X_i$（$R_i$ 上有限表示）で、その $R$ への基底変換が $X$ となるものが存在する。
Limits of Spaces, Lemmas \ref{spaces-limits-lemma-eventually-proper} および
\ref{spaces-limits-lemma-descend-flat} により、$X_i$ は $R_i$ 上固有かつ平坦であると
仮定してよい。Lemma \ref{spaces-more-morphisms-lemma-descend-relatively-perfect} により、$R_i$-完全対象
$E_i$（$D(\mathcal{O}_{X_i})$ のもの）で、その $X$ への引き戻しが $E$ となるものが
存在すると仮定してよい。Derived Categories of Spaces, Lemma
\ref{spaces-perfect-lemma-perfect-direct-image} を $X_i \to \Spec(R_i)$ と $E_i$ に適用し、
すでに示した基底変換性を用いると結果を得る。
\end{proof}

\begin{lemma}
\label{spaces-more-morphisms-lemma-compute-ext-rel-perfect}
$S$ をスキームとする。$f : X \to Y$ を $S$ 上の代数空間の射とする。
$E, K \in D(\mathcal{O}_X)$ とする。次を仮定する：
\begin{enumerate}
\item $Y$ は準コンパクトかつ準分離的である。
\item $f$ は固有、平坦、かつ有限表示である。
\item $E$ は $Y$-完全である。
\item $K$ は擬連接である。
\end{enumerate}
このとき擬連接な $L \in D(\mathcal{O}_Y)$ で
$$
Rf_*R\SheafHom(K, E) = R\SheafHom(L, \mathcal{O}_Y)
$$
を満たすものが存在し、任意の基底変換後にも同じことが成り立つ。すなわち
$$
\vcenter{
\xymatrix{
X' \ar[r]_{g'} \ar[d]_{f'} &
X \ar[d]^f \\
Y' \ar[r]^g &
Y
}
}
\quad\quad
\begin{matrix}
\text{Cartesian ならば } \\
Rf'_*R\SheafHom(L(g')^*K, L(g')^*E) \\
= R\SheafHom(Lg^*L, \mathcal{O}_{Y'})
\end{matrix}
$$
が成り立つ。
\end{lemma}

\begin{proof}
$Y$ は準コンパクトかつ準分離的なので、$X$ についても同じことが成り立つ。
Derived Categories of Spaces, Lemma
\ref{spaces-perfect-lemma-pseudo-coherent-hocolim} により、
$K = \text{hocolim} K_n$ と書ける。ここで $K_n$ は完全で、$K_n \to K$ は切断
$\tau_{\geq -n}$ 上同型を誘導する。$K_n^\vee$ を双対完全複体とする
（Cohomology on Sites, Lemma
\ref{sites-cohomology-lemma-dual-perfect-complex}）。完全対象の逆系
$\ldots \to K_3^\vee \to K_2^\vee \to K_1^\vee$ を得る。Lemma
\ref{spaces-more-morphisms-lemma-perfect-relatively-perfect} により
$K_n^\vee \otimes_{\mathcal{O}_X} E$ は $Y$-完全である。したがって Lemma
\ref{spaces-more-morphisms-lemma-derived-pushforward-rel-perfect} を
$K_n^\vee \otimes_{\mathcal{O}_X} E$ に適用でき、逆系
$$
\ldots \to M_3 \to M_2 \to M_1
$$
を得る。これは $Y$ 上の完全複体の逆系であり、
$$
M_n = Rf_*(K_n^\vee \otimes_{\mathcal{O}_X}^\mathbf{L} E) =
Rf_*R\SheafHom(K_n, E)
$$
を満たす。さらに、これらの複体の構成は任意の基底変換と可換である。すなわち
$Lg^*M_n =
Rf'_*((L(g')^*K_n)^\vee \otimes_{\mathcal{O}_{X'}}^\mathbf{L} L(g')^*E) =
Rf'_*R\SheafHom(L(g')^*K_n, L(g')^*E)$ である。

\medskip\noindent
$K_n \to K$ は $\tau_{\geq -n}$ 上同型を誘導するので、$K_n \to K_{n + 1}$ も
$\tau_{\geq -n}$ 上同型を誘導する。したがって
$K_{n + 1}^\vee \to K_n^\vee$ は $\tau_{\leq n}$ 上同型を誘導する。実際、
$K_n^\vee = R\SheafHom(K_n, \mathcal{O}_X)$ である。
$E$ が $[a, b]$ に Tor 振幅をもつ $f^{-1}\mathcal{O}_Y$-加群の複体であると仮定する。
すると任意の基底変換後にも同じことが成り立つ。Derived Categories of Spaces,
Lemma \ref{spaces-perfect-lemma-tor-independence-and-tor-amplitude} を参照せよ。よって
$$
K_{n + 1}^\vee \otimes_{\mathcal{O}_X} E \to
K_n^\vee \otimes_{\mathcal{O}_X} E
$$
は $\tau_{\leq n + a}$ 上同型を誘導し、任意の基底変換後にも同じことが成り立つ。
右導来関手 $Rf_*$ を適用すると、写像 $M_{n + 1} \to M_n$ は
$\tau_{\leq n + a}$ 上同型を誘導し、任意の基底変換後にも同じことが成り立つ。
区別三角
$$
M_{n + 1} \to M_n \to C_n \to M_{n + 1}[1]
$$
を選ぶ。$y \in |Y|$ とする。基本エタール近傍 $(U, u) \to (Y, y)$ を選ぶ。
これは Decent Spaces, Lemma
\ref{decent-spaces-lemma-decent-space-elementary-etale-neighbourhood} により可能である。
$Y'$ を $u$ における剰余体のスペクトルとする。引き戻すと
$C_n|_U \otimes_{\mathcal{O}_U}^\mathbf{L} \kappa(u)$ は次数 $\geq n + a$ にのみ
非零コホモロジーをもつ。More on Algebra, Lemma
\ref{more-algebra-lemma-lift-perfect-from-residue-field} により、完全複体 $C_n|_U$ は
$[n + a, m_n]$ に Tor 振幅をもつ。ここで $m_n$ はある整数であり、必要なら
$U$ を縮小している。したがって $C_n$ は $[n + a, m_n]$ に Tor 振幅をもつ。
ここで $m_n$ はある整数である
（$Y$ が準コンパクトだから）。特に双対完全複体 $C_n^\vee$ は
$[-m_n, -n - a]$ に Tor 振幅をもつ。

\medskip\noindent
$L_n = M_n^\vee$ を双対完全複体とする。前段落の議論から、
$L_n \to L_{n + 1}$ は $\tau_{\geq -n - a}$ 上同型を誘導する。したがって
$L = \text{hocolim} L_n$ は擬連接である。Derived Categories of Spaces, Lemma
\ref{spaces-perfect-lemma-pseudo-coherent-hocolim} を参照せよ。等式
$$
R\SheafHom(K, E) = R\SheafHom(\text{hocolim} K_n, E) =
R\lim R\SheafHom(K_n, E) = R\lim K_n^\vee \otimes_{\mathcal{O}_X} E
$$
が成り立ち（Cohomology on Sites, Lemma
\ref{sites-cohomology-lemma-colim-and-lim-of-duals}）、$R\lim$ は $Rf_*$ と可換なので、
$$
Rf_*R\SheafHom(K, E) = R\lim M_n = R\lim R\SheafHom(L_n, \mathcal{O}_Y) =
R\SheafHom(L, \mathcal{O}_Y)
$$
を得る。これで $Y$ 上の公式が証明された。$M_n$ の構成は基底変換と可換なので、
% P09-ERR-1230: frozen line 12735 has "base chance" where "base change" is intended.
この公式は任意の基底変換後にも成り立ち続ける。
\end{proof}

\begin{remark}
\label{spaces-more-morphisms-remark-compare-L}
読者は Lemma \ref{spaces-more-morphisms-lemma-compute-ext-rel-perfect} と Derived Categories of Spaces,
Lemma \ref{spaces-perfect-lemma-compute-ext} の類似に気付いたかもしれない。実際、
Lemma \ref{spaces-more-morphisms-lemma-compute-ext-rel-perfect} の擬連接複体 $L$ は、$Y$ 上の一意な
擬連接複体で、関手的同型
$$
\Ext^i_{\mathcal{O}_Y}(L, \mathcal{F}) \longrightarrow
\Ext^i_{\mathcal{O}_X}(K,
E \otimes_{\mathcal{O}_X}^\mathbf{L} Lf^*\mathcal{F})
$$
をもつものとして特徴付けられる。この同型は、$\mathcal{F}$ が
$\QCoh(\mathcal{O}_Y)$ を走るとき境界写像と可換である。これが必要になれば、ここで
精密な結果を定式化し、詳しい証明を与えることにする。
\end{remark}

\begin{lemma}
\label{spaces-more-morphisms-lemma-bounded-on-fibres}
$S$ をスキームとする。$X$ を $S$ 上の代数空間で、構造射 $f : X \to S$ が平坦かつ
局所有限表示であるものとする。$E$ を $D(\mathcal{O}_X)$ の擬連接対象とする。
次は同値である：
\begin{enumerate}
\item $E$ は $S$-完全である。
\item $E$ は局所有界下であり、各点 $s \in S$ に対して対象
$L(X_s \to X)^*E$（$D(\mathcal{O}_{X_s})$ のもの）は局所有界下である。
\end{enumerate}
\end{lemma}

\begin{proof}
すべて局所的なので、直ちに $X$ と $S$ がアフィンである場合へ帰着する。Lemma
\ref{spaces-more-morphisms-lemma-affine-locally-rel-perfect} を参照せよ。この場合は Derived Categories of
Schemes, Lemma \ref{perfect-lemma-bounded-on-fibres} で扱われている。
\end{proof}

\begin{lemma}
\label{spaces-more-morphisms-lemma-characterize-relatively-perfect}
$A$ を環とする。$X$ を $A$ 上分離的、有限表示、かつ平坦な代数空間とする。
$K \in D_\QCoh(\mathcal{O}_X)$ とする。
$R \Gamma (X, E \otimes^\mathbf{L} K)$ が $D(A)$ において完全であることが、
任意の完全な $E \in D(\mathcal{O}_X)$ について成り立つならば、$K$ は
$\Spec(A)$-完全である。
\end{lemma}

\begin{proof}
Lemma \ref{spaces-more-morphisms-lemma-characterize-pseudo-coh-improved} により、$K$ は $A$ に関して
擬連接である。Lemma \ref{spaces-more-morphisms-lemma-relative-pseudo-coherent-is-moot} により、$K$ は
$D(\mathcal{O}_X)$ において擬連接である。Derived Categories of Spaces, Proposition
\ref{spaces-perfect-proposition-detecting-bounded-below} により、$K$ は
$D^-(\mathcal{O}_X)$ に属する。$\mathfrak{p}$ を $A$ の素イデアルとし、
$i : Y \to X$ で $\mathfrak{p}$ 上のスキーム論的ファイバーの包含を表す。すなわち
$Y$ は $\kappa(\mathfrak p)$ 上のスキームである。Lemma
\ref{spaces-more-morphisms-lemma-bounded-on-fibres} により、$Li^*(K)$ が有界下であることを示せば十分である。
$G \in D_{perf} (\mathcal{O}_X)$ を $D_\QCoh (\mathcal{O}_X)$ を生成する完全複体とする。
Derived Categories of Spaces, Theorem
\ref{spaces-perfect-theorem-bondal-van-den-Bergh} を参照せよ。次が成り立つ：
\begin{align*}
R\Hom _{\mathcal{O}_Y}(Li^*(G), Li^*(K))
& =
R\Gamma(Y, Li^*(G ^\vee \otimes ^\mathbf{L} K)) \\
& =
R\Gamma(X, G^\vee \otimes ^{\mathbf{L}} K)
\otimes^\mathbf{L}_A \kappa(\mathfrak{p})
\end{align*}
第一の等式では $Li^*$ が完全対象と双対を保つこと、および Cohomology on Sites,
Lemma \ref{sites-cohomology-lemma-dual-perfect-complex} を用いる。いくつかの詳細は省く。
第二の等式は Derived Categories of Spaces, Lemma
\ref{spaces-perfect-lemma-compare-base-change} から従う。実際、$X$ は $A$ 上平坦である。
仮定により、これは $D(\kappa(\mathfrak{p}))$ の完全対象である。対象
$Li^*(G) \in D_{perf}(\mathcal{O}_Y)$ は $D_\QCoh(\mathcal{O}_Y)$ を生成する。Derived
Categories of Spaces, Remark \ref{spaces-perfect-remark-pullback-generator} を参照せよ。
したがって Derived Categories of Spaces, Proposition
\ref{spaces-perfect-proposition-detecting-bounded-below} により、$Li^*(K)$ は有界下であり、
証明が完了する。
\end{proof}












\section{立方体定理}
\label{spaces-more-morphisms-section-theorem-cube}

\noindent
本節は More on Morphisms, Section
\ref{more-morphisms-section-theorem-cube} の類似である。次の補題は、その仮定のもとで
Picard 関手の対角射が局所閉埋入によって表現可能であることを述べる。

\begin{lemma}
\label{spaces-more-morphisms-lemma-diagonal-picard-flat-proper}
$S$ をスキームとする。$f : X \to Y$ を $S$ 上の代数空間の平坦、固有、有限表示射と
する。$\mathcal{E}$ を有限局所自由 $\mathcal{O}_X$-加群とする。射 $g : Y' \to Y$ に
対して基底変換図式
$$
\xymatrix{
X' \ar[d]_{f'} \ar[r]_{g'} & X \ar[d]^f \\
Y' \ar[r]^g & Y
}
$$
を考える。写像 $\mathcal{O}_{Y'} \to f'_*\mathcal{O}_{X'}$ がすべての
$g : Y' \to Y$ に対して同型であると仮定する。このとき、有限表示の埋入
$j : Z \to Y$ が存在し、射 $g : Y' \to Y$ が $Z$ を経由するための必要十分条件は、
有限局所自由 $\mathcal{O}_{Y'}$-加群 $\mathcal{N}$ で
% P09-ERR-1231: frozen line 12872 uses undefined mathcal L although the lemma data define mathcal E and the proof alternates between the two symbols.
$(f')^*\mathcal{N} \cong (g')^*\mathcal{L}$ を満たすものが存在することである。
\end{lemma}

\begin{proof}
$y : \Spec(k) \to Y$ を体値点とする。$X_y$ を $f$ の $y$ におけるファイバーとする。
これは $H^0(X_y, \mathcal{O}_{X_y}) = k$ という仮定により連結である。したがって
$\mathcal{E}$ の階数はファイバー上一定である。$f$ は開（Morphisms of Spaces, Lemma
\ref{spaces-morphisms-lemma-fppf-open}）かつ閉なので、$Y = \coprod Y_r$ という
$Y$ の開かつ閉な部分空間への分解が存在し、$\mathcal{E}$ は逆像上で階数 $r$ が一定で
ある $Y_r$ をもつ。したがって $\mathcal{E}$ は一定階数 $r$ をもつと仮定してよい。
% P09-ERR-1232: frozen line 12885 omits "by" in the notation declaration for the dual.
$\mathcal{E}^\vee = \SheafHom(\mathcal{E}, \mathcal{O}_X)$ で階数 $r$ の双対加群を表す。

\medskip\noindent
コホモロジーと基底変換（より正確には Derived Categories of Spaces, Lemma
\ref{spaces-perfect-lemma-flat-proper-perfect-direct-image-general}）により、
$E = Rf_*\mathcal{E}$ は $Y$ の導来圏の完全対象であり、その構成は任意の基底変換と
可換である。同様のことが $E' = Rf_*\mathcal{E}^\vee$ にも成り立つ。次数 $< 0$ には
コホモロジーが決してないので、$E$ と $E'$ は（局所的に）
$[0, b]$ に Tor 振幅をもつ。ここで $b$ はある整数である。任意の射 $g : Y' \to Y$ に対して
$f'_*((g')^*\mathcal{E}) = H^0(Lg^*E)$ および
$f'_*((g')^*\mathcal{E}^\vee) = H^0(Lg^*E')$ であることに注意せよ。
$j : Z \to Y$ および $j' : Z' \to Y$ で、Derived Categories of Spaces, Lemma
\ref{spaces-perfect-lemma-locally-closed-where-H0-locally-free} において $E$ および $E'$、
$a = 0$ および $r = r$ に対して構成された局所閉埋入を表す。これらは
$H^0(Lj^*E)$ および $H^0((j')^*E')$ が引き戻しと可換な階数 $r$ の局所自由加群で
あるという性質によって特徴付けられる。

\medskip\noindent
$g : Y' \to Y$ を射とする。補題のような $\mathcal{N}$ が存在するならば、射影公式
Cohomology on Sites, Lemma \ref{sites-cohomology-lemma-projection-formula} により、
{\small
$$
f'_*((g')^*\mathcal{E}) \cong
f'_*((f')^*\mathcal{N}) \cong
\mathcal{N} \otimes_{\mathcal{O}_{Y'}} f'_*\mathcal{O}_{X'} \cong
\mathcal{N}\quad\text{and similarly }\quad
f'_*((g')^*\mathcal{E}^\vee) \cong \mathcal{N}^\vee
$$
}
が成り立つ。これらは階数 $r$ の局所自由加群であり、さらに階数 $r$ の局所自由加群
として任意のさらなる基底変換 $Y'' \to Y'$ 後も保たれる。したがってこの場合、射
$g : Y' \to Y$ は $j$ および $j'$ を経由する。よって $Y$ を
$Z \times_Y Z'$ で置き換え、$f_*\mathcal{E}$ と $f_*\mathcal{E}^\vee$ が
任意の基底変換と可換な局所自由 $\mathcal{O}_Y$-加群であり、階数 $r$ をもつと仮定してよい。

\medskip\noindent
この状況で、射 $g : Y' \to Y$ と補題のような $\mathcal{N}$ が存在するならば、
写像（次数 $0$ の cup 積）
$$
f'_*((g')^*\mathcal{E})
\otimes_{\mathcal{O}_{Y'}}
f'_*((g')^*\mathcal{E}^\vee)
\longrightarrow \mathcal{O}_{Y'}
$$
は完全ペアリングである。逆にこの cup 積写像が完全ペアリングならば、$Y'$ 上局所的に
切断の基底
$\sigma_1, \ldots, \sigma_r$ が $f'_*((g')^*\mathcal{L})$ に、切断の基底
$\tau_1, \ldots, \tau_r$ が $f'_*((g')^*\mathcal{E}^\vee)$ に存在し、それらの積は
$$
\sigma_i \tau_j = \delta_{ij}
$$
を満たす。$\sigma_i$ を $(g')^*\mathcal{L}$ の切断として $X'$ 上で、$\tau_j$ を
$(g')^*\mathcal{L}^\vee$ の切断として $X'$ 上でみなすと、
$$
\sigma_1, \ldots, \sigma_r :
\mathcal{O}_{X'}^{\oplus r}
\longrightarrow
(g')^*\mathcal{E}
$$
は、逆写像
$$
\tau_1, \ldots, \tau_r :
(g')^*\mathcal{E}
\longrightarrow
\mathcal{O}_{X'}^{\oplus r}
$$
をもつ同型である。言い換えれば、
$$
(f')^*f'_*(g')^*\mathcal{E} \cong (g')^*\mathcal{E}
$$
である。しかし cup 積が非退化であるという条件は、レトロコンパクト開部分スキーム
（すなわち適切な行列式が非零となる軌跡）を取り出す。これで証明が完了する。
\end{proof}

\section{複体の有限性性質の降下}
\label{spaces-more-morphisms-section-descent-finiteness}

\noindent
本節は More on Morphisms, Section
\ref{more-morphisms-section-descent-finiteness} および Derived Categories of Schemes,
Section \ref{perfect-section-descent-finiteness} の類似である。

\begin{lemma}
\label{spaces-more-morphisms-lemma-pseudo-coherent-descends-fpqc}
$S$ をスキームとする。$\{f_i : X_i \to X\}$ を $S$ 上の代数空間の fpqc 被覆とする。
$E \in D_\QCoh(\mathcal{O}_X)$ とし、$m \in \mathbf{Z}$ とする。このとき $E$ が
$m$-擬連接であるための必要十分条件は、各 $Lf_i^*E$ が $m$-擬連接であることである。
\end{lemma}

\begin{proof}
引き戻しは常に $m$-擬連接性を保つ。Cohomology on Sites, Lemma
\ref{sites-cohomology-lemma-pseudo-coherent-pullback} を参照せよ。したがって
$Lf_i^*E$ が $m$-擬連接であると仮定して、$E$ が $m$-擬連接であることを示せば十分である。
まずすべての $X_i$（各 $i$ について）がスキームであると仮定してよい。Topologies on Spaces, Lemma
\ref{spaces-topologies-lemma-refine-fpqc-schemes} を参照せよ。次に、スキームである
全射エタール射 $U \to X$ を選ぶ。このとき $U$ はスキームであり、
$U_i = U \times_X X_i$ はスキームであり、
$\{U_i \to U\}$ はスキームの fpqc 被覆となる。Topologies on Spaces, Lemma
\ref{spaces-topologies-lemma-recognize-fpqc-covering} を参照せよ。スキームの場合により、
$\{U_i \to U\}_{i \in I}$ について結果が成り立つことが分かる。Derived Categories of
Schemes, Lemma \ref{perfect-lemma-pseudo-coherent-descends-fpqc} を参照せよ。
一方、制限 $E|_U$ は $D_\QCoh(\mathcal{O}_U)$ の対象（$U$ のザリスキー位相と
「通常の」構造層を用いて定義したもの）から来る。Derived Categories of Spaces, Lemma
\ref{spaces-perfect-lemma-derived-quasi-coherent-small-etale-site} を参照せよ。
エタール位相とザリスキー位相による二つの擬連接性は Derived Categories of Spaces,
Lemma \ref{spaces-perfect-lemma-descend-pseudo-coherent} により一致するので、補題が従う。
\end{proof}

\begin{lemma}
\label{spaces-more-morphisms-lemma-tor-amplitude-descends-fppf}
$S$ をスキームとする。$\{g_i : Y_i \to Y\}$ を $S$ 上の代数空間の fpqc 被覆とする。
$f : X \to Y$ を代数空間の射とし、$X_i = Y_i \times_Y X$ とおく。射影を
$f_i : X_i \to Y_i$ および $g'_i : X_i \to X$ と表す。$E \in D_\QCoh(\mathcal{O}_X)$ とし、
$a, b \in \mathbf{Z}$ とする。このとき次は同値である。
\begin{enumerate}
\item $E$ は $[a, b]$ に Tor 振幅をもち、$D(f^{-1}\mathcal{O}_Y)$ の対象であり、
\item $L(g'_i)^*E$ は $[a, b]$ に Tor 振幅をもち、$D(f_i^{-1}\mathcal{O}_{Y_i})$ の
対象である（すべての $i$ について）。
\end{enumerate}
さらに、「$[a, b]$ に Tor 振幅をもつ」を「局所有限 Tor 次元をもつ」に置き換えても成り立つ。
\end{lemma}

\begin{proof}
引き戻しは Derived Categories of Spaces, Lemma
\ref{spaces-perfect-lemma-tor-independence-and-tor-amplitude} により「$[a, b]$ に
Tor 振幅をもつ」を保つ。$Y_i$ と $X$ は $Y$ 上 Tor 独立である。これは $Y_i \to Y$ が
平坦だからである。(2) を仮定して (1) を示そう。Tor 次元は茎で計算できる。
Cohomology on Sites, Lemma \ref{sites-cohomology-lemma-tor-amplitude-stalk} および
Properties of Spaces, Theorem \ref{spaces-properties-theorem-exactness-stalks} を参照せよ。
$\overline{x}$ を $X$ の幾何点とする。ある $i$ と幾何点を選ぶ。これを
$\overline{x}_i$ とし、$X_i$ にあり、その像 $\overline{x}$ は $X$ にある。このとき
$$
(L(g_i')^*E)_{\overline{x}_i} =
E_{\overline{x}}
\otimes_{\mathcal{O}_{X, \overline{x}}}^\mathbf{L}
\mathcal{O}_{X_, \overline{x}_i}
$$
である。$\overline{y}_i$ を $Y_i$ の点、$\overline{y}$ を $Y$ の点とし、それぞれ
$\overline{x}_i$ および $\overline{x}$ の像とする。$X$ と $Y_i$ は $Y$ 上 Tor 独立なので、More on Algebra, Lemma
\ref{more-algebra-lemma-base-change-comparison} により、表示した右辺は
$$
E_{\overline{x}}
\otimes_{\mathcal{O}_{Y, \overline{y}}}^\mathbf{L}
\mathcal{O}_{Y_i, \overline{y}_i}
$$
$D(\mathcal{O}_{Y_i, \overline{y}_i})$ においてこれと等しい。これの Tor 振幅が
$[a, b]$ にあると仮定したので、More on Algebra, Lemma
\ref{more-algebra-lemma-flat-descent-tor-amplitude} により、$E_{\overline{x}}$ の
$D(\mathcal{O}_{Y, \overline{y}})$ における Tor 振幅も $[a, b]$ にある。したがって (1) が従う。

\medskip\noindent
初等的な位相の議論を用いれば、「局所有限 Tor 次元」の場合も同様に従う。
\end{proof}

\noindent
以下の補題は実際には本節に属するものではない。

\begin{lemma}
\label{spaces-more-morphisms-lemma-thickening-pseudo-coherent}
$S$ をスキームとする。$i : X \to X'$ を代数空間の有限位数肥厚とする。
$K' \in D(\mathcal{O}_{X'})$ を対象とし、$K = Li^*K'$ が擬連接であると仮定する。
このとき $K'$ は擬連接である。
\end{lemma}

\begin{proof}
まず $K'$ が準連接コホモロジー層をもつことを示す。この部分は読み飛ばしてよい。
第一位肥厚の場合へ帰着してよい。Section \ref{spaces-more-morphisms-section-thickenings} を参照せよ。
$\mathcal{I} \subset \mathcal{O}_{X'}$ を $X$ を切り出す準連接イデアル層とする。
短完全列
$$
0 \to \mathcal{I} \to \mathcal{O}_{X'} \to i_*\mathcal{O}_X \to 0
$$
を $K'$ とテンソルすると、区別三角
$$
K' \otimes_{\mathcal{O}_{X'}}^\mathbf{L} \mathcal{I}
\to K' \to
K' \otimes_{\mathcal{O}_{X'}}^\mathbf{L} i_*\mathcal{O}_X
\to
(K' \otimes_{\mathcal{O}_{X'}}^\mathbf{L} \mathcal{I})[1]
$$
を得る。$i_* = Ri_*$ であり、$\mathcal{I}$ を準連接 $\mathcal{O}_X$-加群とみなせる
ので（第一位肥厚だから）、これは
$$
i_*(K \otimes_{\mathcal{O}_X}^\mathbf{L} \mathcal{I})
\to K' \to
i_*K \to
i_*(K \otimes_{\mathcal{O}_X}^\mathbf{L} \mathcal{I})[1]
$$
と書き換えられる。各項の同定には Cohomology of Spaces, Lemma
\ref{spaces-cohomology-lemma-projection-formula-finite} を用いよ。$K$ は
$D_\QCoh(\mathcal{O}_X)$ に属するので、$K'$ は $D_\QCoh(\mathcal{O}_{X'})$ に属する。
ここでは Derived Categories of Spaces, Lemmas
\ref{spaces-perfect-lemma-pseudo-coherent},
\ref{spaces-perfect-lemma-quasi-coherence-tensor-product}, および
\ref{spaces-perfect-lemma-quasi-coherence-direct-image} を用いた。

\medskip\noindent
$K'$ が $D_\QCoh(\mathcal{O}_{X'})$ に属すると仮定する。この問題は $X'$ 上エタール局所的
なので、$X'$ はアフィンと仮定してよい。この場合はスキームの場合、More on Morphisms,
Lemma \ref{more-morphisms-lemma-thickening-pseudo-coherent} から従う。スキームの言葉への
移行には Derived Categories of Spaces, Lemmas
\ref{spaces-perfect-lemma-derived-quasi-coherent-small-etale-site} および
\ref{spaces-perfect-lemma-descend-pseudo-coherent}、ならびに Remark
\ref{spaces-perfect-remark-match-total-direct-images} を用いる。
\end{proof}

\begin{lemma}
\label{spaces-more-morphisms-lemma-thickening-relatively-perfect}
$S$ をスキームとする。代数空間の Cartesian 図式
$$
\xymatrix{
X \ar[r]_i \ar[d]_f & X' \ar[d]^{f'} \\
Y \ar[r]^j & Y'
}
$$
を $S$ 上で考える。$X' \to Y'$ が平坦かつ局所有限表示であり、$Y \to Y'$ が有限位数
肥厚であると仮定する。$E' \in D(\mathcal{O}_{X'})$ とする。$E = Li^*(E')$ が
$Y$-完全ならば、$E'$ は $Y'$-完全である。
\end{lemma}

\begin{proof}
$Y$-完全性について、$E$ は擬連接であり、$E$ は $f^{-1}\mathcal{O}_Y$-加群の複体として
局所有限 Tor 次元をもつことであることを思い出そう
（Definition \ref{spaces-more-morphisms-definition-relatively-perfect}）。Lemma
\ref{spaces-more-morphisms-lemma-thickening-pseudo-coherent} により $E'$ は擬連接である。特に
$E'$ は $D_\QCoh(\mathcal{O}_{X'})$ に属する。Derived Categories of Spaces, Lemma
\ref{spaces-perfect-lemma-pseudo-coherent} を参照せよ。Lemma
\ref{spaces-more-morphisms-lemma-affine-locally-rel-perfect} により、これはスキームの場合へ帰着する。
スキームの場合は More on Morphisms, Lemma
\ref{more-morphisms-lemma-thickening-relatively-perfect} である。
\end{proof}

\begin{lemma}
\label{spaces-more-morphisms-lemma-henselian-relatively-perfect}
$(R, I)$ を環とイデアルからなる組とし、$I$ は Jacobson 根基に含まれるとする。
$S = \Spec(R)$ および $S_0 = \Spec(R/I)$ とおく。$X$ を $R$ 上の代数空間とし、
その構造射 $f : X \to S$ は固有、平坦、有限表示であるとする。$X_0 = S_0 \times_S X$
とおく。$E \in D(\mathcal{O}_X)$ を擬連接とする。$E_0$ を $E$ の導来制限とし、
$X_0$ 上で $S_0$-完全ならば、$E$ は $S$-完全である。
\end{lemma}

\begin{proof}
$U \to X$ を全射エタール射で $U$ がアフィンとなるものとして選ぶ。
閉埋入 $U \to \mathbf{A}^d_S$ を選び、$U_0 = S_0 \times_S U$ とおく。
$E_0|_{U_0}$ は Tor 振幅が $[a, b]$ にあり、ここで $a, b \in \mathbf{Z}$ とする。
$\overline{x}$ を $X$ の幾何点とする。$E_{\overline{x}}$ の Tor 振幅が $R$ 上
$[a - d, b]$ にあることを示す。これで証明は完了する。実際、Tor 振幅は茎から読み取れる。
Cohomology on Sites, Lemma \ref{sites-cohomology-lemma-tor-amplitude-stalk} および
Properties of Spaces, Theorem \ref{spaces-properties-theorem-exactness-stalks} を参照せよ。

\medskip\noindent
$x \in |X|$ を $\overline{x}$ が定める点とする。$|X| \to |S|$ は閉写像であることを
思い出そう（固有射の定義による）。$I$ は Jacobson 根基に含まれるので、$S$ の空でない
任意の閉集合は閉部分スキーム $S_0$ の点を含む。したがって $x \leadsto x_0$ という
$|X|$ における特殊化で、$x_0 \in |X_0|$ となるものを見つけられる。
$u_0 \in U_0$ を選び、$x_0$ へ写るものとする。Decent Spaces, Lemma
\ref{decent-spaces-lemma-generalizations-lift-flat}（またはエタール射に直接適用できる
Decent Spaces, Lemma \ref{decent-spaces-lemma-specialization}）により、$u \leadsto u_0$ という
$U$ における特殊化を見つける。この $u$ は $x$ へ写る。$\overline{x}$ を幾何点
$\overline{u}$ へ持ち上げることができ、これは $U$ の点で $u$ 上にある。このとき
$E_{\overline{x}} = (E|_U)_{\overline{u}}$ である。

\medskip\noindent
$U = \Spec(A)$ と書く。このとき $A$ は平坦な有限表示 $R$-代数であり、$R$-代数の
$d$ 変数多項式環の商である。制限 $E|_U$ は、Derived Categories of Spaces, Lemmas
\ref{spaces-perfect-lemma-pseudo-coherent},
\ref{spaces-perfect-lemma-derived-quasi-coherent-small-etale-site}, および
\ref{spaces-perfect-lemma-descend-pseudo-coherent}、ならびに Derived Categories of
Schemes, Lemma \ref{perfect-lemma-affine-compare-bounded} および
\ref{perfect-lemma-pseudo-coherent-affine} により、$K$（$D(A)$ の擬連接対象）に対応する。
$E_0$ は $K \otimes_A^\mathbf{L} A/IA$ に対応することに注意せよ。
$\mathfrak q \subset \mathfrak q_0 \subset A$ を、$u \leadsto u_0$ に対応する素イデアルと
する。すると
$$
E_{\overline{x}} =
(E|_U)_{\overline{u}} =
E_u \otimes_{\mathcal{O}_{U, u}}^\mathbf{L} \mathcal{O}_{U, \overline{u}} =
K_{\mathfrak q} \otimes_{A_\mathfrak q}^\mathbf{L} A_{\mathfrak q}^{sh}
$$
である（いくつかの詳細は省略する）。$A_\mathfrak q \to A_\mathfrak q^{sh}$ は平坦なので、
これの $R$-加群としての Tor 振幅は、$K_\mathfrak q$ の $R$-加群としての Tor 振幅と
同じである。More on Algebra, Lemma
\ref{more-algebra-lemma-no-change-tor-amplitude} を参照せよ。また、$K_{\mathfrak q}$ は
$K_{\mathfrak q_0}$ の局所化である。したがって、$K_{\mathfrak q_0}$ が $[a - d, b]$ に
Tor 振幅をもつ $R$-加群の複体であることを示せば十分である。

\medskip\noindent
$I \subset \mathfrak p_0 \subset R$ を $f(x_0)$ に対応する素イデアルとする。
このとき
\begin{align*}
K \otimes_R^\mathbf{L} \kappa(\mathfrak p_0)
& =
(K \otimes_R^\mathbf{L} R/I) \otimes_{R/I}^\mathbf{L}
\kappa(\mathfrak p_0) \\
& =
(K \otimes_A^\mathbf{L} A/IA) \otimes_{R/I}^\mathbf{L}
\kappa(\mathfrak p_0)
\end{align*}
である。第二の等式は $R \to A$ が平坦だからである。$a, b$ の選択により、この複体は
区間 $[a, b]$ 内の次数にのみコホモロジーをもつ。最後に More on Algebra, Lemma
\ref{more-algebra-lemma-lift-from-fibre-relatively-perfect} を
$R \to A$、$\mathfrak q_0$、$\mathfrak p_0$、$K$ に適用して結論する。
\end{proof}




\section{節点曲線の族}
\label{spaces-more-morphisms-section-families-nodal}

\noindent
この節は Algebraic Curves, Section \ref{curves-section-families-nodal} の続きである。
用語の選択については、その節も参照せよ。

\medskip\noindent
スキームの射が「相対次元 $1$ の高々節点的」であるという性質は、始域および終域上
エタール局所的である。Descent, Lemma
\ref{descent-lemma-etale-local-source-target}、および
Algebraic Curves, Lemmas
\ref{curves-lemma-nodal-family-etale-local-source}、
\ref{curves-lemma-nodal-family-fpqc-local-target}、および
\ref{curves-lemma-nodal-family-postcompose-etale} を参照せよ。
この性質は基底変換で安定であり、終域上 fpqc 局所的でもある。Algebraic Curves,
Lemmas \ref{curves-lemma-base-change-nodal-family} および
\ref{curves-lemma-nodal-family-fpqc-local-target} を参照せよ。したがって、
Morphisms of Spaces, Lemma \ref{spaces-morphisms-lemma-local-source-target} により、
代数空間について相対次元 $1$ の高々節点的射の概念を次のように定義できる。
射が表現可能であるとき、これは Morphisms of Spaces, Section
\ref{spaces-morphisms-section-representable} で既に定義された概念と一致する。

\begin{definition}
\label{spaces-more-morphisms-definition-nodal-family}
$S$ をスキームとする。$f : X \to Y$ を $S$ 上の代数空間の射とする。
$f$ が{\it 相対次元 $1$ の高々節点的}であるとは、
Morphisms of Spaces, Lemma
\ref{spaces-morphisms-lemma-local-source-target} の同値な条件が
$\mathcal{P} =$「相対次元 $1$ の高々節点的」として成り立つことをいう。
\end{definition}

\begin{lemma}
\label{spaces-more-morphisms-lemma-base-change-nodal}
相対次元 $1$ の高々節点的である性質は基底変換で保たれる。
\end{lemma}

\begin{proof}
Morphisms of Spaces, Remark \ref{spaces-morphisms-remark-base-change-P} および
Algebraic Curves, Lemma \ref{curves-lemma-base-change-nodal-family} を参照せよ。
\end{proof}

\begin{lemma}
\label{spaces-more-morphisms-lemma-nodal-local}
$S$ をスキームとする。$f : X \to Y$ を $S$ 上の代数空間の射とする。
次の条件は同値である。
\begin{enumerate}
\item $f$ は相対次元 $1$ の高々節点的である。
\item 任意のスキーム $Z$ と任意の射 $Z \to Y$ に対して、射
$Z \times_Y X \to Z$ は相対次元 $1$ の高々節点的である。
\item 任意のアフィンスキーム $Z$ と任意の射 $Z \to Y$ に対して、射
$Z \times_Y X \to Z$ は相対次元 $1$ の高々節点的である。
\item スキーム $V$ と全射エタール射 $V \to Y$ が存在し、射
$V \times_Y X \to V$ は相対次元 $1$ の高々節点的である。
\item スキーム $U$ と全射エタール射 $\varphi : U \to X$ が存在し、合成射
$f \circ \varphi$ は相対次元 $1$ の高々節点的である。
\item 次の任意の可換図式
$$
\xymatrix{
U \ar[d] \ar[r] & V \ar[d] \\
X \ar[r] & Y
}
$$
において、$U$, $V$ はスキームで、鉛直射はエタールならば、上の水平射は
相対次元 $1$ の高々節点的である。
\item 次の可換図式
$$
\xymatrix{
U \ar[d] \ar[r] & V \ar[d] \\
X \ar[r] & Y
}
$$
で、$U$, $V$ はスキーム、鉛直射はエタール、$U \to X$ は全射であり、上の水平射が
相対次元 $1$ の高々節点的となるものが存在する。
\item ザリスキー被覆 $Y = \bigcup_{i \in I} Y_i$ と
$f^{-1}(Y_i) = \bigcup X_{ij}$ が存在し、各射 $X_{ij} \to Y_i$ が
相対次元 $1$ の高々節点的である。
\end{enumerate}
\end{lemma}

\begin{proof}
省略する。
\end{proof}

\noindent
次の補題は、射が相対次元 $1$ の高々節点的であるかどうかをファイバー上で判定できる
ことを示す。

\begin{lemma}
\label{spaces-more-morphisms-lemma-locus-where-nodal}
$S$ をスキームとする。$f : X \to Y$ を $S$ 上の代数空間の射で、平坦かつ局所有限表示
とする。このとき、$X' \subset X$ という最大の開部分空間が存在し、
$f|_{X'} : X' \to Y$ は相対次元 $1$ の高々節点的である。さらに、$X'$ の形成は
任意の基底変換と可換である。
\end{lemma}

\begin{proof}
次の可換図式を選ぶ。
$$
\xymatrix{
U \ar[d] \ar[r]_h & V \ar[d] \\
X \ar[r]^f & Y
}
$$
ここで $U$, $V$ はスキーム、鉛直射はエタール、$U \to X$ は全射とする。
スキームの場合の補題（Algebraic Curves, Lemma
\ref{curves-lemma-locus-where-nodal}）により、$U' \subset U$ という最大の開部分スキームが
存在し、$h|_{U'} : U' \to V$ は相対次元 $1$ の高々節点的であり、$U'$ の形成は
基底変換と可換である。
$X' \subset X$ を、点が開集合 $\Im(|U'| \to |X|)$ に対応する開部分空間とする。
Lemma \ref{spaces-more-morphisms-lemma-nodal-local} により、$X' \to Y$ は相対次元 $1$ の高々節点的であり、
$X'$ はこの性質をもつ最大の開部分空間である（このことから $U'$ が $X'$ の $U$ における
逆像であることも従うが、ここでは必要ない）。基底変換後も同じことが成り立つので、
証明が完了する。
\end{proof}




\section{分解性質}
\label{spaces-more-morphisms-section-resolution-property}

\noindent
Derived Categories of Spaces, Section
\ref{spaces-perfect-section-resolution-property} の議論を続ける。

\begin{situation}
\label{spaces-more-morphisms-situation-descend-resolution-property}
$S$ をスキームとする。$X$ を $S$ 上の準コンパクトかつ準分離的な代数空間とする。
$V \to X$ を全射エタール射とし、$V$ はアフィンスキームとする（このようなものは
Properties of Spaces, Lemma
\ref{spaces-properties-lemma-quasi-compact-affine-cover} により存在する）。
次の可換図式を選ぶ。
$$
\xymatrix{
V \ar[rd]_\varphi \ar[rr]_j & & Y \ar[ld]^\pi \\
& X
}
$$
ここで $j$ は開埋入、$\pi$ は代数空間の有限射である（このような図式は
Lemma \ref{spaces-more-morphisms-lemma-quasi-finite-separated-pass-through-finite} により存在する）。
$\mathcal{I} \subset \mathcal{O}_Y$ を $Y$ 上の有限型準連接イデアル層で、
$V(\mathcal{I}) = Y \setminus j(V)$ を満たすものとする（このようなイデアル層は
Limits of Spaces, Lemma
\ref{spaces-limits-lemma-quasi-coherent-finite-type-ideals} により存在する）。
\end{situation}

\begin{lemma}
\label{spaces-more-morphisms-lemma-surjection-in-Noetherian-case}
Situation \ref{spaces-more-morphisms-situation-descend-resolution-property} において $X$ がネーターであると仮定する。
任意の連接 $\mathcal{O}_X$-加群 $\mathcal{F}$ に対し、$r \geq 0$、
$n_1, \ldots, n_r \geq 0$ なる整数、および全射
$$
\bigoplus\nolimits_{i = 1, \ldots, r} \pi_*(\mathcal{I}^{n_i})
\longrightarrow \mathcal{F}
$$
という $\mathcal{O}_X$-加群の射が存在する。
\end{lemma}

\begin{proof}
$\omega_{Y/X}$ を、次の同型をもつ連接 $\mathcal{O}_Y$-加群とする。
$$
\pi_*\omega_{Y/X} \cong
\SheafHom_{\mathcal{O}_X}(\pi_*\mathcal{O}_Y, \mathcal{O}_X)
$$
これは $\pi_*\mathcal{O}_Y$-加群の同型である。Morphisms of Spaces, Lemma
\ref{spaces-morphisms-lemma-affine-equivalence-modules} および
Descent on Spaces, Lemma
\ref{spaces-descent-lemma-finite-over-finite-module} を参照せよ。
標準射 $\mathcal{O}_X \to \pi_*\mathcal{O}_Y$ は、標準射
$$
\text{Tr}_\pi : \pi_*\omega_{Y/X} \longrightarrow \mathcal{O}_X
$$
を誘導する。
$V$ はネーターアフィンなので、次の切断を選べる。
$$
s_1, \ldots, s_r \in
\Gamma(V, \pi^*\mathcal{F} \otimes_{\mathcal{O}_Y} \omega_{Y/X})
$$
これらは連接加群
$\pi^*\mathcal{F} \otimes_{\mathcal{O}_X} \omega_{Y/X}$ を $V$ 上で生成する。
Cohomology of Spaces, Lemma
\ref{spaces-cohomology-lemma-homs-over-open} により、$n_i \geq 0$ なる整数を選び、
$s_i$ が次の射へ拡張するようにできる。
$s_i' : \mathcal{I}^{n_i} \to
\pi^*\mathcal{F} \otimes_{\mathcal{O}_Y} \omega_{Y/X}$.
$X$ へ押し出すと、次の射を得る。
$$
\sigma_i :
\pi_*\mathcal{I}^{n_i}
\xrightarrow{\pi_*s'_i}
\pi_*(\pi^*\mathcal{F} \otimes_{\mathcal{O}_Y} \omega_{Y/X}) =
\mathcal{F} \otimes_{\mathcal{O}_X} \pi_*\omega_{Y/X}
\xrightarrow{\text{Tr}_\pi} \mathcal{F}
$$
ここで等号は Cohomology of Spaces, Lemma
\ref{spaces-cohomology-lemma-finite-rings} による。
証明を終えるため、これらの射の和が全射であることを示す。

\medskip\noindent
$x \in |X|$ を $X$ の点とする。$v \in |V|$ を $x$ へ写る点とする。
次を満たすエタール近傍 $(U, u) \to (X, x)$ を選べる。
$$
U \times_X Y = W \coprod W'
$$
（代数空間の互いに素な和）であり、$W \to U$ は同型で、
その唯一の点 $w \in W$ は $u$ 上にあり、$v$ へ $V \subset Y$ 内で写る。
これを見るには Lemma \ref{spaces-more-morphisms-lemma-etale-splits-off-just-one-quasi-finite-part} および
\ref{etale-lemma-etale-etale-local} を用いればよい。
必要ならさらに $U$ を縮小し、射影により $W$ が $V \subset Y$ に写ると仮定する。
$\omega_{Y/X}$ の形成はエタール局所化と可換なので、
$$
\pi_*\omega_{Y/X}|_U = (\pi|_W)_*\omega_{W/U} \oplus (\pi|_{W'})_*\omega_{W'/U}
$$
を得る。
$(\pi|_W)_*\omega_{W/U} = \mathcal{O}_U$ であり、この同型は
上の分解の第一成分に制限したトレース射 $\text{Tr}_\pi|_U$ によって与えられる。
$W$ が $V$ に写るので、$\mathcal{I}^{n_i}|_W = \mathcal{O}_W$ である。
したがって
$$
\pi_*(\mathcal{I}^{n_i})|_U = \mathcal{O}_U \oplus
(W' \to U)_*(\mathcal{I}^{n_i}|_{W'})
$$
図式を追えば（詳細は省略）、$\sigma_i|_U$ の
$\mathcal{O}_U$ 成分は、$\mathcal{O}_U \to \mathcal{F}$ という射であり、これは次の切断に対応する。
$$
s_i|_W \in
\Gamma(W,\pi^*\mathcal{F} \otimes_{\mathcal{O}_Y} \omega_{Y/X}) =
\Gamma(W, \mathcal{F}|_W \otimes_{\mathcal{O}_W} \omega_{W/U}) =
\Gamma(U, \mathcal{F})
$$
$s_i$ は $\pi^*\mathcal{F} \otimes_{\mathcal{O}_Y} \omega_{Y/X}$ という加群を $V$ 上で
生成し、$W$ は $V$ に写るので、
$\bigoplus \sigma_i$ の $U$ への制限は全射であると結論する。
$x$ は任意の点だったので、証明が完了する。
\end{proof}

\begin{lemma}
\label{spaces-more-morphisms-lemma-resolution-property-in-Noetherian-case}
Situation \ref{spaces-more-morphisms-situation-descend-resolution-property} において
$X$ がネーターであると仮定する。$X$ が分解性質をもつことと、
$\pi_*\mathcal{I}$ が有限局所自由な
$\mathcal{O}_X$-加群の商であることは同値である。
\end{lemma}

\begin{proof}
$\pi_*\mathcal{I}$ は Cohomology of Spaces, Lemma
\ref{spaces-cohomology-lemma-finite-pushforward-coherent} により連接加群である。
したがって $X$ が分解性質をもてば、$\pi_*\mathcal{I}$ は有限局所自由な
$\mathcal{O}_X$-加群の商である。逆に、$\mathcal{E} \to \pi_*\mathcal{I}$ という全射が、
有限局所自由な $\mathcal{O}_X$-加群 $\mathcal{E}$ に対して与えられているとする。
すべての $n \geq 1$ に対して全射
$$
\pi_*\mathcal{I} \otimes_{\mathcal{O}_X} \pi_*\mathcal{I}^n
\longrightarrow
\pi_*\mathcal{I}^{n + 1}
$$
が存在することに注意する。したがって $\mathcal{E}^{\otimes n}$ は
$\pi_*\mathcal{I}^n$ へ、すべての $n \geq 1$ に対して全射する。
Lemma \ref{spaces-more-morphisms-lemma-surjection-in-Noetherian-case} の結果と組み合わせれば、
$X$ が分解性質をもつことが従う。
\end{proof}

\begin{lemma}
\label{spaces-more-morphisms-lemma-resolution-property-one-sheaf}
Situation \ref{spaces-more-morphisms-situation-descend-resolution-property} において、
代数空間 $X$ が分解性質をもつことと、$\pi_*\mathcal{I}$ が
有限局所自由な $\mathcal{O}_X$-加群の商であることは同値である。
\end{lemma}

\begin{proof}
$\pi_*\mathcal{G}$（有限型準連接 $\mathcal{O}_Y$-加群 $\mathcal{G}$ の押し出し）は、Descent on Spaces, Lemma
\ref{spaces-descent-lemma-finite-over-finite-module} により有限型準連接
$\mathcal{O}_X$-加群である。特に、$X$ が分解性質をもつならば、
$\pi_*\mathcal{I}$ は Derived Categories of Spaces, Definition
\ref{spaces-perfect-definition-resolution-property} により有限局所自由な
$\mathcal{O}_X$-加群の商である。

\medskip\noindent
全射 $\mathcal{E} \to \pi_*\mathcal{I}$ が、有限局所自由な
$\mathcal{O}_X$-加群 $\mathcal{E}$ に対して与えられていると仮定する。
証明の残りでは、Lemma \ref{spaces-more-morphisms-lemma-resolution-property-in-Noetherian-case} で扱った
ネーターの場合へ帰着することにより、$X$ が分解性質をもつことを示す。
証明の残りは読み飛ばしてよい。

\medskip\noindent
第一の帰着として、$X$ を $\Spec(\mathbf{Z})$ 上の代数空間とみなしてよい。
これは Spaces, Definition \ref{spaces-definition-base-change} を参照せよ。
（この変更は補題の条件も結論も変えない。）

\medskip\noindent
Limits of Spaces, Lemma
\ref{spaces-limits-lemma-finite-in-finite-and-finite-presentation} により、
$Y = \lim Y_i$ と書ける。ここで各 $Y_i$ は $X$ 上有限かつ有限表示で、
遷移射は閉埋入である。$Z = V(\mathcal{I})$ という $Y$ の閉部分空間を考える。
$\mathcal{I}$ は有限型なので、射 $Z \to Y$ は有限表示である。したがって、ある $i$ と
有限表示射 $Z_i \to Y_i$ で、その $Y$ への基底変換が $Z \to Y$ となるものを見つけられる。
Limits of Spaces, Lemma
\ref{spaces-limits-lemma-descend-finite-presentation} を参照せよ。
$i' \geq i$ に対して $Z_{i'} = Z_i \times_{Y_i} Y_{i'}$ とおく。
$i$ を増やせば、$Z_i \to Y_i$ が（有限表示の）閉埋入であると仮定できる。
Limits of Spaces, Lemma
\ref{spaces-limits-lemma-descend-closed-immersion} を参照せよ。
$\mathcal{I}_i \subset \mathcal{O}_{Y_i}$ を $Z_i$ のイデアル層とし、
$\pi_i : Y_i \to X$ を構造射とする。
$i' \geq i$ についても同様にする。$Z = \lim_{i' \geq i} Z_{i'}$ なので、
$$
\pi_*\mathcal{I} = \colim \pi_{i', *}\mathcal{I}_{i'}
$$
を得る。
この系の遷移射はすべて全射である。これは射
$$
\pi_{i, *}\mathcal{O}_{Y_i} \to \pi_{i', *}\mathcal{O}_{Y_{i'}}
$$
が全射であることと、$Z_{i'} = Z_i \times_{Y_i} Y_{i'}$ であることから分かる。
Cohomology of Spaces, Lemma
\ref{spaces-cohomology-lemma-finite-presentation-quasi-compact-colimit} により、
ある $i' \geq i$ について、射 $\mathcal{E} \to \pi_*\mathcal{I}$ は
射 $\mathcal{E} \to \pi_{i', *}\mathcal{I}_{i'}$ へ持ち上がる。
$i'$ を増やせば、この射
$$
\mathcal{E} \to \pi_{i', *}\mathcal{I}_{i'}
$$
は全射となる（そうでなければ、全射な遷移射をもつ余核の余極限が零でないことに反する）。
これで、次の段落で扱う場合へ帰着される。

\medskip\noindent
$X$ を $\mathbf{Z}$ 上の代数空間とし、$Y \to X$ は有限表示とする。
絶対ネーター近似により、$X = \lim X_i$ と有向極限として書ける。
各 $X_i$ は $\mathbf{Z}$ 上有限型の準分離代数空間であり、遷移射はアフィンである。
Limits of Spaces, Proposition \ref{spaces-limits-proposition-approximate} を参照せよ。
$\pi : Y \to X$ は有限表示なので、ある $i$ と有限表示射
$\pi_i : Y_i \to X_i$ で、その $X$ への基底変換が $\pi$ となるものを見つけられる。
Limits of Spaces, Lemma
\ref{spaces-limits-lemma-descend-finite-presentation} を参照せよ。
$i$ を増やせば $\pi_i$ は有限であると仮定できる。Limits of Spaces, Lemma
\ref{spaces-limits-lemma-descend-finite} を参照せよ。
次に、有限局所自由な $\mathcal{O}_{X_i}$-加群 $\mathcal{E}_i$ で、その
$X$ への引き戻しが $\mathcal{E}$ となるものが存在すると仮定できる。
Limits of Spaces, Lemma \ref{spaces-limits-lemma-descend-invertible-modules} を参照せよ。
さらに、射 $\mathcal{E}_i \to \pi_{i, *}\mathcal{O}_{Y_i}$ で、その
$X$ への引き戻しが合成
$\mathcal{E} \to \pi_*\mathcal{I} \to \pi_*\mathcal{O}_Y$ となるものが存在すると仮定できる。
Limits of Spaces, Lemma
\ref{spaces-limits-lemma-descend-modules-finite-presentation} を参照せよ。
余核
$$
\mathcal{E}_i \to \pi_{i, *}\mathcal{O}_{Y_i} \to \mathcal{Q}_i \to 0
$$
は連接 $\mathcal{O}_{Y_i}$-加群であり、その $X$ への引き戻しは（有限表示な）余核
$\mathcal{Q}$、すなわち射 $\mathcal{E} \to \pi_*\mathcal{O}_Y$ の余核である。
すなわち、$\mathcal{Q} = \pi_*(\mathcal{O}_Y/\mathcal{I})$ である。
次の射を考える。
$$
\mathcal{E}_i
\otimes_{\mathcal{O}_{X_i}}
\pi_{i, *} \mathcal{O}_{Y_i}
\longrightarrow
\pi_{i, *}\mathcal{O}_{Y_i}
\otimes_{\mathcal{O}_{X_i}}
\pi_{i, *} \mathcal{O}_{Y_i}
\to
\pi_{i, *} \mathcal{O}_{Y_i}
\to
\mathcal{Q}_i
$$
ここで第二の射は $\pi_{i, *}\mathcal{O}_{Y_i}$ 上の代数構造で与えられる。
この射を $Y$ へ引き戻すと零になる。なぜなら
$\mathcal{E} \to \pi_*\mathcal{O}_Y$ の像はイデアル $\pi_*\mathcal{I}$ だからである。
したがって Limits of Spaces, Lemma
\ref{spaces-limits-lemma-descend-modules-finite-presentation} により、
$i$ を増やせば、表示した合成は零であると仮定できる。
これはまさに、$\mathcal{E}_i \to \pi_{i, *}\mathcal{O}_{Y_i}$ の像が
$\pi_{i, *}\mathcal{I}_i$ の形であることを意味する。ここで
$\mathcal{I}_i \subset \mathcal{O}_{Y_i}$ はある連接イデアル層である。
射 $\mathcal{E}_i \to \pi_{i, *}\mathcal{O}_{Y_i}$ は
$\mathcal{E} \to \pi_*\mathcal{O}_Y$ へ引き戻されるので、
$\mathcal{I}_i$ の $Y$ への引き戻しは $\mathcal{I}$ を生成する。
$V_i \subset Y_i$ を、その補集合が $V(\mathcal{I}_i) \subset Y_i$ である開部分空間とする。
上の議論により、$V$ は $V_i$ の逆像である。
$i$ を増やせば、$V_i$ はアフィンであり、
$\pi_i|_{V_i} : V_i \to X_i$ はエタールであると仮定できる。Limits of Spaces, Lemmas
\ref{spaces-limits-lemma-limit-is-affine} および
\ref{spaces-limits-lemma-descend-etale} を参照せよ。
以上により、Lemma \ref{spaces-more-morphisms-lemma-resolution-property-in-Noetherian-case} を適用して
$X_i$ が分解性質をもつことが分かる。$X \to X_i$ はアフィンなので、
Derived Categories of Spaces, Lemma
\ref{spaces-perfect-lemma-resolution-property-goes-up-affine} により $X$ も分解性質をもつ。
\end{proof}

\begin{lemma}
\label{spaces-more-morphisms-lemma-resolution-property-descends}
$S$ をスキームとする。$X = \lim X_i$ を、アフィンな遷移射をもつ
$S$ 上の準コンパクトかつ準分離的な代数空間の順系の極限とする。
$X$ が分解性質をもつことと、$X_i$ が分解性質をもつような $i$ が存在することは同値である。
\end{lemma}

\begin{proof}
$X_i$ が分解性質をもてば、Derived Categories of Spaces, Lemma
\ref{spaces-perfect-lemma-resolution-property-goes-up-affine} により $X$ もそうである。
$X$ が分解性質をもつと仮定する。
$i \in I$ を選ぶ。アフィンスキーム $V_i$ と全射エタール射
$V_i \to X_i$ を選べる（Properties of Spaces, Lemma
\ref{spaces-properties-lemma-quasi-compact-affine-cover}）。
埋入 $j : V_i \to Y_i$ で、$Y_i$ が $X_i$ 上有限かつ有限表示となるものを選べる
（Lemma \ref{spaces-more-morphisms-lemma-quasi-finite-separated-pass-through-finite-addendum}）。
有限型準連接イデアル
$\mathcal{I}_i \subset \mathcal{O}_{Y_i}$ で、
$V_i = Y_i \setminus V(\mathcal{I}_i)$ を満たすものを選べる
（Limits of Spaces, Lemma
\ref{spaces-limits-lemma-quasi-coherent-finite-type-ideals}）。
$V \to Y \to X$ を $V_i \to Y_i \to X_i$ の $X$ への基底変換とする。
$\mathcal{I} \subset \mathcal{O}_Y$ をイデアル $\mathcal{I}_i$ の引き戻しとする。
Lemma \ref{spaces-more-morphisms-lemma-resolution-property-one-sheaf} の容易な向きにより、
有限局所自由な $\mathcal{O}_X$-加群 $\mathcal{E}$ と全射
$\mathcal{E} \to \pi_*\mathcal{I}$ が存在する。
$\pi_i : Y_i \to X_i$ は有限かつ有限表示なので、
$\pi : Y \to X$ も有限かつ有限表示である。また、
$\mathcal{O}_{X_i}$-加群
$\pi_{i, *}\mathcal{O}_{Y_i}$ および
$\pi_{i, *}(\mathcal{O}_{Y_i}/\mathcal{I}_i)$ は有限表示であり、
$X$ へ引き戻すと
$\pi_*\mathcal{O}_Y$ および
$\pi_*(\mathcal{O}_Y/\mathcal{I})$ を与える。したがって Limits of Spaces, Lemma
\ref{spaces-limits-lemma-descend-modules-finite-presentation} により、
$i$ を増やせば、有限局所自由な $\mathcal{O}_{X_i}$-加群
$\mathcal{E}_i$ と射
$\mathcal{E}_i \to \pi_{i, *}\mathcal{O}_{Y_i}$ で、その $X$ への基底変換が合成
$\mathcal{E} \to \pi_*\mathcal{I} \to \pi_*\mathcal{O}_Y$ を復元するものを見つけられる。
有限表示 $\mathcal{O}_{X_i}$-加群
$\Coker(\mathcal{E}_i \to \pi_{i, *}\mathcal{O}_{Y_i})$ と
$\pi_{i, *}(\mathcal{O}_{Y_i}/\mathcal{I}_i)$ の $X$ への引き戻しは、
$\pi_*\mathcal{O}_Y$ の商として一致する。
したがって Limits of Spaces, Lemma
\ref{spaces-limits-lemma-descend-modules-finite-presentation} により、
これらが一致すると仮定してよい。言い換えると、
% P09-ERR-1234: frozen source uses O_{X_i} here although the preceding and following maps use pi_{i,*}O_{Y_i}; preserve literally and record for canon review.
$\mathcal{E}_i \to \pi_{i, *}\mathcal{O}_{X_i}$ の像は
$\pi_{i, *}\mathcal{I}_i$ に等しい。
すると Lemma \ref{spaces-more-morphisms-lemma-resolution-property-one-sheaf} により、
$X_i$ が分解性質をもつことが従う。
\end{proof}

\begin{lemma}
\label{spaces-more-morphisms-lemma-resolution-property-affine-diagonal}
$S$ をスキームとする。$X$ を分解性質をもつ準コンパクトかつ準分離的な代数空間とする。
このとき $X$ は $\mathbf{Z}$ 上アフィンな対角射をもつ
（Properties of Spaces, Definition
\ref{spaces-properties-definition-separated} の意味で）。
\end{lemma}

\begin{proof}
スキームの場合と同じように証明できるが、ここではスキームの場合から導く。
まず $S = \Spec(\mathbf{Z})$ と仮定してよいし、そうする。
次に、スキーム $Y$ と全射整射 $f : Y \to X$ を選ぶ。
Decent Spaces, Lemma
\ref{decent-spaces-lemma-there-is-a-scheme-integral-over} を参照せよ。
$f$ はアフィンなので、Derived Categories of Spaces, Lemma
\ref{spaces-perfect-lemma-resolution-property-goes-up-affine} により $Y$ は分解性質をもつ。
したがってスキームの場合により、スキーム $Y$ はアフィンな対角射をもつ。
Derived Categories of Schemes, Lemma
\ref{perfect-lemma-resolution-property-affine-diagonal} を参照せよ。
次の可換図式を考える。
$$
\xymatrix{
Y \ar[d] \ar[rr]_{\Delta_Y} & & Y \times_{\mathbf{Z}} Y \ar[d] \\
X \ar[rr]^{\Delta_X} & & X \times_{\mathbf{Z}} X
}
$$
右の鉛直射は整であり、特にアフィンであることに注意する。
$W \to X \times_{\mathbf{Z}} X$ を、$W$ がアフィンである射とする。このとき
$$
Y \times_{X \times_{\mathbf{Z}} X} W =
Y \times_{\Delta_Y, Y \times_{\mathbf{Z}} Y}
(Y \times_{\mathbf{Z}} Y) \times_{X \times_{\mathbf{Z}} X} W
$$
はアフィンである。一方、$Y \to X$ は整かつ全射なので、
$$
Y \times_{X \times_{\mathbf{Z}} X} W
\longrightarrow
X \times_{X \times_{\mathbf{Z}} X} W
$$
は $Y \to X$ の $W$ への基底変換として整かつ全射である。
したがって Limits of Spaces, Proposition
\ref{spaces-limits-proposition-affine} により、この射の終域はアフィンである。
ゆえに $\Delta_X$ はアフィンであり、望む結論を得る。
\end{proof}







\section{ブローアップと分解性質}
\label{spaces-more-morphisms-section-resolution-property-by-blowing-up}

\noindent
ブローアップ後に分解性質が成り立つことを証明する。

\begin{lemma}
\label{spaces-more-morphisms-lemma-blow-up-resolution-property}
$S$ をスキームとする。$X$ を $S$ 上の準コンパクトかつ準分離的な代数空間とする。
$|X|$ は有限個の既約成分をもつと仮定する。このとき、稠密な準コンパクト開部分空間
$U \subset X$ と $U$-許容ブローアップ $X' \to X$ が存在し、
代数空間 $X'$ は分解性質をもつ。
\end{lemma}

\begin{proof}
Limits of Spaces, Lemma
\ref{spaces-limits-lemma-there-is-a-scheme-finite-over-refined} により、
全射、有限かつ有限表示な射 $f : Y \to X$ で $Y$ がスキームとなるもの、および
準コンパクトな稠密開部分空間 $U \subset X$ で
$f^{-1}(U) \to U$ が有限エタールとなるものが存在する。
More on Morphisms, Lemma
\ref{more-morphisms-lemma-blow-up-ample-family} により、
準コンパクトな稠密開部分空間 $V \subset Y$ と $V$-許容ブローアップ
$Y' \to Y$ が存在し、$Y'$ は可逆加群の豊富族をもつ。
$U$ を縮小すれば $f^{-1}(U) \subset V$ と仮定できる（詳細は省略する）。
したがって $f' : Y' \to X$ は $U$ 上有限エタールであり、特に射
$(f')^{-1}(U) \to U$ は有限局所自由である。
Lemma \ref{spaces-more-morphisms-lemma-finite-after-blowing-up} により、$U$-許容ブローアップ
$X' \to X$ が存在し、$Y''$、すなわち $Y'$ の狭義変換は $X'$ 上有限局所自由である。図式は
$$
\xymatrix{
Y'' \ar[d] \ar[r]_g &
Y' \ar[r] &
Y \ar[d] \\
X' \ar[rr] & & X
}
$$
である。$g : Y'' \to Y'$ は $X' \to X$ の中心の逆像におけるブローアップなので
（Divisors on Spaces, Lemma
\ref{spaces-divisors-lemma-strict-transform}）、$g : Y'' \to Y'$ は射影的であり、
ある $g$-豊富な可逆加群が $Y''$ 上に存在する。
したがって More on Morphisms, Lemma
\ref{more-morphisms-lemma-ample-family-ample-relative} により、
$Y''$ は可逆加群の豊富族をもつ。ゆえに $Y''$ は分解性質をもつ。
Derived Categories of Schemes, Lemma
\ref{perfect-lemma-resolution-property-ample-family} を参照せよ。
最後に Derived Categories of Spaces, Lemma
\ref{spaces-perfect-lemma-resolution-property-goes-down-finite-flat} により、
$X'$ は分解性質をもつと結論する。
\end{proof}

\begin{lemma}
\label{spaces-more-morphisms-lemma-blow-up-resolution-property-filtered}
$S$ をスキームとする。$X$ を $S$ 上の準コンパクトかつ準分離的な代数空間とする。
このとき $t \geq 0$ と閉部分空間
$$
X \supset Z_0 \supset Z_1 \supset \ldots \supset Z_t = \emptyset
$$
で、$Z_i \to X$ が有限表示、$Z_0 \subset X$ が肥厚化となるものが存在する。
さらに各 $i = 0, \ldots t - 1$ に対し、
% P09-ERR-1235: frozen source says Z_i minus Z_{i-1}; i=0 makes this undefined and the proof uses Z_{1,i} minus Z_{1,i+1}; preserve literally and record for canon review.
$(Z_i \setminus Z_{i - 1})$-許容ブローアップ $Z'_i \to Z_i$ で、
$Z'_i$ が分解性質をもつものが存在する。
\end{lemma}

\begin{proof}
この段落では絶対ネーター近似を用い、$\Spec(\mathbf{Z})$ 上有限表示な代数空間の場合へ
帰着する。$X$ を $\Spec(\mathbf{Z})$ 上の代数空間とみなしてよい。
Spaces, Definition \ref{spaces-definition-base-change} および
Properties of Spaces, Definition
\ref{spaces-properties-definition-separated} を参照せよ。
したがって Limits of Spaces, Proposition
\ref{spaces-limits-proposition-approximate} を適用できる。
この結果、アフィン射 $X \to X_0$ で、$X_0$ が $\mathbf{Z}$ 上有限表示となるものを
見つけられる。$X_0$ について補題を証明できれば、層化とブローアップの中心を
$X$ へ引き戻して $X$ についての結果を得る。この議論では、分解性質がアフィン射に沿って
上へ移ること（Derived Categories of Spaces, Lemma
\ref{spaces-perfect-lemma-resolution-property-goes-up-affine}）、および
アフィン射の狭義変換がアフィンであることを用いる。詳細は省略する。
これで次の段落で扱う場合へ帰着される。

\medskip\noindent
$X$ が $\mathbf{Z}$ 上有限表示であると仮定する。
すると $X$ はネーターであり、$|X|$ は有限次元のネーター位相空間
（有限個の既約成分をもつ）である。したがって $\dim(|X|)$ に関する帰納法を使える。
Lemma \ref{spaces-more-morphisms-lemma-blow-up-resolution-property} により、稠密開部分空間
$U \subset X$ と $U$-許容ブローアップ $X' \to X$ が存在し、
$X'$ は分解性質をもつ。
$Z_0 = X$ とおき、$Z_1 \subset X$ を
$|Z_1| = |X| \setminus |U|$ を満たす被約閉部分空間とする。
帰納法により整数 $t \geq 0$ と閉部分空間によるフィルトレーション
$$
Z_1 \supset Z_{1, 0} \supset Z_{1, 1} \supset \ldots
\supset Z_{1, t} = \emptyset
$$
を得る。ここで $Z_{1, 0} \to Z_1$ は肥厚化であり、
$(Z_{1, i} \setminus Z_{1, i + 1})$-許容ブローアップ
$Z'_{1, i} \to Z_{1, i}$ で、$Z'_{1, i}$ が分解性質をもつものが存在する。
$Z_1$ は被約なので $Z_1 = Z_{1, 0}$ である。
したがって $Z_i = Z_{1, i - 1}$ および $Z'_i = Z'_{1, i - 1}$ とおけばよい。
ここで $i \geq 1$ である。これで補題が証明される。
\end{proof}
